\documentclass[11pt,colorlinks]{charlie}
\usepackage{booktabs,array}
\newcolumntype{L}[1]{>{\raggedright\arraybackslash}p{#1}}
\setlength{\emergencystretch}{3em}
\usepackage[capitalize,nameinlink]{cleveref}
\usepackage{aliascnt}
\makeatletter
\def\ch@redefthm#1#2#3{\expandafter\let\csname #1\endcsname\relax
  \expandafter\let\csname end#1\endcsname\relax
  \newaliascnt{#1}{theorem}\newtheorem{#1}[#1]{#2}\aliascntresetthe{#1}%
  \crefname{#1}{#2}{#3}\Crefname{#1}{#2}{#3}}
\crefname{theorem}{Theorem}{Theorems}\Crefname{theorem}{Theorem}{Theorems}
\theoremstyle{chplain}
\ch@redefthm{lemma}{Lemma}{Lemmas}
\ch@redefthm{proposition}{Proposition}{Propositions}
\ch@redefthm{corollary}{Corollary}{Corollaries}
\ch@redefthm{conjecture}{Conjecture}{Conjectures}
\theoremstyle{chdefinition}
\ch@redefthm{definition}{Definition}{Definitions}
\ch@redefthm{example}{Example}{Examples}
\ch@redefthm{hypothesis}{Hypothesis}{Hypotheses}
\makeatother

\newcommand{\R}{\mathbb{R}}
\newcommand{\C}{\mathbb{C}}
\newcommand{\Z}{\mathbb{Z}}
\newcommand{\Q}{\mathbb{Q}}
\newcommand{\N}{\mathbb{N}}
\newcommand{\E}{\mathbb{E}}
\newcommand{\one}{\mathbf{1}}
\newcommand{\dd}{\,\mathrm{d}}

\newif\ifdraftmarkers \draftmarkersfalse
\newcounter{gapcount}
\newcommand{\gap}[1]{\ifdraftmarkers\stepcounter{gapcount}%
  \par\smallskip\noindent\fcolorbox{red!60!black}{red!5}{\parbox{\dimexpr\linewidth-2\fboxsep-2\fboxrule}{%
  \textcolor{red!60!black}{\textbf{Gap \thegapcount.}} #1}}\par\smallskip\fi}
\newcommand{\unchecked}[1]{\ifdraftmarkers%
  \par\smallskip\noindent\fcolorbox{orange!80!black}{orange!6}{\parbox{\dimexpr\linewidth-2\fboxsep-2\fboxrule}{%
  \textcolor{orange!80!black}{\textbf{Unchecked.}} #1}}\par\smallskip\fi}
\newcommand{\source}[1]{\ifdraftmarkers{\textcolor{gray}{\footnotesize[source: \texttt{#1}]}}\fi}
\newcommand{\todo}[1]{\ifdraftmarkers{\textcolor{blue!70!black}{\footnotesize[\textbf{todo:} #1]}}\fi}

\newcommand{\coderepo}{\url{https://github.com/charlieoneill11/maths-releases}}
\newcommand{\aiusestatement}{The mathematical arguments and computer programs in this paper were developed by Astra (GPT-6.2 Astra, OpenAI). Fable (Fable~5.2, Anthropic) worked throughout as a second model: it proposed alternative approaches, reviewed arguments and raised objections. The manuscript was prepared from the project's research records with Claude (Anthropic). All computations were re-run before release, and the programs and certificates are available at \coderepo.}
\newcommand{\responsibilitystatement}{Charles O'Neill has checked the arguments and takes full responsibility for the correctness of the results.}
\newcommand{\mscandkeywords}[2]{\par\medskip\noindent{\small\textit{2020 Mathematics Subject Classification.} #1\par\noindent\textit{Keywords.} #2\par}}

\usepackage{mathrsfs}

\newcommand{\Om}{\Omega}
\newcommand{\ip}[2]{\langle #1,#2\rangle}
\newcommand{\F}{\mathbb{F}}

\title{An improved lower bound for the sixth moment of the Riemann zeta function}
\author{Charles O'Neill, Max Kirkby, Astra and Fable}
\date{24 November 2026}

\begin{document}
\maketitle

\begin{abstract}
Let $M_6(T)=\int_0^T|\zeta(\tfrac12+it)|^6\dd t$. The conjecture of Conrey and Ghosh predicts $M_6(T)\sim 42\,c_3\,T(\log T)^9$, where $c_3=a_3/9!$ and $a_3$ is the usual arithmetic factor. We prove unconditionally that
\[
 \liminf_{T\to\infty}\frac{M_6(T)}{c_3T(\log T)^9}\ \ge\ 41.6406029094\ldots,
\]
improving the previous record $41.008$ of Fraccaro. The argument places every Cauchy--Schwarz test in a common Hilbert space of Schur profiles on the simplex $\{0<x<y<z<1\}$, where the target $\zeta^3$ has mass exactly $42$. A lower bound then becomes the measure of the region on which physical pairings of test Dirichlet polynomials can be shown to agree with the unitary model. We first determine exactly what the two classical test families can give, namely $264076794201/6400000000=41.2619990939\ldots$. We then enlarge the admissible region using the Kuznetsov formula for marked short divisors with the level-averaged spectral large sieves of Deshouillers--Iwaniec and Watt, Heath-Brown's estimate for $d_3$ in arithmetic progressions, and, for moduli up to $T^{3/5}$, a dual-sum estimate for hyper-Kloosterman sums based on the sums-of-products theorem of Fouvry, Kowalski and Michel, combined with the bounds of Deshouillers and Iwaniec for sums of Kloosterman fractions and with the Kuznetsov formula at the level fixed by the denominators. We also show that enlarging the supports within the classes used here (the dominant shapes, the extreme traces, the shapes with a short coordinate, and the balanced and mixed-row regions, with their reflections) cannot take the projection beyond $41.9462895683\ldots<42$. The exact masses, this ceiling, the covering argument for the mixed row and the deduction of the bound from the arithmetic pairings have been checked in Lean~4 with Mathlib.
\end{abstract}
\mscandkeywords{11M06; 11F72, 11L05, 11N37}{Riemann zeta function, sixth moment, mean values, Kuznetsov formula, spectral large sieve}

\section*{AI use and author contributions}
\aiusestatement

Open-weight models, Kimi~K3.1~Pro (Moonshot AI) and DeepSeek~V4.1~Flash (DeepSeek), were able to derive a significant proportion of these results when given guidance. We will report the details separately.

Charles O'Neill proposed the problem because the Riemann hypothesis, now a theorem \cite{Alpoge}, gives Harper's bound $M_6(T)\ll T(\log T)^9$ \cite{Harper}, which controls the error terms of the Cauchy--Schwarz method at the correct order. His technical contribution was to ask for the test-function method to be set up as an optimisation problem and, if it fell short of $42$, for its ceiling to be proved. This led to the exact value in \Cref{thm:twofamily} and to the ceiling in \Cref{prop:ceiling}. Max Kirkby set up and ran the experiments with open-weight models. \responsibilitystatement

\section{Introduction}

\subsection{The sixth moment}
For $k\ge1$ let $M_{2k}(T)=\int_0^T|\zeta(\tfrac12+it)|^{2k}\dd t$. Hardy and Littlewood evaluated $M_2$ and Ingham evaluated $M_4$; no asymptotic is known for $k\ge3$. Conrey and Ghosh conjectured \cite{CG}, and the random-matrix model of Keating and Snaith \cite{KS} and the recipe of Conrey, Farmer, Keating, Rubinstein and Snaith \cite{CFKRS} independently predict,
\begin{equation}\label{eq:conj}
 M_6(T)\sim 42\,c_3\,T(\log T)^9,\qquad
 c_3=\frac{a_3}{9!},
\end{equation}
where
\[
 a_3=\prod_p\Big[(1-p^{-1})^9\sum_{j\ge0}\binom{j+2}{2}^{2}p^{-j}\Big].
\]
Throughout we normalise every coefficient by $c_3T(\log T)^9$, so the conjectured value is $42$.

\subsection{Previous bounds}
The order of magnitude $T(\log T)^9$ is known from below \cite{RS,HS}. Harper's upper bound $M_6(T)\ll T(\log T)^9$ \cite{Harper} rests on the Riemann hypothesis, which is now a theorem of Alp\"oge \cite{Alpoge}, so the order of magnitude is known unconditionally. Ng \cite{Ng} derived the full asymptotic conditionally on a conjecture for ternary additive divisor sums. For the leading constant, the lower-bound method of Conrey and Ghosh pairs $\zeta^3$ with a test Dirichlet polynomial and applies Cauchy--Schwarz. Durkan and Page \cite{DP} obtained $34.4012$. Fraccaro \cite{Fraccaro}, using two families of tests, obtained $41.008$; an exact rational evaluation of Fraccaro's $250$ rows gives $41.0089916904\ldots$, and an expanded $397$-row certificate gives $41.0298094152\ldots$ (see the remark at the end of \Cref{sec:twofamily}).

\subsection{Results}
\begin{theorem}[Two families]\label{thm:twofamily}
Let $\mathcal F$ be the two test families of Fraccaro, with smooth length profiles and centred derivatives of arbitrary fixed polynomial degree. The supremum of the leading Cauchy--Schwarz coefficients over $\mathcal F$ is
\[
 C_*=\frac{264076794201}{6400000000}=41.26199909390625 .
\]
In particular $\liminf_{T\to\infty}M_6(T)/(c_3T(\log T)^9)\ge C_*$, and $C_*$ is an upper bound for the coefficient certified by every finite system of tests from $\mathcal F$. Restricted to derivative degree at most four, the supremum is strictly smaller than $C_*$.
\end{theorem}

\begin{theorem}[Main theorem]\label{thm:main}
Unconditionally,
\[
 \liminf_{T\to\infty}\frac{M_6(T)}{c_3T(\log T)^9}\ \ge\ C=41.64060290940640373015\ldots,
\]
where
\[
 C=\frac{112666832560750258461501594014215737455802061}{2705696476246250148185796472625838239963136}.
\]
Here $C=\mu(U_7)$ is the mass of the explicit region $U_7$ of \Cref{thm:mixedrow}.
\end{theorem}

\begin{proposition}[Limit of the method]\label{prop:ceiling}
Suppose arithmetic pairings were available for every Schur profile supported on the dominant shapes $x+y<z$, on all shapes of trace below $1$ or above $2$, on all shapes with $x<1/4$ or $z>3/4$, on the balanced and mixed-row regions $\mathcal R$ and $\mathcal F^+$ of \Cref{sec:balanced,sec:mixedrow}, and on the reflections of all of these. Projection onto that span would still give at most
\[
 \frac{13860404849860091}{330432202525696}=41.946289568378975\ldots<42 .
\]
\end{proposition}

The structure of the argument is the following. In \Cref{sec:schur} we put every test in a common Hilbert space $L^2(\Om)$ of Schur profiles, in which $\zeta^3$ corresponds to the explicit function $v$ with $9!\|v\|^2=42$. This identification is proved test by test, from arithmetic evaluations; it does not assume the unknown norm of $\zeta^3$. A lower bound is then $9!\|P_{\mathcal H}v\|^2$, the squared projection onto the span $\mathcal H$ of the profiles of admissible tests. When $\mathcal H$ contains all profiles supported on an open reflection-invariant set $U$, the bound is $\mu(U)$ for the measure $d\mu=9!\,v^2\dd x\dd y\dd z$. \Cref{thm:twofamily} computes this exactly for the two classical families. \Cref{thm:main} enlarges $U$ by proving new arithmetic pairings, in five steps:
\begin{center}
\small
\begin{tabular}{@{}L{0.30\linewidth}L{0.40\linewidth}r@{}}
\toprule
New tests & Input & $\mu$ of the region \\
\midrule
Short cubic, unequal lengths (\Cref{thm:third}) & Fourth moment of Dirichlet $L$-functions & $41.3163069650$ \\
Marked short divisors (\Cref{thm:marked}) & Kuznetsov formula, spectral large sieves & $41.5917968750$ \\
Balanced row (\Cref{thm:row}) & $d_3$ in progressions \cite{HB} & $41.5958462794$ \\
Factorable row (\Cref{thm:factorable}) & Hyper-Kloosterman sums, sums of products \cite{FKM} & $41.6117132746$ \\
Mixed row (\Cref{thm:mixedrow}) & Kloosterman fractions \cite{DI}, Kuznetsov at fixed level & $41.6406029094$ \\
\bottomrule
\end{tabular}
\end{center}
The fifth step, \Cref{thm:mixedrow}, keeps the same witnesses as \Cref{thm:factorable} but proves the four-shift row for moduli $q\asymp T^\kappa$ with $\kappa$ up to $\frac35$, by averaging over the moduli instead of bounding each one. At each step the new family must also be paired with all earlier ones.

\emph{Hypotheses.} The proofs of \Cref{thm:twofamily}, \Cref{thm:third} and \Cref{prop:ceiling} use neither the Riemann hypothesis nor the generalised Riemann hypothesis. The proofs of \Cref{thm:marked}, \Cref{thm:row}, \Cref{thm:factorable}, \Cref{thm:mixedrow} and \Cref{thm:main} use Harper's bound $\int_T^{2T}|\zeta(\frac12+it)|^6\dd t\ll T(\log T)^9$ \cite{Harper}, through \cref{lem:sixth} and \cref{lem:sourcenorm}. Harper's bound rests on the Riemann hypothesis, which is a theorem of Alp\"oge \cite{Alpoge}; this is the sense in which \Cref{thm:main} is unconditional. The four-shift estimates \cref{prop:fourshift,prop:fourshiftplus} themselves use neither hypothesis: their inputs are the Weil and Deligne bounds, the theorem of Fouvry, Kowalski and Michel, the Kuznetsov formula and the spectral large sieve, and the Kim--Sarnak bound for exceptional eigenvalues \cite{KimSarnak}; \cref{prop:fourshiftplus} also uses the bound of Deshouillers and Iwaniec for sums of Kloosterman fractions \cite[Theorem~12]{DI}, with the correction of \cite{Lichtman}, and Pascadi's estimate for the exceptional spectrum \cite{Pascadi}. The generalised Riemann hypothesis for Dirichlet $L$-functions is not used anywhere.

\Cref{sec:main} assembles the result, and \Cref{sec:ceiling} shows that enlarging supports in this way cannot reach $42$. \Cref{sec:lean} describes the parts of the proof that have been checked in Lean. The computations quoted in the text are carried out in exact rational arithmetic unless stated otherwise, and the programs are available at \coderepo.

\section{Schur profiles and the projection principle}\label{sec:schur}

\subsection{The chamber and the target}
Throughout, $L=\log T$ and $s=\tfrac12+it$. Let
\begin{gather*}
 \Om=\{(x,y,z)\in\R^3:0<x<y<z<1\},\\ v(x,y,z)=\tfrac12(z-y)(z-x)(y-x),\qquad J(x,y,z)=(1-z,1-y,1-x),
\end{gather*}
and let $\mu$ be the measure $d\mu=9!\,v^2\dd x\dd y\dd z$ on $\Om$. We write $k=x+y$ and $r=x+y+z$ (the \emph{trace}), and $\ip fg=\int_\Om fg$ for the inner product of $L^2(\Om)$. A function $f$ on $\Om$ is \emph{$J$-symmetric} if $f\circ J=f$. For a measurable $A\subseteq\Om$ we write $L^2_J(A)$ for the space of $J$-symmetric $f\in L^2(\Om)$ vanishing almost everywhere off $A$.

\begin{lemma}\label{lem:chamber}
$J$ is a measure-preserving involution of $\Om$, and $v\circ J=v$. Moreover $\int_\Om v^2=42/9!$, that is, $\mu(\Om)=42$.
\end{lemma}
\begin{proof}
$J$ is affine with linear part $-\mathrm{id}$ composed with the transposition of the first and third coordinates, so $|\det J|=1$, $J\circ J=\mathrm{id}$, and $0<1-z<1-y<1-x<1$ shows $J(\Om)=\Om$. The differences $(z-y,z-x,y-x)$ are mapped by $J$ to $(y-x,z-x,z-y)$, so $v\circ J=v$. For the integral, $2v=\Delta(x,y,z):=(z-y)(z-x)(y-x)$ changes at most by a sign under permutations of the coordinates, so $\int_\Om\Delta^2=\frac16\int_{[0,1]^3}\Delta^2$. Expanding both Vandermonde determinants (Heine's identity),
\[
 \int_{[0,1]^3}\Delta^2=3!\,\det\Big(\int_0^1t^{i+j}\dd t\Big)_{i,j=0}^{2}=6\det\Big(\frac1{i+j+1}\Big)_{i,j=0}^{2}=\frac{6}{2160}.
\]
Hence $\int_\Om v^2=\frac14\cdot\frac16\cdot\frac6{2160}=\frac1{8640}=\frac{42}{9!}$.
\end{proof}

\subsection{Symmetric functions on the unitary group}
Let $N\ge1$, let $U$ be Haar-distributed on $\mathrm U(N)$, and put $\ip FG_N=\E\,F(U)\overline{G(U)}$ for class functions $F,G$. Define $e_i(U)$ by $\det(I+uU)=\sum_{i=0}^Ne_i(U)u^i$, with $e_i=0$ for $i\notin[0,N]$, and put $D=\det U$. For a partition $\lambda$ with $\ell(\lambda)\le N$ let $s_\lambda(U)$ be the Schur polynomial in the eigenvalues of $U$, and let $\lambda'$ be the conjugate partition. We use the following standard facts \cite{Mac,BG}.
\begin{enumerate}[label=(S\arabic*)]
\item\label{S:orth} The $s_\lambda$ with $\ell(\lambda)\le N$ are orthonormal for $\ip\cdot\cdot_N$.
\item\label{S:pieri} (Pieri) $e_ms_\nu=\sum_\lambda s_\lambda$, the sum over $\lambda\supseteq\nu$ with $\ell(\lambda)\le N$ such that $\lambda/\nu$ is a vertical strip of $m$ boxes; equivalently $\lambda'_1\ge\nu'_1\ge\lambda'_2\ge\nu'_2\ge\cdots$ and $|\lambda|=|\nu|+m$ \cite[I.(5.17)]{Mac}.
\item\label{S:cauchy} (Dual Cauchy) $\prod_{k=1}^N(1+u_k)^3=\sum_\lambda s_\lambda(u)s_{\lambda'}(1,1,1)$, where $s_{\lambda'}(1,1,1)=0$ unless $\lambda_1\le3$, and for $\lambda'=(Z,Y,X)$ Weyl's dimension formula gives $s_{\lambda'}(1,1,1)=d_\lambda:=\tfrac12(Z-Y+1)(Z-X+2)(Y-X+1)$ \cite[I.(4.3$'$)]{Mac}.
\item\label{S:jt} (Dual Jacobi--Trudi) $s_\lambda=\det(e_{\lambda'_i-i+j})_{1\le i,j\le\ell(\lambda')}$ \cite[I.(3.5)]{Mac}.
\item\label{S:dual} (Duality) $e_i=D\,\overline{e_{N-i}}$ for all $i$. If $\lambda_1\le3$ and $\lambda'=(Z,Y,X)$, then $D^3\overline{s_\lambda}=s_{\lambda^*}$, where $(\lambda^*)'=(N-X,N-Y,N-Z)$.
\item\label{S:stable} (Stable range) If $m,m'\in\Z_{\ge0}^3$ and $|m|=|m'|\le N$, then $\ip{e_{m_1}e_{m_2}e_{m_3}}{e_{m'_1}e_{m'_2}e_{m'_3}}_N$ is the number of $3\times3$ matrices of non-negative integers with row sums $m$ and column sums $m'$.
\end{enumerate}
For \ref{S:dual}, if the eigenvalues are $\omega_1,\dots,\omega_N$ then $e_{N-i}(\omega)=\prod_k\omega_k\cdot e_i(\bar\omega)=D\,\overline{e_i(\omega)}$; the second statement is the identity $\det^3\cdot s_\lambda(\omega^{-1})=s_{(3-\lambda_N,\dots,3-\lambda_1)}(\omega)$, whose conjugate has the stated columns. For \ref{S:stable}, every $\lambda$ with $|\lambda|\le N$ has $\ell(\lambda)\le N$, so by \ref{S:orth} the left side is the Hall inner product $\ip{e_m}{e_{m'}}$ in the ring of symmetric functions; the involution $\omega$ is an isometry with $\omega e_m=h_m$, and $\ip{h_m}{h_{m'}}$ is the stated count \cite[\S\S7.5, 7.9]{St}.

The Schur polynomials occurring below all have $\lambda_1\le3$. We index them by their column lengths $\lambda'=(Z,Y,X)$, $N\ge Z\ge Y\ge X\ge0$, and write $\lambda/N=(X/N,Y/N,Z/N)\in\overline\Om$. By \ref{S:dual}, $\lambda^*/N=J(\lambda/N)$.

\subsection{Cubic weights and their profiles}
For a bounded Borel function $F$ on $[0,1]^3$ put
\[
 E_N[F]=\sum_{m\in\{0,\dots,N\}^3}F(m/N)\,e_{m_1}e_{m_2}e_{m_3},\qquad c_\lambda(F)=\ip{E_N[F]}{s_\lambda}_N .
\]
Thus $E_N[1]=\det(I+U)^3$. For $(x,y,z)\in\overline\Om$ and $c\in\R^3$ define the \emph{Pieri kernel}
\[
 \kappa(x,y,z;c)=\big[\min(z,\,c_1+c_2-x)-\max(y,\,c_1+c_2-y,\,c_1,\,c_2)\big]_+
\]
if $c\in[0,1]^3$ and $c_1+c_2+c_3=x+y+z$, and $\kappa=0$ otherwise. The \emph{profile} of $F$ is
\begin{equation}\label{eq:profile}
 h_F(x,y,z)=\int_0^1\!\!\int_0^1F(c)\,\kappa(x,y,z;c)\dd c_1\dd c_2,\qquad c_3=x+y+z-c_1-c_2 .
\end{equation}
For fixed $c$ the kernel is continuous in $(x,y,z)$ except where $c_3=r-c_1-c_2\in\{0,1\}$, a null set of $(c_1,c_2)$ for each $(x,y,z)$, and $0\le\kappa\le1$. If $F$ is continuous, dominated convergence shows that $h_F$ is continuous on $\overline\Om$. For a general bounded Borel $F$ the profile is bounded and measurable but need not be continuous: for $F(c)=\one_{c_1+c_2+c_3\in\Q}$ one gets $h_F=\one_{x+y+z\in\Q}\,v$ (by \cref{lem:profiles}(ii) below).

\begin{lemma}[Pieri multiplicities]\label{lem:pieri}
Let $K_\lambda(m)$ be the multiplicity of $s_\lambda$ in $e_{m_1}e_{m_2}e_{m_3}$, where $\lambda'=(Z,Y,X)$. Then $K_\lambda(m)=0$ unless $|m|=|\lambda|$ and $0\le m_3\le N$, in which case $K_\lambda(m)$ is the number of integers $u$ with
\[
 \max(Y,\,m_1+m_2-Y,\,m_1,\,m_2)\le u\le\min(Z,\,m_1+m_2-X).
\]
In particular $|K_\lambda(m)-N\kappa(\lambda/N;m/N)|\le1$ for all $m$, and $\sum_mK_\lambda(m)=d_\lambda$.
\end{lemma}
\begin{proof}
By \ref{S:pieri} applied to $e_{m_2}$ and the one-column shape $(1^{m_1})$, $e_{m_1}e_{m_2}$ is the sum of $s_\nu$ over the two-column shapes $\nu'=(u,w)$ with $u\ge m_1\ge w\ge0$ and $u+w=m_1+m_2$, that is, $\max(m_1,m_2)\le u\le m_1+m_2$ and $w=m_1+m_2-u$. Applying \ref{S:pieri} again, $s_\lambda$ occurs in $e_{m_3}s_\nu$ (once) if and only if $|\lambda|=|\nu|+m_3$, $0\le m_3\le N$ and $Z\ge u\ge Y\ge w\ge X$; the conditions $w\le Y$ and $w\ge X$ read $u\ge m_1+m_2-Y$ and $u\le m_1+m_2-X$. This gives the count. The number of integers in a real interval differs from the positive part of its length by at most one, and in scaled units the interval has length $N\kappa(\lambda/N;m/N)$. Finally $\sum_mK_\lambda(m)$ is the coefficient of $s_\lambda$ in $(\sum_ie_i)^3=\det(I+U)^3$, which is $d_\lambda$ by \ref{S:cauchy}.
\end{proof}

\begin{lemma}[Profiles]\label{lem:profiles}
\begin{enumerate}[label=\textup{(\roman*)}]
\item If $F$ is continuous, then $\sup_\lambda|N^{-3}c_\lambda(F)-h_F(\lambda/N)|\to0$ as $N\to\infty$.
\item $h_1=v$. Consequently $|h_F|\le\|F\|_\infty v$ for every bounded $F$, and $F\mapsto h_F$ is linear.
\item If $F,G$ are continuous and real-valued, then as $N\to\infty$
\begin{gather*}
 N^{-9}\ip{E_N[F]}{E_N[G]}_N\to\ip{h_F}{h_G},\\
 N^{-9}\ip{E_N[F]}{D^3\overline{E_N[G]}}_N\to\ip{h_F}{h_G\circ J}.
\end{gather*}
\end{enumerate}
\end{lemma}
\begin{proof}
(i) By \cref{lem:pieri}, $N^{-3}c_\lambda(F)=N^{-2}\sum_{m_1,m_2}F(m/N)\kappa(\lambda/N;m/N)+O(N^{-1})$, the sum over integers $0\le m_1,m_2\le N$ with $m_3=|\lambda|-m_1-m_2$, since only $O(N^2)$ terms are non-zero and each differs from its replacement by $O(N^{-3})$. The main term is a Riemann sum, with mesh $1/N$, for the integral \eqref{eq:profile} at $(x,y,z)=\lambda/N$. Its integrand is bounded by $\|F\|_\infty$; off the two lines $c_3\in\{0,1\}$ it has a modulus of continuity bounded by that of $F$ plus $2|\cdot|$, uniformly in $(x,y,z)$, because $\kappa$ is a positive part of a difference of a minimum and a maximum of affine functions with coefficients in $\{0,\pm1\}$. The $O(N)$ cells of the Riemann sum that meet the two lines contribute $O(N^{-1})$. Hence the error tends to $0$ uniformly in $\lambda$.

(ii) Apply (i) with $F=1$. Since $\sum_mK_\lambda(m)=d_\lambda$ and $N^{-3}d_\lambda=v(\lambda/N)+O(N^{-1})$, we get $h_1(\lambda/N)=v(\lambda/N)+o(1)$ uniformly. Both functions are continuous on $\overline\Om$ and the points $\lambda/N$ become dense, so $h_1=v$. As $\kappa\ge0$, $|h_F|\le\|F\|_\infty h_1$. Linearity is clear.

(iii) By \ref{S:orth}, $\ip{E_N[F]}{E_N[G]}_N=\sum_\lambda c_\lambda(F)c_\lambda(G)$, the sum over $\lambda'=(Z,Y,X)$ with $N\ge Z\ge Y\ge X\ge0$. There are $O(N^3)$ such $\lambda$, and by (i) and (ii) both $N^{-3}c_\lambda(\cdot)$ are bounded and equal $h_\cdot(\lambda/N)+o(1)$ uniformly. Hence $N^{-9}\sum_\lambda c_\lambda(F)c_\lambda(G)=N^{-3}\sum_\lambda h_F(\lambda/N)h_G(\lambda/N)+o(1)$, a Riemann sum for the continuous function $h_Fh_G$ on $\overline\Om$. For the second limit, \ref{S:dual} gives $D^3\overline{E_N[G]}=\sum_\lambda c_\lambda(G)s_{\lambda^*}$, so the inner product is $\sum_\lambda c_{\lambda^*}(F)c_\lambda(G)$, and $\lambda^*/N=J(\lambda/N)$. The same argument gives $\int_\Om h_F(J\cdot)h_G=\ip{h_F}{h_G\circ J}$ by \cref{lem:chamber}. For complex-valued continuous $F,G$ the same proof applies with $c_\lambda(G)$ replaced by $\overline{c_\lambda(G)}$, and the limits are $\int_\Om h_F\overline{h_G}$ and $\int_\Om h_F\,\overline{h_G\circ J}$.
\end{proof}

\begin{remark}
In particular $N^{-9}\|\det(I+U)^3\|_N^2\to\|v\|^2=42/9!$. This is only the motivation for the normalisation: the unitary model enters the proofs below solely through \cref{lem:profiles}, which computes inner products of explicit profiles.
\end{remark}

\subsection{Faithful families and the projection principle}
A \emph{test} is a family $B=(B_T)_{T\ge T_0}$ of continuous functions $B_T:\R\to\C$. Let $\mathcal W$ be the set of non-negative $\Phi\in C_c^\infty((1,2))$, $\Phi\not\equiv0$. For $\Phi\in\mathcal W$ and tests $B,B'$ put
\begin{gather*}
 \mathcal N_T=a_3\,TL^9\!\int\Phi,\qquad
 \mathrm N_T(B)=\int_\R\Phi(t/T)\,\zeta(s)^3\,\overline{B_T(t)}\dd t,\\
 \mathrm G_T(B,B')=\int_\R\Phi(t/T)\,B_T(t)\overline{B'_T(t)}\dd t .
\end{gather*}
\begin{definition}\label{def:faithful}
A \emph{family} is a set $\mathcal T$ of pairs $(B,h_B)$, where $B$ is a test and $h_B\in L^2(\Om)$ is real-valued (the \emph{profile} of $B$). It is \emph{faithful} if for every $\Phi\in\mathcal W$ and all $(B,h_B),(B',h_{B'})\in\mathcal T$,
\[
 \frac{\mathrm N_T(B)}{\mathcal N_T}\to\ip{v}{h_B},\qquad
 \frac{\mathrm G_T(B,B')}{\mathcal N_T}\to\ip{h_B}{h_{B'}}\qquad(T\to\infty).
\]
Its \emph{span} $\mathcal H(\mathcal T)$ is the closed linear span of $\{h_B\}$ in $L^2(\Om)$.
\end{definition}

If $\mathcal T$ is faithful, so is the family of all finite real linear combinations $(\sum c_iB_i,\sum c_ih_{B_i})$, by bilinearity. The union of two faithful families is faithful if and only if all \emph{cross-pairings} $\mathrm G_T(B,B')$, with $B$ from one family and $B'$ from the other, have the limits $\ip{h_B}{h_{B'}}$. This is why, in \Cref{sec:marked,sec:balanced,sec:factorable}, each new family is paired with every earlier one before its support is added.

\begin{proposition}[Projection principle]\label{prop:projection}
Let $\mathcal T$ be faithful and $\mathcal H=\mathcal H(\mathcal T)$. Then
\[
 \liminf_{T\to\infty}\frac{M_6(T)}{c_3TL^9}\ \ge\ 9!\,\|P_{\mathcal H}v\|^2 ,
\]
where $P_{\mathcal H}$ is orthogonal projection. If $\mathcal H\supseteq L^2_J(U)$ for a $J$-invariant measurable set $U\subseteq\Om$, then $9!\|P_{\mathcal H}v\|^2\ge\mu(U)$.
\end{proposition}
\begin{proof}
All profiles and all coefficients of combinations here are real. Put $C=9!\|P_{\mathcal H}v\|^2$. If $C=0$ there is nothing to prove, so assume $C>0$. Let $\epsilon>0$. Since $\ip{v}{h}=\ip{P_{\mathcal H}v}{h}$ for $h\in\mathcal H$, and the map $h\mapsto\ip{v}{h}^2/\|h\|^2$ is continuous at $h=P_{\mathcal H}v\ne0$ with value $\|P_{\mathcal H}v\|^2$, there is a finite real combination $h=\sum_ic_ih_{B_i}$ of profiles of $\mathcal T$ with $h\ne0$ and $9!\ip{v}{h}^2/\|h\|^2>C-\epsilon$. Put $B=\sum_ic_iB_i$. By Cauchy--Schwarz with the weight $\Phi(t/T)\ge0$,
\[
 |\mathrm N_T(B)|^2\le\int\Phi(t/T)|\zeta(s)|^6\dd t\cdot\mathrm G_T(B,B).
\]
By faithfulness $\mathrm N_T(B)/\mathcal N_T\to\ip vh$ and $\mathrm G_T(B,B)/\mathcal N_T\to\|h\|^2>0$. As $\mathcal N_T=9!\,c_3TL^9\int\Phi$, we obtain, for $T\ge T_1(\epsilon,\Phi)$,
\begin{equation}\label{eq:smoothbound}
 \int\Phi(t/T)|\zeta(\tfrac12+it)|^6\dd t\ \ge\ (C-2\epsilon)\,c_3\,TL^9\!\int\Phi .
\end{equation}
Now fix an integer $K\ge1$ and choose $\Phi\in\mathcal W$ with $\Phi\le\one_{(1,2)}$ and $\int\Phi\ge1-\epsilon$. For $T$ large, applying \eqref{eq:smoothbound} at the $K$ heights $T_k=2^{-k}T$,
\[
 M_6(T)\ \ge\ \sum_{k=1}^{K}\int_{T_k}^{2T_k}|\zeta(\tfrac12+it)|^6\dd t\ \ge\ (C-2\epsilon)(1-\epsilon)\,c_3\sum_{k=1}^K2^{-k}T(\log T_k)^9 .
\]
Since $\log T_k/\log T\to1$ for each fixed $k$, the $\liminf$ of $M_6(T)/(c_3TL^9)$ is at least $(C-2\epsilon)(1-\epsilon)(1-2^{-K})$. Letting $\epsilon\to0$ and $K\to\infty$ proves the first claim.

For the second, if $\mu(U)=0$ the claim follows from $\|P_{\mathcal H}v\|^2\ge0$, so assume $\mu(U)>0$. Then $v\one_U$ is $J$-symmetric because $v\circ J=v$ and $J(U)=U$, so $v\one_U\in\mathcal H$ and $\|P_{\mathcal H}v\|\ge\ip{v}{v\one_U}/\|v\one_U\|=\|v\one_U\|$. Thus $9!\|P_{\mathcal H}v\|^2\ge9!\int_Uv^2=\mu(U)$.
\end{proof}

No step uses the size of $\int|\zeta|^6$: the profile of each test is identified from its own arithmetic evaluation, and the target enters only through the numerators $\mathrm N_T(B)$. Every finite combination is fixed before $T\to\infty$.

\section{The two classical families: proof of \Cref{thm:twofamily}}\label{sec:twofamily}

\subsection{The families}
For a real function $f$ and a real polynomial $p$ put
\begin{gather*}
 A_f(s)=\sum_{n\ge1}f\Big(\frac{\log n}{L}\Big)n^{-s},\qquad
 C_p(s)=p\Big(\frac12+\frac1L\frac{d}{ds}\Big)\zeta(s),\\
 \Pi H=\tfrac12\big(H+\chi(s)^3\overline H\big),
\end{gather*}
where $\zeta(s)=\chi(s)\zeta(1-s)$. For $s=\frac12+it$, $\overline{\zeta(s)^3}\chi(s)^3=\zeta(s)^3$, so $\mathrm N_T(\Pi H)=\operatorname{Re}\mathrm N_T(H)$.

\emph{$U$-tests.} Data $(a_0,b_0,f,g,p)$ with $0<a_0\le b_0$, $3a_0+5b_0<2$, real $f\in C_c^\infty((0,a_0))$, $g\in C_c^\infty((0,b_0))$ and a real polynomial $p$. The test is $\Pi U$ with $U=\chi(s)C_p(s)A_f(s)A_g(\bar s)$, its cubic weight is $F_U(c)=f(c_1)\,g(1-c_2)\,p(\tfrac12-c_3)$, and its profile is $\frac12(h_{F_U}+h_{F_U}\circ J)$.

\emph{$V$-tests.} Data $(\eta,f,g,p)$ with $0<\eta<\frac14$, real polynomials $f,g$ on $[0,\eta]$ vanishing at $0$ and $\eta$ and extended by $0$, and a real polynomial $p$. The test is $\Pi V$ with $V=C_p(s)A_f(s)A_g(s)$, its cubic weight is $F_V(c)=f(c_1)g(c_2)p(\tfrac12-c_3)$, and its profile is $\frac12(h_{F_V}+h_{F_V}\circ J)$.

Let $\mathcal T_2$ be the family of all these pairs. The cubic weights record the matrix analogues: by \ref{S:dual}, $E_N[F_U]=\sum_{i,A,B}p(\tfrac12-\tfrac iN)f(\tfrac AN)g(\tfrac BN)\,e_i\,e_A\,D\overline{e_B}$ mirrors $U$, with $D$ in the role of $\chi$, and $\Pi$ corresponds to $h\mapsto\frac12(h+D^3\overline h)$, which acts on profiles by $J$ (\cref{lem:profiles}(iii)).

For $0\le\alpha,\beta\le1$ and a polynomial $p$ let
\[
 H_{\alpha,\beta,p}(x,y,z)=p(\tfrac12-j)\,\kappa(x,y,z;\alpha,\beta,j),\qquad j=x+y+z-\alpha-\beta,
\]
be the \emph{point feature}. By Fubini, the raw profiles are
\[
 h_{F_U}=\iint f(\alpha)g(\beta)H_{\alpha,1-\beta,p}\dd\alpha\dd\beta,\qquad h_{F_V}=\iint f(\alpha)g(\beta)H_{\alpha,\beta,p}\dd\alpha\dd\beta.
\]

\subsection{The arithmetic input}
We define three kernels. Let $X=(X_0,\dots,X_3)$ be uniformly distributed on the simplex $\{X\in[0,1]^4:\sum X_i=1\}$, and for $\varrho\in\R^4$ put $\Delta(\varrho)=\prod_{0\le i<k\le3}(\varrho_i-\varrho_k)$. For $\alpha,\beta\in[0,1]$ put $\varrho_w=(1,\alpha+\beta-w,w,0)$ and
\begin{equation}\label{eq:Kkernel}
 \mathcal K_p(\alpha,\beta)=\int_{\max(0,\alpha+\beta-1)}^{\min(\alpha,\beta)}\frac{\Delta(\varrho_w)}{12}\,\E\,p\Big(\sum_{i=0}^3(\varrho_{w,i}-\tfrac12)X_i\Big)\dd w.
\end{equation}
For $a,b,c,d\ge0$ with $a+b+c+d<1$ and polynomials $p,q$ let
\begin{equation}\label{eq:Bkernel}
 \mathcal B_{p,q}(a,b;c,d)=\int_{\mathcal R_0}\int_{-w}^{w}p(\delta-h)\,q(-\delta-h)\dd h\dd\varrho,
\end{equation}
where $\mathcal R_0=\{\varrho\in[0,\infty)^{2\times2}:\varrho_{11}+\varrho_{12}\le a,\ \varrho_{21}+\varrho_{22}\le b,\ \varrho_{11}+\varrho_{21}\le c,\ \varrho_{12}+\varrho_{22}\le d\}$, $R=a+b-|\varrho|$, $S=c+d-|\varrho|$, $w=\frac12(1-R-S)$ and $\delta=\frac12(R-S)$. For $c\in[0,\infty)^3$ and $0\le\sigma\le|c|<1$ let
\begin{equation}\label{eq:Mkernel}
 \mathcal M_{p,q}(c;\sigma)=\int_{\Sigma}\int_{\prod_k[-\varrho_k/2,\varrho_k/2]}p\Big(\tau-\sum_ky_k\Big)\,q\Big(\tau+\sum_ky_k\Big)\dd y\dd a_1\dd a_2,
\end{equation}
where $\Sigma=\{a\in[0,\infty)^3:a_1+a_2+a_3=\sigma,\ a_k\le c_k\}$, $\varrho_k=c_k-a_k$, and $\tau=\frac12(1-|c|+\sigma)$; put $\mathcal M_{p,q}(c;\sigma)=0$ if $\sigma>|c|$. Write $p^-(X)=p(-X)$.

\begin{theorem}[Fraccaro]\label{thm:fraccaro}
Let $\Phi\in\mathcal W$ and let $\Pi U_i$, $\Pi V_i$ be tests from $\mathcal T_2$ with data $(f_i,g_i,p_i)$, $i=1,2$. For a kernel $\mathcal Q(a,b,c,d)$ write $[\mathcal Q]=\iiiint f_1(a)g_1(b)f_2(c)g_2(d)\,\mathcal Q(a,b,c,d)$. Then, as $T\to\infty$,
\begin{align*}
 \mathrm N_T(\Pi U_1)/\mathcal N_T&\to\textstyle\iint f_1(a)g_1(b)\,\mathcal K_{p_1}(a,1-b)\dd a\dd b,\\
 \mathrm N_T(\Pi V_1)/\mathcal N_T&\to\textstyle\iint f_1(a)g_1(b)\,\mathcal K_{p_1}(a,b)\dd a\dd b,\\
 \mathrm G_T(\Pi U_1,\Pi U_2)/\mathcal N_T&\to\tfrac12\big[\mathcal B_{p_1,p_2}(a,d;b,c)+\mathcal B_{p_1,p_2^-}(a,c;b,d)\big],\\
 \mathrm G_T(\Pi V_1,\Pi V_2)/\mathcal N_T&\to\tfrac12\big[\mathcal B_{p_1,p_2}(a,b;c,d)\big],\\
 \mathrm G_T(\Pi U_1,\Pi V_2)/\mathcal N_T&\to\tfrac12\big[\mathcal M_{p_1,p_2^-}(b,c,d;a)+\mathcal M_{p_1^-,p_2^-}(a,c,d;b)\big].
\end{align*}
\end{theorem}
For $\deg p_i\le4$ this is \cite[Props.~3 and 7--10, and (2.65)]{Fraccaro}; the arbitrary fixed degree is \cref{lem:degree} below.

\begin{lemma}[Fixed derivative degree]\label{lem:degree}
\Cref{thm:fraccaro} holds for real polynomials $p_1,p_2$ of every fixed degree.
\end{lemma}
\begin{proof}
Fraccaro's proofs evaluate the shifted mixed moments with $C_p$ replaced by $\zeta(s+\gamma)$, $\gamma$ in a fixed small complex disc, and obtain the polynomial weights by differentiating in $\gamma$. The asymptotic formulae for the shifted $U$ and $V$ numerators at the end of \cite[\S\S2.2.3--2.2.4]{Fraccaro} and for both diagonal Gram calculations in the proofs of \cite[Props.~7--8]{Fraccaro} hold uniformly for $\gamma$ in a fixed disc of radius $\asymp1/L$, so the Cauchy integral formula on that disc gives every fixed number of derivatives with the same leading forms. In the reflection step, \cite[Lemma~6]{Fraccaro} holds at each fixed order; the sum in \cite[(2.64)]{Fraccaro} then runs up to $\deg p_2$, its squared weighted norm is $O(TL^7)$ by the balanced estimate, and its cross term is $O(TL^8)$. The discarded reflected $VV$ estimate also holds uniformly on the disc before differentiation. The implied constants depend on the tests, which are fixed.
\end{proof}

\subsection{Identification with Schur profiles}
\begin{lemma}[Bridge]\label{lem:bridge}
For all tests in $\mathcal T_2$ the limits in \Cref{thm:fraccaro} equal $\ip{v}{h_B}$ and $\ip{h_B}{h_{B'}}$, where $h_B$ is the profile of $B$. Consequently $\mathcal T_2$ is faithful.
\end{lemma}
\begin{proof}
\emph{Numerators.} Since $v\circ J=v$ and $J$ preserves measure, $\ip{v}{\frac12(h+h\circ J)}=\ip vh$. By Fubini it suffices to show $\ip{v}{H_{\alpha,\beta,p}}=\mathcal K_p(\alpha,\beta)$. In \eqref{eq:profile} write $u$ for the variable of the interval defining $\kappa$, and $w=\alpha+\beta-u$. Then $u$ lies in that interval if and only if
\[
 1\ge z\ge u\ge y\ge w\ge x\ge0,\qquad \max(0,\alpha+\beta-1)\le w\le\min(\alpha,\beta),
\]
and on this set $j=x+y+z-u-w\in[0,1]$ automatically. Hence
\[
 \ip{v}{H_{\alpha,\beta,p}}=\int_{\max(0,\alpha+\beta-1)}^{\min(\alpha,\beta)}\int_{I(\varrho_w)}v\;p\big(\tfrac12-x-y-z+\alpha+\beta\big)\dd x\dd y\dd z\dd w,
\]
where $I(\varrho)=\{\varrho_0>z>\varrho_1>y>\varrho_2>x>\varrho_3\}$ is the interlacing region of $\varrho_w=(1,\alpha+\beta-w,w,0)$. For $\varrho_0>\dots>\varrho_3$ the Dixon--Anderson change of variables (see \cite{Anderson}, \cite[\S4.2]{Forrester}; we include the short proof)
\[
 X_i=\frac{\prod_{\ell\in\{x,y,z\}}(\varrho_i-\ell)}{\prod_{k\ne i}(\varrho_i-\varrho_k)}\qquad(0\le i\le3)
\]
is a bijection from $I(\varrho)$ onto the open simplex $\{X_i>0,\ \sum_iX_i=1\}$, with $\dd X_1\dd X_2\dd X_3=\Delta(z,y,x)\dd x\dd y\dd z/\Delta(\varrho)$, where $\Delta(z,y,x)=2v$ and $\Delta(\varrho)=\prod_{i<k}(\varrho_i-\varrho_k)$. Indeed, by partial fractions $\prod_\ell(t-\ell)/\prod_i(t-\varrho_i)=\sum_iX_i/(t-\varrho_i)$. Exactly $i$ of the $\ell$ and $i$ of the $\varrho_k$ exceed $\varrho_i$, so interlacing makes every $X_i$ positive; comparing the coefficients of $t^{-1}$ and $t^{-2}$ at infinity gives $\sum X_i=1$ and $\sum\varrho_iX_i=\sum\varrho_i-(x+y+z)$. Conversely, given $X$ in the open simplex, $R(t)=\sum_iX_i/(t-\varrho_i)$ is strictly decreasing on each interval $(\varrho_{i+1},\varrho_i)$, from $+\infty$ to $-\infty$, so it has exactly one zero in each of the three intervals; and $R(t)\prod_i(t-\varrho_i)$ is a monic cubic, since $\sum X_i=1$. So it is $\prod_\ell(t-\ell)$ for a unique interlacing triple, and taking residues at $\varrho_i$ recovers $X$. For the Jacobian, $\partial X_i/\partial\ell=-X_i/(\varrho_i-\ell)$, so by Cauchy's determinant
\[
 \Big|\det\Big(\frac{\partial X_i}{\partial\ell}\Big)_{i\in\{1,2,3\},\ \ell\in\{z,y,x\}}\Big|=\prod_{i=1}^3X_i\cdot\frac{\Delta(\varrho_1,\varrho_2,\varrho_3)\,\Delta(z,y,x)}{\prod_{i=1}^3\prod_\ell|\varrho_i-\ell|}=\frac{\Delta(z,y,x)}{\Delta(\varrho)},
\]
after inserting the definition of $X_i$ and using $\prod_{i=1}^3\prod_{k\ne i}|\varrho_i-\varrho_k|=\Delta(\varrho_1,\varrho_2,\varrho_3)^2\prod_{i=1}^3(\varrho_0-\varrho_i)$. Since $\sum_i\varrho_{w,i}=1+\alpha+\beta$, the argument of $p$ equals $\sum_i(\varrho_{w,i}-\frac12)X_i$, and the inner integral is $\frac12\Delta(\varrho_w)\int p(\cdot)\dd X=\frac{\Delta(\varrho_w)}{12}\E\,p(\cdot)$, the simplex having volume $1/6$. This is \eqref{eq:Kkernel}.

\emph{Gram entries.} By \cref{lem:profiles}(iii) and $\ip{\frac12(h+h\circ J)}{\frac12(h'+h'\circ J)}=\frac12\ip h{h'}+\frac12\ip h{h'\circ J}$,
\begin{equation}\label{eq:grammatrix}
 \ip{h_B}{h_{B'}}=\lim_{N\to\infty}\frac1{2N^9}\Big(\ip{E_N[F]}{E_N[F']}_N+\ip{E_N[F]}{D^3\overline{E_N[F']}}_N\Big)
\end{equation}
for the cubic weights $F,F'$ of $B,B'$. We expand each inner product over the short indices and compute the inner sum over the two derivative indices. Throughout, $A,B,C,D$ denote short indices, scaled as $a=A/N$ and so on; for both families the short indices satisfy $a+b+c+d<1$ (for $U$-tests $a_0+b_0\le\frac14(3a_0+5b_0)<\frac12$).

\emph{(VV).} Here $\ip{E_N[F]}{E_N[F']}_N=\sum f_1(a)g_1(b)f_2(c)g_2(d)\,\Sigma_N$ with
\[
 \Sigma_N=\sum_{i,j}p(\tfrac12-\tfrac iN)q(\tfrac12-\tfrac jN)\,\ip{e_ie_Ae_B}{e_je_Ce_D}_N ,
\]
$p=p_1$, $q=p_2$. The inner product vanishes unless $i+A+B=j+C+D$. If $i+j\le N$ the common degree is $\frac12(i+j+A+B+C+D)<N$, so by \ref{S:stable} it counts $3\times3$ tables with row sums $(i,A,B)$ and column sums $(j,C,D)$. Such a table is determined by its short-to-short block $(\varrho_{kl})_{k,l\le2}$, whose row sums are at most $A,B$ and column sums at most $C,D$, together with its corner entry $\ell\ge0$; then $i=\ell+(C+D-|\varrho|)$ and $j=\ell+(A+B-|\varrho|)$. If $i+j>N$, \ref{S:dual} gives $\ip{e_ie_Ae_B}{e_je_Ce_D}_N=\ip{e_{N-j}e_Ae_B}{e_{N-i}e_Ce_D}_N$, and $(N-j)+(N-i)<N$, so the same parametrisation applies to $(N-j,N-i)$. Summing over the two halves and scaling ($\varrho=N\varrho^{\rm sc}$, $\ell=N(w+h)$ in the first half and $\ell=N(w-h)$ in the second) gives
\[
 \Sigma_N=N^5\,\mathcal B_{p,q}(a,b;c,d)+O(N^4)
\]
uniformly in the short indices: both halves are sums of a fixed polynomial over the lattice points of a five-dimensional polytope of bounded complexity, and in the scaled variables the two derivative weights become $p(\delta-h)q(-\delta-h)$ with $h\in[-w,0]$ and $h\in[0,w]$ respectively. The reflected term vanishes identically: by \ref{S:dual} every Schur constituent of $D^3\overline{e_je_Ce_D}$ has degree $3N-(j+C+D)$, and $i+A+B=3N-(j+C+D)$ is impossible because $i,j\le N$ and $A+B+C+D<N$. Summing over the $O(N^4)$ short indices, \eqref{eq:grammatrix} equals the stated $VV$ limit.

\emph{(UU).} With $e_{N-B}=D\overline{e_B}$ and $|D|=1$,
\begin{gather*}
 \ip{e_ie_Ae_{N-B}}{e_je_Ce_{N-D}}_N=\ip{e_ie_Ae_D}{e_je_Ce_B}_N,\\
 \ip{e_ie_Ae_{N-B}}{D^3\overline{e_je_Ce_{N-D}}}_N=\ip{e_ie_Ae_C}{e_{N-j}e_Be_D}_N ,
\end{gather*}
the second because $D\cdot\overline D^{3}\cdot D=\overline D$ and $\overline De_j=\overline{e_{N-j}}$. In the second, substituting $j'=N-j$ turns the weight $q(\frac12-\frac jN)$ into $q^-(\frac12-\frac{j'}N)$. Both are now of the form treated in (VV), which gives $N^5\mathcal B_{p,q}(a,d;c,b)$ and $N^5\mathcal B_{p,q^-}(a,c;b,d)$ up to $O(N^4)$; and $\mathcal B$ is symmetric in its last two arguments.

\emph{(UV).} Similarly
\begin{gather*}
 \ip{e_ie_Ae_{N-B}}{e_je_Ce_D}_N=\ip{e_ie_{N-j}e_A}{e_Be_Ce_D}_N,\\
 \ip{e_ie_Ae_{N-B}}{D^3\overline{e_je_Ce_D}}_N=\ip{e_Ae_Ce_D}{e_{N-i}e_{N-j}e_B}_N .
\end{gather*}
Consider in general $\ip{e_{i}e_{j}e_{S}}{e_{C_1}e_{C_2}e_{C_3}}_N$ with $C_1+C_2+C_3<N$, weighted by $P(\frac12-\frac iN)Q(\frac12-\frac jN)$. A table with rows $(i,j,S)$ and columns $(C_1,C_2,C_3)$ is determined by the split $(a_1,a_2,a_3)$ of the row $S$, with $a_k\le C_k$, and by the entries $v_k\in[0,C_k-a_k]$ of the row $i$; then $j=H-\sum v_k$ with $H=|C|-S$. With $v_k=\frac12(C_k-a_k)+Ny_k$ the weights become $P(\tau-\sum y_k)Q(\tau+\sum y_k)$ in scaled variables, so the weighted sum is $N^5\mathcal M_{P,Q}(c;\sigma)+O(N^4)$, $\sigma=S/N$. In the first identity substitute $j'=N-j$, so $(P,Q)=(p,q^-)$ with columns $(b,c,d)$ and $\sigma=a$. In the second substitute $i'=N-i$, $j'=N-j$, so $(P,Q)=(p^-,q^-)$ with columns $(a,c,d)$ and $\sigma=b$. This gives the stated $UV$ limit.

Combining with \Cref{thm:fraccaro}, the limits of $\mathrm N_T/\mathcal N_T$ and $\mathrm G_T/\mathcal N_T$ are $\ip v{h_B}$ and $\ip{h_B}{h_{B'}}$, so $\mathcal T_2$ is faithful.
\end{proof}

\begin{remark}
The identity of \cref{lem:bridge}, including the overall normalisation, has also been compared numerically with the exact integer evaluation of Fraccaro's kernels. We took twelve point tests from Fraccaro's certificate: its first and last rows and ten rows drawn at random. It integrates their projected point features over $\Om$ by the midpoint rule on the cubes of side $1/120$ and $1/240$ whose centres lie in $\Om$, and combines the two by Richardson extrapolation. The largest relative error of a numerator, $|\text{computed}-\text{exact}|/|\text{exact}|$, is $1.5\cdot10^{-3}$. The largest Gram discrepancy, divided by the largest absolute Gram entry of the sample, is $5.5\cdot10^{-4}$. This is a floating-point consistency check; the proof is the one given above.
\end{remark}

\subsection{Supports}
Put $\varrho=1-z$ and define the open sets
\begin{align*}
 \mathcal C&=\{3\min(x,\varrho)+5\max(x,\varrho)<2\}\cup\{y<\tfrac14\}\cup\{y>\tfrac34\},\\
 \mathcal B&=\{3\min(x,\varrho)+5\max(x,\varrho)>2,\ x+y>\tfrac12,\ y+z<\tfrac32\},\\
 W&=\{0<x<\tfrac14,\ \tfrac14<y<\tfrac12-x,\ y<z<\tfrac35(1+x)\}.
\end{align*}
($J$ exchanges $x$ and $\varrho$ and maps $y$ to $1-y$.) The sets $\mathcal C$ and $\mathcal B$ are $J$-invariant; $W$ is not, and $JW$ is the opposite wing.

\begin{lemma}\label{lem:partition}
Up to a null set, $\Om$ is the disjoint union of $\mathcal C$, $\mathcal B$, $W$ and $JW$.
\end{lemma}
\begin{proof}
Off the null set where one of the defining inequalities is an equality, a point outside $\mathcal C\cup\mathcal B$ has $3\min(x,\varrho)+5\max(x,\varrho)>2$, $\frac14<y<\frac34$, and either $x+y<\frac12$ or $y+z>\frac32$. These two alternatives are exchanged by $J$ and are incompatible, since together they force $z-x>1$. If $x+y<\frac12$, then $y>\frac14$ forces $x<\frac14$ and $y<\frac12-x$; if $\varrho\le x$ we would have $3\varrho+5x\le8x<2$, so $\varrho>x$ and the first inequality reads $3x+5\varrho>2$, that is, $z<\frac35(1+x)$. Conversely the inequalities defining $W$ imply all of these. So the first alternative is $W$ and the second is $JW$.
\end{proof}

\begin{lemma}[Dominance]\label{lem:dominance}
If $\kappa(x,y,z;c)>0$ then $x\le c_k\le z$ for $k=1,2,3$. Consequently:
\begin{enumerate}[label=\textup{(\roman*)}]
\item every $U$-profile vanishes off $\{3\min(x,\varrho)+5\max(x,\varrho)<2\}\subseteq\mathcal C$;
\item every raw $V$-profile $h_{F_V}$ vanishes off $\{x+y<\frac12\}$, so every $V$-profile vanishes off $\{x+y<\frac12\}\cup\{y+z>\frac32\}$;
\item every profile in $\mathcal T_2$ vanishes on $\mathcal B$, and on $W$ the $V$-profile equals $\frac12h_{F_V}$.
\end{enumerate}
\end{lemma}
\begin{proof}
In the notation of the proof of \cref{lem:bridge}, $\kappa>0$ requires $z\ge u\ge\max(c_1,c_2)$ and $x\le w\le\min(c_1,c_2)$, while $c_3=x+(y-w)+(z-u)\ge x$ and $c_3=z-(u-y)-(w-x)\le z$.

(i) For the raw $U$-profile, $c_1=\alpha<a_0$ and $c_2=1-\beta$ with $\beta<b_0$, so $x<a_0$ and $\varrho<b_0$. If $x\le\varrho$ then $3x+5\varrho<3a_0+5b_0<2$; if $\varrho<x$ then $3\varrho+5x<8a_0\le3a_0+5b_0<2$. The set is $J$-invariant, so the reflected part vanishes off it too.

(ii) For $V$, $c_1,c_2<\eta<\frac14$ and $c_3\le z$, so $x+y=c_1+c_2+c_3-z<\frac12$. Reflection maps $\{x+y<\frac12\}$ to $\{y+z>\frac32\}$.

(iii) Both sets in (i) and (ii) miss $\mathcal B$. On $W$ the $U$-profiles vanish by (i), and $y+z<\frac12+\frac34<\frac32$, so the reflected half of a $V$-profile vanishes by (ii).
\end{proof}

\subsection{The span}
\begin{lemma}[Density on $\mathcal C$]\label{lem:density}
$L^2_J(\mathcal C)\subseteq\mathcal H(\mathcal T_2)$.
\end{lemma}
\begin{proof}
Write $\mathcal H=\mathcal H(\mathcal T_2)$.

\emph{Step 1: point features and limits.} Let $\mathcal F$ be the set of bounded Borel $F$ such that $\frac12(h_F+h_F\circ J)\in\mathcal H$. It is a linear space, and it is closed under uniform convergence because $\|h_F\|\le\|F\|_\infty\|v\|$ (\cref{lem:profiles}(ii)) and $\mathcal H$ is closed. It contains $F_U$ and $F_V$ for all admissible data. It also contains $f(c_1)g(c_2)P(c_3)$ whenever $P$ is continuous on $[0,1]$ and either $f\in C_c((0,a_0))$, $g(1-\cdot)\in C_c((0,b_0))$ with $0<a_0\le b_0$, $3a_0+5b_0<2$, or $f,g\in C_c((0,\eta))$ with $\eta<\frac14$. Indeed, $P(c_3)=p(\frac12-c_3)$ for a continuous $p$ on $[-\frac12,\frac12]$, which is a uniform limit of polynomials; continuous compactly supported $f,g$ are uniform limits of smooth ones; and a continuous $f$ with compact support in $(0,\eta)$ is the uniform limit of $u(\eta-u)P_n(u)$, where the polynomials $P_n$ approximate $f(u)/(u(\eta-u))$ uniformly on $[0,\eta]$.

\emph{Step 2: synthesis of a bump.} Fix a point of $\mathcal C$ off the null set where $x=\varrho$, $y=\frac14$, $y=\frac34$ or $3\min(x,\varrho)+5\max(x,\varrho)=2$. Replacing it by its image under $J$ if necessary, we may assume that it lies in $\{y<\frac14\}$, or in $\{3x+5\varrho<2,\ x<\varrho\}$. Let $Q$ be a closed box around it, inside $\Om$ and inside the same one of these two sets, so small that its projections leave room for admissible data: in the first case $Q\subseteq\{y<\eta\}$ for some $\eta<\frac14$; in the second, $x<a_0$ and $\varrho<b_0$ on $Q$ for some $a_0\le b_0$ with $3a_0+5b_0<2$.

Let $\psi\in C_c^\infty(\operatorname{int}Q)$ and write $\hat\psi(m_1,m_2,m_3)=\psi(m_3,m_2,m_1)$ for its expression in descending coordinates. By \ref{S:jt},
\[
 \sum_{\lambda}\psi(\lambda/N)\,s_\lambda=\sum_{m}\beta_N(m)\,e_{m_1}e_{m_2}e_{m_3},\]\[ \beta_N(m)=\sum_{\sigma\in S_3}\operatorname{sgn}(\sigma)\,\hat\psi\Big(\frac{m+\iota-\sigma}{N}\Big),
\]
where $\iota=(1,2,3)$, $\sigma=(\sigma(1),\sigma(2),\sigma(3))$, and $\hat\psi$ is extended by zero (for large $N$ every argument in the support has strictly decreasing coordinates in $(0,1)$, so the substitution $m_i=\lambda'_i-i+\sigma(i)$ is a bijection on the relevant terms). Taylor expansion of $\hat\psi$ at $m/N$ gives
\[
 \beta_N(m)=N^{-3}\Psi(m/N)+O(N^{-4}),\qquad \Psi=\frac16\sum_{\sigma\in S_3}\operatorname{sgn}(\sigma)\big((\iota-\sigma)\cdot\nabla\big)^3\hat\psi,
\]
uniformly in $m$, with $\beta_N(m)=0$ unless $m/N$ is within $O(1/N)$ of $\operatorname{supp}\hat\psi$. The terms of order $0,1,2$ cancel because $\sum_\sigma\operatorname{sgn}(\sigma)R(\sigma)=0$ for every polynomial $R$ of degree at most $2$ in three variables: its alternation is an alternating polynomial of degree at most $2$, hence divisible by the Vandermonde polynomial of degree $3$, hence zero. Taking the coefficient of $s_\lambda$ and using $K_\lambda\ge0$, $\sum_mK_\lambda(m)=d_\lambda=O(N^3)$,
\[
 \psi(\lambda/N)=N^{-3}c_\lambda(\Psi)+O(N^{-1})=h_\Psi(\lambda/N)+o(1)
\]
uniformly, by \cref{lem:profiles}(i). As $\psi$ and $h_\Psi$ are continuous, $h_\Psi=\psi$ on $\Om$.

\emph{Step 3: approximation.} $\Psi$ is continuous with support in $\operatorname{supp}\hat\psi$. In the first case assign the coordinates $(m_3,m_2,m_1)\leftrightarrow(x,y,z)$ of $\Psi$ to $(c_1,c_2,c_3)$; in the second, to $(c_1,c_3,c_2)$, with $c_2=1-\varrho$ the reflected short factor. By the Stone--Weierstrass theorem, $\Psi$ is a uniform limit on $[0,1]^3$ of finite sums $\sum f_k(c_1)g_k(c_2)P_k(c_3)$ of continuous functions. Multiplying the two short factors by continuous cutoffs $\chi_1(c_1)\chi_2(c_2)$ that equal $1$ on the projections of $\operatorname{supp}\Psi$ and are supported where the data of Step 1 are admissible, we keep a uniform approximation of $\Psi=\Psi\chi_1\chi_2$. Since the order of the three factors in $E_N$ is immaterial, $h_{F\circ\pi}=h_F$ for every permutation $\pi$ of the coordinates, so each approximant lies in $\mathcal F$, and hence $\Psi\in\mathcal F$, that is, $\frac12(\psi+\psi\circ J)\in\mathcal H$.

\emph{Step 4.} Every $J$-symmetric $\phi\in C_c^\infty$ supported in $\mathcal C$ off the null set is, by a partition of unity, a finite sum of functions $\frac12(\psi+\psi\circ J)$ as in Step 2. These are dense in $L^2_J(\mathcal C)$, and $\mathcal H$ is closed.
\end{proof}

For a continuous $g$ on $\R^2$ put $\Gamma g(x,y,z)=g(r,x+y)-g(r,x+z)$, where $r=x+y+z$. Let $\mathcal G$ be the closure in $L^2(W)$ of $\{\Gamma g|_W: g\in C(\R^2),\ g(r,u)=0\text{ for }u\ge\frac12-\epsilon\text{ for some }\epsilon>0\}$, and let $\widetilde{\mathcal G}$ be the set of $J$-symmetric functions supported in $W\cup JW$ whose restriction to $W$ lies in $\mathcal G$.

\begin{proposition}[The span]\label{prop:span}
$\mathcal H(\mathcal T_2)=L^2_J(\mathcal C)\oplus\widetilde{\mathcal G}$, an orthogonal sum. In particular every element of $\mathcal H(\mathcal T_2)$ vanishes on $\mathcal B$.
\end{proposition}
\begin{proof}
The two summands are orthogonal since $\mathcal C$ and $W\cup JW$ are disjoint (\cref{lem:partition}).

\emph{$\subseteq$.} Let $h$ be a profile in $\mathcal T_2$. It is $J$-symmetric and vanishes on $\mathcal B$ (\cref{lem:dominance}), so by \cref{lem:partition} $h=h\one_{\mathcal C}+h\one_{W\cup JW}$ with $h\one_{\mathcal C}\in L^2_J(\mathcal C)$. On $W$, $h$ vanishes for $U$-tests, and for a $V$-test $h=\frac12h_{F_V}$. For $\alpha,\beta<\eta<\frac14$ and $(x,y,z)\in W$ we have $y>\frac14>\max(\alpha,\beta,\alpha+\beta-y)$, so
\begin{equation}\label{eq:Wfeature}
 H_{\alpha,\beta,p}=p\big(\tfrac12-r+\alpha+\beta\big)\big[(\alpha+\beta-x-y)_+-(\alpha+\beta-x-z)_+\big]\quad\text{on }W,
\end{equation}
the cutoff $j\in[0,1]$ being automatic where the bracket is non-zero. Integrating against $f(\alpha)g(\beta)$ shows $\frac12h_{F_V}=\Gamma G$ on $W$ with $G(r,u)=\frac12\iint f(\alpha)g(\beta)p(\frac12-r+\alpha+\beta)(\alpha+\beta-u)_+\dd\alpha\dd\beta$, which is continuous and vanishes for $u\ge2\eta$. So $h\one_{W\cup JW}\in\widetilde{\mathcal G}$, and $\mathcal H(\mathcal T_2)$, the closure of the span of such $h$, is contained in the closed space on the right.

\emph{$\supseteq$.} By \cref{lem:density}, $L^2_J(\mathcal C)\subseteq\mathcal H(\mathcal T_2)$. Let $g$ be continuous with $g(r,u)=0$ for $u\ge\frac12-\epsilon$. On the compact set of values $(r,u)$ attained on $\overline W$, $g$ is a uniform limit of finite sums $\sum_kP_k(r)(b_k-u)_+$ with polynomials $P_k$ and knots $b_k\in(0,\frac12)$: interpolate $g(r,\cdot)$ linearly at fixed knots below $\frac12-\epsilon$ (a continuous piecewise linear function of $u$ vanishing beyond its last knot is a combination of the $(b_k-u)_+$), and approximate the continuous coefficients by polynomials. For one term, take the point feature with $\alpha=\beta=b/2$ and $p(t)=P(\frac12+b-t)$; by \eqref{eq:Wfeature}, its projection $h=\frac12(H+H\circ J)$ equals $\frac12P(r)\Gamma[(b-u)_+]$ on $W$. This $h$ lies in $\mathcal H(\mathcal T_2)$: the map $(\alpha,\beta)\mapsto H_{\alpha,\beta,p}\in L^2(\Om)$ is continuous (dominated convergence; for fixed $(x,y,z)$ the feature is continuous in $(\alpha,\beta)$ off the lines $\alpha+\beta\in\{r,r-1\}$), so $h$ is the limit of the profiles of $V$-tests whose short polynomials $f,g$ concentrate at $b/2$, constructed as in Step 1 of \cref{lem:density}. Its restriction to $\mathcal C$ lies in $L^2_J(\mathcal C)\subseteq\mathcal H(\mathcal T_2)$ and it vanishes on $\mathcal B$, so $h\one_{W\cup JW}\in\mathcal H(\mathcal T_2)$. Summing, and using that the maps $(x,y,z)\mapsto(r,x+y)$ and $(r,x+z)$ are linear of rank two on the bounded set $W$ (so uniform convergence of $g$ implies convergence of $\Gamma g$ in $L^2(W)$), we obtain the $J$-symmetric extension of $\Gamma g|_W$ in $\mathcal H(\mathcal T_2)$. Taking closures proves the claim.
\end{proof}

\subsection{The residual least-squares problem}
By \cref{prop:span} and \cref{lem:partition},
\begin{equation}\label{eq:threelosses}
 9!\,\|P_{\mathcal H(\mathcal T_2)}v\|^2=42-\mu(\mathcal B)-2\cdot9!\operatorname{dist}_{L^2(W)}(v,\mathcal G)^2,
\end{equation}
because $P_{\mathcal H}v=v\one_{\mathcal C}+P_{\widetilde{\mathcal G}}v$, and $J$-symmetric functions on $W\cup JW$ have twice their squared norm on $W$. We compute the two losses. All integrals below are of polynomials over polytopes and are evaluated exactly by computer, and an independently organised computation gives the same values.

\emph{The invisible set.} By $J$-symmetry, $\mu(\mathcal B)=2\mu(\mathcal B\cap\{x<\varrho\})$. On $\mathcal B\cap\{x<\varrho\}$ the first inequality forces $\varrho>\frac14$, so $y<z<\frac34$ and $y+z<\frac32$ is automatic. In the coordinates $(x,\varrho,t)$ with $t=y-x$ and $d=1-\varrho-x=z-x$ we have $4v^2=t^2(d-t)^2d^2$, and the region is
\begin{gather*}
 \{0<x<\tfrac14,\ \tfrac{2-3x}5<\varrho<\tfrac12+x,\ \tfrac12-2x<t<d\}\\ \cup\ \{\tfrac14<x<\tfrac12,\ x<\varrho<1-x,\ 0<t<d\}.
\end{gather*}
(For $x>\frac14$ the constraints $3x+5\varrho>2$ and $x+y>\frac12$ are automatic.) Integrating $\frac{9!}2t^2(d-t)^2d^2$ over the two pieces gives $\frac{60361491}{100000000}$ and $\frac{21}{256}$, so
\begin{equation}\label{eq:L0}
 L_0:=\mu(\mathcal B)=\frac{8570577}{12500000}.
\end{equation}
The second piece is the half $\{x<\varrho\}$ of the scaled chamber $\frac14<x<y<z<\frac34$; with the factor $2$ included above, its value $\frac{21}{256}$ is the mass $42/2^9$ of the whole scaled chamber, by scaling.

\emph{The fit on $W$.} Let $g$ be admissible. On $W_>=W\cap\{z>\frac12\}$ we have $x+z>\frac12$, so $\Gamma g=g(r,x+y)$; on $W_<=W\cap\{z<\frac12\}$ both terms may be non-zero. The values of $g$ used on $W_>$ lie in $\{r-u>\frac12\}$, because $r-(x+y)=z$; those used on $W_<$ lie in $\{r-u<\frac12\}$, because $r-(x+y)=z<\frac12$ and $r-(x+z)=y<\frac12$. Hence the two fits are independent, and $\operatorname{dist}(v,\mathcal G)^2=\delta_>+\delta_<$, the two squared distances on $W_>$ and $W_<$ in which $g$ ranges over all bounded measurable functions on the respective half-planes vanishing for $u\ge\frac12$. (Such $g$ are $L^2$-limits of admissible ones on the bounded set $W$, by truncation, smoothing and dominated convergence, the push-forwards of Lebesgue measure under the two rank-two maps being absolutely continuous.)

\emph{The large region.} Write $S=x+y$ and $h=S-\frac14$. On $W_>$, $\frac12<z<\frac34$ and, at fixed $(z,S)$, $x$ ranges over $(\ell,h)$ with $\ell=\max(0,\frac{5z-3}3)$, for $\frac14+\ell<S<\frac12$. Since $r=z+S$, the fits $g(r,S)$ are exactly the functions of $(z,S)$, and the best fit is the mean of
\[
 P(x)=v(x,S-x,z)=\tfrac12(z-S+x)(z-x)(S-2x)
\]
over $x\in(\ell,h)$. Thus, splitting at $z=\frac35$ where $\ell$ changes,
\begin{equation}\label{eq:L1}
\begin{split}
 L_1:=2\cdot9!\,\delta_>&=2\cdot9!\int_{1/2}^{3/4}\int_{1/4+\ell}^{1/2}\Big(\int_\ell^hP^2\dd x-\frac{(\int_\ell^hP\dd x)^2}{h-\ell}\Big)\dd S\dd z\\
 &=\frac{588879}{25600000}+\frac{1114119}{40960000}.
\end{split}
\end{equation}

\emph{The small region.} Put $u=4z-1$, $v_0=4y-1$, $w=4x$ and $m=u+v_0+w$. Then $W_<$ becomes $\{0<v_0<u<1,\ 0<w<1-v_0\}$ (the constraint $z<\frac35(1+x)$ is automatic), and $v=d_m(u,v_0)/128$ with
\[
 d_m(u,v_0)=(u-v_0)(1+u-w)(1+v_0-w),\qquad w=m-u-v_0 .
\]
Since $r=\frac14(m+2)$, $x+y=\frac14(m+1-u)$ and $x+z=\frac14(m+1-v_0)$, the fit is $\frac1{128}(f_m(u)-f_m(v_0))$ with $f_m$ an arbitrary function vanishing for $u\le c_m:=\max(0,m-1)$. At fixed $m\in(0,2)$ the fibre is $\{(u,v_0):0<v_0<u<1,\ u>m-1,\ u+v_0<m\}$. Since $d_m$ is antisymmetric in $(u,v_0)$, the first variation in the direction of a bounded $\phi$ supported in $(c_m,\min(1,m))$ is $-2\int\phi(u)\int_0^{m-u}\big(d_m(u,v_0)-f_m(u)+f_m(v_0)\big)\dd v_0\dd u$, so the normal equation is
\begin{equation}\label{eq:normal}
 (m-u)f_m(u)-\int_{c_m}^{m-u}f_m(v_0)\dd v_0=b_m(u):=\int_0^{m-u}d_m(u,v_0)\dd v_0
\end{equation}
for $c_m<u<\min(1,m)$. Here $b_m(u)=\frac12(m-u)(4u^2+3u-2mu-m)$. The equation is solved by
\begin{equation}\label{eq:fm}
 f_m(u)=3u^2+(1-2m)u+\tfrac12c_m\quad(u>c_m),\qquad f_m(u)=0\quad(u\le c_m).
\end{equation}
Indeed, for $m\le1$ this is a direct polynomial identity; for $m>1$ the unpinned polynomial satisfies the same identity with lower limit $0$, the constant $\frac12c_m$ makes its integral over $(0,c_m)$ vanish, and $m-u\ge c_m$ for all free $u$. A function satisfying the normal equation for every $\phi$ is the orthogonal projection, since the residual is then orthogonal to the closed subspace of fits and $f_m$ itself is admissible. The squared norm of the fit equals its inner product with $d_m$, which by antisymmetry is $\int_{c_m}^{\min(1,m)}f_mb_m$. With
\begin{gather*}
 \iiint_{0<v_0<u<1,\ 0<w<1-v_0}d_m^2=\frac{104}{945},\\
 \int_0^2\int_{c_m}^{\min(1,m)}f_m(u)b_m(u)\dd u\dd m=\frac{41}{2520}+\frac{457}{5040}=\frac{77}{720},
\end{gather*}
and the Jacobian $4^{-3}$ and scaling $128^{-2}$,
\begin{equation}\label{eq:L2}
 L_2:=2\cdot9!\,\delta_<=\frac{9!}{2\cdot4^9}\Big(\frac{104}{945}-\frac{77}{720}\Big)=\frac{141}{65536}.
\end{equation}

\subsection{Proof of \Cref{thm:twofamily}}
By \eqref{eq:threelosses}--\eqref{eq:L2},
\[
 9!\,\|P_{\mathcal H(\mathcal T_2)}v\|^2=42-L_0-L_1-L_2=\frac{264076794201}{6400000000}=C_* .
\]
By \cref{lem:bridge} the family $\mathcal T_2$ is faithful, so \cref{prop:projection} gives the lower bound. Conversely, for any finite system $B_1,\dots,B_n$ from $\mathcal T_2$ with numerator vector $\mathbf k$ and Gram matrix $G$ (the limits in \Cref{thm:fraccaro}), \cref{lem:bridge} gives $k_i=\ip v{h_i}$ and $G_{ij}=\ip{h_i}{h_j}$, so for every real $\mathbf b$ with $\mathbf b^tG\mathbf b>0$,
\[
 \frac{9!(\mathbf b^t\mathbf k)^2}{\mathbf b^tG\mathbf b}=9!\frac{\ip{v}{h}^2}{\|h\|^2}\le9!\,\|P_{\mathcal H(\mathcal T_2)}v\|^2=C_*,\qquad h=\sum b_ih_i .
\]
Since $\sup_h\ip vh^2/\|h\|^2$ over the span of all profiles is $\|P_{\mathcal H}v\|^2$, $C_*$ is the supremum. \qed

\begin{remark}\label{rem:degreefour}
Restricted to derivative degree at most four, the supremum is strictly smaller than $C_*$. Suppose finite systems of degree at most four had Rayleigh quotients $\ip v{h_n}^2/\|h_n\|^2$ tending to $\|P_{\mathcal H}v\|^2$. Replacing $h_n$ by $-h_n$ where necessary, we may assume $\ip v{h_n}\ge0$. Since $\ip{v}{h_n}=\ip{P_{\mathcal H}v}{h_n}$, the normalised profiles then satisfy $\|h_n/\|h_n\|-P_{\mathcal H}v/\|P_{\mathcal H}v\|\|^2=2-2\ip v{h_n}/(\|h_n\|\|P_{\mathcal H}v\|)\to0$, so they converge in $L^2$ to $P_{\mathcal H}v/\|P_{\mathcal H}v\|$. On $W_>$ every degree-four profile is, at fixed $S$, a polynomial of degree at most four in $z$ (by \eqref{eq:Wfeature} with $r=z+S$), and on a rectangle in $(z,S)$ where $h-\ell$ is bounded below this property is closed in $L^2$, by projection onto the five-dimensional space on each slice. But the optimal mean in \eqref{eq:L1} has a kink at $z=\frac35$: its $z$-derivative jumps by $\frac5{3h}\big(\frac1h\int_0^hP-P(0)\big)$, which equals $-\frac{83}{960}$ at $S=\frac38$ and is non-zero for $S$ in an open interval by continuity. The exact degree-four supremum is not determined here.
\end{remark}

\begin{remark}
Exact rational evaluation of Fraccaro's $250$ rows with his integer weights gives $41.0089916904\ldots$, and an expanded $397$-row certificate with $\deg p\le4$ gives $41.0298094152\ldots$. These finite certificates are not used in the proofs.
\end{remark}

\section{Short cubic tests and unequal same-side lengths}\label{sec:third}

We add two kinds of tests to $\mathcal T_2$. Put $\varrho=1-z$ as before and
\begin{gather*}
 \mathcal D_{\rm acc}=\{3\min(x,\varrho)+5\max(x,\varrho)<2\}\cup\{x+y<\tfrac12\}\cup\{y+z>\tfrac32\},\\
 \mathcal A=\{r<1\}\cup\{r>2\},\qquad S_3=\mathcal D_{\rm acc}\cup\mathcal A .
\end{gather*}

\emph{$W$-tests.} For real $\phi\in C_c^\infty((0,1))$ let $D_\phi(s)=\sum_{n\ge1}d_3(n)\phi(\log n/L)n^{-s}$. The test is $\Pi D_\phi$, its cubic weight is $F_\phi(c)=\phi(c_1+c_2+c_3)$, and its profile is $\frac12(\phi(r)+\phi(3-r))v$. (Indeed $h_{F_\phi}=\phi(r)h_1=\phi(r)v$, because $\kappa(x,y,z;c)$ vanishes unless $c_1+c_2+c_3=r$.)

\emph{$V'$-tests.} Data $(a_0,b_0,f,g,p)$ with $a_0,b_0>0$, $a_0+b_0<\frac12$, real $f\in C_c^\infty((0,a_0))$, $g\in C_c^\infty((0,b_0))$ and a real polynomial $p$. The test is $\Pi V$ with $V=C_pA_fA_g$, and its cubic weight and profile are those of a $V$-test.

Let $\mathcal T_3=\mathcal T_2\cup\{W\text{-tests}\}\cup\{V'\text{-tests}\}$.

\begin{theorem}\label{thm:third}
$\mathcal T_3$ is faithful, and $v\one_{S_3}\in\mathcal H(\mathcal T_3)$. Moreover
\[
 \mu(S_3)=42-L_0+\frac1{512}=\frac{8263261393}{200000000}=41.316306965 .
\]
In particular $\liminf_{T\to\infty}M_6(T)/(c_3TL^9)\ge41.316306965$.
\end{theorem}

The $V'$-tests remove the losses $L_1,L_2$ of \Cref{sec:twofamily}: they reach all of $\{x+y<\frac12\}\supseteq W$. The $W$-tests recover the part $\mathcal B\cap\mathcal A$ of the invisible set.

\subsection{The diagonal and the stable range}
For $t\in[0,\infty)^{3\times3}$ write $R(t)\in[0,\infty)^3$ for its row sums and $C(t)$ for its column sums, and $\dd t$ for Lebesgue measure on $\R^9$.

\begin{lemma}[Stable pairings]\label{lem:stable}
Let $F,G$ be continuous on $[0,1]^3$, and suppose $F$ vanishes on $\{|c|\ge1-\delta\}$ for some $\delta>0$. Then
\[
 \ip{h_F}{h_G}=\int_{[0,\infty)^{3\times3}}F(R(t))\,G(C(t))\dd t,\qquad \ip{v}{h_F}=\int_{[0,\infty)^{3\times3}}F(R(t))\dd t .
\]
If moreover $G$ vanishes on $\{|c|\ge2-\delta\}$, then $\ip{h_F}{h_G\circ J}=0$.
\end{lemma}
\begin{proof}
By \cref{lem:profiles}(iii),\[\ip{h_F}{h_G}=\lim_{N\to\infty} N^{-9}\sum_{m,m'}F(m/N)G(m'/N)\ip{e_m}{e_{m'}}_N,\]where $e_m=e_{m_1}e_{m_2}e_{m_3}$. Only $|m|=|m'|<N$ contributes, so by \ref{S:stable} the sum counts the integer tables $t$ with $R(t)=m$, $C(t)=m'$, weighted by $F(R(t)/N)G(C(t)/N)$. This is a Riemann sum over the lattice points of the region $\{t\in[0,\infty)^9:|t|<(1-\delta)N\}$ for a bounded integrand that is continuous there, and it converges to the integral. The case $G=1$ gives the second formula, since $h_1=v$. For the last claim, by \ref{S:dual} every constituent of $D^3\overline{e_{m'}}$ has degree $3N-|m'|>N+\delta N>|m|$, so $\ip{E_N[F]}{D^3\overline{E_N[G]}}_N=0$.
\end{proof}

\begin{lemma}[Arithmetic diagonal]\label{lem:diagonal}
Let $\Psi$ be continuous with compact support in $[0,\infty)^6$, and write $\ell(n)=(\log n_1,\log n_2,\log n_3)/L$ for $n\in\N^3$. Then
\[
 \frac1{L^9}\sum_{n_1n_2n_3=m_1m_2m_3}\frac{\Psi(\ell(n),\ell(m))}{n_1n_2n_3}\ \longrightarrow\ a_3\int_{[0,\infty)^{3\times3}}\Psi(R(t),C(t))\dd t .
\]
The same holds with $\Psi$ replaced by $\Psi_L$, if $\Psi_L\to\Psi$ uniformly with supports in a fixed compact set; and the left side with $|\Psi|$ in place of $\Psi$ is $O_\Psi(1)$.
\end{lemma}
\begin{proof}
For $u,w\in\C^3$ with positive real parts, the Euler product gives
\[
 \sum_{n_1n_2n_3=m_1m_2m_3}\frac{\prod_in_i^{-u_i/L}\prod_jm_j^{-w_j/L}}{n_1n_2n_3}=\mathcal A\Big(\frac uL,\frac wL\Big)\prod_{i,j=1}^3\zeta\Big(1+\frac{u_i+w_j}L\Big),
\]
where $\mathcal A$ is an Euler product that converges absolutely and uniformly for all real parts in $(-\frac1{10},\frac1{10})$, because the terms of order $p^{-1}$ cancel; at the origin its local factor is $(1-X)^9\sum_{r\ge0}\binom{r+2}2^2X^r$ with $X=1/p$, so $\mathcal A(0,0)=a_3$. Since $\zeta(1+\epsilon)\sim1/\epsilon$, the left side divided by $L^9$ tends to $a_3\prod_{i,j}(u_i+w_j)^{-1}$ for fixed $u,w$, locally uniformly. Since
\[
 \prod_{i,j}(u_i+w_j)^{-1}=\int_{[0,\infty)^{3\times3}}e^{-\sum_iu_iR_i(t)-\sum_jw_jC_j(t)}\dd t,
\]
the measures $\nu_L=L^{-9}\sum(n_1n_2n_3)^{-1}\delta_{(\log n/L,\log m/L)}$ on $[0,\infty)^6$ have Laplace transforms converging on the positive orthant to that of the push-forward of $a_3\dd t$ under $t\mapsto(R(t),C(t))$. By the continuity theorem for Laplace transforms of positive measures (after an exponential tilt, which makes them finite), $\nu_L$ converges vaguely, which is the claim. The variant with $\Psi_L$ and the bound follow from vague convergence and a positive Laplace majorant.
\end{proof}

\begin{lemma}[Smooth expansion]\label{lem:expansion}
Fix $1\le K\le4$ and a compact set $\Gamma\subset\C^K$. Let $\omega\in C^\infty([0,\infty))$ equal $1$ on $[0,1]$ and $0$ on $[2,\infty)$. For $T\le t\le2T$ and $\gamma\in\Gamma$,
\[
 \prod_{k=1}^K\zeta\Big(s+\frac{\gamma_k}L\Big)=\sum_{m\ge1}\frac{d_{\gamma/L}(m)}{m^{s}}\,\omega\Big(\frac m{T^8}\Big)+O_A(T^{-A}),
\]
where $d_{\gamma/L}(m)$ are the coefficients of $\prod_k\zeta(s+\gamma_k/L)$, uniformly in $\gamma$, and similarly for $t$-derivatives of fixed order and for $\partial_\gamma$-derivatives (so for $C_p$ in place of a shifted zeta factor).
\end{lemma}
\begin{proof}
By Mellin inversion the sum equals $\frac1{2\pi i}\int_{(2)}\prod_k\zeta(s+w+\gamma_k/L)\widetilde\omega(w)T^{8w}\dd w$, where $\widetilde\omega(w)=\int_0^\infty\omega(x)x^{w-1}\dd x$ has a simple pole at $w=0$ with residue $1$ and decays faster than any power of $|\Im w|$ in vertical strips. Move the line to $\Re w=-B$. The residue at $w=0$ is the product of zeta values; the poles at $w=1-s-\gamma_k/L$ have $|\Im w|\asymp T$ and contribute $O_A(T^{-A})$. On $\Re w=-B$ the functional equation gives $|\zeta(s+w+\gamma_k/L)|\ll T^{B+1/2}$, so the integral is $O(T^{K(B+1/2)-8B})$, which is $O(T^{-A})$ for $B$ large in terms of $A$, since $K\le4$. (For $K>4$ the same argument applies with $T^{2K}$ in place of $T^8$; we only use $K\in\{1,3,4\}$.) Derivatives are treated in the same way.
\end{proof}

\begin{lemma}[Short-side diagonal rule]\label{lem:shortside}
Let $\theta<1$ and let $B(s)=\sum_{n\le T^\theta}b(n)n^{-s}$ with $|b(n)|\ll n^{\epsilon}$. Let $B'_t(s)=\sum_{m\le T^{K}}\beta_t(m)m^{-s}$, where each $\beta_t(m)$ is smooth in $t\in[T,2T]$ with $\partial_t^j\beta_t(m)\ll_jT^{-j}m^{\epsilon}$ and $K$ is fixed. Then for $\Phi\in\mathcal W$
\[
 \int\Phi(t/T)\,B'_t(s)\,\overline{B(s)}\dd t=\sum_{n\le T^\theta}\frac{\overline{b(n)}}n\int\Phi(t/T)\beta_t(n)\dd t+O_A(T^{-A}).
\]
\end{lemma}
\begin{proof}
The off-diagonal terms are $\overline{b(n)}(mn)^{-1/2}\int\Phi(t/T)\beta_t(m)(n/m)^{it}\dd t$ with $m\ne n$, $n\le T^\theta$, so $|\log(m/n)|\ge\log(1+1/n)\gg T^{-\theta}$. Integrating by parts $j$ times gives $O_j(T(T^{1-\theta})^{-j}m^{\epsilon})$. Summing over the $O(T^{K+1})$ pairs and choosing $j$ large proves the claim.
\end{proof}

\begin{lemma}[Dual short factor]\label{lem:dualfactor}
Let $g\in C_c^\infty((0,b))$ with $b<1$. Uniformly for $T\le t\le2T$,
\[
 \chi(s)\overline{A_g(s)}=\sum_{k\ge1}k^{-s}\,g\Big(\frac{\log(t/(2\pi k))}{L}\Big)+O_g(T^{-1/2}\log T).
\]
\end{lemma}
\begin{proof}
$\overline{A_g(s)}=\sum_ng(\log n/L)n^{-1/2+it}$. Split $x^{-1/2}g(\log x/L)e^{it\log x}$ by a smooth dyadic partition and apply Poisson summation on each block $x\asymp X$. The frequency $k$ contributes $\int x^{-1/2}g(\log x/L)e^{it\log x-2\pi ikx}\dd x$. For $k\ge1$ with $k\asymp T/X$ the phase has a non-degenerate stationary point at $x_0=t/(2\pi k)$, where the second derivative is $\asymp T/X^2$; after rescaling $x=Xy$ the amplitude has derivatives $O_j(X^{-1/2})$, and one-dimensional stationary phase gives
\[
 k^{-1/2}e^{it\log(t/(2\pi k))-it-i\pi/4}g(\log x_0/L)+O_g(X^{1/2}T^{-3/2}),
\]
including the partition factor at $x_0$; the constant phase is $-\pi/4$ because the second derivative of $t\log x-2\pi kx$ is negative. Away from a fixed enlargement of the stationary range, including $k\le0$, repeated integration by parts gives $O_{g,A}(X^{1/2}(T+|k|X)^{-A})$. Summing over $k$ and the $O(L)$ blocks gives the error $O_g(T^{-1/2}\log T)$. Stirling's formula $\chi(s)=e^{-it\log(t/2\pi)+it+i\pi/4}(1+O(1/T))$ removes the phase, and the dual amplitudes have absolute sum $O(T^{1/2})$, so the Stirling error is admissible.
\end{proof}

\subsection{The unequal-length numerator}
\begin{proposition}\label{prop:unequal}
For every $V'$-test, $\mathrm N_T(\Pi V)/\mathcal N_T\to\iint f(a)g(b)\,\mathcal K_p(a,b)\dd a\dd b$.
\end{proposition}
\begin{proof}
Put $\varkappa=a_0+b_0<\frac12$. By linearity and the Cauchy integral formula in a shift variable (as in \cref{lem:degree}) it suffices to treat $p(X)=e^{\zeta_0X}$ for $\zeta_0$ in a fixed compact set, for which $C_p(s)=e^{\zeta_0/2}\zeta(s+\zeta_0/L)$, with all estimates uniform in $\zeta_0$. Since $\zeta(s)^3\overline{\zeta(s+\zeta_0/L)}=\chi(1-s+\bar\zeta_0/L)\zeta(s)^3\zeta(s-\bar\zeta_0/L)$ on the critical line, and $e^{\bar\zeta_0/2}\chi(1-s+\bar\zeta_0/L)/\chi(1-s)\to e^{-\bar\zeta_0/2}$ in $C^j$ of $t/T\in[1,2]$ for each $j$ (differentiated Stirling), it suffices to evaluate, for four scaled shifts $\boldsymbol\sigma=(\sigma_1,\dots,\sigma_4)$ in a compact set and $\boldsymbol\alpha=\boldsymbol\sigma/L$,
\[
 I(\boldsymbol\alpha)=\int\Phi(t/T)\,\chi(1-s)\prod_{i=1}^4\zeta(s+\alpha_i)\,\overline{A_f(s)A_g(s)}\dd t,
\]
uniformly for $\Phi$ in a bounded subset of $C_c^\infty((1,2))$; the case needed is $\boldsymbol\sigma=(0,0,0,-\bar\zeta_0)$.

\emph{Stationary phase.} Write $f_h=f(\log h/L)$, $g_k=g(\log k/L)$, $N=hk\le T^{\varkappa}$, and let $d_{\boldsymbol\alpha}(m)$ be the coefficients of $\prod_i\zeta(s+\alpha_i)$. Insert the expansion of \cref{lem:expansion} (with a cutoff $e^{-m/T^{K}}$ instead, for $K$ large; the difference is immaterial). The phase of $\chi(1-s)(N/m)^{it}$ is $t\log(t/2\pi)-t-\frac\pi4+t\log(N/m)$, stationary at $t_0=2\pi m/N$, and the amplitude of $\chi(1-s)$ is $1+O(T^{-1})$ with symbol bounds. For $m\asymp NT$ stationary phase gives $2\pi\sqrt{m/N}\,e(-m/N)\Phi(t_0/T)+O(T^{-1/2})$, where $e(x)=e^{2\pi ix}$; outside a fixed range of $m/(NT)$ integration by parts gives any power saving. Including $(mN)^{-1/2}$ and summing over $m\asymp NT$ gives an error $O(T^\epsilon)$ per pair $(h,k)$, so
\begin{equation}\label{eq:saddle}
 I(\boldsymbol\alpha)=2\pi\sum_{h,k}\frac{f_hg_k}{N}\sum_md_{\boldsymbol\alpha}(m)e(-m/N)\Phi\Big(\frac{2\pi m}{NT}\Big)+O_\epsilon(T^{\varkappa+\epsilon}).
\end{equation}

\emph{Characters.} For $q=N/d$,
\[
 e(-m/N)=\sum_{d\mid(m,N)}\frac1{\varphi(q)}\sum_{\xi\bmod q}\tau(\bar\xi)\,\xi(-m/d).
\]
After Mellin inversion of $\Phi$, the Dirichlet series attached to the pair $(d,\xi)$ is $\sum_nd_{\boldsymbol\alpha}(dn)\xi(n)(dn)^{-w}$. On $\Re w=\frac12+\delta$, $0<\delta<\frac1{10}$, it equals $d^{-w}\prod_iL(w+\alpha_i,\xi)$ times local corrections at the primes dividing $d$, whose product is $O_\epsilon(d^\epsilon)$ by the divisor bound. We use the fourth-moment bound
\begin{equation}\label{eq:fourth}
 \frac1{\varphi(q)}\sum_{\xi\bmod q}|L(w,\xi)|^4\ll_{\delta,\epsilon}q^\epsilon(1+|\Im w|)^{C_{\delta,\epsilon}}\qquad(\Re w=\tfrac12+\delta+O(1/L)).
\end{equation}
To prove it, group the characters by their primitive inducing character of conductor $r\mid q$ and parity; removing the Euler factors at primes dividing $q$ costs $O_\epsilon(q^\epsilon)$. The approximate functional equation for $L(w,\xi)^2$, $\xi$ primitive modulo $r$, has two smoothed sums of length $r^{1+\epsilon}(1+|\Im w|)^{C}$ with coefficients $d_2(n)n^{-w}$ and $d_2(n)n^{-(1-w)}$, the latter multiplied by a factor of absolute value $\ll(r(1+|\Im w|))^{1-2\Re w}$. The large sieve for characters modulo $r$ bounds the mean square of a block of length $Y$ by $(r+Y)\sum|a_n|^2$; the prefactor absorbs the growth of the dual coefficients. Dyadic summation gives $O(r^{1+\epsilon}(1+|\Im w|)^C)$ for the sum over characters modulo $r$, and summing over $r\mid q$ gives \eqref{eq:fourth}; H\"older's inequality then bounds the mean of the product of four shifted $L$-functions. For non-principal $\xi$ move the Mellin contour to $\Re w=\frac12+\delta$; there are no poles. Since $|\tau(\xi)|\le\sqrt q$, the contribution for fixed $N$, before the factor $1/N$ in \eqref{eq:saddle}, is
\[
 \ll_\epsilon(NT)^{1/2+\delta}\sum_{d\mid N}d^{-1/2-\delta+\epsilon}(N/d)^{1/2+\epsilon}\ll_\epsilon T^{1/2+\delta+\epsilon}N^{1+\delta+\epsilon}.
\]
Summing over $h,k$ with weight $1/N$ gives $O(T^{1/2+\delta+\epsilon+\varkappa(1+\delta+\epsilon)})=O(T^{1-\eta})$ for some $\eta>0$, on choosing $\delta,\epsilon$ small in terms of $\frac12-\varkappa$. For principal characters the same contour shift crosses only the poles at $w=1-\alpha_i$, and the new line obeys the same bound.

\emph{Residues.} For the principal characters, write $w=1+\omega$ and $D_p(w)=\prod_i(1-p^{-w-\alpha_i})^{-1}$. After summing over $d\mid N$, the local factor at $p^\nu\Vert N$ is
\begin{gather*}
 B_p(0,w)=D_p(w),\\ B_p(\nu,w)=\sum_{r\ge\nu}d_{\boldsymbol\alpha}(p^r)p^{-rw}-\frac{d_{\boldsymbol\alpha}(p^{\nu-1})p^{-(\nu-1)w}}{p-1}\quad(\nu\ge1),
\end{gather*}
because the principal Gauss sum modulo $q$ is $\mu(q)$, so only $v_p(d)\in\{\nu,\nu-1\}$ survive. With Mellin variables $u,v'$ for the profiles $f,g$, the Euler product of the local factors $\sum_{j,k\ge0}B_p(j+k,1+\omega)p^{(j+k)\omega-ju-kv'}$ factors as
\begin{equation}\label{eq:euler}
 G(\boldsymbol\alpha,\omega,u,v')\prod_i\zeta(1+\omega+\alpha_i)\prod_{h\in\{u,v'\}}\frac{\prod_i\zeta(1+h+\alpha_i)}{\zeta(1+h-\omega)},
\end{equation}
with $G$ absolutely and locally uniformly convergent when all real parts are at most a small $\delta'$ in absolute value: indeed $p^\omega B_p(1,1+\omega)/D_p(1+\omega)=\sum_ip^{-1-\alpha_i}-p^{-1+\omega}+O(p^{-2+O(\delta')})$ and the terms with $j+k\ge2$ are $O(p^{-2+O(\delta')})$. At the origin the unfactored local series is
\[
 \sum_{r\ge0}\binom{r+3}3\binom{r+2}2X^r-\frac X{1-X}\sum_{r\ge0}(r+2)\binom{r+3}3X^r=\frac{1+4X+X^2}{(1-X)^6}
\]
with $X=1/p$,
as one sees by applying $X\partial_X$ to $(1-X)^{-4}$. Removing the ten net polar factors leaves $(1-X)^4(1+4X+X^2)$, so $G(0,0,0,0)=a_3$.

First let the $\sigma_i$ be pairwise distinct. At $\omega=-\alpha_j$ the two denominators in \eqref{eq:euler} cancel the $j$-th numerator factors. With
\[
 F_j(x)=\mathcal L^{-1}_{u\to x}\prod_{i\ne j}(u+\sigma_i)^{-1},\qquad
 S(\boldsymbol\sigma;x,y)=\sum_{j=1}^4\frac{e^{-\sigma_j}F_j(x)F_j(y)}{\prod_{i\ne j}(\sigma_i-\sigma_j)},
\]
the residue sum divided by $a_3TL^9$ tends to $\int\Phi\cdot\iint f(x)g(y)S(\boldsymbol\sigma;x,y)\dd x\dd y$: the three zeta factors at the shift differences give $L^3/\prod_{i\ne j}(\sigma_i-\sigma_j)$, the remaining six give $L^6$ times the inverse Laplace transforms (dominated on contours $\Re u=\Re v'=c/L$ by the rapid decay of the transforms of $f,g$), $(T/2\pi)^{-\alpha_j}\to e^{-\sigma_j}$ and $\widetilde\Phi(1-\alpha_j)\to\widetilde\Phi(1)=\int\Phi$. To remove the separation, write $Q(u)=\prod_i(u+\sigma_i)$ and $F(x,v)=\mathcal L^{-1}_{u\to x}((u-v)/Q(u))$; then $F(x,-\sigma_j)=F_j(x)$ and the residue theorem gives
\[
 S(\boldsymbol\sigma;x,y)=\frac1{2\pi i}\oint\frac{e^vF(x,v)F(y,v)}{Q(v)}\dd v
\]
over a circle enclosing the roots of $Q$. This is holomorphic in $\boldsymbol\sigma$, and the convergence extends from separated shifts to all shifts in a compact set by the maximum principle on polydiscs.

\emph{Identification.} For $\boldsymbol\sigma=(0,0,0,-z)$, $z=\bar\zeta_0$, we have $Q(u)=u^3(u-z)$, and an exact computation gives, for $0<x\le y$ with $x+y<1$,
\[
 e^{-z/2}S(\boldsymbol\sigma;x,y)=\mathcal K_{\exp(z\,\cdot)}(x,y),
\]
both sides being computed in closed form; the right side is evaluated by the Hermite--Genocchi formula $\E\exp(z\sum_ic_iX_i)=6z^{-3}\,[c_0,c_1,c_2,c_3]\exp(z\,\cdot)$ for the third divided difference. Both sides are entire in $z$ and symmetric in $(x,y)$. Collecting the factor $e^{-\bar\zeta_0/2}$ from the reduction, taking real parts and differentiating in $\zeta_0$ at $0$ gives the claim for every polynomial $p$.
\end{proof}

\begin{remark}
This recovers, with a different organisation, the four-residue calculation of Durkan and Page \cite[\S5.2]{DP} for same-side amplifiers, with the product of the two lengths, rather than their maximum, controlling the error. Dirichlet GRH would not improve the exponent of $N$ in the non-principal estimate, so the argument is the same with or without it.
\end{remark}

\subsection{Proof of \Cref{thm:third}}
\emph{Faithfulness.} Pairs within $\mathcal T_2$ are covered by \cref{lem:bridge}. We use $\ip{\Pi H}{\Pi K}=\frac12\operatorname{Re}\ip HK+\frac12\operatorname{Re}\ip H{\mathcal JK}$ for real-weighted pairings, where $\mathcal JK=\chi^3\overline K$, which holds because $\ip{\mathcal JH}{\mathcal JK}=\overline{\ip HK}$ and $\ip{\mathcal JH}{K}=\overline{\ip H{\mathcal JK}}$; on the Schur side the corresponding identity is $\ip{\frac12(h+h\circ J)}{\frac12(h'+h'\circ J)}=\frac12\ip h{h'}+\frac12\ip h{h'\circ J}$.

(a) \emph{Pairings with a $W$-test.} Let $\theta<1$ with $\operatorname{supp}\phi\subset(0,\theta)$, so $D_\phi$ has length $\le T^\theta$. Every other test $K$ in $\mathcal T_3$, and $\zeta^3$, is, after \cref{lem:expansion} for its zeta factors and \cref{lem:dualfactor} for a reflected short factor $\chi\overline{A_g}$, of the form $B'_t$ in \cref{lem:shortside}, up to an error whose weighted $L^1$-norm against $|D_\phi|$ is $o(TL^9)$. For $U$ this uses $A_f\ll T^{a_0/2}$, $C_p\ll_\epsilon T^{1/4+\epsilon}$ (convexity and Cauchy's formula), $\int\Phi|D_\phi|\ll TL^{9/2}$ and $-\frac12+\frac{a_0}2+\frac14<0$, since $a_0<\frac14$. The diagonal of \cref{lem:shortside} imposes $n=m_1m_2m_3$ with $n\le T^\theta$, and on it the dual profile $g(\log(t/(2\pi k))/L)$ may be replaced by $g(1-\log k/L)$ at a cost $O(1/L)$ per term, hence $O(TL^8)$ in total by \cref{lem:diagonal}. Then \cref{lem:diagonal} evaluates the direct pairings $\ip{D_\phi}{K}$, $\ip{\zeta^3}{D_\phi}$ as $a_3TL^9\int\Phi$ times $\int\phi(|R(t)|)G(C(t))\dd t$, where $G$ is the cubic weight of $K$ (respectively $\int\phi(|R(t)|)\dd t$), and by \cref{lem:stable} these are $\ip{h_{F_\phi}}{h_G}$ and $\ip{v}{h_{F_\phi}}$. The reflected pairings: $\ip{D_\phi}{\mathcal JD_\psi}$ has no stationary point, since the product of the two indices is $<T^2$ while the phase of $\chi^3$ requires order $T^3$, so it is $O(T^{-A})$; $\ip{D_\phi}{\mathcal JV}$ is negligible because after $\chi\overline{C_p}=C_{p^-}+O(\cdot)$ (the centred functional equation, with squared weighted error $O(TL^7)$ as in \cite[\S2.3]{Fraccaro}) the remaining phase has derivative $-2\log(t/2\pi)-\log m+\log(nhk)+O(1/T)$ with $nhk\le T^{\theta+1/2}$; and $\mathcal JU=\chi^2\overline{C_p}\,\overline{A_f}A_g$ is, by the same two replacements, a $U$-type product with $(f,g,p)$ replaced by $(g,f,p^-)$, to which the direct argument applies. On the Schur side \cref{lem:stable} gives $\ip{h_{F_\phi}}{h_G\circ J}=0$ whenever $G$ vanishes on $|c|\ge2-\delta$, which covers $F_\psi$ and $F_V$, and for $U$ it gives the profile of the reflected data, in agreement.

(b) \emph{$V'$-tests.} The numerator is \cref{prop:unequal}. For Gram entries among $U$, $V$ and $V'$-tests, Fraccaro's evaluation \cite[\S2.3]{Fraccaro} uses only that the four short exponents of the two tests sum to less than $1$, which holds for every pair here (for $U$ with $V'$: $a_0+b_0+\varkappa<\frac12+\frac12$); the limits are therefore those of \Cref{thm:fraccaro}, and the proof of \cref{lem:bridge}, which used only this condition, identifies them with Schur inner products.

\emph{The span.} The proof of \cref{lem:density} applies verbatim on boxes in $\{x+y<\frac12\}$, assigning the short factors of a $V'$-test to the coordinates $x,y$ (with $a_0,b_0$ slightly above their ranges and $a_0+b_0<\frac12$) and the derivative to $z$; with \cref{lem:density} and reflection this gives $L^2_J(\mathcal D_{\rm acc})\subseteq\mathcal H(\mathcal T_3)$. Let $0\le\phi_n\le1$ in $C_c^\infty((0,1))$ increase to $\one_{(0,1)}$. The profiles $\frac12(\phi_n(r)+\phi_n(3-r))v$ of the $W$-tests converge in $L^2$ to $\frac12v\one_{\mathcal A}$, so $v\one_{\mathcal A}\in\mathcal H(\mathcal T_3)$. Since $v\one_{\mathcal A\cap\mathcal D_{\rm acc}}\in L^2_J(\mathcal D_{\rm acc})$, we get $v\one_{S_3}=v\one_{\mathcal D_{\rm acc}}+v\one_{\mathcal A}-v\one_{\mathcal A\cap\mathcal D_{\rm acc}}\in\mathcal H(\mathcal T_3)$.

\emph{The measure.} Up to null sets $\Om\setminus\mathcal D_{\rm acc}=\mathcal B$, by the definitions, so $\Om\setminus S_3=\mathcal B\setminus\mathcal A$ and $\mu(S_3)=42-\mu(\mathcal B)+\mu(\mathcal B\cap\mathcal A)$. On $\{r<1\}$, $\mathcal B$ is $E=\{r<1,\ x+y>\frac12\}$: there $z<\frac12$, so $\varrho>\frac12>x$, $3x+5\varrho>2$ and $y+z<\frac32$ hold automatically. By $J$-symmetry $\mu(\mathcal B\cap\mathcal A)=2\mu(E)$. The substitution $b_1=y+z-\frac12$, $b_2=x+z-\frac12$, $b_3=x+y-\frac12$ maps $E$ onto $\{b_1>b_2>b_3>0,\ \sum b_i<\frac12\}$, preserves the three pairwise differences up to order, and has Jacobian $2$. Hence $\mu(E)=\frac12\cdot2^{-9}\mu(\{r<1\})$, and $\mu(\{r<1\})=9!\int_{|t|<1}\dd t=1$ by \cref{lem:stable} with $F=\one_{\{|c|<1\}}$ (approximated by continuous functions). So $2\mu(E)=\frac1{512}$, which is also checked by direct integration. Together with \eqref{eq:L0} this gives $\mu(S_3)$, and \cref{prop:projection} gives the lower bound. \qed

\begin{remark}
One can also adjoin, on top of $\mathcal T_3$, three-factor tests just past total degree one, evaluated with Heath-Brown's theorem on $d_3$ in progressions \cite{HB}. Their gain lies, inside the region $N_0$ of \Cref{thm:marked}, in a strip $\{40r+21k<51\}$ that is contained in the balanced region $\mathcal R$ of \Cref{thm:row}. They therefore do not affect \Cref{thm:main}, and we omit them.
\end{remark}

\section{Marked tests and the spectral estimate}\label{sec:marked}

In this section we enlarge the support from $S_3$ to $S_3\cup\mathcal Z$, where
\[
 \mathcal Z=\{(x,y,z)\in\Om:\ z-x>\tfrac12\}
\]
is the \emph{spread region}. It is $J$-invariant, since $J$ preserves $z-x$. Write $\mathcal D_+=\{x+y<z\}$ for the dominant shapes. Every point of $\mathcal Z$ lies in $\mathcal D_+\cup J\mathcal D_+$: as $J\mathcal D_+=\{y>1-(z-x)\}$ and $1-(z-x)<\frac12<z-x$, a point with $y\ge z-x$ has $y>1-(z-x)$.

We use the coordinates $(r,k,x)$, $r=x+y+z$, $k=x+y$, on $\Om$; the map $(x,y,z)\mapsto(r,k,x)$ has Jacobian $\pm1$. Let
\[
 \mathcal Y=\Big\{(r,k,x):\ 1<r<2,\ \max(r-1,\tfrac12)<k<\tfrac r2,\ 0<x<\tfrac k2,\ r-k-x>\tfrac12\Big\},
\]
an open set. For a function $\psi$ on $\mathcal Y$ put $\psi^\flat(x,y,z)=\psi(x+y+z,x+y,x)$, extended by $0$. The conditions defining $\mathcal Y$ say exactly that $(x,y,z)=(x,k-x,r-k)\in\Om\cap\mathcal D_+\cap\mathcal Z$ with $1<r<2$ and $k>\frac12$: indeed $x<k/2$ is $x<y$, $k<r/2$ is $x+y<z$ (and implies $y<z$), $k>r-1$ is $z<1$, and $r-k-x>\frac12$ is $z-x>\frac12$.

\begin{theorem}\label{thm:marked}
For real $\psi\in C_c^\infty(\mathcal Y)$ let $W_\psi$ be the marked test of \cref{def:marked} below. Let
\[
 \mathcal T_{\mathcal M}=\big\{(W_\psi,\psi^\flat),\ (\mathcal JW_\psi,\psi^\flat\circ J):\ \psi\in C_c^\infty(\mathcal Y)\text{ real}\big\},\qquad \mathcal T_4=\mathcal T_3^\infty\cup\mathcal T_{\mathcal M},
\]
where $\mathcal JB=\chi(s)^3\overline{B}$ and $\mathcal T_3^\infty$ is the set of $U$-, $V'$- and $W$-tests in $\mathcal T_3$ (all of which have smooth data). Then $\mathcal H(\mathcal T_3^\infty)=\mathcal H(\mathcal T_3)$, $\mathcal T_4$ is faithful, $v\one_{S_3\cup\mathcal Z}\in\mathcal H(\mathcal T_4)$, and
\[
 \mu(S_3\cup\mathcal Z)=42-2\mu(N_0)=42-\frac{209}{512}=41.591796875,
\]
where $N_0=\{(x,y,z)\in\Om:\ 1<r<\tfrac32,\ k>\tfrac12,\ z-x<\tfrac12\}$.
\end{theorem}

The proof occupies the rest of the section; it is completed in \Cref{sec:markedproof}. Its analytic core is a mixed pairing (\cref{prop:mixed}) between a marked test and an arbitrary smooth cubic source, which rests on the signed estimate \cref{thm:signed} for the ternary divisor function on rational arcs. That estimate is proved with the Kuznetsov formula at an Atkin--Lehner cusp of $\Gamma_0(M)$, where the numerator phase disappears and the shifted divisor coefficients factor through three standard $L$-functions; the levels $M$ are then averaged with the large sieves for exceptional eigenvalues of Deshouillers--Iwaniec and Watt.

Throughout this section, $\Phi$ denotes a smooth function supported in a fixed compact subset of $(1,2)$, possibly complex-valued, and statements are uniform for $\Phi$ in bounded subsets of $C_c^\infty$ (with fixed support). We write
\[
 \ip B{B'}_\Phi=\int_\R\Phi(t/T)\,B(\tfrac12+it)\,\overline{B'(\tfrac12+it)}\dd t,\qquad\|B\|_\Phi^2=\ip BB_{|\Phi|}.
\]
For $\Phi\in\mathcal W$ this is $\mathrm G_T(B,B')$. All implied constants may depend on $\psi$, the support of $\Phi$, and fixed smooth weights; $\epsilon>0$ is arbitrarily small and may change between occurrences.

\subsection{Marked witnesses}
\begin{definition}\label{def:marked}
Let $\psi\in C_c^\infty(\mathcal Y)$ be real and put $m_\psi(r,k,u)=-\partial_3\psi(r,k,u)$. For $q\ge1$ and $r\in\R$ let
\[
 \beta_\psi(q;r)=\sum_{lde=q}\frac{\mu(l)}l\,m_\psi\Big(r,\frac{\log q}L,\frac{\log d}L\Big),
\]
and for $n\ge1$
\[
 M_\psi(n)=L^2\sum_{q\ge1}\frac{c_q(n)}q\,\beta_\psi\Big(q;\frac{\log n}L\Big),\qquad W_\psi(s)=\sum_{n\ge1}M_\psi(n)n^{-s},
\]
where $c_q(n)=\sum_{a\bmod q}^*e(an/q)$ is Ramanujan's sum.
\end{definition}

The sums are finite: $m_\psi(r,\cdot,\cdot)=0$ unless $r$ lies in a compact subset $I_\psi\subset(1,2)$, so $n\le T^2$, and then $q\le T^{r/2-\delta}$ for some $\delta=\delta_\psi>0$. Since $|\beta_\psi(q;r)|\le\|m_\psi\|_\infty\tau_3(q)$ and $|c_q(n)|\le(q,n)$, we have $M_\psi(n)\ll_\epsilon n^\epsilon$. More precisely, on the support of $\beta_\psi(q;\log n/L)$,
\begin{equation}\label{eq:qrange}
 T^{r-1+\delta}\le q\le T^{r/2-\delta},\qquad r=\frac{\log n}L,
\end{equation}
because $k=\log q/L$ ranges over a compact subset of $(r-1,r/2)$. With $X=T^r$ and $H=X/T$ this says $HX^{\delta}\le q\le X^{1/2-\delta}$.

The \emph{cubic weight} of $W_\psi$ is
\begin{equation}\label{eq:cpsi}
 c_\psi(a_1,a_2,a_3)=(\partial_u\partial_vf_\psi)(a_1+a_2+a_3,a_1,a_2),\quad f_\psi(r,u,v)=m_\psi(r,u+v,u),
\end{equation}
the derivatives being taken with $r$ fixed. Since $m_\psi(r,k,u)$ vanishes unless $0<u<k/2$ and $r-1<k<r/2<1$, the function $c_\psi$ is smooth with compact support in $(0,1)^3$: on its support $a_1=u$, $a_2=k-u$ and $a_3=r-k$ lie in $(0,1)$.

\begin{lemma}[Marked profile]\label{lem:markedprofile}
$h_{c_\psi}=\psi^\flat$.
\end{lemma}
\begin{proof}
Fix $(x,y,z)\in\Om$ and put $r=x+y+z$, $F(a_1,a_2)=f_\psi(r,a_1,a_2)=m(r,a_1+a_2,a_1)$, writing $m=m_\psi$. Since $c_\psi$ has compact support in $(0,1)^3$, \eqref{eq:profile} reads $h_{c_\psi}(x,y,z)=\iint\partial_1\partial_2F(a_1,a_2)\,\kappa(x,y,z;a)\dd a_1\dd a_2$ with $a_3=r-a_1-a_2$. By definition $\kappa$ is the length of the set of $u\in[y,z]$ with $w:=a_1+a_2-u\in[x,y]$ and $\max(a_1,a_2)\le u$; the last condition is equivalent to $w\le a_1\le u$ given $u+w=a_1+a_2$. Changing variables $(a_1,a_2,u)\mapsto(a_1,w,u)$, which has Jacobian $1$,
\[
 h_{c_\psi}(x,y,z)=\int_y^z\!\!\int_x^y\Lambda(u,w)\dd w\dd u,\quad\Lambda(u,w)=\int_w^u(\partial_1\partial_2F)(a,u+w-a)\dd a .
\]
Now $\partial_1\partial_2F(a_1,a_2)=(m_{kk}+m_{kx})(r,a_1+a_2,a_1)$, where subscripts denote partial derivatives in the second and third slots, so with $k'=u+w$,
\[
 \Lambda(u,w)=\partial_k^2\Big[\int_w^um(r,k,a)\dd a\Big]_{k=k'}+m_k(r,k',u)-m_k(r,k',w).
\]
For $u>w$ we have $u>k'/2$, and $m(r,k,a)=0$ for $a\ge k/2$; hence for $k$ near $k'$ the bracket equals $\int_w^{k/2}m(r,k,a)\dd a=\psi(r,k,w)$, and $m_k(r,k',u)=0$. Thus $\Lambda(u,w)=(\psi_{kk}+\psi_{kx})(r,u+w,w)=\partial_w\partial_u G(u,w)$ with $G(u,w)=\psi(r,u+w,w)$, and
\[
 h_{c_\psi}(x,y,z)=G(z,y)-G(z,x)-G(y,y)+G(y,x).
\]
Here $G(z,y)=\psi(r,y+z,y)$ and $G(z,x)=\psi(r,x+z,x)$ vanish because $y+z>r/2$ and $x+z>r/2$, while $G(y,y)=\psi(r,2y,y)$ vanishes because $\psi(r,k,\cdot)$ vanishes near $k/2$. The remaining term is $G(y,x)=\psi(r,x+y,x)=\psi^\flat(x,y,z)$.
\end{proof}

\noindent The identity has also been checked by numerical quadrature.

We shall also need the effect of reflection on profiles. For a bounded Borel $F$ on $[0,1]^3$ put $F^J(c)=F(1-c_1,1-c_2,1-c_3)$.

\begin{lemma}\label{lem:reflectedprofile}
$h_{F^J}=h_F\circ J$ for continuous $F$.
\end{lemma}
\begin{proof}
By \ref{S:dual}, $D\overline{e_i}=e_{N-i}$, so $D^3\overline{E_N[F]}=E_N[F^J]$ up to the factor $F^J(m/N)$ versus $F(1-m/N)$, which agree. Hence $c_\lambda(F^J)=\ip{D^3\overline{E_N[F]}}{s_\lambda}_N=\overline{\ip{E_N[F]}{D^3\overline{s_\lambda}}_N}=c_{\lambda^*}(F)$, using that $D^3\overline{\,\cdot\,}$ is antiunitary and $D^3\overline{s_\lambda}=s_{\lambda^*}$. Since $\lambda^*/N=J(\lambda/N)$, \cref{lem:profiles}(i) gives the claim.
\end{proof}

\subsection{Ramanujan polynomials}
The norms and pairings of marked tests with one another are explicit; no estimate for $d_3$ is involved. We first record the arithmetic of the denominators.

\begin{lemma}[Local factors]\label{lem:localfactor}
For complex $\alpha,\beta$ with positive real parts, $q\ge1$ and $n\ge1$,
\begin{equation}\label{eq:ramanujaninv}
 \sum_{d,e\ge1:\ de\mid n}d^{-\alpha}e^{-\beta}=\zeta(1+\alpha)\zeta(1+\beta)\sum_{q\ge1}\frac{c_q(n)}q\,\mathcal D_q(\alpha,\beta),
\end{equation}
where $\mathcal D_q(\alpha,\beta)=\sum_{lde=q}\mu(l)l^{-1-\alpha-\beta}d^{-\alpha}e^{-\beta}$, the right side converging absolutely. Moreover
\begin{equation}\label{eq:Dlocal}
 \sum_{j\ge0}\mathcal D_{p^j}(\alpha,\beta)Y^j=\frac{1-p^{-1-\alpha-\beta}Y}{(1-p^{-\alpha}Y)(1-p^{-\beta}Y)},
\end{equation}
and $q^{-1}\mathcal D_q(\alpha,\beta)\zeta(1+\alpha)\zeta(1+\beta)=\sum_{d,e\ge1:\,q\mid de}d^{-1-\alpha}e^{-1-\beta}$.
\end{lemma}
\begin{proof}
\eqref{eq:Dlocal} is immediate from the definition, since $\mu(l)\ne0$ only for $l\in\{1,p\}$ at $p$. For the last identity, both sides are multiplicative in $q$ after dividing by $\zeta(1+\alpha)\zeta(1+\beta)$: the right side divided by it is $\prod_p\Pr(p^{v_p(q)}\mid UV)$, where $U,V$ are independent random positive integers with $\Pr(U=d)=d^{-1-\alpha}/\zeta(1+\alpha)$ (as formal weights), whose valuations are independent and geometric with parameters $P=p^{-1-\alpha}$ and $R=p^{-1-\beta}$. Then $\sum_{j\ge0}(pY)^j\Pr(v_p(UV)\ge j)$ equals $[1-pY(1-P)(1-R)/((1-pPY)(1-pRY))]/(1-pY)$, which is \eqref{eq:Dlocal} (a rational-function identity, also checked by computer). Finally, $\one_{b\mid n}=b^{-1}\sum_{q\mid b}c_q(n)$ gives
\[
 \sum_{de\mid n}\frac1{d^{\alpha}e^{\beta}}=\sum_{d,e}\frac{1}{d^{1+\alpha}e^{1+\beta}}\sum_{q\mid de}c_q(n)=\sum_qc_q(n)\!\sum_{q\mid de}\frac{1}{d^{1+\alpha}e^{1+\beta}},
\]
the interchange being justified by the bound $|c_q(n)|\le(q,n)$ and the convergence of $\sum_b\tau(b)^2b^{-1-\min(\Re\alpha,\Re\beta)}$.
\end{proof}

\begin{lemma}[Four-edge limit]\label{lem:fouredge}
Let $m_0\in\{\varphi,\mathrm{id}\}$, and let $w$, $\boldsymbol a=(a_1,a_2)$, $\boldsymbol b=(b_1,b_2)$ be complex with real parts in $(-\frac1{10},\frac1{10})$. If moreover $\Re(w+a_i+b_j)>0$ for all $i,j$, the left side of \eqref{eq:Gfactor} converges absolutely and
\begin{equation}\label{eq:Gfactor}
 \sum_{q\ge1}\frac{m_0(q)}{q^{2+w}}\,\mathcal D_q(a_1,a_2)\,\mathcal D'_q(b_1,b_2)=\mathcal G_{m_0}(w,\boldsymbol a,\boldsymbol b)\prod_{i,j=1}^2\zeta(1+w+a_i+b_j),
\end{equation}
where $\mathcal D'_q(b_1,b_2)=\sum_{lde=q}\mu(l)l^{-1}d^{-b_1}e^{-b_2}$ and $\mathcal G_{m_0}$ is given by an Euler product converging absolutely and uniformly when all real parts are at most $\frac1{10}$ in absolute value, uniformly in the imaginary parts, and $\mathcal G_\varphi(0,0,0)=a_3$. The right side of \eqref{eq:Gfactor} continues the left side meromorphically to the whole region where all real parts lie in $(-\frac1{10},\frac1{10})$; below, \eqref{eq:Gfactor} is only used where the left side converges absolutely. The same holds with $\mathcal D'$ in place of $\mathcal D$ in the first factor, with the same value at the origin.

Consequently, let $\Psi$ be smooth with compact support in $(0,\infty)^3$, let $a_1,a_2\in i\R$ be fixed, and write, for $t=(t_{ij})\in[0,\infty)^{2\times2}$, $R_i=t_{i1}+t_{i2}$, $C_j=t_{1j}+t_{2j}$ and $|t|=\sum t_{ij}$. Then, as $L\to\infty$,
\begin{align}
 &\frac1{L^4}\sum_{q\ge1}\frac{\varphi(q)}{q^2}\,\mathcal D_q\Big(\frac{a_1}L,\frac{a_2}L\Big)\sum_{lde=q}\frac{\mu(l)}l\,\Psi\Big(\frac{\log q}L,\frac{\log d}L\Big)\notag\\
 &\hspace{12em}\longrightarrow a_3\int_{[0,\infty)^{2\times2}}e^{-a_1R_1-a_2R_2}\,\Psi(|t|,C_1)\dd t,\label{eq:fouredge}\\
 &\frac1{L^4}\sum_{q\ge1}\frac{\varphi(q)}{q^2}\sum_{l'd'e'=q}\frac{\mu(l')}{l'}\sum_{lde=q}\frac{\mu(l)}l\,\Psi\Big(\frac{\log q}L,\frac{\log d'}L,\frac{\log d}L\Big)\notag\\
 &\hspace{12em}\longrightarrow a_3\int_{[0,\infty)^{2\times2}}\Psi(|t|,R_1,C_1)\dd t,\label{eq:fouredge2}
\end{align}
where in \eqref{eq:fouredge} $\Psi$ is a function of two variables. The left sides with absolute values inside are $O_\Psi(1)$.
\end{lemma}
\begin{proof}
At a prime $p$ the left side of \eqref{eq:Gfactor} has degree-one coefficient $m_0(p)p^{-2-w}(p^{-a_1}+p^{-a_2}-p^{-1-a_1-a_2})\times{}$ $(p^{-b_1}+p^{-b_2}-p^{-1})=\sum_{i,j}p^{-1-w-a_i-b_j}+O(p^{-2+4/10})$, which matches the degree-one coefficient of $\prod_{i,j}\zeta(1+w+a_i+b_j)$; by \eqref{eq:Dlocal} the higher coefficients are $O((j+1)^2p^{-j(1-4/10)})$ in modulus, uniformly in the imaginary parts. Hence the local factors of $\mathcal G_{m_0}$ are $1+O(p^{-6/5})$ uniformly, the second-degree terms giving the main error; this is summable over $p$, which proves the first claim. At the origin $\mathcal D_{p^j}(0,0)=\mathcal D'_{p^j}(0,0)$ is $1$ for $j=0$ and $j+1-j/p$ for $j\ge1$ (from \eqref{eq:Dlocal}), and with $P=1/p$ the local factor of the left side of \eqref{eq:Gfactor} for $m_0=\varphi$ is $1+(1-P)\sum_{j\ge1}P^j(1+j(1-P))^2$; multiplying by $(1-P)^4$ gives $(1-P)^9\sum_{j\ge0}\binom{j+2}2^2P^j$ (an identity of power series, both sides being $(1-P)^4(1+4P+P^2)$; also checked by computer). So $\mathcal G_\varphi(0,0,0)=a_3$.

For \eqref{eq:fouredge2} fix $w_0>0$ and write
\[
 \Psi(k,\varrho,\upsilon)e^{w_0k}=(2\pi)^{-3}\int\widehat\Psi_{w_0}(\xi)e^{i(\xi_1k+\xi_2\varrho+\xi_3\upsilon)}\dd\xi
\]
with $\widehat\Psi_{w_0}$ a Schwartz function. Then the left side of \eqref{eq:fouredge2} equals
\[
 (2\pi)^{-3}\int\widehat\Psi_{w_0}(\xi)\,L^{-4}\,\Sigma\Big(\frac{w_0-i\xi_1}L;\Big(\frac{-i\xi_2}L,0\Big),\Big(\frac{-i\xi_3}L,0\Big)\Big)\dd\xi,
\]
where $\Sigma(w;\boldsymbol a,\boldsymbol b)$ is the left side of \eqref{eq:Gfactor} with $m_0=\varphi$ and $\mathcal D'$ in both places; the interchange is justified by absolute convergence for $\Re w=w_0/L>0$. By \eqref{eq:Gfactor} the integrand is bounded by $|\widehat\Psi_{w_0}(\xi)|\sup|\mathcal G_\varphi|\,(L^{-1}\zeta(1+w_0/L))^4\ll|\widehat\Psi_{w_0}(\xi)|$, since $|\zeta(1+\sigma+i\tau)|\le\zeta(1+\sigma)$ and all shifts other than $w$ are purely imaginary. For fixed $\xi$ it converges, as $L\to\infty$, to $\widehat\Psi_{w_0}(\xi)\,a_3\prod_{i,j}(w_0-i\xi_1+\tilde a_i+\tilde b_j)^{-1}$ with $\tilde{\boldsymbol a}=(-i\xi_2,0)$ and $\tilde{\boldsymbol b}=(-i\xi_3,0)$, by continuity of $\mathcal G_\varphi$ at the origin and $L^{-1}\zeta(1+z/L)\to1/z$. Since
\[
 \prod_{i,j}(w+\tilde a_i+\tilde b_j)^{-1}=\int_{[0,\infty)^{2\times2}}e^{-w|t|-\sum_i\tilde a_iR_i-\sum_j\tilde b_jC_j}\dd t\qquad(\Re w>0),
\]
dominated convergence and Fourier inversion give \eqref{eq:fouredge2}. The proof of \eqref{eq:fouredge} is the same, with $\mathcal D_q(a_1/L,a_2/L)$ in the first factor, $\tilde{\boldsymbol a}=(a_1,a_2)$, and only the two Fourier variables $\xi_1,\xi_3$. The bounds with absolute values follow from the same Euler products with all shifts real and positive.
\end{proof}

\begin{lemma}[Ramanujan polynomials]\label[lemma]{lem:ramanujan}
Let $\delta>0$ and let $\beta(q;r)$, $\beta'(q;r)$ be functions, smooth in $r$, with $\partial_r^j\beta(q;r)\ll_j\tau_3(q)^{C}$ uniformly, vanishing unless $r$ lies in a fixed compact $I\subset(1,2)$ and
\begin{equation}\label{eq:qrange2}
 T^{r-1+\delta}\le q\le T^{r/2-\delta}.
\end{equation}
Put $M(n)=L^2\sum_qq^{-1}c_q(n)\beta(q;\log n/L)$, and similarly $M'$. Then
\[
 \ip{P[M]}{P[M']}_\Phi=TL^5\Big(\int\Phi\Big)\sum_{q\ge1}\frac{\varphi(q)}{q^2}\int\beta(q;r)\overline{\beta'(q;r)}\dd r+O(T^{1-\delta/2}L^{C}),
\]
where $P[M](s)=\sum_nM(n)n^{-s}$. In particular, taking for $\Phi$ a fixed smooth majorant $\Phi^+\ge|\Phi|$ with support in a fixed compact subset of $(1,2)$,
\[
 \|P[M]\|_\Phi^2\le\ip{P[M]}{P[M]}_{\Phi^+}\ll TL^5\sum_q\frac{\varphi(q)}{q^2}\sup_r|\beta(q;r)|^2+O(T^{1-\delta/2}L^C).
\]
\emph{One-sided bound.} Suppose instead that $\beta$ vanishes unless $r\in I$ and $q\le T^{r/2-\delta}$, with no lower bound on $q$. Then there is $c=c(I,\delta)>0$ such that
\[
 \|P[M]\|_\Phi^2\le\ip{P[M]}{P[M]}_{\Phi^+}\ll TL^5\sum_q\frac1q\sup_r|\beta(q;r)|^2+O(T^{1-c}),
\]
where the implied constants depend only on $I$, $\delta$, $\Phi^+$ and the constants in the bounds for $\partial_r^j\beta$.
\end{lemma}
\begin{proof}
Write $\ip{P[M]}{P[M']}_\Phi=T\sum_{n,m}M(n)\overline{M'(m)}(nm)^{-1/2}\widehat\Phi\big(\frac T{2\pi}\log\frac nm\big)$ with $\widehat\Phi(\xi)=\int\Phi(x)e(-x\xi)\dd x$. Insert a smooth partition $\sum_jV(m/2^j)=1$, $V\in C_c^\infty([\frac12,2])$, and consider one block $m\asymp X=2^j$, with $X=T^{r_0}$, $r_0\in I+[-\frac1L,\frac1L]$; put $H=X/T$. As $\widehat\Phi$ is Schwartz and $M,M'\ll n^\epsilon$, pairs with $|n-m|>HX^{\epsilon}$ contribute $O(T^{-A})$.

Since $\overline{M'(m)}=L^2\sum_qq^{-1}c_q(m)\overline{\beta'(q;\log m/L)}$, we need, for fixed $n$ and $q$, the sum $\sum_mc_q(m)G_n(m)$ with
\[
 G_n(y)=\overline{\beta'(q;\tfrac{\log y}L)}\,V\Big(\frac yX\Big)(ny)^{-1/2}\,\widehat\Phi\Big(\frac T{2\pi}\log\frac ny\Big).
\]
As $c_q(m)=\sum_{a\bmod q}^*e(am/q)$ and $\sum_me(am/q)G(m)=\sum_{\nu\in\Z}\widehat G(\nu-a/q)$ by Poisson summation, where $\widehat G(\xi)=\int G(y)e(-\xi y)\dd y$, we get
\[
 \sum_mc_q(m)G_n(m)=\sum_{a\in\Z,\ (a,q)=1}\widehat{G_n}\Big(\frac aq\Big).
\]
Substitute $y=n-u$. For $n\asymp X$ and $|u|\le HX^\epsilon$ we have $\frac T{2\pi}\log\frac n{n-u}=\frac{Tu}{2\pi n}+O(X^{2\epsilon}/T)$, and Fourier inversion of $\widehat\Phi$ on the scale $u\asymp n/T\asymp H$ gives
\begin{equation}\label{eq:packet}
 \begin{split}
 T\,\widehat{G_n}\Big(\frac aq\Big)={}&2\pi\,e\Big(-\frac{an}q\Big)\overline{\beta'(q;\tfrac{\log n}L)}\,V\Big(\frac nX\Big)\,\Phi\Big(\frac{2\pi an}{qT}\Big)\\
 &+O\Big(\frac{X^\epsilon}T\Big(1+\frac{|a|H}q\Big)^{-A}\Big),
 \end{split}
\end{equation}
the error coming from the Taylor expansions of $\beta'$, $V$ and $(ny)^{-1/2}$ in $u$, each of which changes by $O(H/X)=O(1/T)$ over the range $|u|\le HX^\epsilon$. In particular only $|a|\ll X^\epsilon q/H$ contributes, and since $\Phi$ is supported in $(1,2)$, the main term has $a>0$ and $a\asymp q/H$.

Now expand $M(n)=L^2\sum_{q'}q'^{-1}\sum_{b\bmod q'}^*e(bn/q')\beta(q';\log n/L)$. For each $(q',b)$ and $(q,a)$ the $n$-sum is $\sum_ne((b/q'-a/q)n)\Theta(n)$ with $\Theta$ smooth, supported on $n\asymp X$, with $\Theta^{(j)}\ll X^{-j+\epsilon}$. By Poisson summation it is $O(X^{-A})$ unless $\|b/q'-a/q\|\le X^{-1+\epsilon}$. On the block, $n,m\le2X$, so \eqref{eq:qrange2} gives $q,q'\le Q:=(2X)^{1/2}T^{-\delta}$. If $b/q'\ne a/q$ in $\Q/\Z$ then $\|b/q'-a/q\|\ge(qq')^{-1}\ge T^{2\delta}/(2X)$, which exceeds $X^{-1+\epsilon}$ once $\epsilon<\delta/2$, since $X\le T^2$. So only $q'=q$, $b\equiv a$ survive, and then the $n$-sum is $\int\Theta+O(X^{-A})$. The error term in \eqref{eq:packet}, summed over $a$, is $\ll X^\epsilon T^{-1}(1+q/H)\le2X^\epsilon T^{-1}q/H$, because $q\ge HT^{\delta}/2$ by \eqref{eq:qrange2}; with the weight $L^2|\beta'|/q$, summed over $q\le Q$ and over $n\asymp X$ with $|M(n)|\ll X^\epsilon$, it is $\ll X^{O(\epsilon)}QX/(HT)=X^{O(\epsilon)}Q\le T^{1-\delta+O(\epsilon)}$. Collecting,
\begin{multline*}
 \ip{P[M]}{P[M'_X]}_\Phi\\
 =2\pi L^4\sum_q\frac1{q^2}\int\beta(q;\tfrac{\log x}L)\overline{\beta'(q;\tfrac{\log x}L)}V\Big(\frac xX\Big)\sum_{(a,q)=1}\Phi\Big(\frac{2\pi ax}{qT}\Big)\dd x\\
 +O(T^{1-\delta+O(\epsilon)}),
\end{multline*}
where $M'_X$ is the block. By M\"obius inversion and Poisson summation, $\sum_{(a,q)=1}F(a)=\frac{\varphi(q)}q\int F+O(\tau(q)X^\epsilon)$ for $F(a)=\Phi(2\pi ax/(qT))$, whose scale is $qT/x\asymp q/H\ge T^{\delta}/4$ by \eqref{eq:qrange2}. The main term is $\frac{\varphi(q)}q\cdot\frac{qT}{2\pi x}\int\Phi\asymp q/H$, so the relative error is $O(\tau(q)X^{\epsilon}T^{-\delta})$. Summing over the blocks and writing $x=T^r$, $\dd x/x=L\dd r$, gives the claim: the error terms total $O(T^{1-\delta+O(\epsilon)}L^{C})$, using $\sum_q\varphi(q)q^{-2}\tau_3(q)^{2C}\tau(q)\ll L^{C'}$, and this is $O(T^{1-\delta/2}L^C)$ for $\epsilon$ small.

For the one-sided bound take $\beta'=\beta$ and $\Phi^+$ in place of $\Phi$; let $\operatorname{supp}\Phi^+\subseteq[c_1,c_2]\subset(1,2)$. The first part of the proof, up to and including \eqref{eq:packet}, did not use the lower cutoff. On a block the upper cutoff alone gives $q,q'\le Q=(2X)^{1/2}T^{-\delta}$, so distinct reduced fractions are separated exactly as above and only the diagonal $q'=q$, $b\equiv a$ survives. The error in \eqref{eq:packet}, summed over $a$, is $\ll X^\epsilon T^{-1}(1+q/H)$; with the weight $L^2|\beta|/q$, summed over $q\le Q$ and over $n\asymp X$ with $|M(n)|\ll X^\epsilon$, it is $\ll X^{O(\epsilon)}\big((X/T)\log Q+QX/(HT)\big)=X^{O(\epsilon)}(H+Q)$. Here $H=X/T\le T^{\max I-1+o(1)}$ and $Q\le T^{1-\delta+o(1)}$, so this is $O(T^{1-c})$ for $c=\frac12\min(2-\max I,\delta)$ and $\epsilon$ small. The diagonal is at most
\[
 2\pi L^4\|\Phi^+\|_\infty\|V\|_\infty\sum_{q\le Q}\frac{\sup_r|\beta(q;r)|^2}{q^2}\int_{X/2}^{2X}\#\mathcal A_q(x)\dd x,
\]
where $\mathcal A_q(x)$ is the set of integers $a\ge1$ with $2\pi ax/(qT)\in[c_1,c_2]$. If $\mathcal A_q(x)$ is non-empty, it contains some $a\ge1$ with $a\le c_2qT/(2\pi x)$, so $qT/x\ge2\pi/c_2>\pi$, and then $\#\mathcal A_q(x)\le1+(c_2-c_1)qT/(2\pi x)\le qT/x$. If $\mathcal A_q(x)$ is empty, the bound $\#\mathcal A_q(x)\le qT/x$ holds trivially. So in every case $\#\mathcal A_q(x)\le2q/H$ for $x\ge X/2$, and the integral is at most $3qX/H=3qT$. Hence the diagonal is $\ll TL^4\sum_qq^{-1}\sup_r|\beta(q;r)|^2$. Summing over the $O(L)$ blocks that meet the support gives the one-sided bound.
\end{proof}

\begin{proposition}[Gram limits]\label{prop:markedgram}
For real $\psi,\eta\in C_c^\infty(\mathcal Y)$ and $\Phi$ as above,
\[
 \ip{W_\psi}{W_\eta}_\Phi=a_3TL^9\Big(\int\Phi\Big)\ip{\psi^\flat}{\eta^\flat}+o(TL^9).
\]
\end{proposition}
\begin{proof}
By \cref{lem:ramanujan} with $\beta=\beta_\psi$, $\beta'=\beta_\eta$, it suffices to show that for each fixed $r$
\[
 L^{-4}\sum_q\frac{\varphi(q)}{q^2}\beta_\psi(q;r)\beta_\eta(q;r)\to a_3\iint\psi(r,k,x)\eta(r,k,x)\dd k\dd x,
\]
with a bound $O(1)$ uniform in $r$ (dominated convergence in $r$ then finishes the proof). The left side is that of \eqref{eq:fouredge2} with $\Psi(k,\varrho,\upsilon)=m_\psi(r,k,\varrho)m_\eta(r,k,\upsilon)$, so its limit is $a_3\int m_\psi(r,|t|,R_1)m_\eta(r,|t|,C_1)\dd t$. For fixed $k=|t|$ and fixed $R_1=\varrho$, $C_1=\upsilon$ in $(0,k)$, the remaining free entry $t_{11}$ ranges over an interval of length
\begin{align*}
 \min(\varrho,\upsilon)-\max(0,\varrho+\upsilon-k)&=\min(\varrho,\upsilon,k-\varrho,k-\upsilon)\\
 &=\int_0^{k/2}\one_{x<\varrho<k-x}\one_{x<\upsilon<k-x}\dd x,
\end{align*}
and $(k,\varrho,\upsilon,t_{11})\mapsto t$ has Jacobian $1$. Hence the limit is
\[
 a_3\int_0^\infty\!\!\int_0^{k/2}\Big(\int_x^{k-x}m_\psi(r,k,\varrho)\dd\varrho\Big)\Big(\int_x^{k-x}m_\eta(r,k,\upsilon)\dd\upsilon\Big)\dd x\dd k .
\] Since $m_\psi(r,k,\cdot)$ vanishes on $[k/2,\infty)$, $\int_x^{k-x}m_\psi(r,k,\varrho)\dd\varrho=\psi(r,k,x)$, and similarly for $\eta$. Integrating over $r$ and changing variables to $(x,y,z)$ gives $a_3\ip{\psi^\flat}{\eta^\flat}$.
\end{proof}

\subsection{Cubic realisation}
For $g\in C_c^\infty(\R)$ let $A_g(s)=\sum_ng(\log n/L)n^{-s}$, and for $c\in C_c^\infty(\R^3)$ let $P_c(s)=\sum_{n_1,n_2,n_3}c(\ell(n))(n_1n_2n_3)^{-s}$, $\ell(n)=(\log n_j/L)_j$.

\begin{lemma}[Sixth-moment bounds]\label{lem:sixth}
Let $\chi\in C_c^\infty(\R)$ and, for $\xi\in\R$, let $g_\xi(a)=\chi(a)e^{i\xi a}$.
\begin{enumerate}[label=(\roman*)]
\item Uniformly for $|\xi|\le LT/4$, $\int_{T}^{2T}|A_{g_\xi}(\frac12+it)|^6\dd t\ll_\chi TL^9$.
\item Let $K\subset\R^3$ be compact. There is $C_K$ such that for every $c\in C_c^\infty(\R^3)$ supported in $K$,
\[
 \|P_c\|_\Phi\ll_K(TL^9)^{1/2}\int_{\R^3}|\widehat c(\xi)|\dd\xi+T^{C_K}\int_{\max_j|\xi_j|>LT/4}|\widehat c(\xi)|\dd\xi ,
\]
where $\widehat c(\xi)=\int c(a)e^{-i\xi\cdot a}\dd a$.
\end{enumerate}
\end{lemma}
\begin{proof}
(i) Put $g=g_\xi$ and $s=\frac12+it$ with $T\le t\le2T$. By Mellin inversion $A_g(s)=\frac1{2\pi i}\int_{(2)}\zeta(s+w)\widetilde g_L(w)\dd w$, where
\[
 \widetilde g_L(w)=\int_0^\infty g\Big(\frac{\log x}L\Big)x^{w-1}\dd x=L\int_\R\chi(a)e^{La\Re w}e^{ia(\xi+L\Im w)}\dd a .
\]
For fixed $\Re w$ this is rapidly decreasing in $\Im w$, so we may move the line of integration to $\Re w=0$. We cross the pole at $w=1-s$, with residue $\widetilde g_L(\frac12-it)$. Here the frequency is $|\xi-Lt|\ge\frac34LT$, while the $j$-th derivative of $\chi(a)e^{La/2}$ is $\ll_{\chi,j}L^jT^{C_\chi}$; integrating by parts $j$ times shows that the residue is $O_{\chi,A}(T^{-A})$, uniformly in $|\xi|\le LT/4$. On $\Re w=0$ we have $|\widetilde g_L(iv)|=L|\widehat\chi(-\xi-Lv)|$, so by Minkowski's inequality
\begin{multline*}
 \Big(\int_T^{2T}|A_g(\tfrac12+it)|^6\dd t\Big)^{1/6}\\
 \le\frac1{2\pi}\int_\R L|\widehat\chi(-\xi-Lv)|\Big(\int_T^{2T}|\zeta(\tfrac12+i(t+v))|^6\dd t\Big)^{1/6}\dd v+O(T^{-A}).
\end{multline*}
For $|v|\le T/2$ the inner integral is at most $\int_{T/2}^{3T}|\zeta(\frac12+i\tau)|^6\dd\tau\ll TL^9$, by Harper's theorem \cite{Harper}, which is unconditional by \cite{Alpoge}; and $\int_\R L|\widehat\chi(-\xi-Lv)|\dd v=\int|\widehat\chi|$. For $|v|>T/2$ we have $|\xi+Lv|\ge L|v|-LT/4\ge L|v|/2$, so $|\widehat\chi(-\xi-Lv)|\ll_A(L|v|)^{-A}$, while the convexity bound $\zeta(\frac12+i\tau)\ll(2+|\tau|)^{1/4}$ bounds the inner integral by $\ll T^{1/6}(T+|v|)^{1/4}$; this range contributes $O(T^{-A})$.

(ii) Choose $\chi\in C_c^\infty(\R)$ equal to $1$ on the three coordinate projections of $K$; it depends only on $K$. By Fourier inversion $c(a)=c(a)\prod_j\chi(a_j)=(2\pi)^{-3}\int\widehat c(\xi)\prod_j\chi(a_j)e^{i\xi_ja_j}\dd\xi$, so for each $s$ (the sums being finite)
\[
 P_c(s)=(2\pi)^{-3}\int_{\R^3}\widehat c(\xi)\prod_{j=1}^3A_{g_{\xi_j}}(s)\dd\xi .
\]
By Minkowski's inequality $\|P_c\|_\Phi\le(2\pi)^{-3}\int|\widehat c(\xi)|\,\|\prod_jA_{g_{\xi_j}}\|_\Phi\dd\xi$. If $\max_j|\xi_j|\le LT/4$, H\"older's inequality and (i) give $\|\prod_jA_{g_{\xi_j}}\|_\Phi^2\le\|\Phi\|_\infty\prod_j(\int_T^{2T}|A_{g_{\xi_j}}|^6)^{1/3}\ll TL^9$. Otherwise we use the trivial bound $|A_{g_{\xi_j}}(\frac12+it)|\le\sum_n|\chi(\log n/L)|n^{-1/2}\ll T^{C_K}$.
\end{proof}

\begin{proposition}[Realisation]\label[proposition]{prop:realisation}
$\|W_\psi-P_{c_\psi}\|_\Phi=o\big((TL^9)^{1/2}\big)$.
\end{proposition}
\begin{proof}
Let $f=f_\psi$ and define
\[
 M^*(n)=L^2\sum_q\frac{c_q(n)}q\sum_{lde=q}\frac{\mu(l)}lf\Big(r,\frac{\log ld}L,\frac{\log le}L\Big),\qquad r=\frac{\log n}L .
\]
Since $f(r,\frac{\log ld}L,\frac{\log le}L)=m_\psi(r,\frac{\log lq}L,\frac{\log ld}L)$ (as $l^2de=lq$), the coefficient of $L^2c_q(n)/q$ in $M^*(n)$ is
\[
 \beta^*(q;r)=\sum_{lde=q}\frac{\mu(l)}l\,m_\psi\Big(r,\frac{\log lq}L,\frac{\log ld}L\Big):
\]
the mark is evaluated at $lq,ld$ in place of $q,d$. By compactness of $\operatorname{supp}\psi\subset\mathcal Y$ there is $\delta>0$ such that $r-1+\delta\le k\le r/2-\delta$ on the support of $m_\psi(r,k,u)$. So $\beta_\psi$ satisfies \eqref{eq:qrange2}, while each term of $\beta^*$ has $T^{r-1+\delta}\le lq\le T^{r/2-\delta}$. The upper cutoff $q\le lq\le T^{r/2-\delta}$ therefore holds for $\beta^*$, but the lower cutoff can fail. For example, for a prime $p$ let $T=p^{9/2}$, $n=p^7$, $q=p^2$ and $(l,d,e)=(p,1,p)$. Then $r=\frac{14}9$, the shifted mark is evaluated at $(r,\frac{\log lq}L,\frac{\log ld}L)=(\frac{14}9,\frac23,\frac29)$, which lies in $\mathcal Y$, and $\log q/L=\frac49<r-1=\frac59$. Such terms have large $l$ (here $l=T^{2/9}$), and are covered by the one-sided bound below.

Put $D(q)=\sum_{lde=q}l^{-1/2}\le\tau_3(q)$. Fix $0<\eta_0\le\delta/2$, split $\beta_\psi=\beta^{\le}+\beta^{>}$ and $\beta^*=\beta^{*\le}+\beta^{*>}$ according as $l\le T^{\eta_0}$ or $l>T^{\eta_0}$, and put $\Delta^{\le}=\beta^{\le}-\beta^{*\le}$ and $\Delta^{>}=\beta^{>}-\beta^{*>}$. Then $W_\psi-P[M^*]=P[M_{\Delta^\le}]+P[M_{\Delta^>}]$, where $M_\Delta(n)=L^2\sum_qq^{-1}c_q(n)\Delta(q;\log n/L)$. Let $\Phi^+\ge|\Phi|$ be a fixed smooth majorant as in \cref{lem:ramanujan}.

\emph{Small $l$.} For $l\le T^{\eta_0}$ we have $\log q/L\ge\log lq/L-\eta_0$, so every term of $\beta^{*\le}$ has $T^{r-1+\delta/2}\le q\le T^{r/2-\delta}$: both cutoffs of \eqref{eq:qrange2} hold for $\Delta^\le$, with $\delta/2$ in place of $\delta$. Moreover both mark arguments move by $\log l/L\le\eta_0$, so by the mean value theorem $|\partial_r^j\Delta^\le(q;r)|\ll_{\psi,j}\eta_0D(q)$. The two-sided norm bound of \cref{lem:ramanujan} gives
\[
 \|P[M_{\Delta^\le}]\|_\Phi^2\ll\eta_0^2\,TL^5\sum_{q\le T}\frac{D(q)^2}q+O(T^{1-\delta/4}L^C).
\]
(The lower cutoff is not needed for this estimate, since the one-sided bound would also apply to $\Delta^\le$; it gives the better error term.)

\emph{Large $l$.} For $l>T^{\eta_0}$ we use $l^{-1}\le T^{-\eta_0/2}l^{-1/2}$, which bounds the M\"obius tails uniformly: $|\partial_r^j\Delta^>(q;r)|\ll_{\psi,j}T^{-\eta_0/2}D(q)$ for all $q$ and $r$. As $\Delta^>$ satisfies the upper cutoff $q\le T^{r/2-\delta}$, the one-sided bound of \cref{lem:ramanujan} gives
\[
 \|P[M_{\Delta^>}]\|_\Phi^2\ll T^{1-\eta_0}L^5\sum_{q\le T}\frac{D(q)^2}q+O(T^{1-c}).
\]

We have $\sum_{q\le T}D(q)^2/q\ll L^4$. Indeed, $D$ is multiplicative with $D(p)=2+p^{-1/2}$ and $D(p^j)\le(1-2^{-1/2})^{-1}(j+1)$ for all $p,j$. At $s=1+1/L$ the Euler product of $\sum_qD(q)^2q^{-s}$ is $\ll\zeta(s)^4$, and $q^{1/L}\le e$ for $q\le T$. The implied constants above depend only on $\psi$ and $\Phi^+$, not on $\eta_0$. Hence
\[
 \limsup_{T\to\infty}\frac{\|W_\psi-P[M^*]\|_\Phi^2}{TL^9}\ll\eta_0^2 .
\]
Next we invert $M^*$ exactly. Write $f(r,u,v)=\iint\widehat f_r(A,B)e^{-Au-Bv}\dd A\dd B$ on the lines $\Re A=\Re B=1$, with $\widehat f_r$ rapidly decreasing on them, uniformly for $r\in I_\psi$ with all $r$-derivatives. Applying \eqref{eq:ramanujaninv} with $(\alpha,\beta)=(A/L,B/L)$ and $1/\zeta(1+\alpha)=\sum_i\mu(i)i^{-1-\alpha}$ gives
\[
 M^*(n)=\sum_{def_0=n}F_L\Big(r,\frac{\log d}L,\frac{\log e}L\Big),
\]
where
\[
 F_L(r,a,b)=\iint\widehat f_r(A,B)\frac{L}{\zeta(1+A/L)}\frac{L}{\zeta(1+B/L)}e^{-Aa-Bb}\dd A\dd B,
\]
all interchanges being absolutely convergent for fixed $n$. On $\Re A=1$ one has $L/\zeta(1+A/L)\to A$ and $|L/\zeta(1+A/L)|\ll(1+|A|)^2$ uniformly in $L$ (use the Laurent expansion of $\zeta$ for $|A|\le L$, and $1/\zeta(1+\sigma+i\tau)\ll\log(2+|\tau|)$ for $\sigma\ge0$ otherwise). By dominated convergence $F_L\to\partial_a\partial_bf$ in every $C^j$ norm on compact subsets of $\R^3$. Hence $P[M^*]=P_{c_L}$ with $c_L(a)=F_L(a_1+a_2+a_3,a_1,a_2)$, and $c_L\to c_\psi$ in $C^j$ on compacts. Only the values of $c_L$ and $c_\psi$ on $[0,\infty)^3$ enter $P_{c_L}$ and $P_{c_\psi}$, and there both vanish unless $a_1+a_2+a_3\in I_\psi$, since $f(r,\cdot,\cdot)=0$ for $r\notin I_\psi$. Let $\omega\in C_c^\infty(\R^3)$ be equal to $1$ on the compact set $\{a\in[0,\infty)^3:\ a_1+a_2+a_3\in I_\psi\}$, and let $K$ be its support. Then $P_{c_L}-P_{c_\psi}=P_{h_L}$ with $h_L=\omega(c_L-c_\psi)$, and $h_L\to0$ in every $C^j$ norm with support in $K$. Integrating by parts, $|\widehat h_L(\xi)|\ll_j\|h_L\|_{C^j}(1+|\xi|)^{-j}$; with $j=4$ this gives $\int|\widehat h_L|\to0$, and with $j=C_K+4$ it gives $\int_{\max|\xi_j|>LT/4}|\widehat h_L|=O(T^{-C_K-1})$, because the norms $\|h_L\|_{C^j}$ are bounded. Hence \cref{lem:sixth}(ii) gives $\|P[M^*]-P_{c_\psi}\|_\Phi=o((TL^9)^{1/2})$. Letting $\eta_0\to0$ proves the proposition.
\end{proof}

\begin{remark}
In this section Harper's bound $\int_{T/2}^{3T}|\zeta(\frac12+it)|^6\dd t\ll TL^9$ \cite{Harper} enters through \cref{lem:sixth}, which is used in \cref{prop:realisation} and, below, in \cref{lem:reflfactor,lem:reflcube,lem:factors}, and through \cref{lem:sourcenorm}, which is used in \cref{prop:mixed}. Harper's bound rests on the Riemann hypothesis, which is a theorem \cite{Alpoge}.
\end{remark}

\subsection{Sources and the mixed pairing}
For $\boldsymbol u\in\C^3$ and $\phi\in C_c^\infty((1,2))$ the \emph{elementary source} is
\[
 P_{\boldsymbol u,\phi}(s)=\sum_{n\ge1}b_{\boldsymbol u}(n)\,\phi\Big(\frac{\log n}L\Big)n^{-s},\qquad b_{\boldsymbol u}(n)=\sum_{n_1n_2n_3=n}\prod_{j=1}^3n_j^{-u_j/L},
\]
with cubic weight $c_{\boldsymbol u,\phi}(a)=\phi(a_1+a_2+a_3)e^{-\boldsymbol u\cdot a}$. Thus $P_{\boldsymbol 0,\phi}$ is a smoothed $d_3$-polynomial. Let $c\in C_c^\infty(\R^3)$ satisfy $\phi(a_1+a_2+a_3)=1$ for every $a\in[0,\infty)^3\cap\operatorname{supp}c$. Only the values of $c$ on $[0,\infty)^3$ enter $P_c$, and there $c(a)=\phi(a_1+a_2+a_3)c(a)=(2\pi)^{-3}\int\widehat c(\xi)c_{-i\xi,\phi}(a)\dd\xi$, so $P_c$ is a superposition of sources:
\begin{equation}\label{eq:sourcesuperposition}
 P_c=(2\pi)^{-3}\int_{\R^3}\widehat c(\xi)\,P_{-i\xi,\phi}\dd\xi .
\end{equation}

\begin{lemma}[Source norms]\label{lem:sourcenorm}
Uniformly for $\boldsymbol u\in(i\R)^3$ and $\Phi$ in bounded sets, $\|P_{\boldsymbol u,\phi}\|_\Phi\ll(1+|\boldsymbol u|)^{3}(TL^9)^{1/2}$.
\end{lemma}
\begin{proof}
Write $s=\frac12+it$ with $T\le t\le2T$, and recall that $\phi$ is supported in $(1,2)$, so that only $n\le T^2$ occur. Since $|b_{\boldsymbol u}(n)|\le\tau_3(n)$ for $\boldsymbol u\in(i\R)^3$, we have $|P_{\boldsymbol u,\phi}(s)|\le\sum_{n\le T^2}\tau_3(n)n^{-1/2}\ll TL^2$, hence $\|P_{\boldsymbol u,\phi}\|_\Phi\ll T^{3/2}L^2$. This proves the lemma when $|\boldsymbol u|\ge T^{1/2}$, so assume $|\boldsymbol u|<T^{1/2}$.

By Mellin inversion $P_{\boldsymbol u,\phi}(s)=\frac1{2\pi i}\int_{(2)}Z(w)\widetilde\phi_L(w)\dd w$, where $Z(w)=\prod_j\zeta(s+w+u_j/L)$ and $\widetilde\phi_L(w)=L\int\phi(a)e^{Law}\dd a$. Move the line of integration to $\Re w=0$. The poles crossed are at $w_j=\frac12-it-u_j/L$; they may nearly coincide, so instead of computing residues we bound their total contribution, which is $\frac1{2\pi i}\oint_{\partial R}Z\widetilde\phi_L$, where $R$ is the rectangle
\[
 \big|\Re w-\tfrac12\big|\le\tfrac1L,\qquad\big|\Im w+t\big|\le\tfrac{|\boldsymbol u|+1}L .
\]
It contains every $w_j$, and every point of $\partial R$ has distance at least $1/L$ from every $w_j$. For $w\in\partial R$ put $z_j=w-w_j$, so that $s+w+u_j/L=1+z_j$ with $\Re z_j\ge-1/L$, $|z_j|\ge1/L$ and $|\Im z_j|\le2(|\boldsymbol u|+1)/L\le T^{1/2}$. If $|\Im z_j|\le1$ then $\zeta(1+z_j)=z_j^{-1}+O(1)\ll L$; otherwise the approximation $\zeta(w)=\sum_{n\le x}n^{-w}-x^{1-w}/(1-w)+O(x^{-\Re w})$ with $x=|\Im w|$ \cite[Theorem~4.11]{Titchmarsh} gives $\zeta(1+z_j)\ll\log(2+|\Im z_j|)\ll L$, because $n^{-\Re z_j}\le n^{1/L}\le e$ for $n\le T^{1/2}$. Thus $Z\ll L^3$ on $\partial R$. On $\partial R$ we also have $|\Re w|\le1$ and $|\Im w|\ge t-T^{1/2}\ge T/2$, so repeated integration by parts in $\widetilde\phi_L(w)=L\int\phi(a)e^{La\Re w}e^{iLa\Im w}\dd a$ gives $\widetilde\phi_L(w)\ll_AT^{-A}$. As $\partial R$ has length $O(T^{1/2})$, the poles contribute $O(T^{-A})$, uniformly in $t$ and $\boldsymbol u$.

On $\Re w=0$ the three zeta factors are at heights $t+v+\Im u_j/L$ on the critical line, and $|\widetilde\phi_L(iv)|=L|\widehat\phi(-Lv)|$. By Minkowski's inequality
\[
 \|P_{\boldsymbol u,\phi}\|_\Phi\le\frac1{2\pi}\int_\R L|\widehat\phi(-Lv)|\,\Big\|\prod_j\zeta\big(\tfrac12+i(\cdot+v+\Im u_j/L)\big)\Big\|_\Phi\dd v+O(T^{-A}).
\]
For $|v|\le T/4$ all heights lie in $[T/2,3T]$, and H\"older's inequality with Harper's bound \cite{Harper} (unconditional by \cite{Alpoge}) bounds the norm by $\ll(TL^9)^{1/2}$; since $\int L|\widehat\phi(-Lv)|\dd v=\int|\widehat\phi|$, this range contributes $\ll(TL^9)^{1/2}$. For $|v|>T/4$ the convexity bound gives a norm $\ll T^{1/2}(T+|v|)^{3/4}$, and $L|\widehat\phi(-Lv)|\ll_A(L|v|)^{-A}$, so this range contributes $O(T^{-A})$. Hence $\|P_{\boldsymbol u,\phi}\|_\Phi\ll(TL^9)^{1/2}$ for $|\boldsymbol u|<T^{1/2}$.
\end{proof}

\begin{lemma}[Polar part]\label{lem:polar}
Let $(a,q)=1$, let $W\in C_c^\infty((0,\infty))$, and let $\boldsymbol u\in\C^3$ have distinct coordinates with $|\Re u_j|\le1$. Then
\[
 \sum_{n\ge1}b_{\boldsymbol u}(n)e\Big(\frac{an}q\Big)W(n)=\mathcal P_{\boldsymbol u}(q;W)+\mathcal E_{\boldsymbol u}(a,q;W),
\]
where, for $\{i,j,k\}=\{1,2,3\}$, $\alpha_i=(u_j-u_i)/L$ and $\beta_i=(u_k-u_i)/L$,
\begin{equation}\label{eq:polarterm}
 \mathcal P_{\boldsymbol u}(q;W)=\sum_{i=1}^3\frac{\mathcal D_q(\alpha_i,\beta_i)}q\,\zeta(1+\alpha_i)\zeta(1+\beta_i)\int_0^\infty W(x)x^{-u_i/L}\dd x,
\end{equation}
and $\mathcal E_{\boldsymbol u}(a,q;W)$ is given by the dual expansion \eqref{eq:dualexpansion} below. The polar term $\mathcal P_{\boldsymbol u}$ does not depend on $a$, and extends holomorphically to all $\boldsymbol u$ (as a sum of residues).
\end{lemma}
\begin{proof}
Splitting $n_j$ into residue classes $x_j\bmod q$ and using Mellin inversion of $W$,
\begin{equation}\label{eq:hurwitz}
 \begin{split}
 &\sum_nb_{\boldsymbol u}(n)e\Big(\frac{an}q\Big)W(n)\\
 &=\frac1{2\pi i}\int_{(2)}\widetilde W(w)\,q^{-3w-\sum_ju_j/L}\sum_{x\in(\Z/q)^3}e\Big(\frac{ax_1x_2x_3}q\Big)\prod_{j=1}^3\zeta\Big(w+\frac{u_j}L,\frac{x_j}q\Big)\dd w,
 \end{split}
\end{equation}
where $\zeta(w,\theta)=\sum_{n\ge0}(n+\theta)^{-w}$ is the Hurwitz zeta function, $x_j\in\{1,\dots,q\}$, and $\widetilde W(w)=\int W(x)x^{w-1}\dd x$. The only poles of the integrand are at $w=1-u_i/L$, and they are simple because the $u_i$ are distinct. At such a pole the residue of $\zeta(w+u_i/L,x_i/q)$ is $1$, and $\sum_{x_i\bmod q}e(ax_1x_2x_3/q)=q\one_{q\mid x_jx_k}$ because $(a,q)=1$. Hence, if $\Re\alpha_i>0$ and $\Re\beta_i>0$, so that the Hurwitz series converge absolutely and may be regrouped as sums over $d\equiv x_j$ and $e\equiv x_k\pmod q$, the residue is
\begin{multline*}
 \widetilde W(1-\tfrac{u_i}L)\,q^{-2-\alpha_i-\beta_i}\!\!\sum_{\substack{x_j,x_k\bmod q\\ q\mid x_jx_k}}\!\!\zeta(1+\alpha_i,\tfrac{x_j}q)\zeta(1+\beta_i,\tfrac{x_k}q)\\
 =\widetilde W(1-\tfrac{u_i}L)\sum_{\substack{d,e\ge1\\ q\mid de}}d^{-1-\alpha_i}e^{-1-\beta_i},
\end{multline*}
which is the $i$th term of \eqref{eq:polarterm} by \cref{lem:localfactor}. Both the first line of this display (a finite sum of products of Hurwitz zeta values) and the $i$th term of \eqref{eq:polarterm} (in which $\mathcal D_q$ is a finite sum) are meromorphic in $(\alpha_i,\beta_i)$, so the identity extends by analytic continuation to all $\boldsymbol u$ with distinct coordinates. Moving the contour in \eqref{eq:hurwitz} to $\Re w=-\delta$ leaves $\mathcal E_{\boldsymbol u}$, which is the integral over that line. Finally, for $\boldsymbol u$ in a compact set and $L$ large, $\mathcal P_{\boldsymbol u}(q;W)$ is $\frac1{2\pi i}$ times the integral of the integrand of \eqref{eq:hurwitz} over a fixed small circle about $w=1$ enclosing the three poles, which is holomorphic in $\boldsymbol u$; this treats collisions of the coordinates.
\end{proof}

\begin{lemma}[Exponential profiles]\label{lem:expprofile}
For $\boldsymbol u\in\C^3$ with distinct coordinates and $(x,y,z)\in\Om$,
\[
 h_{e^{-\boldsymbol u\cdot a}}(x,y,z)=\frac{\det\big(e^{-u_j\lambda_i}\big)_{i,j=1}^3}{(u_1-u_2)(u_1-u_3)(u_2-u_3)},\qquad(\lambda_1,\lambda_2,\lambda_3)=(x,y,z).
\]
Moreover, with $k=x+y$, $r=x+y+z$ and $\alpha_i,\beta_i$ as in \cref{lem:polar} but without the factor $1/L$,
\begin{equation}\label{eq:residuesum}
 h_{e^{-\boldsymbol u\cdot a}}(x,y,z)=\sum_{i=1}^3\frac{e^{-ru_i}}{(u_j-u_i)(u_k-u_i)}\int_x^{k-x}e^{-(u_j-u_i)\varrho-(u_k-u_i)(k-\varrho)}\dd\varrho .
\end{equation}
\end{lemma}
\begin{proof}
As in the proof of \cref{lem:markedprofile},
\[
 h_{e^{-\boldsymbol u\cdot a}}(x,y,z)=\int_y^z\!\!\int_x^y\!\!\int_w^{u'}e^{-u_1a-u_2(u'+w-a)-u_3(r-u'-w)}\dd a\dd w\dd u' .
\]
The inner integral is $(e^{-u_1w-u_2u'}-e^{-u_1u'-u_2w})e^{-u_3(r-u'-w)}/(u_1-u_2)$. The remaining integrand factorises in $(u',w)$; integrating, the two terms in which both variables are evaluated at $y$ cancel, and since $r=x+y+z$ each of the six remaining terms is $\pm e^{-\sum_iu_{\sigma(i)}\lambda_i}/((u_1-u_2)(u_1-u_3)(u_2-u_3))$ for a permutation $\sigma$, with the sign of $\sigma$. This is the determinant formula. For \eqref{eq:residuesum}, the $i$th term on the right equals $\big(e^{-u_iz-u_jx-u_ky}-e^{-u_iz-u_jy-u_kx}\big)/\big((u_j-u_i)(u_k-u_i)(u_j-u_k)\big)$, and the sum over $i$ is the expansion of the determinant along the row $\lambda_3=z$. (Both identities have also been checked numerically.)
\end{proof}

\begin{lemma}[Arcs]\label{lem:arcs}
Let $\psi\in C_c^\infty(\mathcal Y)$ be real, and let $b(n)\ll n^\epsilon$ be supported on $n\le T^{10}$. Let $\sum_jV(n/2^j)=1$ be a smooth dyadic partition with $V\in C_c^\infty([\frac12,2])$. Then
\[
 \begin{split}
 \ip{P[b]}{W_\psi}_\Phi={}&2\pi L^2\sum_{X}\sum_{q\ge1}\frac1q\sum_{\substack{a\ge1\\(a,q)=1}}\sum_nb(n)e\Big(-\frac{an}q\Big)\beta_\psi\Big(q;\frac{\log n}L\Big)\\
 &\times V\Big(\frac nX\Big)\Phi\Big(\frac{2\pi an}{qT}\Big)+O(T^{1-\delta}),
 \end{split}
\]
where $X$ runs over powers of $2$ and $\delta=\delta_\psi>0$.
\end{lemma}
\begin{proof}
This is the first half of the proof of \cref{lem:ramanujan}, with $M'=M_\psi$ and $b$ in place of $M$: Poisson summation in the witness variable $m$ gives \eqref{eq:packet}, with $\beta'=\beta_\psi$ (which is real). The error term in \eqref{eq:packet}, summed over $n\asymp X$, $q\le X^{1/2-\delta}$ and $a\ll X^\epsilon q/H$ with weight $|b(n)||\beta_\psi|/q$, is $O(X^{1+\epsilon}T^{-1}(L+X^{1/2-\delta}/H))=O(X^{\epsilon}(HL+X^{1/2}))$, and $H\le T^{1-2\delta}$ and $X^{1/2}\le T^{1-\delta}$ on the support since $r<2-2\delta$; after shrinking $\delta$ this is $O(T^{1-\delta})$. Pairs with $|n-m|>HX^\epsilon$ contribute $O(T^{-A})$.
\end{proof}

By \cref{lem:polar}, for $b=b_{\boldsymbol u}\phi(\log\cdot/L)$ the inner $n$-sum in \cref{lem:arcs} splits as $\mathcal P_{\boldsymbol u}(q;\mathcal W)+\mathcal E_{\boldsymbol u}(-a,q;\mathcal W)$, with
\begin{equation}\label{eq:arcweight}
 \mathcal W(y)=\mathcal W_{X,q,a}(y)=\phi\Big(\frac{\log y}L\Big)\beta_\psi\Big(q;\frac{\log y}L\Big)V\Big(\frac yX\Big)\Phi\Big(\frac{2\pi ay}{qT}\Big).
\end{equation}
The \emph{polar main term} is the contribution of $\mathcal P_{\boldsymbol u}$, and the \emph{signed error} the contribution of $\mathcal E_{\boldsymbol u}$.

\begin{lemma}[Polar limit]\label{lem:polarlimit}
Let $\boldsymbol u\in(i\R)^3$ have distinct coordinates and let $\phi=1$ on a neighbourhood of $I_\psi$. The polar main term equals $a_3TL^9(\int\Phi)\int_\Om h_{e^{-\boldsymbol u\cdot a}}\psi^\flat+o(TL^9)$.
\end{lemma}
\begin{proof}
Since $\mathcal P_{\boldsymbol u}(q;\mathcal W)$ is linear in $\mathcal W$, we first sum over $a$. For fixed $y\asymp X$ the function $a\mapsto\Phi(2\pi ay/(qT))$ varies on the scale $qT/y\asymp q/H\ge X^\delta$ by \eqref{eq:qrange}, so by M\"obius inversion and Poisson summation $\sum_{(a,q)=1}\Phi(2\pi ay/(qT))=\frac{\varphi(q)}q\frac{qT}{2\pi y}\int\Phi+O(\tau(q)X^\epsilon)$. The error contributes $O(T^{1-\delta/2})$ after summing over $q$ and $X$, and $\sum_XV(y/X)=1$. Using \eqref{eq:polarterm}, the substitution $y=T^r$ and $\phi=1$ on the support of $\beta_\psi$, the polar main term is therefore
\begin{multline*}
 TL^3\Big(\int\Phi\Big)\sum_{i=1}^3\zeta(1+\alpha_i)\zeta(1+\beta_i)\int\!e^{-ru_i}\sum_q\frac{\varphi(q)}{q^2}\mathcal D_q(\alpha_i,\beta_i)\beta_\psi(q;r)\dd r\\
 +O(T^{1-\delta/2}L^C).
\end{multline*}
Here $L^{-2}\zeta(1+\alpha_i)\zeta(1+\beta_i)\to((u_j-u_i)(u_k-u_i))^{-1}$, and by \eqref{eq:fouredge} with $\Psi(k,\upsilon)=m_\psi(r,k,\upsilon)$ and $(a_1,a_2)=(u_j-u_i,u_k-u_i)$, the $q$-sum divided by $L^4$ tends to $a_3\int e^{-(u_j-u_i)R_1-(u_k-u_i)R_2}m_\psi(r,|t|,C_1)\dd t$, boundedly in $r$. As in the proof of \cref{prop:markedgram}, fixing $k=|t|$, $C_1=\upsilon$ and $R_1=\varrho$ and integrating out the free entry,
\[
 \begin{split}
 &\int e^{-(u_j-u_i)R_1-(u_k-u_i)R_2}m_\psi(r,|t|,C_1)\dd t\\
 &\qquad=\int_0^\infty\!\!\int_0^{k/2}\psi(r,k,x)\int_x^{k-x}e^{-(u_j-u_i)\varrho-(u_k-u_i)(k-\varrho)}\dd\varrho\dd x\dd k .
 \end{split}
\]
Summing over $i$ with the factors $e^{-ru_i}/((u_j-u_i)(u_k-u_i))$ and using \eqref{eq:residuesum} at $(x,y,z)=(x,k-x,r-k)$ gives $a_3TL^9(\int\Phi)\iiint\psi(r,k,x)h_{e^{-\boldsymbol u\cdot a}}(x,k-x,r-k)\dd x\dd k\dd r+o(TL^9)$, which is the claim.
\end{proof}

\begin{proposition}[Mixed pairing]\label{prop:mixed}
Let $\psi\in C_c^\infty(\mathcal Y)$ be real and let $\phi=1$ on a neighbourhood of $I_\psi$. For every $\boldsymbol u\in(i\R)^3$ with distinct coordinates,
\[
 \ip{P_{\boldsymbol u,\phi}}{W_\psi}_\Phi=a_3TL^9\Big(\int\Phi\Big)\int_\Om h_{c_{\boldsymbol u,\phi}}\psi^\flat+o(TL^9),
\]
and $|\ip{P_{\boldsymbol u,\phi}}{W_\psi}_\Phi|\ll(1+|\boldsymbol u|)^3TL^9$ uniformly in $\boldsymbol u\in(i\R)^3$, all uniformly for $\Phi$ in bounded sets. Consequently, let $c\in C_c^\infty(\R^3)$ be such that $1<a_1+a_2+a_3<2$ for every $a\in[0,\infty)^3\cap\operatorname{supp}c$. Then $\ip{P_c}{W_\psi}_\Phi=a_3TL^9(\int\Phi)\ip{h_c}{\psi^\flat}+o(TL^9)$, where $h_c$ is the profile of the restriction of $c$ to $[0,1]^3$.
\end{proposition}
\begin{proof}
The uniform bound follows from Cauchy--Schwarz with respect to a smooth non-negative majorant of $|\Phi|$, \cref{lem:sourcenorm} and \cref{prop:markedgram}. For the asymptotic, fix $\eta_0>0$ and split $\beta_\psi=\beta^{\le}+\beta^{>}$ according as $l\le T^{\eta_0}$ or $l>T^{\eta_0}$ in \cref{def:marked}, and correspondingly $W_\psi=W^{\le}+W^{>}$. As a part of $\beta_\psi$, $\beta^>$ satisfies both cutoffs \eqref{eq:qrange2}, and $|\partial_r^j\beta^{>}(q;r)|\ll_jT^{-\eta_0/2}\sum_{lde=q}l^{-1/2}$. So the two-sided norm bound of \cref{lem:ramanujan} (the statement beginning ``In particular'', with a smooth majorant of $|\Phi|$) gives $\|W^{>}\|_{\Phi}^2\ll T^{1-\eta_0}L^{C}+O(T^{1-\delta/2}L^C)\ll T^{1-\min(\eta_0,\delta/2)}L^C$, and by Cauchy--Schwarz and \cref{lem:sourcenorm} the pairing with $W^>$ is $o(TL^9)$. For $W^{\le}$ apply \cref{lem:arcs} and \cref{lem:polar}. The polar main term for $W^{\le}$ differs from that for $W_\psi$ by $O(T^{1-\eta_0/2}L^{C})$, by the absolute-value version of \eqref{eq:fouredge}, so \cref{lem:polarlimit} evaluates it. It remains to show that the signed error
\[
 2\pi L^2\sum_X\sum_q\frac1q\sum_{(a,q)=1}\mathcal E_{\boldsymbol u}(-a,q;\mathcal W^{\le}_{X,q,a})
\]
is $O(T^{1-c})$ for some $c>0$, where $\mathcal W^{\le}$ is \eqref{eq:arcweight} with $\beta^{\le}$. This follows from \cref{thm:signed} below, applied on each of the $O(L^3)$ dyadic boxes $X$, $q\asymp Q$, $d\asymp D$ and for each $l\le T^{\eta_0}$, once $\eta_0$ is smaller than the saving of that theorem: on the support of $m_\psi(r,\log q/L,\log d/L)$ the exponents $h=1-1/r$, $s=\log Q/\log X$ and $t=\log D/\log X$ lie in a compact set $\mathcal K_\psi$ of the form required there. Indeed, with $(r,k,u)$ in the support of $m_\psi$ we have $h=1-\frac1r\in(0,\frac12)$, $s=k/r<\frac12$, $s>h$ because $k>r-1$, $s>\frac13$ because $k>\frac12>\frac r3$ when $r\le\frac32$ and $k>r-1\ge\frac r3$ when $r\ge\frac32$, and
\begin{equation}\label{eq:pcriterion}
 r\Big(\frac{1+h}2-s-t\Big)=r-\frac12-k-u=(r-k)-u-\frac12>0,
\end{equation}
because $(r,k,u)\in\operatorname{supp}m_\psi\subseteq\operatorname{supp}\psi\subset\mathcal Y$ and $r-k-x>\frac12$ on $\mathcal Y$. Finally, given $c$, choose $\phi\in C_c^\infty((1,2))$ equal to $1$ on a neighbourhood of $I_\psi$ and on the compact set $\{a_1+a_2+a_3:\ a\in[0,\infty)^3\cap\operatorname{supp}c\}$. By \eqref{eq:sourcesuperposition}, $\ip{P_c}{W_\psi}_\Phi=(2\pi)^{-3}\int\widehat c(\xi)\ip{P_{-i\xi,\phi}}{W_\psi}_\Phi\dd\xi$. The integrand is $\ll|\widehat c(\xi)|(1+|\xi|)^3TL^9$, which is integrable, and for almost every $\xi$ (those with distinct coordinates) it is $a_3TL^9(\int\Phi)\widehat c(\xi)\ip{h_{c_{-i\xi,\phi}}}{\psi^\flat}+o(TL^9)$. By dominated convergence and Fubini's theorem in \eqref{eq:profile}, the limit is $a_3(\int\Phi)\ip{h_c}{\psi^\flat}$, because $(2\pi)^{-3}\int\widehat c(\xi)c_{-i\xi,\phi}\dd\xi=c$ on $[0,1]^3$.
\end{proof}

\begin{remark}
The criterion \eqref{eq:pcriterion} is where the spread region comes from: the Kuznetsov estimate below saves a power exactly when $p=\frac{1+h}2-s-t>0$, and in the coordinates of the chamber this is $z-x>\frac12$ for the marked divisor $d\approx T^x$.
\end{remark}

\subsection{The signed estimate}
Put $\theta=\frac7{64}$, the Kim--Sarnak exponent \cite{KimSarnak}: every exceptional Laplace eigenvalue $\frac14+t^2$ of a Hecke--Maass cusp form on $\Gamma_0(M)$ has $t=i\tau$ with $0<\tau\le\theta$.

\begin{theorem}[Signed estimate]\label{thm:signed}
Let $\mathcal K$ be a compact subset of
\[
 \Big\{(h,s,t):\ 0<h<\tfrac12,\ \max(h,\tfrac13)<s<\tfrac12,\ t\ge0,\ p:=\tfrac{1+h}2-s-t>0\Big\}.
\]
There is $c=c(\mathcal K)>0$ with the following property. Let $X$ be large, $(h,s,t)\in\mathcal K$, $H=X^h$, $Q=X^s$, $A=Q/H$, $D=X^t$, let $L\asymp\log X$, and let $1\le l\le X^{c}$. Let $V\in C_c^\infty((0,\infty)^3)$, let $W(y;\alpha,\kappa)$ be smooth with compact support in $(0,\infty)^3$, and let $\boldsymbol u\in(i\R)^3$ be fixed. Then
\[
 \Sigma:=\sum_{\substack{q,a\ge1\\(a,q)=1}}\ \sum_{de=q/l}\frac1q\,V\Big(\frac aA,\frac qQ,\frac dD\Big)\,\mathcal E_{\boldsymbol u}\Big(\pm a,q;W\Big(\frac\cdot X;\frac aA,\frac qQ\Big)\Big)\ll(X/H)X^{-c},
\]
for either sign, where $\mathcal E_{\boldsymbol u}$ is defined in \cref{lem:polar}. The implied constant depends on $\mathcal K$, $\boldsymbol u$ and finitely many seminorms of $V,W$; $W$ may in addition depend on $L$, provided its seminorms, taken in $\log y$, are bounded uniformly.
\end{theorem}

The weights in \cref{prop:mixed} are of this form: on each dyadic box, $\mathcal W_{X,q,a}(Xy)$ in \eqref{eq:arcweight} is a smooth function of $(y,a/A,q/Q,d/D)$ times $\mu(l)/l$, because $X/(HT)=1$, and the dependence on $\log y/L$ through $\phi$ and $m_\psi$ has bounded derivatives in $\log y$. We write $N=Q^3/X$ and $K=Q(A/N)^{1/2}=X^{(1-h)/2}$, and note that $Q\le X^{1/2}$ gives $N\le Q X^{-2(1/2-s)}<Q$.

\subsubsection{Reduction to Kloosterman sums}
\begin{lemma}[Dual expansion]\label{lem:voronoi}
For $(a,q)=1$ and $\boldsymbol u\in(i\R)^3$,
\begin{equation}\label{eq:dualexpansion}
 \mathcal E_{\boldsymbol u}(a,q;W(\cdot/X))=\sum_{\boldsymbol\varepsilon\in\{\pm1\}^3}\ \sum_{\boldsymbol h\in\N^3}G(a,\boldsymbol\varepsilon\boldsymbol h;q)\prod_{j=1}^3\Big(\frac{h_j}q\Big)^{u_j/L}\frac X{q^3}\,k_{\boldsymbol\varepsilon}\Big(\frac{Xh_1h_2h_3}{q^3}\Big),
\end{equation}
with $G(a,\boldsymbol v;q)=\sum_{x\in(\Z/q)^3}e\big((ax_1x_2x_3+v_1x_1+v_2x_2+v_3x_3)/q\big)$ and
\begin{gather*}
 k_{\boldsymbol\varepsilon}(Y)=\frac1{2\pi i}\int_{(-\delta)}\widetilde W(w)\,Y^{w-1}\prod_{j=1}^3\gamma_{\varepsilon_j}\Big(w+\frac{u_j}L\Big)\dd w,\\
 \gamma_\pm(w)=\Gamma(1-w)(2\pi)^{w-1}e^{\mp\pi i(1-w)/2}.
\end{gather*}
For every $j\ge0$, $Y^jk^{(j)}_{\boldsymbol\varepsilon}(Y)\ll_{j}Y^{-\epsilon}$ for all $Y>0$, and $\ll_{j,B}Y^{-B}$ for $Y\ge1$; the same holds for derivatives in the parameters of $W$.
\end{lemma}
\begin{proof}
On $\Re w=-\delta$ every $w+u_j/L$ has negative real part, so Hurwitz's formula $\zeta(w,\theta)=\sum_\pm\gamma_\pm(w)\sum_{h\ge1}e(\pm h\theta)h^{w-1}$ \cite[(12.10)]{Apostol} converges absolutely; inserting it in the integral over $\Re w=-\delta$ that defines $\mathcal E_{\boldsymbol u}$ (\cref{lem:polar}) and using $\widetilde{W(\cdot/X)}(w)=X^w\widetilde W(w)$ gives \eqref{eq:dualexpansion}. By Stirling's formula $|\gamma_\pm(\sigma+i\tau)|\ll(1+|\tau|)^{1/2-\sigma}$ for $\sigma$ in compact sets away from the poles $w\in\{1,2,\dots\}$, while $\widetilde W$ decays faster than any power of $|\Im w|$. Moving the contour to $\Re w=1-\epsilon$ (no poles are crossed, the poles of $\gamma_\pm(w+u_j/L)$ being at $\Re w\ge1$) gives the bound $Y^{-\epsilon}$, and moving it to $\Re w=-B$ gives $Y^{-B-1}$. Derivatives bring down polynomial factors in $w$.
\end{proof}

Consequently the terms with $h_1h_2h_3>NX^{\nu}$ contribute $O(X^{-B})$ for any fixed $\nu>0$, and we may restrict to $h_1h_2h_3\le NX^\nu$.

\begin{lemma}[The complete sum]\label{lem:complete}
For $(a,q)=1$ and $\boldsymbol h\in\N^3$,
\[
 G(a,\boldsymbol\varepsilon\boldsymbol h;q)=q\sum_{f\mid q}\ \sum_{\substack{gj=f\\ g\mid(h_2,h_3),\ j\mid(h_1,h_2h_3/g^2)}}f\,S\Big(1,-\varepsilon_1\varepsilon_2\varepsilon_3\,\overline a\,\frac{h_1h_2h_3}{f^2};\frac qf\Big),
\]
where $S(m,n;c)$ is the classical Kloosterman sum and $\overline a$ is an inverse of $a$ modulo $q/f$.
\end{lemma}
\begin{proof}
Summing over $x_3$ gives $q\one[ax_1x_2\equiv-\varepsilon_3h_3\ (q)]$. Group $x_1$ by $g=(x_1,q)$, so $g\mid h_3$; write $x_1=gx'$, $q=gq'$ with $(x',q')=1$. The congruence becomes $x_2\equiv-\varepsilon_3\overline a\,\overline{x'}h_3/g\pmod{q'}$, and summing $e(\varepsilon_2h_2x_2/q)$ over the $g$ lifts of $x_2$ modulo $q$ gives $g\one_{g\mid h_2}e(-\varepsilon_2\varepsilon_3\overline a\,\overline{x'}h_2h_3/(g^2q'))$. Hence $G=q\sum_{g\mid(q,h_2,h_3)}g\,S(\varepsilon_1h_1,-\varepsilon_2\varepsilon_3\overline a h_2h_3/g^2;q/g)$. Selberg's identity $S(m,n;c)=\sum_{j\mid(m,n,c)}j\,S(1,mn/j^2;c/j)$ then gives the claim, with $f=gj$. For completeness we prove the identity. If $c=c_1c_2$ with $(c_1,c_2)=1$, then $S(m,n;c)=S(m\bar c_2,n\bar c_2;c_1)S(m\bar c_1,n\bar c_1;c_2)$, where $\bar c_2c_2\equiv1\ (c_1)$ and $\bar c_1c_1\equiv1\ (c_2)$; the right side factors in the same way, because $j\mid(m,n,c)$ splits uniquely as $j_1j_2$ with $j_i\mid(m,n,c_i)$, and $S(a,b;c')=S(1,ab;c')$ whenever $(a,c')=1$. So we may take $c=p^k$. If $p\nmid m$ (or $p\nmid n$) only $j=1$ occurs, and $S(m,n;p^k)=S(1,mn;p^k)$. Let $m=pm'$, $n=pn'$. For $k=1$ the left side is $p-1$, and the right side is $S(1,mn;p)+p=-1+p$. For $k\ge2$ the left side is $pS(m',n';p^{k-1})$. On the right, the term $j=1$ is $S(1,N;p^k)$ with $p\mid N$, which vanishes: replacing $x$ by $x(1+p^{k-1}t)$ multiplies each term by $e(t(x-N\bar x)/p)=e(tx/p)$, and $\sum_{t\bmod p}e(tx/p)=0$ as $p\nmid x$. The terms $j=pj'$ give $p\sum_{j'\mid(m',n',p^{k-1})}j'S(1,m'n'/j'^2;p^{k-1}/j')$, which is $pS(m',n';p^{k-1})$ by induction on $k$.
\end{proof}

Write $h_1=ju$, $h_2=gb$, $h_3=gc$ and $bc=jw$. Then $h_1h_2h_3=f^2n$ with $n=uw$, and for fixed $(g,j,n)$ the number of $(u,w,b,c)$ is $\alpha_j(n)=\sum_{uw=n}\tau(jw)\le\tau(j)d_3(n)$. Inserting \cref{lem:complete} in \eqref{eq:dualexpansion}, $\Sigma$ becomes, up to $O(X^{-B})$,
\begin{equation}\label{eq:Sigmakloosterman}
 \begin{split}
 \Sigma=\sum_{\boldsymbol\varepsilon}\sum_{f\le X}\ \sum_{gj=f}\ \sum_{n\le NX^\nu/f^2}&b_{g,j}(n)\sum_{\substack{a,q,\ (a,q)=1\\ f\mid q,\ de=q/l}}\frac{fX}{q^3}\,V(\cdots)\\
 &\times S\Big(1,\mp\varepsilon\,\overline an;\frac qf\Big)k_{\boldsymbol\varepsilon}\Big(\frac{Xf^2n}{q^3}\Big)q^{-\sum u_j/L},
 \end{split}
\end{equation}
where $\varepsilon=\varepsilon_1\varepsilon_2\varepsilon_3$ and $b_{g,j}(n)=\sum_{uw=n}\sum_{bc=jw}(ju)^{u_1/L}(gb)^{u_2/L}(gc)^{u_3/L}$, so that $|b_{g,j}(n)|\le\alpha_j(n)$. The factor $q^{-\sum u_j/L}$ is unimodular and smooth in $q/Q$. Since $X/q^3\asymp1/N$, the claim of \cref{thm:signed} is that the triple sum over $(a,q,d)$ in \eqref{eq:Sigmakloosterman}, summed over $n$ and $f$, saves $X^{c}$ against $AQ^2$, as $(X/H)/(1/N)=AQ^2$. Weil's bound gives $X^\epsilon AQ^{3/2}N$ for it, which suffices only when $s<\frac25$; for larger $s$ we need the spectral theory of $\Gamma_0(M)$.

\subsubsection{The Atkin--Lehner cusp}
Fix $\boldsymbol\varepsilon$, $f=gj$, a dyadic range $n\asymp B$ with $1\le B\le NX^\nu/f^2$, and $a,d$. Put $G=(ld,f)$ and $m=ld/G$. The conditions $f\mid q$ and $ld\mid q$ are equivalent to $\mathrm{lcm}(ld,f)=mf\mid q$, that is, $q=fc$ with $m\mid c$; then $e=q/(ld)$ is determined. The condition $(a,q)=1$ is $(a,f)=1$ together with $(a,c)=1$, and implies $(a,m)=1$. Put $M=am$.

\begin{lemma}\label{lem:ALcusp}
Let $M=am$ with $(a,m)=1$, choose $\beta,\delta'\in\Z$ with $a\delta'-m\beta=1$, and let
\[
 \sigma_{\mathfrak a}=W_a=a^{-1/2}\begin{pmatrix}a&\beta\\ am&a\delta'\end{pmatrix}\in\mathrm{SL}_2(\R).
\]
Then $\sigma_{\mathfrak a}$ is a scaling matrix for the cusp $\mathfrak a=1/m$ of $\Gamma_0(M)$, it normalises $\Gamma_0(M)$ (it is the Atkin--Lehner operator at the exact divisor $a$), the set of moduli of the pair $(\infty,\mathfrak a)$ is $\{c\sqrt a:\ c\ge1,\ m\mid c,\ (c,a)=1\}$, and
\[
 S_{\infty\mathfrak a}(1,n;c\sqrt a)=S(1,n\overline a;c)\qquad(m\mid c,\ (c,a)=1).
\]
\end{lemma}
\begin{proof}
$W_a$ has determinant $1$, maps $\infty$ to $1/m$, and is the classical Atkin--Lehner matrix \cite{AtkinLehner}; the stabiliser computation is standard \cite[\S2.3]{IwSpec}. A product $\gamma=\begin{pmatrix}A_0&B_0\\ MC_0&D_0\end{pmatrix}W_a\in\Gamma_0(M)W_a$ has lower row $\sqrt a\big(m(aC_0+D_0),\ (MC_0\beta+a\delta'D_0)/a\big)$ and upper-left entry $\sqrt a(A_0+mB_0)$. Put $c=m(aC_0+D_0)$; then $m\mid c$, and $(c,a)=(D_0,a)=1$ because $A_0D_0\equiv1\ (M)$. With $x=A_0+mB_0$ and $y=(MC_0\beta+a\delta'D_0)/a$ the Kloosterman phase is $e((x+ny)/c)$. Now $ay=MC_0\beta+a\delta'D_0=D_0+\beta c$ since $a\delta'=1+m\beta$, so $ay\equiv D_0\ (c)$, and $xD_0=1+M B_0C_0+mB_0D_0\equiv1\ (c)$. Hence $y\equiv\overline a\,\overline x\ (c)$, and as $\gamma$ runs over the double cosets with lower-left entry $c\sqrt a$, $x$ runs once over $(\Z/c)^*$. This gives the identity; the shift in $\sigma_{\mathfrak a}$ by $\begin{pmatrix}1&\beta/a\\0&1\end{pmatrix}$ is exactly what removes the phase $e(n\overline m/a)$ that a scaling matrix without it would produce.
\end{proof}

Thus, for fixed $(a,d,f,g,j,B)$, the $q$-sum in \eqref{eq:Sigmakloosterman} is a sum over the moduli $c\sqrt a$ of the cusp pair $(\infty,\mathfrak a)$ of $\Gamma_0(M)$, of $S_{\infty\mathfrak a}(1,\mp\varepsilon n;c\sqrt a)$ against a smooth weight. By \cref{lem:ALcusp} the Fourier coefficients of a form $u$ at $\mathfrak a$ are those of $u|W_a$ at $\infty$, and since $a$ is a positive power of $X$ the two cusps are inequivalent, so the Kuznetsov formula has no diagonal term.

\subsubsection{The Kuznetsov formula}
Write the weight in \eqref{eq:Sigmakloosterman} as $(c\sqrt a)^{-1}F(4\pi\sqrt n/(c\sqrt a))$. Then $F$ is supported on $\xi\asymp\xi_B:=K_B^{-1}$, where
\[
 K_B=\frac{Q\sqrt A}{f\sqrt B}\ \ge\ \frac{Q\sqrt A}{\sqrt{NX^\nu}}=KX^{-\nu/2},
\]
and $\|F\|_\infty\ll X^\epsilon f\cdot(Q/f)\sqrt A=X^\epsilon Q\sqrt A$, with $\xi^iF^{(i)}(\xi)\ll_iX^{\epsilon}Q\sqrt A$; the factor $f$ is the one in \eqref{eq:Sigmakloosterman}. The dependence of $F$ on $n$ (through $\sqrt n$ in the conversion) is separated by Mellin inversion in $n/B$ on a fixed enlarged interval: $F=\int F_y(\xi)(n/B)^{iy}\dd y$ with $F_y$ of the same type and rapidly decaying in $y$. We apply the Kuznetsov formula for the cusp pair $(\infty,\mathfrak a)$ of $\Gamma_0(M)$ in the form of \cite[Theorem~1]{DI} (see also \cite[Theorems~9.4,~9.8]{IwSpec}): the sum is
\begin{equation}\label{eq:kuznetsov}
 \begin{split}
 &\sum_j\overline{\rho_{j\infty}(1)}\,\rho_{j\mathfrak a}(\pm n)\,\widetilde F(t_j)+\sum_{k,j}\widetilde F(k)\overline{\psi_{jk\infty}(1)}\psi_{jk\mathfrak a}(n)\\
 &\qquad+\frac1{4\pi}\sum_{\mathfrak c}\int\overline{\rho_{\mathfrak c\infty}(1,t)}\rho_{\mathfrak c\mathfrak a}(\pm n,t)\widetilde F(t)\dd t,
 \end{split}
\end{equation}
over the Maass, holomorphic (only for the $+$ sign), and Eisenstein spectra, with the transforms $\widetilde F$ of \cite{DI}. We use the following Fourier coefficient convention throughout. Maass forms are normalised by $\int_{\Gamma_0(M)\backslash\mathbb H}|u|^2\,y^{-2}\dd x\dd y=1$, the Petersson inner product not divided by the volume, as in \cite{DI}. For such a Maass form $u$ of parameter $t$ at a cusp $\mathfrak b$, write
\[
 u(\sigma_{\mathfrak b}z)=\sqrt y\sum_{n\ne0}\rho^{\rm raw}_{u\mathfrak b}(n)K_{it}(2\pi|n|y)e(nx),
 \qquad \rho_{u\mathfrak b}(n)=\frac{\rho^{\rm raw}_{u\mathfrak b}(n)}{\sqrt{\cosh(\pi t)}} .
\]
Use the same convention for the nonconstant coefficients of Eisenstein series. In the Whittaker expansion with $W_{0,it}(4\pi|n|y)=2\sqrt{|n|y}\,K_{it}(2\pi|n|y)$ the raw $K$-coefficient is $2\sqrt{|n|}$ times the Whittaker coefficient. Thus the two Maass coefficients in \eqref{eq:kuznetsov} together carry $1/\cosh(\pi t)$, including at an exceptional parameter $t=i\tau$, where $\cosh(\pi t)=\cos(\pi\tau)\ge\cos(7\pi/64)>0.94$ for $0<\tau\le\theta$. In \cite[Theorem~1]{DI} the Maass and Eisenstein terms are products of two raw coefficients with the weight $1/\cosh(\pi t)$, up to an absolute constant factor, which we include in $\widetilde F$; in \eqref{eq:kuznetsov} the weight is absorbed into the coefficients, and otherwise $\widetilde F$ is the transform of \cite{DI}. For a holomorphic form of weight $k$, $\psi_{jk\mathfrak b}(n)$ denotes the coefficient of $e(nz)$ at that cusp multiplied by $\sqrt{(k-1)!}/(4\pi n)^{(k-1)/2}$, the Petersson large-sieve convention of \cite[(1.28)]{DI}. In \cite[Theorem~1]{DI} the holomorphic term is the product of the raw coefficients at $m$ and $n$ times $\Gamma(k)(4\pi\sqrt{mn})^{1-k}$, up to a factor depending only on $k$, which we include in $\widetilde F(k)$ (a factor polynomial in $k$ does not affect the bound for $\widetilde F(k)$ in \eqref{eq:transforms} below); and $\Gamma(k)(4\pi\sqrt{n})^{1-k}$ times the raw coefficients at $1$ and $n$ is exactly $\overline{\psi_{jk\infty}(1)}\psi_{jk\mathfrak a}(n)$.

For $\xi_B\le1$ the transforms satisfy \cite[Lemma~7.1]{DI}, \cite[Lemma~3.4]{Pascadi}
\begin{equation}\label{eq:transforms}
 \begin{gathered}
 \widetilde F(t)\ll_AX^\epsilon Q\sqrt A\,(1+|t|)^{-A}\ (t\in\R),\qquad\widetilde F(k)\ll_AX^\epsilon Q\sqrt A\,k^{-A},\\
 \widetilde F(i\tau)\ll X^\epsilon Q\sqrt A\,K_B^{2\tau}\ (0<\tau\le\theta),
 \end{gathered}
\end{equation}
the logarithm at $t=0$ being absorbed in $X^\epsilon$.

\subsubsection{Factorisation through three $L$-functions}
For a primitive Hecke form $\pi$ (Maass or holomorphic) of conductor $M_0\mid M$ with Hecke eigenvalues $\lambda_\pi(n)$, let $V_\pi(M)$ be the span of its old translates $u(\delta z)$, $\delta\mid M/M_0$. It is preserved by $W_a$, which acts unitarily on it \cite{AtkinLehner}.

\begin{lemma}[Oldform bases]\label{lem:oldforms}
$V_\pi(M)$ has an orthonormal basis $\{u_{\pi,\iota}\}$, $\iota\le\tau(M/M_0)$, whose Fourier coefficients at $\infty$ are
\[
 \rho_{u_{\pi,\iota}}(n)=\kappa_{\pi,M}\sum_{\delta\mid M/M_0}\xi_{\iota\delta}\sqrt\delta\,\lambda_\pi(n/\delta)\one_{\delta\mid n},\qquad\xi_{\iota\delta}\ll_\epsilon M^\epsilon,
\]
where the vector with $\iota=1$ is the newvector scaled to have $L^2$ norm one at level $M$ (for the measure $y^{-2}\dd x\dd y$ on $\Gamma_0(M)\backslash\mathbb H$), and $\xi_{1\delta}=\one_{\delta=1}$. The factor $\sqrt\delta$ is the same in both conventions: $u(\delta z)$ has raw $K$-Bessel coefficient $\sqrt\delta\,\rho^{\rm raw}_u(n/\delta)$ at $n$, and for a holomorphic form the translate $\delta^{k/2}f(\delta z)$ has normalised coefficient $\delta^{k/2}\delta^{-(k-1)/2}=\sqrt\delta$ times that of $f$ at $n/\delta$. For Maass forms put $\kappa^{\rm raw}_{\pi,M}=\rho^{\rm raw}_{u_{\pi,1}\infty}(1)$ and $\kappa_{\pi,M}=\kappa^{\rm raw}_{\pi,M}/\sqrt{\cosh(\pi t_\pi)}$. For holomorphic forms $\rho$ and $\kappa_{\pi,M}$ denote the Petersson-normalised coefficients and first coefficient, respectively.
\end{lemma}
\begin{proof}
This is \cite[Lemma~3.15 and (3.16)]{Humphries} for Maass forms and \cite[Theorem~8]{SPY} for holomorphic forms; see also \cite[\S3.5]{KY}. The bound $\xi_{\iota\delta}\ll M^\epsilon$ follows from the explicit local factors there, the Kim--Sarnak bound and multiplicativity: the normalised local factors are $1+O(p^{-1/2+\theta})$ for large $p$ and bounded at each fixed prime uniformly in the exponent.
\end{proof}

\begin{lemma}[Euler numerator]\label{lem:euler}
Let $\boldsymbol u\in(i\R)^3$, $j,\delta\ge1$, and $b_{g,j}$ as in \eqref{eq:Sigmakloosterman}. For $\Re z>1$,
\[
 \sum_{n\ge1}b_{g,j}(\delta n)\lambda_\pi(n)n^{-z}=g^{(u_2+u_3)/L}\,\mathcal E_{\pi,j,\delta}(z)\prod_{i=1}^3L\Big(z-\frac{u_i}L,\pi\Big),
\]
where the factor $g^{(u_2+u_3)/L}$ has modulus one, and $\mathcal E_{\pi,j,\delta}$ is holomorphic on $\Re z>\frac12$ and satisfies $\mathcal E_{\pi,j,\delta}(z)\ll_\epsilon(j\delta M)^\epsilon$ on $\Re z\ge\frac12+\epsilon$, uniformly in the imaginary parts. At primes $p\nmid j\delta M$ its local factor is
\[
 1-(z_1z_2+z_1z_3+z_2z_3)p^{-2z}+\lambda_\pi(p)z_1z_2z_3p^{-3z},\qquad z_i=p^{u_i/L}.
\]
\end{lemma}
\begin{proof}
By definition $b_{g,j}(n)=g^{(u_2+u_3)/L}b_{1,j}(n)$, so it suffices to treat $g=1$. At $p\nmid j\delta$ the local coefficients of $b_{1,j}(\delta\,\cdot)$ are, up to a unimodular constant, $b(p^v)=h_v(z_1,z_2,z_3)$, the complete homogeneous symmetric polynomial, so that $\sum_vb(p^v)Y^v=\prod_i(1-z_iY)^{-1}$. With Satake parameters $\alpha\beta=1$ and $\lambda_\pi(p^v)=(\alpha^{v+1}-\beta^{v+1})/(\alpha-\beta)$, the series $\sum_vh_v(\boldsymbol z)\lambda_\pi(p^v)Y^v$ equals $(1-e_2(\boldsymbol z)Y^2+\lambda_\pi(p)e_3(\boldsymbol z)Y^3)/\prod_i(1-\alpha z_iY)(1-\beta z_iY)$, an identity of rational functions (also checked by computer; repeated parameters follow by continuity). The displayed numerator is $1+O(p^{-1-2\epsilon})+O(p^{\theta-3/2-3\epsilon})$ on $\Re z\ge\frac12+\epsilon$, so the product over these primes converges absolutely, uniformly. At the primes dividing $j\delta M$ we keep the local series of the left side and multiply by the reciprocals of the local $L$-factors, which are polynomials of degree at most two in $p^{-z}$; the coefficient of $p^{-vz}$ in the left local series is at most $(e+1)(b+1)^2(v+1)^2p^{\theta v}$ in modulus, where $e=v_p(j)$ and $b=v_p(\delta)$, so the local factor is bounded by a constant power of $(e+1)(b+1)$, and the product over $p\mid j\delta M$ is $\ll(j\delta M)^\epsilon$.
\end{proof}

Now consider the cuspidal part of \eqref{eq:kuznetsov}, summed over $n$ against $b_{g,j}(n)(n/B)^{iy}V_0(n/B)$. Since the sum over an orthonormal basis of $V_\pi(M)$ of $\overline{\rho_u(1)}\rho_{u|W_a}(n)$ does not depend on the basis, Cauchy--Schwarz within $V_\pi(M)$, unitarity of $W_a$ and \cref{lem:oldforms} bound the contribution of the class of $\pi$ by
\[
 |\widetilde F(t_\pi)|\,M^\epsilon|\kappa_{\pi,M}|\Big(\sum_\iota\Big|\sum_nb_{g,j}(n)\Big(\frac nB\Big)^{iy}V_0\Big(\frac nB\Big)\rho_{u_{\pi,\iota}}(n)\Big|^2\Big)^{1/2}.
\]
For each $\delta\mid M/M_0$, Mellin inversion of $V_0$, a shift of the contour to $\Re z=\frac12+\epsilon$ (there are no poles, $\pi$ being cuspidal) and \cref{lem:euler} give
\[
 \begin{split}
 &\sqrt\delta\sum_{n'}b_{g,j}(\delta n')\Big(\frac{\delta n'}B\Big)^{iy}V_0\Big(\frac{\delta n'}B\Big)\lambda_\pi(n')\\
 &\qquad\ll X^\epsilon\sqrt B\int_\R\varpi(v)\prod_{i=1}^3\Big|L\Big(\tfrac12+\epsilon+iv-\tfrac{u_i}L,\pi\Big)\Big|\dd v,
 \end{split}
\]
with $\varpi$ rapidly decreasing, using $\sqrt\delta(B/\delta)^{1/2+\epsilon}\le B^{1/2+\epsilon}$. Hence the class of $\pi$ contributes at most
\begin{equation}\label{eq:classbound}
 X^\epsilon\sqrt B\,|\widetilde F(t_\pi)|\,|\kappa_{\pi,M}|^2\int\varpi(v)\prod_{i=1}^3\Big|L\Big(\tfrac12+\epsilon+iv-\tfrac{u_i}L,\pi\Big)\Big|\dd v .
\end{equation}

\subsubsection{Spectral moments}
\begin{proposition}[Moments of $L$-functions]\label[proposition]{prop:moments}
Let $\frac12\le\sigma\le\frac58$, $v\in\R$, $U\ge1$ and $R\ge1$.
\begin{enumerate}[label=\textup{(\roman*)}]
\item For each $M$,
\[
 \sum_{\substack{\pi:\ \mathrm{cond}(\pi)\mid M\\ t_\pi\in\R,\ |t_\pi|\le U}}
 \frac{|\kappa^{\rm raw}_{\pi,M}|^2}{\cosh(\pi t_\pi)}|L(\sigma+iv,\pi)|^{2k}
 \ll_\epsilon M^\epsilon(U+|v|)^{C}\qquad(k=1,2),
\]
the sum over primitive Maass forms. The same bound holds for holomorphic forms of weight at most $U$ with the Petersson weight $|\kappa_{\pi,M}|^2$.
\item Summing over the exceptional primitive forms $\pi$, with $t_\pi=i\tau_\pi$, at all levels in a dyadic range,
\[
 \begin{split}
 &\sum_{M\asymp R}\sum_{\pi\ \mathrm{exc}}|\kappa_{\pi,M}|^2\big(R^{3\tau_\pi}|L(\sigma+iv,\pi)|^2+R^{\tau_\pi}|L(\sigma+iv,\pi)|^4\big)\\
 &\qquad\ll_\epsilon R^{1+\epsilon}(2+|v|)^C .
 \end{split}
\]
\item Consequently, for $K\ge1$ and $\sigma_i,v_i$ as above,
\[
 \begin{split}
 &\sum_{M\asymp R}\sum_{\pi\ \mathrm{exc}}|\kappa_{\pi,M}|^2K^{2\tau_\pi}\prod_{i=1}^3|L(\sigma_i+iv_i,\pi)|\\
 &\qquad\ll_\epsilon R^{1+\epsilon}\max(1,K/R)^{2\theta}(2+\max|v_i|)^C .
 \end{split}
\]
\end{enumerate}
\end{proposition}
\begin{proof}
(i) Fix $M_0\mid M$. The approximate functional equation for $L(\frac12+iv,\pi)^k$ has two smooth sums of effective length $M_0^{k/2}(2+U+|v|)^{C}M^\epsilon$, with weights whose scaled derivatives are bounded by $(U+|v|)^C M^\epsilon$ (Stirling; the gamma factors are separated from their poles by the Kim--Sarnak bound) \cite[Theorem~5.3]{IK}. For $k=2$ use $L(z,\pi)^2=\zeta^{(M_0)}(2z)\sum_n\tau(n)\lambda_\pi(n)n^{-z}$, which holds at every prime (at ramified primes $\lambda_\pi(p^j)=\lambda_\pi(p)^j$), and remove the $\zeta^{(M_0)}(2z)$ by Minkowski over $m=b^2n$, at the cost of a logarithm. The $\iota=1$ vectors of distinct classes are orthonormal at level $M$. At infinity their raw Maass coefficients are $\kappa^{\rm raw}_{\pi,M}\lambda_\pi(n)$ for $n>0$. For a block $n\asymp Y$ the spectral large sieve \cite[Theorem~2]{DI} gives
\[
 \sum_{\substack{j:\ t_j\in\R\\ |t_j|\le U}}
 \frac{1}{\cosh(\pi t_j)}
 \Big|\sum_{n\asymp Y}a_n\rho^{\rm raw}_{j\infty}(n)\Big|^2
 \ll_\epsilon(MYU)^\epsilon(U^2+Y/M)\sum_{n\asymp Y}|a_n|^2 .
\]
Restricting this positive sum to the $\iota=1$ vectors supplies precisely $|\kappa^{\rm raw}_{\pi,M}|^2/\cosh(\pi t_\pi)=|\kappa_{\pi,M}|^2$ for each class. The holomorphic large sieve gives the same bound with the Petersson-normalised coefficients. Mellin inversion in $n/Y$ separates the smooth weights before applying either sieve. Their derivative bounds give an integrable majorant on the Mellin line, uniformly in $\pi$ at fixed $M_0$. In both cases the squared coefficient norm of the separated block is $\ll\sum_{n\asymp Y}\tau(n)^2/n$ (up to polynomial factors in $U+|v|$ from the smooth weights). Here $Y\le M^{1+\epsilon}(U+|v|)^C$, which gives the bound on the critical line after summing over the $\tau(M)$ divisors $M_0$. To the right of the line the direct sum gains $n^{1/2-\sigma}$. The dual sum loses at most $Y^{\sigma-1/2}\le(M_0^{k/2}(2+U+|v|)^CM^\epsilon)^{\sigma-1/2}$, but it is multiplied by the factor from the functional equation of $L(\cdot,\pi)^k$ (root number, conductor and gamma factors), of modulus $\ll M_0^{k(1/2-\sigma)}(2+U+|v|)^C$. The net factor is $\ll M_0^{-k(\sigma-1/2)/2}(2+U+|v|)^CM^\epsilon\le(2+U+|v|)^CM^\epsilon$, which gives the same bound.

(ii) Partition the conductors into dyadic ranges $M_0\asymp R_0$. As in (i), the approximate functional equation reduces $L(\sigma+iv,\pi)^k$ to dyadic blocks of the shape $\sum_{n\asymp Y}\lambda_\pi(n)n^{-1/2-iv}V(n/Y)$ ($k=1$, $Y\le R_0^{1/2+\eta}(2+|v|)^C$) or $\sum_{h\asymp Y_1,k'\asymp Y_2}\lambda_\pi(hk')(hk')^{-1/2-iv}V_1(h/Y_1)V_2(k'/Y_2)$ ($k=2$, $Y_1Y_2\le R_0^{1+\eta}(2+|v|)^C$), up to tails decaying faster than any power and a dual sum treated as in (i); the weight depending on $\pi$ and $M_0$ is removed by Mellin inversion with uniformly rapidly decaying coefficients, after which Minkowski's inequality leaves weights $V_i$ independent of $M$ and $\pi$. Enlarge the sum to all $M\asymp R$ and to a full orthonormal basis of the exceptional spectrum at each level containing the $\iota=1$ vectors; all restrictions on the conductor disappear in this positive majorant. We then apply the level-averaged large sieve for exceptional eigenvalues of Watt \cite{Watt}, in the form of \cite[Proposition~3.5]{PascadiExp}: for $Y_1\le Y_2$ and smooth $V_1,V_2$ whose scaled derivatives of order $i$ are $O(Z^i)$,
\begin{equation}\label{eq:watt}
 \begin{split}
 \sum_{M\asymp R}\sum_{j\ \mathrm{exc}}X_0^{2\tau_j}\Big|\sum_{h,k'}V_1\Big(\frac h{Y_1}\Big)V_2\Big(\frac{k'}{Y_2}\Big)\rho_{j\infty}(hk')\Big|^2&\ll_\epsilon R^\epsilon Z^5(R+Y_1Y_2)Y_1Y_2\\
 &\qquad\Big(X_0\ll\frac{R^2}{Y_1^2Y_2}\Big).
 \end{split}
\end{equation}
For $k=1$ take $Y_1=1$, $Y_2=Y\le R^{1/2+\eta}(2+|v|)^C$, which permits $X_0=R^{3/2-O(\eta)}(2+|v|)^{-C}$; for $k=2$ take $Y_1\le(Y_1Y_2)^{1/2}$, which permits $X_0=R^{1/2-O(\eta)}(2+|v|)^{-C}$. After the normalisation $(Y_1Y_2)^{-1/2}$ both give $R^{1+O(\eta)}(2+|v|)^C$, and $\tau_\pi\le\theta$ restores the exact weights $R^{3\tau}$, $R^{\tau}$ at the cost of $R^{O(\eta)}$. Let $\eta\to0$.

(iii) H\"older's inequality with exponents $(2,4,4)$ over the joint sum, using (ii) with $k=1$ in the first factor and $k=2$ in the other two, gives the weight $R^{(\frac32+\frac14+\frac14)\tau_\pi}=R^{2\tau_\pi}$; then $K^{2\tau}\le R^{2\tau}\max(1,K/R)^{2\theta}$.
\end{proof}

\subsubsection{The continuous spectrum}
\begin{lemma}\label{lem:eisenstein}
For $M=am$ with $(a,m)=1$ let $M^\flat=\prod_{p^e\Vert M}p^{\lceil e/2\rceil}$. With the Eisenstein coefficients normalised as in \eqref{eq:kuznetsov}, for $n\ne0$,
\[
 \Big|\sum_{\mathfrak c}\overline{\rho_{\mathfrak c\infty}(1,t)}\rho_{\mathfrak c\mathfrak a}(n,t)\Big|\ll_\epsilon\frac{(n,M^\flat)^{1/2}}{M^\flat}(M|n|(2+|t|))^\epsilon .
\]
\end{lemma}
\begin{proof}
In \cite{BM} the Eisenstein series are expanded in the Whittaker functions $W_{0,it}(4\pi|n|y)$. Since $W_{0,it}(4\pi|n|y)=2\sqrt{|n|y}\,K_{it}(2\pi|n|y)$, a Whittaker coefficient $\rho^W(n,t)$ corresponds to the raw coefficient $2\sqrt{|n|}\,\rho^W(n,t)$, and to the coefficient $2\sqrt{|n|}\,\rho^W(n,t)/\sqrt{\cosh(\pi t)}$ in the normalisation of \eqref{eq:kuznetsov}. (At level one, for example, $|\rho_{\infty\infty}(n,t)|^2=4|\sigma_{-2it}(|n|)|^2/|\zeta(1+2it)|^2$, using $|\Gamma(\frac12+it)|^2=\pi/\cosh(\pi t)$.) After this conversion, \cite[Lemma~1, (3.10), (4.6)]{BM} give
\[
 \sum_{\mathfrak c}|\rho_{\mathfrak c\infty}(n,t)|^2\ll(n,M^\flat)(M^\flat)^{-1}(M|n|(2+|t|))^\epsilon .
\] Replacing the observation cusp $\infty$ by $\mathfrak a=W_a\infty$ permutes the Eisenstein series and acts unitarily on the continuous fibre, so the same bound holds at $\mathfrak a$. Cauchy--Schwarz gives the claim.
\end{proof}

By \eqref{eq:transforms} and \cref{lem:eisenstein}, the continuous part of \eqref{eq:kuznetsov}, summed over $n\asymp B$ with $|b_{g,j}(n)|\le\alpha_j(n)$, is
\[
 \ll X^{\epsilon}Q\sqrt A\,\frac1{M^\flat}\sum_{n\asymp B}\alpha_j(n)(n,M^\flat)^{1/2}\ll X^\epsilon Q\sqrt A\,\frac B{a^\flat m^\flat},
\]
using $M^\flat=a^\flat m^\flat$ (as $(a,m)=1$) and $\sum_{n\asymp B}(n,M^\flat)^{1/2}\ll B\tau(M)$. Since $d\mid mG$ and $(uv)^\flat\mid u^\flat v^\flat$, we have $1/m^\flat\le f/d^\flat$, and $\sum_{n\le Z}1/n^\flat\ll(\log 2Z)^2$ (write $n=u^2v$ with $v$ squarefree, so $n^\flat=uv$). Summing over $a\asymp A$, $d\asymp D$, the dyadic $B\le NX^\nu/f^2$, and then over $g,j,f$ with $\sum_{f\le X}\tau(f)/f\ll(\log X)^2$, the continuous contribution to the triple sum in \eqref{eq:Sigmakloosterman} is $\ll X^{\epsilon+O(\nu)}Q\sqrt A\,N$. Relative to $AQ^2$ this is
\begin{equation}\label{eq:continuousbound}
 X^{\epsilon+O(\nu)}\frac N{Q\sqrt A}=X^{(3s+h)/2-1+\epsilon+O(\nu)},
\end{equation}
which is a power saving on $\mathcal K$ since $s,h<\frac12$.

\subsubsection{Proof of \texorpdfstring{\cref{thm:signed}}{the signed estimate}}
Fix $\boldsymbol\varepsilon,f=gj,B$ and $G\mid f$, and sum \eqref{eq:classbound} over the levels. For fixed $(l,G)$ the level is $M=a\,ld/G\asymp R:=AlD/G$, and each $M$ arises from at most $\tau(M)$ pairs $(a,d)$, since $a\mid M$ determines $d=MG/(al)$. We may drop the remaining restrictions on $(a,d)$ in this positive majorant. The three $L$-values in \eqref{eq:classbound} have arguments independent of $(a,d)$ at each fixed $v$, and all other factors are bounded by $X^\epsilon$ times integrable majorants independent of the pair. By \cref{prop:moments}(i) (with H\"older's inequality as in (iii), and rapid decay of $\widetilde F$ in $t$ and $k$) for the regular spectrum, and \cref{prop:moments}(iii) with $K=K_B$ and \eqref{eq:transforms} for the exceptional spectrum, the cuspidal contribution summed over $(a,d)$ is
\[
 \ll X^{\epsilon+O(\nu)}Q\sqrt{AB}\,R\big(1+(K_B/R)^{2\theta}\big).
\]
As $2\theta<1$, every power of $B$ here is positive, so $B\le NX^\nu/f^2$ (which turns $K_B$ into $K$ up to $X^{O(\nu)}$) gives $X^{\epsilon+O(\nu)}Q\sqrt{AN}f^{-1}R(1+(K/R)^{2\theta})$. Dividing by $AQ^2$, with $R=AlD/G$,
\[
 X^{\epsilon+O(\nu)}\frac{D\sqrt{AN}}{Qf}\cdot\frac lG\Big(1+\Big(\frac{KG}{AlD}\Big)^{2\theta}\Big)\le X^{\epsilon+O(\nu)}\frac{l}{f}\cdot\frac{D\sqrt{AN}}Q\Big(1+\Big(\frac K{AD}\Big)^{2\theta}\Big),
\]
using $G,l\ge1$ and $2\theta<1$. By the exponent identities (also checked by computer)
\[
 \frac{D\sqrt{AN}}Q=X^{-p},\qquad\frac K{AD}=X^{p},
\]
this is $X^{\epsilon+O(\nu)}l f^{-1}X^{-(1-2\theta)p}$. Summing over $G\mid f$, $g,j$ with $gj=f$, and $f\le X$ costs $\sum_f\tau(f)^2/f\ll X^\epsilon$; the $O(\log X)$ dyadic ranges of $B$ and the eight sign vectors cost $X^\epsilon$. With \eqref{eq:continuousbound}, the triple sum in \eqref{eq:Sigmakloosterman} is
\[
 \ll AQ^2\big(l\,X^{-(1-2\theta)p}+X^{(3s+h)/2-1}\big)X^{\epsilon+O(\nu)} .
\]
Choose $\nu,\epsilon$ small and $c<\frac12\min_{\mathcal K}\min\big((1-2\theta)p,\,1-\frac{3s+h}2\big)$; since $l\le X^c$, this proves \cref{thm:signed}. \qed

\begin{remark}
The exceptional spectrum is where the level average matters. At a single level, \cite[Lemma~3.8]{Pascadi} only gives the factor $(1+K_B/M)^{2\theta}$ with $M=am$ for each pair $(a,d)$, and the extra factor $(1+\sqrt{N/(Ald)})$ from the large sieve on the $n$-side; summing over $d$ then shrinks the region to $(2k-1-x)_+<2(1-2\theta)(z-\frac12-x)$, which, added in both orientations to $S_3$, gives only the coefficient $15953272998824493456684483/384777072096051200000000=41.4610800792\ldots$. The factorisation through $L$-functions avoids the large sieve on the $n$-side altogether, and averaging over the levels $M=a\,ld/G$ gives \cref{prop:moments}(iii) with $R$ in place of $M$.
\end{remark}

\subsection{Reflection and localisation}
For a function $g$ on $\R$ put $g^\vee(a)=\overline{g(1-a)}$, and for $T\le t\le2T$ put
\[
 \ell_t=\frac1L\log\frac t{2\pi T}\in\Big[-\frac2L,0\Big].
\]
Write $\|B\|_6=(\int_T^{2T}|B(\frac12+it)|^6\dd t)^{1/6}$. By H\"older's inequality, $\int|\Phi(t/T)B_1B_2B_3\overline{B_4}|\dd t\le\|\Phi\|_\infty^{1/2}\|B_1\|_6\|B_2\|_6\|B_3\|_6\|B_4\|_\Phi$, and $\|\mathcal JB\|_\Phi=\|B\|_\Phi$ since $|\chi(\frac12+it)|=1$. We use \cref{lem:sixth}(i) in the form $\|A_{\chi_0e^{i\xi\cdot}}\|_6\ll_{\chi_0}(TL^9)^{1/6}$ for $|\xi|\le LT/4$; the case $\xi=0$ applies to every $A_h$ with $h\in C_c^\infty(\R)$.

\begin{lemma}[Reflected factors]\label{lem:reflfactor}
Let $\chi_0\in C_c^\infty((0,1))$ be real and $g_\xi=\chi_0e^{i\xi\cdot}$. Uniformly for $|\xi|\le T^{1/4}$ and $T\le t\le2T$,
\begin{equation}\label{eq:reflpointwise}
 \chi(s)\overline{A_{g_\xi}(s)}=B_\xi(t)+O\big((1+|\xi|)^2T^{-1/2}L\big),\quad B_\xi(t)=\sum_{k\ge1}k^{-s}\,\overline{g_\xi}\Big(\frac1L\log\frac t{2\pi k}\Big),
\end{equation}
and
\[
 \big\|B_\xi-A_{g_\xi^\vee}\big\|_6\ll\frac{1+|\xi|}{L}(TL^9)^{1/6},\qquad \big\|\chi\overline{A_{g_\xi}}-A_{g_\xi^\vee}\big\|_6\ll\frac{(1+|\xi|)^2}{L}(TL^9)^{1/6}.
\]
\end{lemma}
\begin{proof}
Put $t'=t-\xi/L$ and $s'=\frac12+it'$. Since $A_{g_\xi}(s)=A_{\chi_0}(s')$, and the proof of \cref{lem:dualfactor} applies verbatim for $t'\in[T/2,3T]$,
\[
 \chi(s')\overline{A_{g_\xi}(s)}=\sum_kk^{-s'}\chi_0\Big(\frac1L\log\frac{t'}{2\pi k}\Big)+O(T^{-1/2}L).
\]
Write $\chi(\frac12+i\tau)=e^{-i\vartheta(\tau)}$; by Stirling's formula $\vartheta'(\tau)=\log(\tau/2\pi)+O(\tau^{-2})$ and $\vartheta''(\tau)\ll1/\tau$. Hence $\chi(s)/\chi(s')=\exp(-i(\xi/L)\log(t/2\pi))(1+O((|\xi|+\xi^2)/T))$, the linear term coming from the remainder $O(\tau^{-2})$ in $\vartheta'$, and as $k^{-s'}=k^{-s}k^{i\xi/L}$,
\[
 \frac{\chi(s)}{\chi(s')}k^{-s'}=k^{-s}e^{-i\xi\log(t/2\pi k)/L}\big(1+O((|\xi|+\xi^2)/T)\big).
\]
Moreover $\chi_0(\log(t'/2\pi k)/L)=\chi_0(\log(t/2\pi k)/L)+O(|\xi|/(LT))$. Only $k\le2T$ occur, and $\sum_{k\le2T}k^{-1/2}\ll T^{1/2}$; multiplying by $\chi(s)/\chi(s')$, which has modulus $1$, gives \eqref{eq:reflpointwise}.

For the second claim let $G=\overline{g_\xi}$, so that $\widehat G(\eta)=\widehat\chi_0(\eta+\xi)$, and note $\frac1L\log\frac t{2\pi k}=1-\frac{\log k}L+\ell_t$. Choose $\omega\in C_c^\infty((0,1))$ equal to $1$ on a neighbourhood of $1-\operatorname{supp}\chi_0$; for $T$ large, $\omega(a)=1$ whenever $G(1-a+\ell)\ne0$ or $G(1-a)\ne0$ with $|\ell|\le2/L$. By Fourier inversion,
\[
 B_\xi(t)-A_{g_\xi^\vee}(s)=\frac1{2\pi}\int_\R\widehat G(\eta)e^{i\eta}\big(e^{i\eta\ell_t}-1\big)A_{\omega e^{-i\eta\cdot}}(s)\dd\eta .
\]
By Minkowski's inequality and $|e^{i\eta\ell_t}-1|\le2|\eta|/L$, the left side has $\|\cdot\|_6$ at most $\frac1\pi L^{-1}\int|\eta||\widehat\chi_0(\eta+\xi)|\,\|A_{\omega e^{-i\eta\cdot}}\|_6\dd\eta$. For $|\eta|\le LT/4$ the last norm is $\ll(TL^9)^{1/6}$, and $\int|\eta||\widehat\chi_0(\eta+\xi)|\dd\eta\ll1+|\xi|$. For $|\eta|>LT/4$ we use $\|A_{\omega e^{-i\eta\cdot}}\|_6\ll T^{2/3}$ and $|\widehat\chi_0(\eta+\xi)|\ll_A|\eta|^{-A}$, valid as $|\xi|\le T^{1/4}$. The third claim follows from the first two, since the error in \eqref{eq:reflpointwise} has $\|\cdot\|_6\ll(1+|\xi|)^2T^{-1/3}L$.
\end{proof}

For a function $F$ on $[0,1]^3$ recall $F^J(c)=F(1-c_1,1-c_2,1-c_3)$.

\begin{lemma}[Reflected cubes]\label{lem:reflcube}
Let $c\in C_c^\infty((0,1)^3)$ be real. Then $\|\mathcal JP_c-P_{c^J}\|_\Phi\ll_cL^{-1}(TL^9)^{1/2}$. Consequently, for real $\psi\in C_c^\infty(\mathcal Y)$,
\[
 \|\mathcal JW_\psi-P_{c_\psi^J}\|_\Phi=o\big((TL^9)^{1/2}\big),\qquad h_{c_\psi^J}=\psi^\flat\circ J .
\]
\end{lemma}
\begin{proof}
Let $\chi_0\in C_c^\infty((0,1))$ be real and equal to $1$ on the three coordinate projections of $\operatorname{supp}c$, and $g_\xi=\chi_0e^{i\xi\cdot}$. As in the proof of \cref{lem:sixth}(ii), $P_c=(2\pi)^{-3}\int\widehat c(\xi)\prod_jA_{g_{\xi_j}}\dd\xi$, so
\[
 \mathcal JP_c=(2\pi)^{-3}\int\overline{\widehat c(\xi)}\prod_{j=1}^3R_{\xi_j}\dd\xi,\qquad R_\xi=\chi\overline{A_{g_\xi}} .
\]
Put $Q_\xi=A_{g_\xi^\vee}$. Then $(2\pi)^{-3}\int\overline{\widehat c(\xi)}\prod_jQ_{\xi_j}\dd\xi=P_{c'}$ with
\[
 c'(a)=\prod_j\chi_0(1-a_j)\overline{\frac1{(2\pi)^3}\int\widehat c(\xi)e^{i\xi\cdot(1-a)}\dd\xi}=\prod_j\chi_0(1-a_j)\,c(1-a)=c^J(a),
\]
since $c$ is real and $\chi_0(1-a_j)=1$ wherever $c(1-a)\ne0$. By Minkowski's inequality
\[
 \|\mathcal JP_c-P_{c^J}\|_\Phi\le(2\pi)^{-3}\int|\widehat c(\xi)|\,\Big\|\prod_jR_{\xi_j}-\prod_jQ_{\xi_j}\Big\|_\Phi\dd\xi .
\]
If $\max_j|\xi_j|\le T^{1/4}$, write the difference of products as $(R_1-Q_1)R_2R_3+Q_1(R_2-Q_2)R_3+Q_1Q_2(R_3-Q_3)$ with $R_j=R_{\xi_j}$, $Q_j=Q_{\xi_j}$. Here $\|R_j\|_6=\|A_{g_{\xi_j}}\|_6\ll(TL^9)^{1/6}$; $Q_j=e^{-i\xi_j}A_{\chi_0(1-\cdot)e^{i\xi_j\cdot}}$, so also $\|Q_j\|_6\ll(TL^9)^{1/6}$; and $\|R_j-Q_j\|_6\ll(1+|\xi|)^2L^{-1}(TL^9)^{1/6}$ by \cref{lem:reflfactor}. H\"older's inequality bounds the norm of the difference by $\ll(1+|\xi|)^2L^{-1}(TL^9)^{1/2}$. If $\max_j|\xi_j|>T^{1/4}$ we use $|R_j|,|Q_j|\ll T^{1/2}$, and $\int_{\max|\xi_j|>T^{1/4}}|\widehat c|\ll_AT^{-A}$. This proves the first claim, as $\int(1+|\xi|)^2|\widehat c(\xi)|\dd\xi<\infty$.

For the second, $c_\psi\in C_c^\infty((0,1)^3)$ is real, and
\[
 \begin{split}
 \|\mathcal JW_\psi-P_{c_\psi^J}\|_\Phi&\le\|\mathcal J(W_\psi-P_{c_\psi})\|_\Phi+\|\mathcal JP_{c_\psi}-P_{c_\psi^J}\|_\Phi\\
 &=\|W_\psi-P_{c_\psi}\|_\Phi+O\big(L^{-1}(TL^9)^{1/2}\big),
 \end{split}
\]
which is $o((TL^9)^{1/2})$ by \cref{prop:realisation}. Finally $h_{c_\psi^J}=h_{c_\psi}\circ J=\psi^\flat\circ J$ by \cref{lem:reflectedprofile,lem:markedprofile}.
\end{proof}

\begin{lemma}[Centred functional equation]\label{lem:centredfe}
Let $p$ be a real polynomial and $p_t(X)=p(-X-\ell_t)$. Uniformly for $T\le t\le2T$, $\chi(s)\overline{C_p(s)}=C_{p_t}(s)+O_p(T^{-3/4})$.
\end{lemma}
\begin{proof}
By linearity we may take $p=X^d$. Since $e^{\gamma(\frac12+\frac1L\frac d{ds})}\zeta(s)=e^{\gamma/2}\zeta(s+\gamma/L)$, the polynomial $C_{X^d}(s)$ is $d!$ times the coefficient of $\gamma^d$ in $e^{\gamma/2}\zeta(s+\gamma/L)$. On the critical line $\overline{\zeta(w)}=\zeta(\bar w)$ and $\bar s=1-s$, so for real $\gamma$ the functional equation $\zeta(1-w)=\chi(1-w)\zeta(w)$ and $\chi(w)\chi(1-w)=1$ give
\[
 \chi(s)\overline{e^{\gamma/2}\zeta(s+\gamma/L)}=e^{\gamma/2}\,\frac{\chi(s)}{\chi(s-\gamma/L)}\,\zeta(s-\gamma/L)=:G_1(\gamma).
\]
Similarly $C_{p_t}(s)$ is $d!$ times the coefficient of $\gamma^d$ in
\[
 G_2(\gamma)=e^{-\gamma(\frac12+\frac1L\frac d{ds})-\gamma\ell_t}\zeta(s)=e^{-\gamma/2-\gamma\ell_t}\zeta(s-\gamma/L)
\]. Both $G_1$ and $G_2$ are holomorphic for $|\gamma|\le1$, and as $G_1$ is real-analytic in real $\gamma$ with Taylor coefficients $\chi\overline{C_{X^d}}/d!$, we have $\chi\overline{C_{X^d}}=G_1^{(d)}(0)$. By Stirling's formula $\chi'/\chi(\sigma+i\tau)=-\log(\tau/2\pi)+O(1/\tau)$ uniformly for $|\sigma|\le2$, $\tau\ge1$; integrating along the segment from $s-\gamma/L$ to $s$ gives $\chi(s)/\chi(s-\gamma/L)=(t/2\pi)^{-\gamma/L}(1+O(1/(LT)))$, and $(t/2\pi)^{-\gamma/L}=e^{-\gamma(1+\ell_t)}$. Hence on $|\gamma|=1$ we have $|G_1-G_2|\ll|\zeta(s-\gamma/L)|/(LT)\ll T^{-3/4}$, by the convexity bound. Cauchy's formula on $|\gamma|=1$ bounds $G_1^{(d)}(0)-G_2^{(d)}(0)$ by $O_d(T^{-3/4})$.
\end{proof}

\begin{lemma}[Degree localisation]\label{lem:localise}
Let $\psi\in C_c^\infty(\mathcal Y)$ be real and $\varepsilon_0>0$. Let $B(t)=\sum_{N\le T^C}b_t(N)N^{-1/2-it}$, where each $b_t(N)$ is smooth in $t\in[T/2,3T]$ with $\partial_t^jb_t(N)\ll_jT^{C-j}$, and $b_t(N)=0$ whenever $\operatorname{dist}(\log N/L,I_\psi)<\varepsilon_0$. Then $\ip B{W_\psi}_\Phi\ll_AT^{-A}$.
\end{lemma}
\begin{proof}
We have $\ip B{W_\psi}_\Phi=\sum_{N,n}\overline{M_\psi(n)}(Nn)^{-1/2}\int\Phi(t/T)b_t(N)(n/N)^{it}\dd t$, and $M_\psi(n)\ne0$ only if $\log n/L\in I_\psi$, so every non-zero term has $|\log(n/N)|\ge\varepsilon_0L$. Integrating by parts $j$ times bounds each integral by $\ll_jT^{1+C-j}$. There are $O(T^{C+2})$ pairs and $M_\psi(n)\ll n^\epsilon$; take $j$ large.
\end{proof}

\begin{lemma}[Factor products]\label{lem:factors}
Let $f,g\in C_c^\infty((0,\frac12))$ be real, let $p$ be a real polynomial, and let $X=X_1X_2X_3$ with
\[
 X_1\in\{A_f,\ \chi\overline{A_f}\},\qquad X_2\in\{A_g,\ \chi\overline{A_g}\},\qquad X_3\in\{C_p,\ \chi\overline{C_p}\}.
\]
Put $F(a)=F_1(a_1)F_2(a_2)F_3(a_3)$, where $F_1=f$ if $X_1=A_f$ and $F_1=f^\vee$ otherwise, $F_2=g$ or $g^\vee$ in the same way, and $F_3(a)=p(\frac12-a)$ if $X_3=C_p$ and $F_3(a)=p(a-\frac12)$ otherwise. Then for real $\psi\in C_c^\infty(\mathcal Y)$
\[
 \ip X{W_\psi}_\Phi=a_3TL^9\Big(\int\Phi\Big)\ip{h_F}{\psi^\flat}+o(TL^9).
\]
\end{lemma}
\begin{proof}
Let $\Phi^+$ be a fixed smooth majorant of $|\Phi|$ with the same kind of support; by \cref{prop:markedgram}, $\|W_\psi\|_\Phi^2\le\ip{W_\psi}{W_\psi}_{\Phi^+}\ll TL^9$. Throughout, an error $E(t)$ is admissible if $\int|\Phi(t/T)E\,\overline{W_\psi}|\dd t=o(TL^9)$.

\emph{Step 1: explicit factors.} Replace each reflected short factor $\chi\overline{A_f}$ by $B_f(t)=\sum_kk^{-s}f(\log(t/2\pi k)/L)$, using \eqref{eq:reflpointwise} with $\xi=0$, and $\chi\overline{C_p}$ by $C_{p_t}$, using \cref{lem:centredfe}. The errors are admissible: $|A_f|,|A_g|\ll T^{1/4}$ because $f,g$ are supported in $(0,\frac12)$, hence also $|B_f|,|B_g|\ll T^{1/4}$ by \eqref{eq:reflpointwise}, $\int_T^{2T}|C_p|^2\ll_pTL^C$, and so, for instance, an error $O(T^{-1/2}L)$ in one short factor contributes $\ll T^{-1/2}L\cdot T^{1/4}(TL^C)^{1/2}\|W_\psi\|_\Phi\ll T^{3/4}L^C$, while the error $O(T^{-3/4})$ in the third factor contributes $\ll T^{-3/4}T^{1/2}\cdot T^{1/2}\|W_\psi\|_\Phi$. Next expand $C_p$, respectively $C_{p_t}=\sum_{j\le\deg p}(-\ell_t)^j(j!)^{-1}C_{(p^{(j)})^-}$, by \cref{lem:expansion}: up to $O(T^{-A})$ the third factor is $\sum_m\pi_t(\log m/L)\,\omega(m/T^8)m^{-s}$ with $\pi_t(a)=p(\frac12-a)$, respectively $\pi_t(a)=p(a-\frac12-\ell_t)$.

\emph{Step 2: localisation.} Let $S_1,S_2\subset(0,1)$ be the supports of $F_1,F_2$, fix $\varepsilon_0$ with $0<\varepsilon_0<\min(\frac13,\operatorname{dist}(I_\psi,\R\setminus(1,2)))$, so that the sets $K_3$ below and $I_\psi+[-\varepsilon_0,\varepsilon_0]$ of Step~4 lie in $(-2,3)$ and $(1,2)$ respectively, and let $K_3=I_\psi-S_1-S_2+[-3\varepsilon_0,3\varepsilon_0]$. For $T$ large a term $k^{-s}$ of $B_f$ has $\log k/L\in1-\operatorname{supp}f+[-\varepsilon_0,\varepsilon_0]$, so the first two factors only involve $n_1,n_2$ with $\log n_i/L\in S_i+[-\varepsilon_0,\varepsilon_0]$. Choose $\phi_3\in C_c^\infty((-2,3))$ equal to $1$ on $K_3$ and replace $\pi_t(\log m/L)\omega(m/T^8)$ by $\pi_t(\log m/L)\phi_3(\log m/L)$. The removed terms are products $n_1n_2m=N$ with $\operatorname{dist}(\log N/L,I_\psi)\ge\varepsilon_0$, and their coefficients are smooth in $t$ with the derivative bounds of \cref{lem:localise}; so the change is $O(T^{-A})$.

\emph{Step 3: removing the $t$-dependence.} After Step 2 each factor has $\|\cdot\|_6\ll(TL^9)^{1/6}$: for $A_f$, $A_g$ and $A_{\pi\phi_3}$ by \cref{lem:sixth}(i), and for $B_f$ by \cref{lem:reflfactor}. By that lemma $\|B_f-A_{f^\vee}\|_6\ll L^{-1}(TL^9)^{1/6}$, and in the reflected third factor the terms with $j\ge1$ carry $|\ell_t|^j\le(2/L)^j$ and have $\|\cdot\|_6\ll(TL^9)^{1/6}$. Telescoping and H\"older's inequality replace the product by $A_{F_1}A_{F_2}A_{F_3\phi_3}$ with an error that is $O(L^{-1}TL^9)$ after pairing with $W_\psi$. This product is $P_c$ with $c(a)=F(a)\phi_3(a_3)$, and $c\in C_c^\infty(\R^3)$.

\emph{Step 4: evaluation.} Choose $\phi_1\in C_c^\infty((1,2))$ equal to $1$ on $I_\psi+[-\varepsilon_0,\varepsilon_0]$. By \cref{lem:localise}, $\ip{P_c}{W_\psi}_\Phi=\ip{P_{c_1}}{W_\psi}_\Phi+O(T^{-A})$ with $c_1(a)=\phi_1(a_1+a_2+a_3)c(a)$, and \cref{prop:mixed} applies to $c_1$. Its profile is $h_{c_1}=\phi_1(r)h_{F\phi_3(c_3)}$, because $\kappa(x,y,z;c)$ vanishes unless $c_1+c_2+c_3=r$. On $\{r\in I_\psi\}$, which contains the support of $\psi^\flat$, the factor $\phi_1(r)$ is $1$, and every $c$ with $\kappa(x,y,z;c)F(c)\ne0$ has $c_i\in S_i$ ($i=1,2$) and $c_3=r-c_1-c_2\in K_3$, where $\phi_3=1$. Hence $\ip{h_{c_1}}{\psi^\flat}=\ip{h_F}{\psi^\flat}$.
\end{proof}

\subsection{Proof of \Cref{thm:marked}}\label{sec:markedproof}
Fix $\Phi\in\mathcal W$. All tests in $\mathcal T_3^\infty$ have the form $\Pi K$, with $K$ one of
\[
 D_\phi,\qquad U=\chi C_pA_f\overline{A_g}=A_f\cdot\chi\overline{A_g}\cdot C_p,\qquad V=C_pA_fA_g ,
\]
the last with $f\in C_c^\infty((0,a_0))$, $g\in C_c^\infty((0,b_0))$ and $a_0+b_0<\frac12$; for $U$, $a_0<\frac14$ and $b_0<\frac25$ since $3a_0+5b_0<2$. Recall $\Pi K=\frac12(K+\mathcal JK)$, $h_{\Pi K}=\frac12(h_{F_K}+h_{F_K}\circ J)$, and that $\mathcal J$ is an involution with $\ip{\mathcal JH}{\mathcal JK}_\Phi=\overline{\ip HK_\Phi}$ for real $\Phi$; in particular $\mathcal J\Pi K=\Pi K$ and $\mathcal J\zeta^3=\zeta^3$.

\emph{The span.} The profiles of $V$-tests are uniform limits of profiles of $V'$-tests: if $f,g$ are the polynomial data of a $V$-test on $[0,\eta]$, $\eta<\frac14$, choose $f_n,g_n\in C_c^\infty((0,\eta))$ converging uniformly to $f,g$, which is possible because $f,g$ vanish at $0$ and $\eta$; with $a_0=b_0=\eta$ these are $V'$-data, and $|h_F|\le\|F\|_\infty v$ (\cref{lem:profiles}(ii)). Hence $\mathcal H(\mathcal T_3^\infty)=\mathcal H(\mathcal T_3)$, and $v\one_{S_3}\in\mathcal H(\mathcal T_4)$ by \Cref{thm:third}. Let $\mathcal A^\flat=\{(x,k-x,r-k):(r,k,x)\in\mathcal Y\}$, which by the description of $\mathcal Y$ is $\Om\cap\mathcal D_+\cap\mathcal Z\cap\{1<r<2,\ k>\frac12\}$. The functions $\psi^\flat$, $\psi\in C_c^\infty(\mathcal Y)$ real, are dense in $L^2(\mathcal A^\flat)$, and the $\psi^\flat\circ J$ are dense in $L^2(J\mathcal A^\flat)$; so $\mathcal H(\mathcal T_4)\supseteq L^2(\mathcal A^\flat\cup J\mathcal A^\flat)$. We claim that $\mathcal Z\setminus S_3\subseteq\mathcal A^\flat\cup J\mathcal A^\flat$ up to a null set. Let $(x,y,z)\in\mathcal Z\setminus S_3$ lie off the null set where one of the defining inequalities is an equality. As it is not in $\mathcal A$, $1<r<2$; as it is not in $\mathcal D_{\rm acc}$, $x+y>\frac12$ and $y+z<\frac32$. If $x+y<z$ it lies in $\mathcal A^\flat$. Otherwise it lies in $J\mathcal D_+$, and its image under $J$ lies in $\mathcal D_+\cap\mathcal Z$, has trace $3-r\in(1,2)$ and has $(1-z)+(1-y)=2-y-z>\frac12$, so lies in $\mathcal A^\flat$. Hence $v\one_{\mathcal Z\setminus S_3}\in\mathcal H(\mathcal T_4)$, and so $v\one_{S_3\cup\mathcal Z}=v\one_{S_3}+v\one_{\mathcal Z\setminus S_3}\in\mathcal H(\mathcal T_4)$.

\emph{Faithfulness.} Pairs within $\mathcal T_3^\infty$ are covered by \Cref{thm:third}. Let $\psi,\eta\in C_c^\infty(\mathcal Y)$ be real.

(a) \emph{Numerators.} By \cref{lem:expansion}, $\zeta(s)^3=\sum_md_3(m)\omega(m/T^8)m^{-s}+O(T^{-A})$. With $\phi_1$ as in Step 4 of the proof of \cref{lem:factors}, \cref{lem:localise} removes the terms with $\phi_1(\log m/L)\ne1$, so $\mathrm N_T(W_\psi)=\ip{P_{c_0}}{W_\psi}_\Phi+O(T^{-A})$, where $c_0(a)=\phi_1(a_1+a_2+a_3)\omega_0(a)$ and $\omega_0\in C_c^\infty(\R^3)$ is $1$ on $[0,2]^3$. By \cref{prop:mixed}, $\mathrm N_T(W_\psi)/\mathcal N_T\to\ip{h_{c_0}}{\psi^\flat}=\ip{\phi_1(r)v}{\psi^\flat}=\ip v{\psi^\flat}$, using $h_1=v$. Next, $\mathrm N_T(\mathcal JW_\psi)=\ip{\mathcal J\zeta^3}{\mathcal JW_\psi}_\Phi=\overline{\mathrm N_T(W_\psi)}$, whose normalised limit is $\ip v{\psi^\flat}=\ip v{\psi^\flat\circ J}$ because $v\circ J=v$ and $J$ preserves measure.

(b) \emph{Marked Gram entries.} $\mathrm G_T(W_\psi,W_\eta)/\mathcal N_T\to\ip{\psi^\flat}{\eta^\flat}$ by \cref{prop:markedgram}, and $\mathrm G_T(\mathcal JW_\psi,\mathcal JW_\eta)=\overline{\mathrm G_T(W_\psi,W_\eta)}$ has the limit $\ip{\psi^\flat}{\eta^\flat}=\ip{\psi^\flat\circ J}{\eta^\flat\circ J}$. For the mixed orientation, \cref{lem:reflcube} (which rests on \cref{prop:realisation}) gives $\|\mathcal JW_\eta-P_{c_\eta^J}\|_\Phi=o((TL^9)^{1/2})$, so by Cauchy--Schwarz and the bound $\|W_\psi\|_\Phi\ll(TL^9)^{1/2}$, $\mathrm G_T(\mathcal JW_\eta,W_\psi)=\ip{P_{c_\eta^J}}{W_\psi}_\Phi+o(TL^9)$. Every $a\in\operatorname{supp}c_\eta^J$ has trace in $3-I_\eta\subset(1,2)$, so \cref{prop:mixed} gives the limit $\ip{h_{c_\eta^J}}{\psi^\flat}=\ip{\eta^\flat\circ J}{\psi^\flat}$; and $\mathrm G_T(W_\psi,\mathcal JW_\eta)$ is its complex conjugate.

(c) \emph{Cross-pairings with $\mathcal T_3^\infty$.} Since $\mathrm G_T(\Pi K,\mathcal JW_\psi)=\overline{\mathrm G_T(\mathcal J\Pi K,W_\psi)}=\overline{\mathrm G_T(\Pi K,W_\psi)}$ and $h_{\Pi K}$ is $J$-symmetric, it suffices to show
\[
 \frac{\mathrm G_T(K,W_\psi)}{\mathcal N_T}\to\ip{h_{F_K}}{\psi^\flat},\qquad\frac{\mathrm G_T(\mathcal JK,W_\psi)}{\mathcal N_T}\to\ip{h_{F_K}\circ J}{\psi^\flat}
\]
for each $K$; then $\mathrm G_T(\Pi K,W_\psi)/\mathcal N_T\to\ip{h_{\Pi K}}{\psi^\flat}$.

For $K=U$ and $K=V$ these are \cref{lem:factors} with
\[
 \begin{aligned}
 U&=A_f\cdot\chi\overline{A_g}\cdot C_p,&\qquad\mathcal JU&=\chi\overline{A_f}\cdot A_g\cdot\chi\overline{C_p},\\
 V&=A_f\cdot A_g\cdot C_p,&\qquad\mathcal JV&=\chi\overline{A_f}\cdot\chi\overline{A_g}\cdot\chi\overline{C_p},
 \end{aligned}
\]
whose weights $F$ in \cref{lem:factors} are $F_U$, $F_U^J$, $F_V$ and $F_V^J$ respectively; and $h_{F^J}=h_F\circ J$ by \cref{lem:reflectedprofile}.

For $K=D_\phi$ with $\operatorname{supp}\phi\subset(0,\theta)$, $\theta<1$, the profiles $h_{F_\phi}=\phi(r)v$ and $h_{F_\phi}\circ J=\phi(3-r)v$ vanish on $1<r<2$, so both limits must be $0$. All coefficients of $D_\phi$ have $\log N/L<\theta$, at a fixed distance from $I_\psi$, so $\mathrm G_T(D_\phi,W_\psi)=O(T^{-A})$ by \cref{lem:localise}. For the reflected entry, $\mathrm G_T(\mathcal JD_\phi,W_\psi)=\int\Phi(t/T)\chi(s)^3\overline{D_\phi(s)W_\psi(s)}\dd t$, and $D_\phi W_\psi=\sum_Ne(N)N^{-s}$ with $e(N)\ll N^\epsilon$ supported on $N\le T^{\theta+\rho}$, where $\rho=\max I_\psi<2$. The phase of $\chi(s)^3N^{it}$ is $-3\vartheta(t)+t\log N$, whose derivative is at most $-(3-\theta-\rho)L+O(1)$, while its higher derivatives are $O(T^{1-j})$. Repeated integration by parts gives $O_A(T^{-A})$ for each $N$, hence in total.

All limits hold for every $\Phi\in\mathcal W$, so $\mathcal T_4$ is faithful.

\emph{The measure.} Up to null sets $\Om\setminus S_3=\mathcal B\setminus\mathcal A$ (see the proof of \Cref{thm:third}), so
\[
 \begin{split}
 \Om\setminus(S_3\cup\mathcal Z)=\big\{&1<r<2,\ x+y>\tfrac12,\ y+z<\tfrac32,\ z-x<\tfrac12,\\
 &3\min(x,\varrho)+5\max(x,\varrho)>2\big\}.
 \end{split}
\]
The last condition is implied by the others, since $3\min(x,\varrho)+5\max(x,\varrho)\ge4(x+\varrho)=4(1-(z-x))>2$. If $r<\frac32$ then $y+z=r-x<\frac32$ automatically, and if $r>\frac32$ then $x+y=r-z>\frac12$ automatically. Hence the set is the disjoint union, up to a null set, of $N_0$ and $JN_0=\{\frac32<r<2,\ y+z<\frac32,\ z-x<\frac12\}$, and $\mu(\Om\setminus(S_3\cup\mathcal Z))=2\mu(N_0)$ by \cref{lem:chamber}. Finally $\mu(N_0)=\frac{209}{1024}$: $N_0$ is a convex polytope, and integrating the polynomial $9!\,v^2$ exactly over a triangulation of it gives this value (numerical quadrature confirms it). \qed

\section{The balanced region}\label{sec:balanced}

In this section we add profiles supported on a region of shapes of trace just above $1$ whose two shorter coordinates have sum just above $\frac12$. Put $|a|=a_1+a_2+a_3$ for $a\in\R^3$, and let
\begin{gather*}
 \mathcal R=\{(x,y,z)\in\Om:\ 1<r<\tfrac{15}{14},\ 21k+7r<18\},\\
 \mathcal R^\sharp=\{a\in(0,1)^3:\ 1<|a|<\tfrac{15}{14},\ 21(a_1+a_2)+7|a|<18\}.
\end{gather*}
Thus $\mathcal R=\mathcal R^\sharp\cap\Om$, and $a_1+a_2<\frac{18-7}{21}=\frac{11}{21}$ on $\mathcal R^\sharp$. The sets $\mathcal R$ and $J\mathcal R$ are disjoint, since the trace is below $\frac{15}{14}$ on the first and above $\frac{27}{14}$ on the second. Write $\partial_j$ for the derivative in the $j$th coordinate and
\[
 \mathsf V=(\partial_2-\partial_1)(\partial_3-\partial_1)(\partial_3-\partial_2).
\]
For real $\psi\in C_c^\infty(\mathcal R)$ the function $\mathsf V\psi$ is real and smooth with support in $\operatorname{supp}\psi\subset\mathcal R^\sharp$, and we put
\[
 B_\psi=P_{\mathsf V\psi},\qquad I_\psi=\{x+y+z:\ (x,y,z)\in\operatorname{supp}\psi\}\subset(1,\tfrac{15}{14}),
\]
in the notation $P_c(s)=\sum_{n_1,n_2,n_3}c(\ell(n))(n_1n_2n_3)^{-s}$ of \Cref{sec:marked}. So $B_\psi$ is a product-type Dirichlet polynomial in three factors whose lengths $T^{a_1}\le T^{a_2}\le T^{a_3}$ are distributed according to $\mathsf V\psi$.

\begin{theorem}[Balanced row]\label{thm:row}
Let
\[
 \mathcal T_{\mathcal R}=\big\{(B_\psi,\psi),\ (\mathcal JB_\psi,\psi\circ J):\ \psi\in C_c^\infty(\mathcal R)\text{ real}\big\},\qquad\mathcal T_5=\mathcal T_4\cup\mathcal T_{\mathcal R}.
\]
Then $\mathcal T_5$ is faithful, $v\one_{S_3\cup\mathcal Z\cup\mathcal R\cup J\mathcal R}\in\mathcal H(\mathcal T_5)$, and
\begin{align*}
 \mu(S_3\cup\mathcal Z\cup\mathcal R\cup J\mathcal R)&=42-2\mu(N_0)+2\mu(N_0\cap\mathcal R)\\
 &=C_0:=\frac{2578241177999}{61983140352}=41.5958462794\ldots .
\end{align*}
\end{theorem}

The profile of $B_\psi$ is $\psi$ by a Jacobi--Trudi identity (\cref{lem:synthesis}). The analytic input is \cref{prop:balanced}: the pairing of any $P_c$ with $c$ supported in $\mathcal R^\sharp$ against an arbitrary smooth cubic source is given by the Schur inner product. To prove it we group the two shorter factors of $P_c$ into a modulus $q\asymp T^{a_1+a_2}$, write the time integral as a sum over $h=n-qm$, and replace each progression $n\equiv h\pmod q$ by the mean over its gcd class. This uses an estimate for $d_3$ in progressions with smooth weights (\cref{prop:smoothdisc}) deduced from Heath-Brown's work \cite{HB}, and it is here that the conditions defining $\mathcal R$ enter (\cref{lem:exponents}). The gcd-class means are then evaluated by a contour shift and Poisson summation in $h$. The result is a four-shift Euler product, whose limit we identify with the Schur inner product by a four-variable Haar identity (\cref{lem:haar4,lem:balkernel}).

\subsection{Synthesis}
\begin{lemma}[Synthesis]\label{lem:synthesis}
For $\psi\in C_c^\infty(\Om)$ we have $h_{\mathsf V\psi}=\psi$ on $\Om$.
\end{lemma}
\begin{proof}
Step~2 of the proof of \cref{lem:density} uses only that $\psi\in C_c^\infty(\Om)$; the box $Q$ there plays no role in that step. It shows that $h_\Psi=\psi$, where
\[
 \Psi=\frac16\sum_{\sigma\in S_3}\operatorname{sgn}(\sigma)\big((\iota-\sigma)\cdot\nabla\big)^3\hat\psi,\qquad\hat\psi(m_1,m_2,m_3)=\psi(m_3,m_2,m_1).
\]
For $t\in\R^3$, expanding a determinant,
\[
 \sum_{\sigma\in S_3}\operatorname{sgn}(\sigma)e^{(\iota-\sigma)\cdot t}=e^{\iota\cdot t}\det\big(e^{-jt_i}\big)_{i,j=1}^3=e^{\iota\cdot t-t_1-t_2-t_3}\prod_{1\le i<l\le3}\big(e^{-t_l}-e^{-t_i}\big).
\]
The left side has no terms of degree below three (as noted in Step~2), and the right side has cubic part $\prod_{i<l}(t_i-t_l)$. Comparing cubic parts gives $\frac16\sum_\sigma\operatorname{sgn}(\sigma)((\iota-\sigma)\cdot t)^3=\prod_{i<l}(t_i-t_l)$, so $\Psi=(\partial_{m_1}-\partial_{m_2})(\partial_{m_1}-\partial_{m_3})(\partial_{m_2}-\partial_{m_3})\hat\psi$. Since $\partial_{m_1},\partial_{m_2},\partial_{m_3}$ act on $\hat\psi$ as $\partial_3,\partial_2,\partial_1$ act on $\psi$, this says $\Psi(c_1,c_2,c_3)=(\mathsf V\psi)(c_3,c_2,c_1)$. Finally $h_{F\circ\pi}=h_F$ for every permutation $\pi$ of the coordinates (Step~3 of the proof of \cref{lem:density}), so $h_{\mathsf V\psi}=h_\Psi=\psi$.
\end{proof}
\noindent The identity has also been checked by numerical quadrature.

\subsection{The divisor function in progressions}
For $q\ge1$, a unit $a\bmod q$ and a box $\mathcal B=I_1\times I_2\times I_3$ of integer intervals put
\[
 \Delta_q(\mathcal B;a)=\#\{n\in\mathcal B:\ n_1n_2n_3\equiv a\ (q)\}-\frac1{\varphi(q)}\#\{n\in\mathcal B:\ (n_1n_2n_3,q)=1\}.
\]
For $U,V,W\ge1$ let $\Delta^*_q(U,V,W)$ be the supremum of $|\Delta_q(\mathcal B;a)|$ over units $a$ and boxes with $I_1\subseteq[1,4U]$, $I_2\subseteq[1,4V]$, $I_3\subseteq[1,4W]$. It is nondecreasing in each of $U,V,W$.

\begin{lemma}[Heath-Brown]\label{lem:HB}
Let $\epsilon>0$, $Y\ge2$ and $U,V,W\ge1$ with $UVW\le Y$. For $1\le q\le Y$,
\begin{enumerate}[label=\textup{(\roman*)}]
\item $\Delta^*_q(U,V,W)\ll_\epsilon Y^\epsilon\big(q^{5/6}+q^{1/4}(UV)^{1/2}+q^{1/2}W^{1/2}\big)$;
\item $\Delta^*_q(U,V,W)\ll_\epsilon Y^\epsilon q^{1/2}U$;
\end{enumerate}
and for $1\le Q\le Y$,
\begin{enumerate}[label=\textup{(\roman*)},resume]
\item $\sum_{Q<q\le2Q}\Delta^*_q(U,V,W)\ll_\epsilon Y^\epsilon Q^{3/4}U^{3/4}V^{1/2}(W+U)^{1/4}(V^2+Q)^{1/4}$.
\end{enumerate}
\end{lemma}
\begin{proof}[Source]
In \cite[Lemmas~5 and 6]{HB} the number of $n$ in a box with $n_1n_2n_3$ in a given unit class is compared with an approximant that does not depend on the class, with errors bounded by the right sides of (i) and (ii). The proofs use, for an interval of length $I$ and a frequency $0<|j|\le q/2$, only the bound $\min(I,q/|j|)$ for the corresponding geometric sum, so they apply to every box in the stated range, and the bounds increase with the side lengths. If $|S(a)-M|\le E$ for every unit $a$, then also $|S(a)-\varphi(q)^{-1}\sum_bS(b)|\le2E$; this replaces the approximant by the class mean. For (iii), the proof of \cite[Lemma~8]{HB}, through \cite[(6.2)--(6.5)]{HB}, bounds the discrepancy for a fixed modulus by quantities that depend on the box only through the factor $H(I_2,q)^{1/4}$, where $H(I,q)$ counts quadruples of units in $I$ with equal sums of reciprocal pairs modulo $q$. Since $H(I_2,q)\le H([1,4V],q)$, the supremum over classes and boxes may be taken inside the sum over $q$, and \cite[(6.14)]{HB} bounds $\sum_{Q<q\le2Q}H([1,4V],q)$ by $Y^\epsilon(V^4+QV^2)$. H\"older's inequality as in \cite[Lemma~8]{HB} gives (iii).
\end{proof}

\begin{proposition}[Smooth discrepancy]\label{prop:smoothdisc}
Let $Q\ge1$, $U,V,W>0$ and $h\in\Z\setminus\{0\}$. For each $Q<q\le2Q$ let $w_q\in C^\infty(\R^3)$ be supported in $[U,4U]\times[V,4V]\times[W,4W]$, with $|(y_1\partial_1)(y_2\partial_2)(y_3\partial_3)w_q(y)|\le M$ for all $y$. Put $g=(h,q)$ and
\[
 D_q(h;w)=\sum_{\substack{n\in\N^3\\ n_1n_2n_3\equiv h\ (q)}}w(n)-\frac1{\varphi(q/g)}\sum_{\substack{n\in\N^3\\(n_1n_2n_3,q)=g}}w(n).
\]
Write $U^+=\max(U,1)$ and similarly $V^+,W^+$, and let $Y=2+QU^+V^+W^+|h|$. Then
\[
 \sum_{Q<q\le2Q}|D_q(h;w_q)|\ll_\epsilon MY^\epsilon\,Q\min_\pi\min\big(E_1,E_2,E_3\big)\big(\pi(U^+,V^+,W^+)\big),
\]
where $\pi$ runs over the permutations of the three arguments and
\begin{gather*}
 E_1(U,V,W)=Q^{5/6}+Q^{1/4}(UV)^{1/2}+Q^{1/2}W^{1/2},\qquad E_2(U,V,W)=Q^{1/2}U,\\
 E_3(U,V,W)=U^{3/4}V^{1/2}(W+U)^{1/4}\big(1+V^{1/2}Q^{-1/4}\big).
\end{gather*}
\end{proposition}
\begin{proof}
\emph{(a) Unit classes.} Let $q_1\ge1$, let $a$ be a unit modulo $q_1$, and let $w$ be smooth with support in $[U',4U']\times[V',4V']\times[W',4W']$ and $|(y_1\partial_1)(y_2\partial_2)(y_3\partial_3)w|\le M$. Put
\[
 D^\circ_{q_1}(a;w)=\sum_{n_1n_2n_3\equiv a\ (q_1)}w(n)-\frac1{\varphi(q_1)}\sum_{(n_1n_2n_3,q_1)=1}w(n).
\]
Since $w$ vanishes with all its derivatives on the lower faces of the box, $w(y)=\int_{U'}^{y_1}\int_{V'}^{y_2}\int_{W'}^{y_3}\partial_1\partial_2\partial_3w(t)\dd t$ for $y$ in the box. Inserting this and interchanging the finite sum with the integral,
\[
 D^\circ_{q_1}(a;w)=\int_{[U',4U']\times[V',4V']\times[W',4W']}\partial_1\partial_2\partial_3w(t)\,\Delta_{q_1}(\mathcal B_t;a)\dd t,
\]
where $\mathcal B_t$ is the box of integer points $n$ with $t_j\le n_j\le4U'_j$, writing $(U'_1,U'_2,U'_3)=(U',V',W')$. Here $\mathcal B_t\subseteq[1,4U'^+]\times[1,4V'^+]\times[1,4W'^+]$ because $n_j\ge1$, and $|\partial_1\partial_2\partial_3w(t)|\le M/(t_1t_2t_3)\le M/(U'V'W')$. The box has volume $27U'V'W'$, so
\begin{equation}\label{eq:unitsmooth}
 |D^\circ_{q_1}(a;w)|\le27M\,\Delta^*_{q_1}(U'^+,V'^+,W'^+).
\end{equation}

\emph{(b) Reduction to units.} Fix $q$ with $Q<q\le2Q$ and put $g=(h,q)$, $q=gq_1$, $h=gh_1$, so that $(h_1,q_1)=1$. If $n_1n_2n_3\equiv h\pmod q$ or $(n_1n_2n_3,q)=g$, then $g\mid n_1n_2n_3$. For such $n$ put
\[
 e_1=(n_1,g),\qquad e_2=(n_2,g/e_1),\qquad e_3=g/(e_1e_2).
\]
For each prime one of $n_1/e_1$ and $g/e_1$ is not divisible by it, so $(n_1/e_1,g/e_1)=1$; as $g/e_1$ divides $(n_1/e_1)n_2n_3$, it divides $n_2n_3$. In the same way $(n_2/e_2,e_3)=1$ and $e_3\mid n_3$. The map $n\mapsto(e,n')$, $n'_j=n_j/e_j$, is a bijection from the $n$ with $g\mid n_1n_2n_3$ onto the pairs with $e_1e_2e_3=g$, $(n'_1,e_2e_3)=1$ and $(n'_2,e_3)=1$; the inverse is $n_j=e_jn'_j$, and indeed $(e_1n'_1,g)=e_1(n'_1,e_2e_3)=e_1$ and $(e_2n'_2,e_2e_3)=e_2$. In these variables $n_1n_2n_3\equiv h\pmod q$ becomes $n'_1n'_2n'_3\equiv h_1\pmod{q_1}$, and $(n_1n_2n_3,q)=g$ becomes $(n'_1n'_2n'_3,q_1)=1$; the $\varphi(q/g)=\varphi(q_1)$ residue classes modulo $q$ with gcd $g$ correspond to the units modulo $q_1$. Remove the conditions $(n'_1,e_2e_3)=1$ and $(n'_2,e_3)=1$ by M\"obius inversion over $t_1\mid e_2e_3$ and $t_2\mid e_3$, writing $n'_1=t_1n''_1$ and $n'_2=t_2n''_2$. If $(t_1t_2,q_1)>1$ both terms of the discrepancy vanish, because $h_1$ is a unit. Otherwise $t_1t_2n''_1n''_2n'_3\equiv h_1$ is the unit class $n''_1n''_2n'_3\equiv h_1\overline{t_1t_2}$, and the coprimality condition is unchanged. Hence
\[
 D_q(h;w_q)=\sum_{e_1e_2e_3=g}\ \sum_{\substack{t_1\mid e_2e_3,\ t_2\mid e_3\\(t_1t_2,q_1)=1}}\mu(t_1)\mu(t_2)\,D^\circ_{q_1}\big(h_1\overline{t_1t_2};w_q^{e,t}\big),
\]
with $w_q^{e,t}(y)=w_q(e_1t_1y_1,e_2t_2y_2,e_3y_3)$. This weight is supported in a box of the form in (a), with $U'=U/(e_1t_1)$, $V'=V/(e_2t_2)$, $W'=W/e_3$, and satisfies the same derivative bound because each $y_j\partial_j$ is invariant under dilation. Since $U'^+\le U^+$, $V'^+\le V^+$, $W'^+\le W^+$, \eqref{eq:unitsmooth} and monotonicity of $\Delta^*$ give $|D_q(h;w_q)|\le27M\,\tau_3(g)\tau(g)^2\,\Delta^*_{q_1}(U^+,V^+,W^+)$.

\emph{(c) Summation.} Each $q$ determines $g\mid h$, and $\tau_3(g)\tau(g)^2\ll Y^\epsilon$. Thus
\[
 \sum_{Q<q\le2Q}|D_q(h;w_q)|\ll MY^\epsilon\sum_{g\mid h}\ \sum_{Q/g<q_1\le2Q/g}\Delta^*_{q_1}(U^+,V^+,W^+).
\]
The term $q_1=1$ vanishes, since there is only one class modulo $1$, so we may assume $2Q/g\ge2$. By \cref{lem:HB}(i), (ii), summed over $q_1$, and (iii), the inner sum is $\ll Y^\epsilon(Q/g)E_i$ for $i=1,2,3$, evaluated at $Q/g$ in place of $Q$ and at any permutation of $(U^+,V^+,W^+)$. Each of $QE_1$, $QE_2$, $QE_3$ is nondecreasing in $Q$ (for the last, $QE_3\asymp Q^{3/4}U^{3/4}V^{1/2}(W+U)^{1/4}(V^2+Q)^{1/4}$), so we may replace $Q/g$ by $Q$. Summing over the $\tau(h)\ll Y^\epsilon$ divisors $g$ proves the proposition.
\end{proof}

\begin{remark}
For $h=0$ the discrepancy vanishes identically: the class $0\bmod q$ is the only class with gcd $q$, and $\varphi(1)=1$.
\end{remark}

\begin{lemma}[Exponents]\label{lem:exponents}
For $k\ge0$ and $a\in[0,\infty)^3$ put
\begin{gather*}
 \theta_1(k;a)=\max\big(\tfrac{11k}6,\ \tfrac{5k}4+\tfrac{a_1+a_2}2,\ \tfrac{3k}2+\tfrac{a_3}2\big),\qquad\theta_2(k;a)=\tfrac{3k}2+a_1,\\
 \theta_3(k;a)=\max\big(k+\tfrac{3a_1}4+\tfrac{a_2}2+\tfrac{\max(a_1,a_3)}4,\ \tfrac{3k}4+\tfrac{3a_1}4+a_2+\tfrac{\max(a_1,a_3)}4\big),
\end{gather*}
and $\theta(k;a)=\min_\pi\min_i\theta_i(k;\pi a)$ over permutations $\pi$ of the coordinates. Let
\[
 \mathcal O=\{(k,r):\ k\ge0,\ 1<r<\tfrac{15}{14},\ 21k+7r<18\},
\]
and let $K\subset\mathcal O$ be compact. There are $\eta,\delta_0>0$ such that $\theta(\max(k,\frac12);a)\le1-\eta$ whenever $(k,r)\in K$, $a\in[0,\infty)^3$ and $|a_1+a_2+a_3-r|\le\delta_0$.
\end{lemma}
For $Q=T^k$ and $(U,V,W)=(T^{a_1},T^{a_2},T^{a_3})$ we have $QE_i(U,V,W)\le3T^{\theta_i(k;a)}$ (for $i=1$ each of the three summands is at most $T^{\theta_1}$, and for $i=3$ the factor is at most $2\cdot2^{1/4}$), and each $\theta_i$ is nondecreasing in $k$.
\begin{proof}
If $k<\frac12$ then $(\frac12,r)\in\mathcal O$, because $\frac{21}2+7r<18$ is $r<\frac{15}{14}$. So it suffices to prove $\theta(k;a)<1$ when $(k,r)\in\mathcal O$, $k\ge\frac12$, $a\ge0$ and $|a|=r$. The lemma then follows because $(k,r,a)\mapsto\theta(\max(k,\frac12);a)$ is continuous, the set of $(k,r,a)$ with $(k,r)\in K$, $a\ge0$, $|a|=r$ is compact, and $\theta$ is uniformly continuous on a compact neighbourhood of it.

So let $k\ge\frac12$, $1<r<\frac{15}{14}$, $21k+7r<18$ and $|a|=r$. Put $P=r+\frac52k-2$ and $U_*=2-3k$. Substituting $k<(18-7r)/21$ and using $1<r<\frac{15}{14}$ one finds
\[
 k<\tfrac{11}{21},\qquad P<\tfrac r3,\qquad P<U_*,\qquad r+\tfrac{15}2k<5,\qquad 19k+6r<16 ;
\]
for instance $\frac23r+\frac52k<\frac{90-7r}{42}<2$ gives $P<\frac r3$, and $r+\frac{11}2k<\frac{198-35r}{42}<4$ gives $P<U_*$.

\emph{Case 1: some coordinate $w$ of $a$ lies in $(P,U_*)$.} Put it last. Then $\theta_1=\max\big(\frac{11k}6,\frac{5k}4+\frac{r-w}2,\frac{3k}2+\frac w2\big)$. Here $\frac{11k}6<\frac{121}{126}<1$, $\frac{5k}4+\frac{r-w}2<\frac{5k}4+\frac{r-P}2=1$ and $\frac{3k}2+\frac w2<\frac{3k}2+\frac{U_*}2=1$.

\emph{Case 2: every coordinate is at most $P$ or at least $U_*$.} Order them as $\alpha\le\beta\le\gamma$. Since $\gamma\ge\frac r3>P$, we have $\gamma\ge U_*$. If $\beta\ge U_*$, then $\alpha=r-\beta-\gamma\le r-2U_*$, and $\theta_2$ with $\alpha$ in the first slot is $\frac{3k}2+\alpha\le r+\frac{15}2k-4<1$. Otherwise $\alpha\le\beta\le P$, and we use $\theta_3$ with the slots $(\beta,\alpha,\gamma)$. As $\gamma\ge\beta$, its first term is
\[
 k+\tfrac{3\beta}4+\tfrac\alpha2+\tfrac\gamma4=k+\tfrac r4+\tfrac{\alpha+\beta}4+\tfrac\beta4\le k+\tfrac r4+\tfrac{r-U_*}4+\tfrac P4=\tfrac{19k+6r}8-1<1,
\]
using $\alpha+\beta=r-\gamma\le r-U_*$ and $\beta\le P$. Its second term is
\[
 \tfrac{3k}4+\tfrac{3\beta}4+\alpha+\tfrac\gamma4=\tfrac{3k}4+\tfrac r4+\tfrac\beta2+\tfrac{3\alpha}4\le\tfrac{3k}4+\tfrac r4+\tfrac58(\alpha+\beta)\le\tfrac{21k+7r}8-\tfrac54<1,
\]
using $\alpha\le\beta$ and again $\alpha+\beta\le r-U_*=r-2+3k$.
\end{proof}
\noindent The lemma has also been checked numerically on a grid.

\subsection{The balanced pairing}
\begin{lemma}[Localisation]\label{lem:localiseP}
Let $c\in C_c^\infty(\R^3)$ and $I_c=\{|a|:\ a\in[0,\infty)^3\cap\operatorname{supp}c\}$. Then \cref{lem:localise} holds with $W_\psi$ and $I_\psi$ replaced by $P_c$ and $I_c$.
\end{lemma}
\begin{proof}
The coefficient of $P_c$ at $n$ is $\sum_{n_1n_2n_3=n}c(\ell(n))\ll\tau_3(n)$, and it vanishes unless $\log n/L\in I_c$. These are the only properties of $M_\psi$ used in the proof of \cref{lem:localise}.
\end{proof}

\begin{proposition}[Balanced pairing]\label{prop:balanced}
Let $c\in C_c^\infty(\R^3)$ be real with $\operatorname{supp}c\subset\mathcal R^\sharp$, and let $b\in C_c^\infty(\R^3)$. Uniformly for $\Phi$ in bounded subsets of $C_c^\infty$ with a fixed compact support in $(1,2)$,
\[
 \ip{P_b}{P_c}_\Phi=a_3TL^9\Big(\int\Phi\Big)\int_\Om h_bh_c+o(TL^9),
\]
where $h_b$ is the profile of the restriction of $b$ to $[0,1]^3$.
\end{proposition}
Here $b$ may be complex-valued. \Cref{lem:profiles} and its proof hold verbatim for complex-valued continuous $F$, with $c_\lambda(F)$ complex; in part (iii) the first limit is then $\int_\Om h_F\overline{h_G}$.

\emph{Reduction to elementary sources.} Fix $c$; implied constants may depend on $c$ and on the support of $\Phi$. Let $\varphi_1,\phi\in C_c^\infty((1,2))$ with $\varphi_1=1$ on a neighbourhood of $I_c$ and $\phi=1$ on a neighbourhood of $\operatorname{supp}\varphi_1$. Put $b_1(a)=\varphi_1(|a|)b(a)$. The coefficients of $P_b-P_{b_1}$ vanish at every $N$ with $\log N/L$ in a neighbourhood of $I_c$, so $\ip{P_b-P_{b_1}}{P_c}_\Phi=O(T^{-A})$ by \cref{lem:localiseP}. On $[0,\infty)^3\cap\operatorname{supp}b_1$ we have $\phi(|a|)=1$, so \eqref{eq:sourcesuperposition} gives $P_{b_1}=(2\pi)^{-3}\int\widehat{b_1}(\xi)P_{-i\xi,\phi}\dd\xi$. For $\boldsymbol u\in(i\R)^3$, Cauchy--Schwarz with respect to a smooth non-negative majorant of $|\Phi|$, \cref{lem:sourcenorm} and \cref{lem:sixth}(ii) give
\[
 |\ip{P_{\boldsymbol u,\phi}}{P_c}_\Phi|\ll(1+|\boldsymbol u|)^3TL^9 .
\]
We shall show that for every $\boldsymbol u\in(i\R)^3$ with distinct coordinates
\begin{equation}\label{eq:balelementary}
 \ip{P_{\boldsymbol u,\phi}}{P_c}_\Phi=a_3TL^9\Big(\int\Phi\Big)\int_\Om h_{e^{-\boldsymbol u\cdot a}}h_c+o(TL^9).
\end{equation}
Granting this, dominated convergence in $\xi$ (the integrand is $\ll|\widehat{b_1}(\xi)|(1+|\xi|)^3TL^9$, and almost every $\xi$ has distinct coordinates) and Fubini's theorem in \eqref{eq:profile}, with $(2\pi)^{-3}\int\widehat{b_1}(\xi)e^{i\xi\cdot a}\dd\xi=b_1(a)$, give $\ip{P_{b_1}}{P_c}_\Phi=a_3TL^9(\int\Phi)\int_\Om h_{b_1}h_c+o(TL^9)$. Finally $h_c$ vanishes unless $r\in I_c$, because $\kappa(x,y,z;a)=0$ unless $|a|=r$, and there $h_{b_1}=\varphi_1(r)h_b=h_b$. This proves \cref{prop:balanced}, given \eqref{eq:balelementary}.

The rest of this subsection proves \eqref{eq:balelementary}. Fix $\boldsymbol u\in(i\R)^3$ with distinct coordinates, put $\boldsymbol\alpha=\boldsymbol u/L$ and $r_0=1+\max_k|u_k|$, and let
\begin{gather*}
 \rho_-=\min I_c>1,\qquad\rho_+=\max I_c<\tfrac{15}{14},\\
 \kappa_+=\max\{a_1+a_2:\ a\in\operatorname{supp}c\}<\tfrac{11}{21}.
\end{gather*}
For $q\ge1$ and $a_3\in\R$ put $\gamma_q(a_3)=\sum_{d_1d_2=q}c(\log d_1/L,\log d_2/L,a_3)$, so that $P_c(s)=\sum_{q,m}\gamma_q(\log m/L)(qm)^{-s}$. Each $\gamma_q$ is real and smooth, $\partial^j\gamma_q\ll_j\tau(q)$, and $\gamma_q(a_3)\ne0$ implies $\log q/L\le\kappa_+$ and $\log q/L+a_3\in I_c$. With $\widehat\Phi(\xi)=\int\Phi(x)e(-x\xi)\dd x$ let
\[
 W_{q,h}(y)=\frac1y\,\gamma_q\Big(\frac{\log(y/q)}L\Big)\,\widehat\Phi\Big(\frac{Th}{2\pi y}\Big)\qquad(q\ge1,\ h\in\Z,\ y>0),
\]
a smooth function supported in $[T^{\rho_-},T^{\rho_+}]$, and for $g\mid q$ let $Z_{q,g}(s)=\sum_{(n,q)=g}b_{\boldsymbol u}(n)n^{-s}$. We write $\mathcal I=\ip{P_{\boldsymbol u,\phi}}{P_c}_\Phi$, and $\eta>0$ denotes a small constant depending on $c$, which may decrease from one lemma to the next.

\begin{lemma}[Truncation]\label{lem:baltrunc}
$\displaystyle\mathcal I=T\sum_{q\ge1}\sum_{h\in\Z}S_{q,h}+O(T^{1-\eta})$, where $S_{q,h}=\sum_{n\equiv h\ (q)}b_{\boldsymbol u}(n)W_{q,h}(n)$.
\end{lemma}
\begin{proof}
Expanding both polynomials and integrating in $t$,
\[
 \mathcal I=T\sum_n\sum_{q,m}b_{\boldsymbol u}(n)\phi\Big(\frac{\log n}L\Big)\gamma_q\Big(\frac{\log m}L\Big)(nqm)^{-1/2}\,\widehat\Phi\Big(\frac T{2\pi}\log\frac n{qm}\Big),
\]
a finite sum with $n\le T^2$ and $qm\le T^{\rho_+}$. Fix $\epsilon>0$. Since $\widehat\Phi(\xi)\ll_A|\xi|^{-A}$, the terms with $|\log(n/qm)|>T^{-1+\epsilon}$ contribute $O(T^{-A})$. In the remaining terms put $h=n-qm$, so $|h|\le2nT^{-1+\epsilon}$; as $\log(qm)/L\in I_c$ we have $\phi(\log n/L)=1$. Then $(nqm)^{-1/2}=n^{-1}(1+O(T^{-1+\epsilon}))$, $\frac T{2\pi}\log\frac n{qm}=\frac{Th}{2\pi n}+O(T^{-1+2\epsilon})$ and $\gamma_q(\log m/L)=\gamma_q(\log(n/q)/L)+O(\tau(q)T^{-1+\epsilon})$. For each $n\ll T^{\rho_+}$ and each $h$ there are at most $\tau(n-h)\ll T^\epsilon$ moduli $q\mid n-h$, so these replacements cost
\[
 \ll T\cdot T^{-1+3\epsilon}\sum_{n\ll T^{\rho_+}}\frac{\tau_3(n)}n\cdot nT^{-1+\epsilon}\ll T^{\rho_+-1+5\epsilon}.
\]
Since $q\le T^{\kappa_+}<n/2$, every pair $(q,h)$ with $q\mid n-h$ and $|h|\le2nT^{-1+\epsilon}$ comes from a unique $m=(n-h)/q\ge1$. Finally, adding the terms with $q\mid n-h$ and $|h|>2nT^{-1+\epsilon}$ costs $O(T^{-A})$, by the decay of $\widehat\Phi$; in the new terms $\gamma_q(\log(n/q)/L)\ne0$ still forces $\log n/L\in I_c$. As $\rho_+-1<1$, this proves the lemma after rearranging the absolutely convergent sum.
\end{proof}

\begin{lemma}[Equidistribution]\label{lem:balequi}
$\displaystyle T\sum_q\sum_hS_{q,h}=T\sum_q\sum_h\mathfrak m_{q,h}+O(T^{1-\eta})$, where, with $g=(h,q)$,
\[
 \mathfrak m_{q,h}=\frac1{\varphi(q/g)}\sum_{(n,q)=g}b_{\boldsymbol u}(n)W_{q,h}(n).
\]
\end{lemma}
\begin{proof}
For $h=0$, $S_{q,0}=\mathfrak m_{q,0}$. Let $h\ne0$. Fix $V_0\in C_c^\infty([\frac12,2])$ with $\sum_{i\in\Z}V_0(y/2^i)=1$ for $y>0$; for integers $n_j\ge1$ only $i\ge-1$ contribute. Then $S_{q,h}-\mathfrak m_{q,h}=\sum_{U,V,W}D_q(h;w^{UVW}_{q,h})$, in the notation of \cref{prop:smoothdisc}, where $U,V,W$ run over $\{2^i:\ i\ge-1\}$ and
\[
 w^{UVW}_{q,h}(y)=V_0\Big(\frac{y_1}U\Big)V_0\Big(\frac{y_2}V\Big)V_0\Big(\frac{y_3}W\Big)\prod_{j=1}^3y_j^{-u_j/L}\;W_{q,h}(y_1y_2y_3),
\]
which is supported in $[\frac U2,2U]\times[\frac V2,2V]\times[\frac W2,2W]$. As $W_{q,h}$ is supported in $[T^{\rho_-},T^{\rho_+}]$, only $O(L^3)$ triples contribute, each with $X:=UVW\in[T^{\rho_-}/8,8T^{\rho_+}]$; and only $q\le T^{\kappa_+}$ contribute, in $O(L)$ dyadic ranges $Q<q\le2Q$.

Since $(y_j\partial_j)F(y_1y_2y_3)=((y\partial_y)F)(y_1y_2y_3)$, the Leibniz rule gives $|(y_1\partial_1)(y_2\partial_2)(y_3\partial_3)w^{UVW}_{q,h}|\ll(1+|\boldsymbol u|)^3\max_{i\le3}\sup_{y\asymp X}|(y\partial_y)^iW_{q,h}(y)|$. Now $(y\partial_y)^j\widehat\Phi(\lambda/y)$ is a combination of the functions $(\lambda/y)^l\widehat\Phi^{(l)}(\lambda/y)$, $l\le j$, and $\widehat\Phi$ is a Schwartz function; hence, with $H=X/T$, \cref{prop:smoothdisc} applies with
\[
 M\ll_B(1+|\boldsymbol u|)^3T^\epsilon X^{-1}(1+|h|/H)^{-B}.
\]
It gives $\sum_{Q<q\le2Q}|D_q(h;w^{UVW}_{q,h})|\ll(1+|\boldsymbol u|)^3T^{2\epsilon}X^{-1}(1+|h|/H)^{-B}(2+|h|)^\epsilon\,Q\min_\pi\min_iE_i$. As $H\ge T^{\rho_--1}/8\ge1$, summing over $h\ne0$ gives a factor $\ll H^{1+\epsilon}$, so the contribution of one dyadic choice of $(Q,U,V,W)$ to $T\sum_q\sum_h(S_{q,h}-\mathfrak m_{q,h})$ is
\[
 \ll(1+|\boldsymbol u|)^3T^{3\epsilon}\,Q\min_\pi\min_iE_i(\pi(U^+,V^+,W^+))\le3(1+|\boldsymbol u|)^3T^{3\epsilon}T^{\theta(k;a)},
\]
where $Q=T^k$ and $(U^+,V^+,W^+)=(T^{a_1},T^{a_2},T^{a_3})$. If this choice contributes at all, there are $Q<q\le2Q$ and $n\asymp X$ with $\gamma_q(\log(n/q)/L)\ne0$, so $(\log q/L,\log n/L)=(a'_1+a'_2,|a'|)$ for some $a'\in\operatorname{supp}c$. Hence $(k,\log X/L)$ lies within $O(1/L)$ of the compact set $K_c=\{(a_1+a_2,|a|):\ a\in\operatorname{supp}c\}\subset\mathcal O$; and $|a_1+a_2+a_3-\log X/L|\le3\log2/L$, because $U\le U^+\le2U$. Applying \cref{lem:exponents} with $K$ a compact neighbourhood of $K_c$ in $\mathcal O$, and using that $\theta$ is nondecreasing in $k$, we get $\theta(k;a)\le1-\eta_0$ for $L$ large. Summing over the $O(L^4)$ dyadic choices and taking $\epsilon<\eta_0/4$ proves the lemma.
\end{proof}

\begin{lemma}[Polar part]\label{lem:balpolar}
Let $\mathcal C$ be the positively oriented circle $|s-1|=r_0/L$, and $\widetilde W(s)=\int_0^\infty W(y)y^{s-1}\dd y$. Then $T\sum_q\sum_h\mathfrak m_{q,h}=\mathcal M+O(T^{1-\eta})$, where
\[
 \mathcal M=T\sum_q\sum_h\frac1{\varphi(q/g)}\cdot\frac1{2\pi i}\oint_{\mathcal C}\widetilde W_{q,h}(s)Z_{q,g}(s)\dd s .
\]
\end{lemma}
\begin{proof}
For $p\mid q$ let $J_p=\{v_p(g)\}$ if $v_p(g)<v_p(q)$ and $J_p=\{j:\ j\ge v_p(q)\}$ if $v_p(g)=v_p(q)$; then $(n,q)=g$ if and only if $v_p(n)\in J_p$ for every $p\mid q$. Hence
\begin{gather*}
 Z_{q,g}(s)=\prod_{k=1}^3\zeta(s+\alpha_k)\,E_{q,g}(s),\\
 E_{q,g}(s)=\prod_{p\mid q}\Big[\prod_{k=1}^3\big(1-p^{-s-\alpha_k}\big)\sum_{j\in J_p}b_{\boldsymbol u}(p^j)p^{-js}\Big].
\end{gather*}
Each factor of $E_{q,g}$ is entire: if $J_p$ is infinite the sum equals $\prod_k(1-p^{-s-\alpha_k})^{-1}$ minus a finite sum. As $|b_{\boldsymbol u}(p^j)|\le\binom{j+2}2$, for $\Re s\ge\frac34$ each factor is at most $C(v_p(g)+1)^2p^{-3v_p(g)/4}$, so $|E_{q,g}(s)|\ll_\epsilon q^\epsilon g^{-3/4}$.

By Mellin inversion $\sum_{(n,q)=g}b_{\boldsymbol u}(n)W_{q,h}(n)=\frac1{2\pi i}\int_{(2)}\widetilde W_{q,h}(s)Z_{q,g}(s)\dd s$. Write $W_{q,h}(y)=y^{-1}w(y)$. For $j\ge0$, integration by parts gives $\widetilde W_{q,h}(s)=(1-s)^{-j}\int(y\partial_y)^jw(y)\,y^{s-2}\dd y$, and $|(y\partial_y)^jw(y)|\ll_{j,B}\tau(q)(1+|h|/H_+)^{-B}$ with $H_+=T^{\rho_+-1}$, because $T|h|/y\ge|h|/H_+$ on the support. Taking $j=0$ for $|s-1|\le\frac18$ and $j=2$ otherwise, and using $|s|\le9|s-1|$ when $|s-1|\ge\frac18$, we get for $\frac34\le\Re s\le2$
\begin{equation}\label{eq:Wtildebound}
 |\widetilde W_{q,h}(s)|\ll_B\tau(q)(1+|s|)^{-2}(1+|h|/H_+)^{-B}\int_{T^{\rho_-}}^{T^{\rho_+}}y^{\Re s-2}\dd y .
\end{equation}
The integral is at most $4T^{-\rho_-/4}$ on $\Re s=\frac34$ and at most $T^{\rho_+}$ for $\frac34\le\Re s\le2$. In particular $|\widetilde W_{q,h}(\frac34+it)|\ll\tau(q)(1+|t|)^{-2}(1+|h|/H_+)^{-B}T^{-\rho_-/4}$, while $Z_{q,g}(s)$ grows at most polynomially in $|\Im s|$, like $|\Im s|^{3/4+\epsilon}$, uniformly in $\frac34\le\Re s\le2$ away from the poles (convexity bound and $|E_{q,g}|\ll q^\epsilon$). So we may move the line of integration to $\Re s=\frac34$. The poles crossed are at $s=1-\alpha_k$, inside $\mathcal C$ as $|\alpha_k|<r_0/L$, and $\mathcal C$ lies in $\Re s>\frac34$; so the residues sum to the integral over $\mathcal C$. By the convexity bound $\prod_k|\zeta(\frac34+it+\alpha_k)|\ll(1+|t|)^{1/2}(1+|\boldsymbol u|)$. The integrals on $\Re s=\frac34$ therefore contribute to $T\sum_q\sum_h\mathfrak m_{q,h}$ at most
\[
 \ll(1+|\boldsymbol u|)T^{1-\rho_-/4+\epsilon}\sum_{q\le T^{\kappa_+}}\sum_{g\mid q}\frac{g^{-3/4}}{\varphi(q/g)}\sum_{(h,q)=g}(1+|h|/H_+)^{-B}.
\]
The innermost sum is $\ll1+H_+/g$, and $\varphi(q/g)^{-1}\ll T^\epsilon g/q$, so the sum over $g\mid q$ is $\ll T^{2\epsilon}(q^{-3/4}+H_+q^{-1})$ and the sum over $q$ is $\ll T^{3\epsilon}(T^{\kappa_+/4}+H_+)$. The total is
\[
 \ll(1+|\boldsymbol u|)T^{5\epsilon}\big(T^{1-(\rho_--\kappa_+)/4}+T^{\rho_+-\rho_-/4}\big),
\]
and both exponents are below $1$, because $\kappa_+<1<\rho_-$ and $\rho_+-\frac{\rho_-}4<\frac{15}{14}-\frac14$.
\end{proof}

We now evaluate $\mathcal M$. For a prime $p$ and integers $a,j,m\ge0$ put
\[
 f_p(a,j,m)=\begin{cases}1,&a\le j+m,\\-1/(p-1),&a=j+m+1,\\0,&\text{otherwise},\end{cases}
\]
\[
 \varepsilon_p(a,j;s)=\prod_{k=1}^3\big(1-p^{-s-\alpha_k}\big)\sum_{m\ge0}d_{\boldsymbol\alpha}(p^m)p^{-ms}f_p(a,j,m),
\]
where $d_{\boldsymbol\alpha}(n)=\sum_{n_1n_2n_3=n}\prod_kn_k^{-\alpha_k}=b_{\boldsymbol u}(n)$. Note that $\varepsilon_p(a,j;s)=1$ when $a\le j$. For $\boldsymbol v=(v_1,v_2)$ put $\sigma_{\boldsymbol v}(n)=\sum_{d_1d_2=n}d_1^{-v_1}d_2^{-v_2}$, and for $(\omega,\boldsymbol v,w)\in\C^4$ let
\[
 \mathcal S_p=\sum_{a,j\ge0}\varepsilon_p(a,j;1+\omega)\,\sigma_{\boldsymbol v}(p^a)\,p^{a\omega-j(1+\omega+w)},\qquad\mathcal S=\prod_p\mathcal S_p .
\]

\begin{lemma}[Euler product]\label{lem:baleuler}
Suppose that $\omega,v_1,v_2,w$ and $\alpha_1,\alpha_2,\alpha_3$ have real parts in $[-\frac1{20},\frac1{20}]$.
\begin{enumerate}[label=\textup{(\roman*)}]
\item The series $\mathcal S_p$ converges absolutely, and
\[
 \mathcal S=\frac{\zeta(1+\omega+w)\prod_{i=1}^2\zeta(1+v_i+w)\prod_{i=1}^2\prod_{k=1}^3\zeta(1+v_i+\alpha_k)}{\prod_{i=1}^2\zeta(1-\omega+v_i)}\,\mathfrak S,
\]
where $\mathfrak S=\prod_p\mathfrak S_p$ with $\mathfrak S_p=1+O(p^{-3/2})$ uniformly; so $\mathfrak S$ is holomorphic and bounded, and the identity holds wherever the left side converges.
\item Let $0<\delta\le\frac1{20}$ and suppose in addition that $\Re(v_i-\omega)$, $\Re(\omega+w)$, $\Re(v_i+w)$ and $\Re(v_i+\alpha_k)$ are at least $\delta$ for all $i,k$. Then
\[
 \sum_{q,J\ge1}\Big|\sigma_{\boldsymbol v}(q)\,q^{\omega}J^{-1-\omega-w}\prod_{p}\varepsilon_p\big(v_p(q),v_p(J);1+\omega\big)\Big|\ll\delta^{-11},
\]
and the sum without absolute values equals $\mathcal S$.
\item $\mathfrak S(0)=a_3$.
\end{enumerate}
\end{lemma}
\begin{proof}
Expand $\varepsilon_p(a,j;1+\omega)\sigma_{\boldsymbol v}(p^a)p^{a\omega-j(1+\omega+w)}$, using $(p-1)^{-1}=\sum_{l\ge1}p^{-l}$, as a sum of monomials $p^{-d}p^{-\ell}$, where $d\ge0$ is an integer, called the degree, and $\ell$ is a linear form in the shifts. The term $(a,j)$ only produces monomials of degree at least $\max(a,j)$: for $a\le j$ this is clear, and for $a>j$ the factor $f_p$ forces $m\ge a-j-1$, with the extra factor $(p-1)^{-1}$ if $m=a-j-1$. Each monomial of degree $d$ has $|\Re\ell|\le\frac{4d}{20}$, since $a\le d$ and $j+m\le d$, We expand completely before collecting terms. Then every monomial carries a coefficient of modulus at most an absolute constant: a product of signs from $\mu$ and of the finitely many coefficients of the numerator of $f_p$, all other expansions (the geometric series for $(p-1)^{-1}$, the divisor-type sums $\sigma_{\boldsymbol v}$ and $d_{\boldsymbol\alpha}$) having coefficients $1$. A monomial of degree $d$ is determined by finitely many choices in the numerator together with a composition of $d$ into a bounded number of non-negative exponents, so there are $O((1+d)^{C})$ of them. Hence the monomials of degree $d$ contribute $\ll(1+d)^Cp^{-4d/5}$ in absolute value, and $\mathcal S_p$ converges absolutely. The monomials of degree $0$ give $1$, and those of degree $1$ come from $(a,j)\in\{(0,1),(1,1),(1,0)\}$:
\[
 F_p=p^{-1-\omega-w}+\sum_ip^{-1-v_i-w}+\sum_{i,k}p^{-1-v_i-\alpha_k}-\sum_ip^{-1+\omega-v_i};
\]
here $\varepsilon_p(1,0;1+\omega)=\sum_kp^{-1-\omega-\alpha_k}-p^{-1}+(\text{degree}\ge2)$ and $\sigma_{\boldsymbol v}(p)p^\omega=\sum_ip^{\omega-v_i}$. So $\mathcal S_p=1+F_p+O(p^{-8/5})$, while the local factor of the displayed quotient of zeta functions is $1+F_p+O(p^{-8/5})$ as well. Hence $\mathfrak S_p$, defined as $\mathcal S_p$ divided by that local factor, is $1+O(p^{-8/5})$, which proves (i).

For (ii), the summand is multiplicative in the pair $(q,J)$: it is the product over all primes $p$ of the term $(a,j)=(v_p(q),v_p(J))$ of $\mathcal S_p$, and the term $(0,0)$ equals $1$. So the sum of absolute values is at most $\prod_p(1+\sum_{\text{monomials}\ne1}|p^{-d-\ell}|)$. The monomials of degree $1$ are the eleven terms of $F_p$, whose absolute values are $p^{-1-\Re(\omega+w)}$, $p^{-1-\Re(v_i+w)}$, $p^{-1-\Re(v_i+\alpha_k)}$ and $p^{-1-\Re(v_i-\omega)}$, all at most $p^{-1-\delta}$. The monomials of degree $d\ge2$ contribute $\ll(1+d)^Cp^{-4d/5}$, whose sum over $d\ge2$ and $p$ converges. So the product is $\ll\exp(11\sum_pp^{-1-\delta})\ll\delta^{-11}$, since $\sum_pp^{-1-\delta}\le\log(1/\delta)+O(1)$. Absolute convergence allows the multiplicative rearrangement, so the sum is $\mathcal S$.

For (iii), at zero shifts $\sigma_{\boldsymbol 0}(p^a)=a+1$ and $d_{\boldsymbol 0}(p^m)=\binom{m+2}2$. With $X=1/p$ and $r=j+m$, since $\sum_{m\le r}\binom{m+2}2=\binom{r+3}3$,
\[
 (1-X)^{-3}\mathcal S_p(0)=\sum_{r\ge0}\binom{r+3}3X^r\Big[\binom{r+2}2-\frac X{1-X}(r+2)\Big]=\frac{1+4X+X^2}{(1-X)^6},
\]
as in the proof of \cref{prop:unequal}. Removing the local factors of the seven net zeta factors, $\mathfrak S_p(0)=(1-X)^7\mathcal S_p(0)=(1-X)^4(1+4X+X^2)$, which is the local factor of $a_3$ by the computation in the proof of \cref{lem:fouredge}.
\end{proof}
\noindent The zero-shift identity in (iii) has also been checked by computer as an identity of power series in $X$ to order $16$, and the first-order statement in (i) exactly: as polynomials in $p^{-1}$ whose coefficients are Laurent polynomials in $p^{-\omega}$, $p^{-v_i}$, $p^{-w}$ and $p^{-\alpha_k}$, the local factors $\mathcal S_p$ and that of the zeta quotient have the same terms of degree $0$ and $1$ (the eleven terms of $F_p$). It also evaluates their ratio at $p=101$ and $p=409$ at one point with all real parts in $[-\frac1{20},\frac1{20}]$, in floating point.

\begin{lemma}[Poisson summation]\label{lem:balpoisson}
For $\beta\in\C^3$ let
\[
 \mathcal Lc(\beta)=\int_{\R^3}c(a)e^{\beta_1a_1+\beta_2a_2+\beta_3(1-a_3)}\dd a,
\]
and let $\widetilde\Phi(s)=\int\Phi(x)x^{s-1}\dd x$. Then
\begin{multline*}
 \mathcal M=\frac T{(2\pi i)^4}\oint_{|\omega|=r_0/L}\int_{(2r_0)^3}\mathcal Lc(\beta)e^{-\beta_3}\Big(\frac T{2\pi}\Big)^{\omega+w}\\
 \times\widetilde\Phi(1+\omega+w)\prod_{k=1}^3\zeta(1+\omega+\alpha_k)\,\mathcal S\dd\beta\dd\omega,
\end{multline*}
where $\int_{(2r_0)^3}$ denotes integration over $\Re\beta_1=\Re\beta_2=\Re\beta_3=2r_0$, $\mathcal S=\mathcal S(\omega;\boldsymbol v,w)$ with $\boldsymbol v=(\beta_1,\beta_2)/L$ and $w=\beta_3/L$, and the integral converges absolutely.
\end{lemma}
\begin{proof}
\emph{Poisson summation in $h$.} The indicator of $(h,q)=g$ is $q$-periodic in $h$, and $\sum_{h\bmod q,\,(h,q)=g}e(-hJ/q)=c_{q/g}(J)$; hence $\one_{(h,q)=g}=q^{-1}\sum_{J\bmod q}c_{q/g}(J)e(hJ/q)$. For $y>0$ let $F(x)=\widehat\Phi(Tx/(2\pi y))$, so that $\widehat F(\xi)=\frac{2\pi y}T\Phi(-\frac{2\pi y\xi}T)$. By Poisson summation, and since $c_{q/g}$ is even and $q$-periodic,
\[
 \sum_{(h,q)=g}\widehat\Phi\Big(\frac{Th}{2\pi y}\Big)=\frac1q\sum_{J\in\Z}c_{q/g}(J)\widehat F\Big(\frac Jq\Big)=\frac{2\pi y}{qT}\sum_{J\ge1}c_{q/g}(J)\,\Phi\Big(\frac{2\pi yJ}{qT}\Big),
\]
the terms with $J\le0$ vanishing because $\Phi$ is supported in $(1,2)$. The sum over $h$ in $\mathcal M$ is absolutely convergent uniformly for $s\in\mathcal C$ and $y$ in the support, so we may perform it inside the integrals. This gives
\[
 \mathcal M=\sum_q\sum_{J\ge1}\frac{2\pi}q\cdot\frac1{2\pi i}\oint_{\mathcal C}\int_0^\infty y^{s-1}\gamma_q\Big(\frac{\log(y/q)}L\Big)\Phi\Big(\frac{2\pi yJ}{qT}\Big)\dd y\;\mathcal Z_{q,J}(s)\dd s,
\]
where
\[
 \mathcal Z_{q,J}(s)=\sum_{g\mid q}\frac{c_{q/g}(J)}{\varphi(q/g)}Z_{q,g}(s),
\]
with only finitely many $(q,J)$, because $\Phi(2\pi yJ/(qT))\ne0$ forces $1<2\pi yJ/(qT)<2$, so $J<qT/(\pi y)<2qT^{1-\rho_-}$.

\emph{The Ramanujan factor.} By H\"older's evaluation
\[
 c_l(J)=\mu\big(l/(l,J)\big)\varphi(l)/\varphi\big(l/(l,J)\big)
\]
\cite[\S8.4]{Apostol}. If $(n,q)=g$ and $l=q/g$, then $(q,Jn)=g(l,J)$, because $(l,n/g)=1$; so $c_{q/g}(J)/\varphi(q/g)=f(q,Jn)$ with $f(q,N)=\mu(q/(q,N))/\varphi(q/(q,N))$. Hence, for $\Re s>1$ and then by continuation,
\[
 \mathcal Z_{q,J}(s)=\sum_{n\ge1}b_{\boldsymbol u}(n)n^{-s}f(q,Jn)=\prod_{k=1}^3\zeta(s+\alpha_k)\prod_p\varepsilon_p\big(v_p(q),v_p(J);s\big),
\]
because $f(q,Jn)=\prod_pf_p(v_p(q),v_p(J),v_p(n))$; the product over $p$ is finite, as $\varepsilon_p(0,j;s)=1$.

\emph{Change of variables.} Substitute $y=qTx/(2\pi J)$ and $s=1+\omega$. Then $\frac{2\pi}qy^{s-1}\dd y=T(\frac T{2\pi})^\omega q^\omega J^{-1-\omega}x^\omega\dd x$, and
\begin{multline*}
 \mathcal M=\frac T{2\pi i}\oint_{|\omega|=r_0/L}\Big(\frac T{2\pi}\Big)^\omega\sum_{q,J}q^\omega J^{-1-\omega}\mathcal Z_{q,J}(1+\omega)\\
 \times\int x^\omega\Phi(x)\gamma_q\Big(\frac{\log(Tx/(2\pi J))}L\Big)\dd x\dd\omega .
\end{multline*}

\emph{Laplace inversion.} Since $c\in C_c^\infty$, $\mathcal Lc$ is entire and $|\mathcal Lc(\beta)|\ll_N(1+|\beta|)^{-N}$ on $\Re\beta_i=2r_0$, and $c(a)=(2\pi i)^{-3}\int_{(2r_0)^3}\mathcal Lc(\beta)e^{-\beta_1a_1-\beta_2a_2-\beta_3(1-a_3)}\dd\beta$. With $a_3=\log(Tx/(2\pi J))/L$ this gives
\[
 \gamma_q\Big(\frac{\log(Tx/(2\pi J))}L\Big)=\frac1{(2\pi i)^3}\int_{(2r_0)^3}\mathcal Lc(\beta)\,\sigma_{\boldsymbol v}(q)\,e^{-\beta_3}\Big(\frac{Tx}{2\pi J}\Big)^{w}\dd\beta .
\]
Insert this and move the sums over $q,J$ inside the $\beta$-integral. On the contours $\Re\omega\in[-r_0/L,r_0/L]$ and $\Re v_i=\Re w=2r_0/L$, while the $\alpha_k$ are imaginary; so the hypotheses of \cref{lem:baleuler}(ii) hold with $\delta=r_0/L$, for $L$ large. Together with $|\zeta(1+\omega+\alpha_k)|\ll L$ on $\mathcal C$ (as $|\omega+\alpha_k|\ge1/L$) and $|x^w|\le4$, this justifies the interchange, and \cref{lem:baleuler}(ii) identifies $\sum_{q,J}\sigma_{\boldsymbol v}(q)q^\omega J^{-1-\omega-w}\mathcal Z_{q,J}(1+\omega)$ with $\prod_k\zeta(1+\omega+\alpha_k)\mathcal S$. The $x$-integral is $\widetilde\Phi(1+\omega+w)$.
\end{proof}

\begin{lemma}[Limit]\label{lem:ballimit}
$\mathcal M=a_3TL^9(\int\Phi)\Lambda+o(TL^9)$, where, with $\boldsymbol\sigma=(u_1,u_2,u_3,\beta_3)$,
\[
 \Lambda=\frac1{(2\pi i)^4}\oint_{|\nu|=r_0}\int_{(2r_0)^3}\mathcal Lc(\beta)\,\frac{e^\nu\prod_{i=1}^2(\beta_i-\nu)}{\prod_{l=1}^4(\nu+\sigma_l)\prod_{i=1}^2\prod_{l=1}^4(\beta_i+\sigma_l)}\dd\beta\dd\nu .
\]
\end{lemma}
\begin{proof}
Substitute $\omega=\nu/L$ in \cref{lem:balpoisson}. Since $\zeta(1+s)-1/s$ is entire, $L^{-1}\zeta(1+z/L)\to1/z$ for each $z\ne0$. Using \cref{lem:baleuler}(i), the integrand of \cref{lem:balpoisson}, multiplied by $\dd\omega=\dd\nu/L$ and divided by $TL^9$, converges pointwise to $a_3(\int\Phi)$ times the integrand of $\Lambda$: the zeta factors give $L^{3+1+2+6-2}=L^{10}$ times the corresponding linear factors, $\mathfrak S\to\mathfrak S(0)=a_3$, $(\frac T{2\pi})^{\omega+w}e^{-\beta_3}\to e^\nu$ and $\widetilde\Phi(1+\omega+w)\to\int\Phi$; and $\prod_k(\nu+u_k)(\nu+\beta_3)=\prod_l(\nu+\sigma_l)$, $\prod_i(\beta_i+\beta_3)\prod_{i,k}(\beta_i+u_k)=\prod_{i,l}(\beta_i+\sigma_l)$. For dominated convergence we bound the factors on the contours, uniformly in $L$. On $|\nu|=r_0$ we have $|\nu+u_k|\ge1$, so $|L^{-1}\zeta(1+(\nu+u_k)/L)|\le2$. The other zeta factors in the numerator have arguments $1+z/L$ with $\Re z\ge r_0$, so $|L^{-1}\zeta(1+z/L)|\le L^{-1}\zeta(1+\Re z/L)\le2/r_0$. For the denominator, $\Re(\beta_i-\nu)\ge r_0$ and $|L/\zeta(1+z/L)|\ll(1+|z|)^2$ for $\Re z\ge r_0$, as in the proof of \cref{prop:realisation}. Finally $\mathfrak S$ is bounded, $|(\frac T{2\pi})^{\omega+w}e^{-\beta_3}|\le e^{2r_0}$ and $|\widetilde\Phi(1+\omega+w)|\le4\int|\Phi|$. So the normalised integrand is $\ll|\mathcal Lc(\beta)|(1+|\beta|)^4$, which is integrable.
\end{proof}

\begin{lemma}[Four-variable Haar identity]\label{lem:haar4}
Let $A,B\ge0$ be integers with $A+B\le N$. For $X\in\C^4$ with distinct coordinates,
\[
 \Big\langle\prod_{i=1}^4\det(I+X_iU),\ D\,e_Ae_B\Big\rangle_N=\sum_{i=1}^4\frac{X_i^{N+3}h_A(X_{\ne i})h_B(X_{\ne i})}{\prod_{l\ne i}(X_i-X_l)},
\]
where $h_A$ is the complete homogeneous symmetric polynomial of degree $A$ and $X_{\ne i}\in\C^3$ omits $X_i$. If $|X_i|\le1$ for all $i$, the left side has absolute value at most $(N+1)(N+3)^6$.
\end{lemma}
\begin{proof}
By the dual Cauchy identity \cite[I.(4.3$'$)]{Mac}, $\prod_{i=1}^4\det(I+X_iU)=\sum_\lambda s_{\lambda'}(X)s_\lambda(U)$, over the $\lambda$ with $\ell(\lambda)\le N$ and $\lambda_1\le4$. By \ref{S:pieri}, $e_Ae_B=\sum_{y=0}^{\min(A,B)}s_\nu$ with $\nu'=(A+B-y,y)$, and $Ds_\nu=s_{\nu+(1^N)}$, a Schur polynomial whose column lengths are $(N,A+B-y,y)$; this is a partition because $A+B\le N$. By \ref{S:orth} the left side is $\sum_{y}s_{(N,A+B-y,y,0)}(X)$. For a partition $\mu$ with at most four parts, $s_\mu(X)=\det(X_i^{\mu_l+4-l})/\det(X_i^{4-l})$. For $\mu=(N,j,A+B-j,0)$ expand the numerator along the column $l=1$, of exponent $N+3$. The minor omitting row $i$ is $V(X_{\ne i})s_{(j,A+B-j)}(X_{\ne i})$, where $V(Y)=\prod_{a<b}(Y_a-Y_b)$, and $\det(X_i^{4-l})=V(X)=(-1)^{i-1}V(X_{\ne i})\prod_{l\ne i}(X_i-X_l)$. Hence
\[
 s_\mu(X)=\sum_{i=1}^4\frac{X_i^{N+3}s_{(j,A+B-j)}(X_{\ne i})}{\prod_{l\ne i}(X_i-X_l)},
\]
and summing over $\max(A,B)\le j\le A+B$ with Pieri's rule $h_Ah_B=\sum_js_{(j,A+B-j)}$ \cite[I.(5.16)]{Mac} gives the identity. For the bound, $s_\mu$ has non-negative coefficients, so $|s_\mu(X)|\le s_\mu(1,1,1,1)=\prod_{a<b}\frac{\mu_a-\mu_b+b-a}{b-a}\le(N+3)^6$, and there are at most $N+1$ values of $y$.
\end{proof}
\noindent As a finite consistency check, we compared the two sides in exact rational arithmetic for five triples $(N,A,B)$ at one rational point $X\in\Q^4$; the proof above gives the general identity.

\begin{lemma}[Kernel identity]\label{lem:balkernel}
Let $\boldsymbol u\in(i\R)^3$ have distinct coordinates, let $c\in C_c^\infty((0,1)^3)$ be real with $a_1+a_2<1$ on $\operatorname{supp}c$, and define $\Lambda$ as in \cref{lem:ballimit}. Then $\Lambda=\int_\Om h_{e^{-\boldsymbol u\cdot a}}h_c$.
\end{lemma}
\begin{proof}
For $\Re\beta_3>0$ put $\boldsymbol\sigma=(u_1,u_2,u_3,\beta_3)$; its coordinates are distinct. For $l\le4$ and $x\ge0$ let $F_l(x)=\int e^{-\sum_{l'\ne l}\sigma_{l'}t_{l'}}\dd t$, over $t\in[0,\infty)^{\{1,2,3,4\}\setminus\{l\}}$ with $\sum t_{l'}=x$ (two-dimensional measure), and
\begin{gather*}
 S(\boldsymbol\sigma;x_1,x_2)=\sum_{l=1}^4\frac{e^{-\sigma_l}F_l(x_1)F_l(x_2)}{\prod_{l'\ne l}(\sigma_{l'}-\sigma_l)},\\
 \mathcal L_3c(x;\beta_3)=\int c(x_1,x_2,a_3)e^{\beta_3(1-a_3)}\dd a_3 .
\end{gather*}
Since $\Re\sigma_l\ge0$, $|F_l(x)|\le x^2/2$; and on $\Re\beta_3=2r_0$ each $|\sigma_{l'}-\sigma_l|$ is bounded below by $\min(2r_0,\min_{k\ne k'}|u_k-u_{k'}|)$, so $|S(\boldsymbol\sigma;x)|\ll(1+|x|)^4$ uniformly in $\beta_3$. Also $|\mathcal L_3c(x;\beta_3)|\ll_N(1+|\beta_3|)^{-N}$ uniformly in $x$, and $\mathcal L_3c(\cdot;\beta_3)$ vanishes off a compact subset of $(0,1)^2$. We show that both sides equal
\[
 \Lambda'=\int_{\R^2}\frac1{2\pi i}\int_{(2r_0)}\mathcal L_3c(x;\beta_3)S(\boldsymbol\sigma;x)\dd\beta_3\dd x .
\]

\emph{(a) $\Lambda=\Lambda'$.} In the $\nu$-integral the poles inside $|\nu|=r_0$ are $\nu=-u_k$, $k\le3$, while $\nu=-\beta_3$ lies outside as $|\beta_3|\ge2r_0$. The residue at $\nu=-\sigma_l$ is
\[
 R_l(\beta)=\frac{e^{-\sigma_l}}{\prod_{l'\ne l}(\sigma_{l'}-\sigma_l)\prod_{i=1}^2\prod_{l'\ne l}(\beta_i+\sigma_{l'})},
\]
so $\Lambda=(2\pi i)^{-3}\int\mathcal Lc(\beta)\sum_{l\le3}R_l(\beta)\dd\beta$. We may add $l=4$: $\mathcal Lc(\beta)R_4(\beta)=\int c(a)e^{\beta_1a_1+\beta_2a_2-\beta_3a_3}\dd a\cdot\prod_k(u_k-\beta_3)^{-1}\prod_{i,k}(\beta_i+u_k)^{-1}$ is holomorphic in $\Re\beta_3>0$ and, integrating by parts in $a_3$, is $\ll_N e^{-a_*\Re\beta_3}(1+|\beta_3|)^{-N}$ there, uniformly in $\Re\beta_3\ge2r_0$, where $a_*>0$ is the least value of $a_3$ on $\operatorname{supp}c$. Moving the $\beta_3$-line to $\Re\beta_3=R\to\infty$ shows that its $\beta_3$-integral vanishes. Next, $\int_0^\infty F_l(x)e^{-\beta_ix}\dd x=\prod_{l'\ne l}(\beta_i+\sigma_{l'})^{-1}$ for $\Re\beta_i>0$, so $\sum_lR_l(\beta)=\iint_{[0,\infty)^2}S(\boldsymbol\sigma;x)e^{-\beta_1x_1-\beta_2x_2}\dd x$. By the bounds above Fubini's theorem applies, and Laplace inversion in $(\beta_1,\beta_2)$, applied to $\mathcal Lc(\beta)=\iint\mathcal L_3c(x;\beta_3)e^{\beta_1x_1+\beta_2x_2}\dd x$, gives $\Lambda=\Lambda'$.

\emph{(b) $\int_\Om h_{e^{-\boldsymbol u\cdot a}}h_c=\Lambda'$.} Let $F(a)=e^{-\boldsymbol u\cdot a}$ on $[0,1]^3$. By \cref{lem:profiles}(iii), in its complex form, the left side is $\lim N^{-9}\ip{E_N[F]}{E_N[c]}_N$. Here $E_N[F]=\prod_{k=1}^3\det(I+X_kU)$ with $X_k=e^{-u_k/N}$. By \ref{S:dual}, $\overline{e_j}=\overline De_{N-j}$, so $\ip{E_N[F]}{e_Ae_Be_j}_N=\ip{E_N[F]e_{N-j}}{De_Ae_B}_N$. By Laplace inversion in $a_3$, $c(A/N,B/N,j/N)=(2\pi i)^{-1}\int_{(2r_0)}\mathcal L_3c(\frac AN,\frac BN;\beta_3)X_4^{N-j}\dd\beta_3$ with $X_4=e^{-\beta_3/N}$, and $\sum_{j=0}^NX_4^{N-j}e_{N-j}=\det(I+X_4U)$. As $c$ is real,
\[
 \ip{E_N[F]}{E_N[c]}_N=\sum_{A,B}\frac1{2\pi i}\int_{(2r_0)}\mathcal L_3c\Big(\frac AN,\frac BN;\beta_3\Big)H_N(A,B;\beta_3)\dd\beta_3,
\]
where $H_N=\big\langle\prod_{l=1}^4\det(I+X_lU),De_Ae_B\big\rangle_N$.
Only $(A,B)$ with $A+B<N$ contribute, so \cref{lem:haar4} applies with $|X_l|\le1$; in particular $|H_N|\le(N+1)(N+3)^6$. For fixed $\beta_3$, uniformly in $A,B$: $X_l^{N+3}=e^{-\sigma_l}+O(N^{-1})$, $X_l-X_{l'}=N^{-1}(\sigma_{l'}-\sigma_l)(1+O(N^{-1}))$, and $N^{-2}h_A(X_{\ne l})=N^{-2}\sum_{m\in\Z_{\ge0}^3,\,|m|=A}e^{-\sum\sigma_{l'}m_{l'}/N}=F_l(A/N)+O(N^{-1})$, a Riemann sum over the simplex; so $N^{-7}H_N(A,B;\beta_3)=S(\boldsymbol\sigma;\frac AN,\frac BN)+o(1)$. By dominated convergence in $\beta_3$, with the majorant $\sup_x|\mathcal L_3c(x;\beta_3)|$,
\[
 \frac{N^{-7}}{2\pi i}\int\mathcal L_3c\Big(\frac AN,\frac BN;\beta_3\Big)H_N\dd\beta_3=G\Big(\frac AN,\frac BN\Big)+o(1)
\]
uniformly in $(A,B)$, where $G(x)=\frac1{2\pi i}\int_{(2r_0)}\mathcal L_3c(x;\beta_3)S(\boldsymbol\sigma;x)\dd\beta_3$. The function $G$ is continuous with compact support in $(0,1)^2$, so $N^{-9}\ip{E_N[F]}{E_N[c]}_N=N^{-2}\sum_{A,B}G(\frac AN,\frac BN)+o(1)\to\int G=\Lambda'$.
\end{proof}
\noindent The kernel identity has also been checked numerically at two points.

\begin{proof}[Proof of \eqref{eq:balelementary}]
By \cref{lem:baltrunc,lem:balequi,lem:balpolar,lem:ballimit}, $\mathcal I=a_3TL^9(\int\Phi)\Lambda+o(TL^9)$, and $\Lambda=\int_\Om h_{e^{-\boldsymbol u\cdot a}}h_c$ by \cref{lem:balkernel}, which applies since $\operatorname{supp}c\subset\mathcal R^\sharp\subset(0,1)^3$ and $a_1+a_2<\frac{11}{21}$ there. All bounds depend on $\Phi$ only through its support and finitely many seminorms.
\end{proof}

\subsection{Proof of \Cref{thm:row}}
Fix $\Phi\in\mathcal W$, and let $\psi,\eta$ be real. For $\psi\in C_c^\infty(\mathcal R)$, \cref{lem:synthesis} gives $h_{\mathsf V\psi}=\psi$, and \cref{lem:sixth}(ii) gives $\|B_\psi\|_\Phi\ll(TL^9)^{1/2}$. The coefficients of $B_\psi$ live at $\log N/L\in I_\psi$, and those of $P_{(\mathsf V\psi)^J}$ at $\log N/L\in3-I_\psi\subset(\frac{27}{14},2)$. Recall that $\ip H{\mathcal JK}_\Phi=\overline{\ip{\mathcal JH}K_\Phi}$ for real $\Phi$.

\emph{(a) Numerators.} As in part (a) of \Cref{sec:markedproof}, \cref{lem:expansion} and \cref{lem:localiseP} give $\mathrm N_T(B_\psi)=\ip{P_{c_0}}{B_\psi}_\Phi+O(T^{-A})$ with $c_0(a)=\varphi_1(|a|)\omega_0(a)$, where $\varphi_1\in C_c^\infty((1,2))$ is $1$ near $I_\psi$. By \cref{prop:balanced}, $\mathrm N_T(B_\psi)/\mathcal N_T\to\ip{h_{c_0}}{\psi}=\ip{\varphi_1(r)v}\psi=\ip v\psi$. Further $\mathrm N_T(\mathcal JB_\psi)=\overline{\mathrm N_T(B_\psi)}$ has the limit $\ip v\psi=\ip v{\psi\circ J}$.

\emph{(b) Gram entries within $\mathcal T_{\mathcal R}$.} By \cref{prop:balanced}, $\mathrm G_T(B_\psi,B_\eta)/\mathcal N_T\to\ip{h_{\mathsf V\psi}}{h_{\mathsf V\eta}}=\ip\psi\eta$, and $\mathrm G_T(\mathcal JB_\psi,\mathcal JB_\eta)=\overline{\mathrm G_T(B_\psi,B_\eta)}$ has the limit $\ip\psi\eta=\ip{\psi\circ J}{\eta\circ J}$. By \cref{lem:reflcube} and the norm bound, $\mathrm G_T(\mathcal JB_\eta,B_\psi)=\ip{P_{(\mathsf V\eta)^J}}{B_\psi}_\Phi+o(TL^9)$, and this is $O(T^{-A})$ by \cref{lem:localiseP}, as $3-I_\eta$ and $I_\psi$ are disjoint. This agrees with $\ip{\eta\circ J}\psi=0$, since $J\mathcal R\cap\mathcal R=\emptyset$. The entry $\mathrm G_T(B_\psi,\mathcal JB_\eta)$ is the complex conjugate.

\emph{(c) Cross-pairings with $\mathcal T_3^\infty$.} As in part (c) of \Cref{sec:markedproof}, since $\mathrm G_T(\Pi K,\mathcal JB_\psi)=\overline{\mathrm G_T(\Pi K,B_\psi)}$ and $h_{\Pi K}$ is $J$-symmetric, it suffices to show $\mathrm G_T(K,B_\psi)/\mathcal N_T\to\ip{h_{F_K}}\psi$ and $\mathrm G_T(\mathcal JK,B_\psi)/\mathcal N_T\to\ip{h_{F_K}\circ J}\psi$ for $K\in\{U,V,D_\phi\}$. For $U$ and $V$ this is \cref{lem:factors} with $W_\psi$, $I_\psi$ and $\psi^\flat$ replaced by $B_\psi$, $I_\psi$ and $\psi$. Its proof applies verbatim: it uses $\|W_\psi\|_\Phi\ll(TL^9)^{1/2}$, which for $B_\psi$ is \cref{lem:sixth}(ii); \cref{lem:localise}, which for $B_\psi$ is \cref{lem:localiseP}; and \cref{prop:mixed} for the pairing of $W_\psi$ with a smooth compact source, which for $B_\psi$ is \cref{prop:balanced}; and in Step~4 it uses that $\psi^\flat$ vanishes unless $r\in I_\psi$, which holds for $\psi$. For $K=D_\phi$ the profiles $\phi(r)v$ and $\phi(3-r)v$ vanish on $1<r<2$. The coefficients of $D_\phi$ live at $\log N/L<\theta<1$, so $\mathrm G_T(D_\phi,B_\psi)=O(T^{-A})$ by \cref{lem:localiseP}; and $\mathrm G_T(\mathcal JD_\phi,B_\psi)=O(T^{-A})$ by the integration by parts in part (c) of \Cref{sec:markedproof}, with $\rho=\max I_\psi<\frac{15}{14}$, since $\theta+\rho<3$.

\emph{(d) Cross-pairings with marked tests.} Let $\eta\in C_c^\infty(\mathcal Y)$. Every $a\in\operatorname{supp}\mathsf V\psi$ has $|a|\in I_\psi\subset(1,2)$, so \cref{prop:mixed} gives $\mathrm G_T(B_\psi,W_\eta)/\mathcal N_T\to\ip{\psi}{\eta^\flat}$; $\mathrm G_T(W_\eta,B_\psi)$ is the conjugate, with the same real limit. Next $\mathrm G_T(B_\psi,\mathcal JW_\eta)=\overline{\ip{\mathcal JB_\psi}{W_\eta}_\Phi}=\overline{\ip{P_{(\mathsf V\psi)^J}}{W_\eta}_\Phi}+o(TL^9)$ by \cref{lem:reflcube} and $\|W_\eta\|_\Phi\ll(TL^9)^{1/2}$. Every $a\in\operatorname{supp}(\mathsf V\psi)^J$ has $|a|\in3-I_\psi\subset(1,2)$, so \cref{prop:mixed} and \cref{lem:reflectedprofile} give the limit $\ip{\psi\circ J}{\eta^\flat}=\ip\psi{\eta^\flat\circ J}$. The entries $\mathrm G_T(\mathcal JB_\psi,\mathcal JW_\eta)=\overline{\mathrm G_T(B_\psi,W_\eta)}$ and $\mathrm G_T(\mathcal JB_\psi,W_\eta)=\overline{\mathrm G_T(B_\psi,\mathcal JW_\eta)}$ follow.

Pairs within $\mathcal T_4$ are covered by \Cref{thm:marked}. All limits hold for every $\Phi\in\mathcal W$, so $\mathcal T_5$ is faithful.

\emph{The span.} The real $\psi\in C_c^\infty(\mathcal R)$ are dense in $L^2(\mathcal R)$, and the $\psi\circ J$ are dense in $L^2(J\mathcal R)$, so $\mathcal H(\mathcal T_5)\supseteq L^2(\mathcal R\cup J\mathcal R)$. With $v\one_{S_3\cup\mathcal Z}\in\mathcal H(\mathcal T_4)$ from \Cref{thm:marked}, $v\one_{S_3\cup\mathcal Z\cup\mathcal R\cup J\mathcal R}=v\one_{S_3\cup\mathcal Z}+v\one_{(\mathcal R\cup J\mathcal R)\setminus(S_3\cup\mathcal Z)}\in\mathcal H(\mathcal T_5)$.

\emph{The measure.} By the proof of \Cref{thm:marked}, $\Om\setminus(S_3\cup\mathcal Z)$ is, up to a null set, the disjoint union of $N_0$ and $JN_0$, which have traces in $(1,\frac32)$ and $(\frac32,2)$. Hence $\mathcal R\cap JN_0=\emptyset=J\mathcal R\cap N_0$, the complement of $S_3\cup\mathcal Z\cup\mathcal R\cup J\mathcal R$ is $(N_0\setminus\mathcal R)\cup J(N_0\setminus\mathcal R)$ up to a null set, and its mass is $2\mu(N_0)-2\mu(N_0\cap\mathcal R)$. In the coordinates $(r,k,x)$, which have Jacobian one, $N_0\cap\mathcal R$ is
\begin{gather*}
 1<r<\tfrac{15}{14},\qquad\tfrac12<k<\tfrac{18-7r}{21},\qquad x_{\min}<x<\tfrac k2,\\
 x_{\min}=\max\big(2k-r,\ r-k-\tfrac12\big),
\end{gather*}
because $x<y$ is $x<\frac k2$, $y<z$ is $x>2k-r$, $z-x<\frac12$ is $x>r-k-\frac12$, and the conditions $z<1$, $r<\frac32$ and $x>0$ are implied (if $r-k-\frac12\le0$ then $2k-r\ge r-1>0$). Writing $y-x=s$, one has $z-x=\frac{A+s}2$ and $z-y=\frac{A-s}2$ with $A=2r-3k$, so $9!v^2=\frac{9!}{64}s^2(A^2-s^2)^2$, and the $x$-integral is
\[
 \frac{9!}{128}\int_0^t s^2(A^2-s^2)^2\dd s=\frac{9!}{128}\Big(\frac{A^4t^3}3-\frac{2A^2t^5}5+\frac{t^7}7\Big),\qquad t=k-2x_{\min}.
\]
The two expressions for $x_{\min}$ cross at $k=\frac{4r-1}6$, which lies in the range of $k$ exactly when $r<\frac{43}{42}$. Integrating the resulting piecewise polynomial exactly gives
\[
 \mu(N_0\cap\mathcal R)=\frac{62748701}{30991570176}=\frac{62748701}{2^8\cdot3\cdot7^9}
\]
(computed by iterated polynomial integration and, independently, by rational triangulation). With $\mu(N_0)=\frac{209}{1024}$ this gives $C_0$. By \cref{prop:projection}, $\liminf_{T\to\infty}M_6(T)/(c_3TL^9)\ge C_0$. \qed

\section{Factorable moduli just past half length}\label{sec:factorable}

The balanced region $\mathcal R$ reaches shapes whose two shorter coordinates have sum $k=x+y$ up to $\frac{11}{21}$, but only for traces $r<\frac{15}{14}$. In this section we add a region of all traces $1<r<\frac32$ in which $k$ lies just above $\frac12$. Put
\begin{gather*}
 \mathcal P=\Big\{(a_1,a_2)\in(0,1)^2:\ \tfrac12<k<\tfrac{13}{25},\ \ 5(2k-1)<a_1<\tfrac k2-\tfrac32(2k-1)\Big\},\\
 \mathcal F=\{(x,y,z)\in\Om:\ 1<r<\tfrac32,\ (x,y)\in\mathcal P\},\\
 \mathcal F^\sharp=\{a\in(0,1)^3:\ 1<|a|<\tfrac32,\ (a_1,a_2)\in\mathcal P\}.
\end{gather*}
where $k=a_1+a_2$. On $\mathcal P$ we have $a_1<\frac k2<a_2$, and the two bounds for $a_1$ meet exactly at $k=\frac{13}{25}$. In the coordinates
\[
 \varsigma=2-\frac1k,\qquad\mathfrak b=\frac{a_1}k ,
\]
which we use for the arithmetic, $\mathcal P$ becomes $\mathcal P^*=\{0<\varsigma<\frac1{13},\ 5\varsigma<\mathfrak b<\frac12-\frac32\varsigma\}$. Thus $\mathcal F=\mathcal F^\sharp\cap\Om$; and $\mathcal F\cap J\mathcal F=\emptyset$, since the traces on $J\mathcal F$ exceed $\frac32$. For real $\psi\in C_c^\infty(\mathcal F)$ the function $\mathsf V\psi$ is real and smooth with support in $\operatorname{supp}\psi\subset\mathcal F^\sharp$, and, with $\mathsf V$ and $I_\psi$ as in \Cref{sec:balanced}, we put $F_\psi=P_{\mathsf V\psi}$.

\begin{theorem}[Factorable row]\label{thm:factorable}
Let
\[
 \mathcal T_{\mathcal F}=\big\{(F_\psi,\psi),\ (\mathcal JF_\psi,\psi\circ J):\ \psi\in C_c^\infty(\mathcal F)\text{ real}\big\},\qquad\mathcal T_6=\mathcal T_5\cup\mathcal T_{\mathcal F}.
\]
Then $\mathcal T_6$ is faithful, $v\one_{U_6}\in\mathcal H(\mathcal T_6)$ with $U_6=S_3\cup\mathcal Z\cup\mathcal R\cup J\mathcal R\cup\mathcal F\cup J\mathcal F$, and
\begin{align*}
 \mu(U_6)&=C_0+2\mu\big((N_0\cap\mathcal F)\setminus\mathcal R\big)\\
 &=C_0+\frac{840731426179152778718787}{52986177566328125000000000}\\
 &=C_6:=\frac{3307268442615739142376085993}{79479266349492187500000000}\\
 &=41.61171327466424317998\ldots .
\end{align*}
\end{theorem}

\emph{Overview.} The proof follows that of \Cref{thm:row}; only the analytic input changes. In \Cref{sec:balanced} the pairing of a witness $P_c$ with a source was computed by grouping the two shorter factors of $P_c$ into a modulus $q$ and using the distribution of the \emph{source} in progressions modulo $q$; this needed the trace of the witness to be below $\frac{15}{14}$. Here we instead treat the longest factor of the witness as a shifted zeta function, by Fourier inversion and the functional equation, so that every pairing reduces to the four-shift integral $I(\boldsymbol\alpha)$ of \cref{prop:unequal}, now with moduli $q=hk$ of size $T^\kappa$ with $\frac12<\kappa<\frac{13}{25}$. Its principal part is evaluated exactly as in \cref{prop:unequal}. For the non-principal characters we apply the functional equations of the four Dirichlet $L$-functions. The resulting dual sum has length about $q^{1+\varsigma}$ and involves, at primes, the hyper-Kloosterman sums $\mathrm{Kl}_3$; it must be shown to be smaller than its trivial bound by slightly more than $q^\varsigma$. We do this by splitting the conductor as $HK$, where $H$ collects the primes of the shorter factor $h$ and $K$ is squarefree, shifting the dual variables by multiples of $H$ to remove the $H$-part, and completing modulo $K$, where correlation sums of $\mathrm{Kl}_3$ have square-root cancellation \cite{FKM}. The inequalities defining $\mathcal P$ are exactly the conditions under which one of three such estimates applies to every configuration of the four dual variables. The identification with Schur profiles then reuses \cref{lem:balkernel}.

\subsection{The four-shift row}
For $F\in C_c^\infty(\R^2)$, possibly complex-valued, put
\[
 A_F(s)=\sum_{h,k\ge1}F\Big(\frac{\log h}L,\frac{\log k}L\Big)(hk)^{-s},
\]
and let $S(\boldsymbol\sigma;x)$, $x\in\R^2$, be the kernel of \cref{prop:unequal} (equivalently of \cref{lem:balkernel}).

\begin{proposition}[Four-shift row]\label{prop:fourshift}
Let $F\in C_c^\infty(\R^2)$ have support in a compact subset of $\mathcal P$, and let $\boldsymbol\sigma\in(i\R)^4$. For $\Phi\in\mathcal W$,
\[
 \begin{split}
 I_F(\boldsymbol\sigma)&:=\int\Phi(t/T)\,\chi(1-s)\prod_{i=1}^4\zeta\Big(s+\frac{\sigma_i}L\Big)\overline{A_F(s)}\dd t\\
 &=a_3TL^9\Big(\int\Phi\Big)\int_{\R^2}\overline{F(x)}\,S(\boldsymbol\sigma;x)\dd x+o(TL^9).
 \end{split}
\]
\end{proposition}
\begin{proof}
We follow the proof of \cref{prop:unequal} from the paragraph \emph{Stationary phase} on, with $f_hg_k$ replaced by $\overline{F(\log h/L,\log k/L)}$ and $\varkappa$ by $\kappa_+=\max\{a_1+a_2:\ a\in\operatorname{supp}F\}<\frac{13}{25}$. The stationary-phase formula \eqref{eq:saddle} uses only $\varkappa<1$, so it holds with error $O(T^{\kappa_++\epsilon})$. The character expansion is unchanged, and we treat its three parts as follows.

\emph{Non-principal characters.} Their total contribution is $O(T^{1-\eta})$ for some $\eta>0$, by \cref{prop:nonprincipal} below. This replaces the use of \eqref{eq:fourth}, which suffices only for $\varkappa<\frac12$.

\emph{Principal characters off the poles.} For the pair $(d,\xi_0)$, with $\xi_0$ principal modulo $q=N/d$, the contour shift to $\Re w=\frac12+\delta$ crosses only the poles at $w=1-\alpha_i$. On the new line the Dirichlet series is $d^{-w}\prod_i\zeta(w+\alpha_i)$ times local factors at the primes dividing $N$, of total size $O_\epsilon(N^\epsilon)$, and $|\tau(\xi_0)|=|\mu(q)|\le1$. So the new line contributes $\ll_\epsilon(NT)^{\frac12+\delta}N^\epsilon\sum_{d\mid N}d^{-\frac12-\delta}\varphi(N/d)^{-1}\ll T^{\frac12+\delta}N^{2\epsilon}$ for fixed $N$, and $O(T^{\frac12+2\delta})$ after summing over $h,k$ with the weight $1/N$. This replaces the bound used in \cref{prop:unequal} for these terms, which was adequate only for $\varkappa<\frac12$.

\emph{Residues.} The residue computation of \cref{prop:unequal} does not use $\varkappa<\frac12$. In it, replace the Mellin transforms of $f$ and $g$ by the two-variable Laplace transform of $\overline F$, which decays rapidly on the lines $\Re u=\Re v'=c/L$. The local factors are unchanged, and the limit is $\int\Phi\iint\overline{F(x)}S(\boldsymbol\sigma;x)\dd x$, first for distinct $\sigma_i$ and then for all $\boldsymbol\sigma$ by holomorphy, as there.
\end{proof}

\subsection{The dual sum}
Fix $F$ and $\boldsymbol\sigma$ as in \cref{prop:fourshift}, and put $\boldsymbol\alpha=\boldsymbol\sigma/L$. Let $E\subset\mathcal P$ be a compact set containing $\operatorname{supp}F$. For a modulus $q\ge1$ put $X=qT/(2\pi)$ and
\[
 \mathcal S(q)=\sum_{m\ge1}d_{\boldsymbol\alpha}(m)\,e(-m/q)\,\Phi(m/X),
\]
which is the inner sum in \eqref{eq:saddle}. For $g\mid q$ put $q'=q/g$. As in \cref{prop:unequal},
\begin{equation}\label{eq:charexp}
 \mathcal S(q)=\sum_{g\mid q}\frac1{\varphi(q')}\sum_{\xi\bmod q'}\tau(\bar\xi)\xi(-1)\sum_{n\ge1}d_{\boldsymbol\alpha}(gn)\xi(n)\Phi(gn/X).
\end{equation}
Each $\xi$ is induced by a unique primitive character $\psi$ modulo some $d\mid q'$, and $\tau(\bar\xi)=0$ unless $r=q'/d$ is squarefree and coprime to $d$, in which case $\tau(\bar\xi)=\mu(r)\bar\psi(r)\tau(\bar\psi)$ \cite[Lemma~3.1]{IK}. Let $\mathcal S^\flat(q)$ be the part of \eqref{eq:charexp} with $d>1$. For $d\ge1$ and $v\in\Z$ define
\[
 \mathscr K_d(v)=\frac1{d\varphi(d)}\sum_{\psi\bmod d}^{*}\tau(\psi)^3\,\bar\psi(v),
\]
the sum being over primitive characters; so $\mathscr K_d(v)=0$ unless $(v,d)=1$. For $\mathfrak a\in\{0,1\}$ let
\[
 V_{\mathfrak a}(y)=\frac1{2\pi i}\int_{(-\frac1{10})}\widetilde\Phi(w)\,y^{w}\prod_{i=1}^4\frac{\Gamma\big(\frac{1-w-\alpha_i+\mathfrak a}2\big)}{\Gamma\big(\frac{w+\alpha_i+\mathfrak a}2\big)}\dd w\qquad(y>0),
\]
where $\widetilde\Phi(w)=\int_0^\infty\Phi(x)x^{w-1}\dd x$.

\begin{lemma}[Dual sum]\label{lem:dualsum}
Let $q\ge1$. For each triple $(g,d,r)$ with $gdr=q$, $d>1$, $r$ squarefree and $(r,d)=1$ there are complex numbers $b_\ell=b_\ell(g,d,r)$, supported on the divisors $\ell$ of $g^3r^4$ that are coprime to $d$, with $\sum_\ell|b_\ell|\ll_\epsilon q^\epsilon$, such that
\[
 \begin{split}
 \mathcal S^\flat(q)=\sum_{(g,d,r)}\mu(r)\frac{d^2\varphi(d)}{\varphi(dr)}\,\varpi_d\sum_\ell b_\ell\sum_{\mathfrak a\in\{0,1\}}\sum_{n\ge1}&\frac{d_{-\boldsymbol\alpha}(n)}n\,V_{\mathfrak a}\Big(\frac nY\Big)\\
 &\times\frac{\mathscr K_d(rn\bar\ell)+(-1)^{\mathfrak a}\mathscr K_d(-rn\bar\ell)}2,
 \end{split}
\]
where $\bar\ell$ is an inverse of $\ell$ modulo $d$, $Y=Y(g,d,\ell)=g\ell d^4/(\pi^4X)$, $\varpi_d=\pi^{-2+\sum\alpha_i}d^{-\sum\alpha_i}$, and $d_{-\boldsymbol\alpha}(n)=\sum_{n_1n_2n_3n_4=n}\prod_in_i^{\alpha_i}$. Moreover $q^{1+\varsigma}\ll Y/(\ell g^{-3}r^{-4})\ll q^{1+\varsigma}$, where $\varsigma=2-\log T/\log q$.
\end{lemma}
\begin{proof}
Fix $(g,d,r)$ and a primitive $\psi$ modulo $d$, with induced character $\xi$ modulo $q'=dr$. The Dirichlet series $D(w)=\sum_nd_{\boldsymbol\alpha}(gn)\xi(n)n^{-w}$ has an Euler product. At $p\nmid q$ its factor is $\prod_i(1-\psi(p)p^{-w-\alpha_i})^{-1}$. If $p\mid q'$ then $\xi(p)=0$ and the factor is the constant $d_{\boldsymbol\alpha}(p^{v_p(g)})$. If $p\mid g$ and $p\nmid q'$, the factor is $\sum_{j\ge0}d_{\boldsymbol\alpha}(p^{v+j})\psi(p)^jp^{-jw}$ with $v=v_p(g)$. Since $d_{\boldsymbol\alpha}(p^j)=h_j(p^{-\alpha_1},\dots,p^{-\alpha_4})$ is a complete homogeneous polynomial, $\prod_i(1-x_iX)\sum_{j\ge0}h_{v+j}(x)X^j$ is a polynomial in $X$ of degree at most $3$, with coefficients bounded by $C(v+1)^3$ for $|x_i|=1$. Hence
\[
 D(w)=\prod_{i=1}^4L(w+\alpha_i,\psi)\sum_\ell b_\ell\psi(\ell)\ell^{-w},
\]
where $\sum_\ell b_\ell\psi(\ell)\ell^{-w}$ is the product over $p$ of the local polynomials in $\psi(p)p^{-w}$ described next (so $b_\ell$ need not be multiplicative: for zero shifts, $g=2$, $d=3$ and $r=1$ give $b_1=4$), $b_\ell$ is supported on $\ell\mid g^3r^4$ with $(\ell,d)=1$, and it is independent of $\psi$: the primes $p\mid r$ contribute $d_{\boldsymbol\alpha}(p^{v_p(g)})\prod_i(1-\psi(p)p^{-w-\alpha_i})$, of degree $4$, and the primes $p\mid g$, $p\nmid q'$ contribute the polynomial just described. The divisor bound gives $\sum|b_\ell|\ll_\epsilon q^\epsilon$.

By Mellin inversion
\[
 \sum_nd_{\boldsymbol\alpha}(gn)\xi(n)\Phi(gn/X)=\frac1{2\pi i}\int_{(2)}\widetilde\Phi(w)(X/g)^wD(w)\dd w .
\]
As $\psi$ is non-principal, $D$ is entire; $\widetilde\Phi$ decays rapidly on vertical lines and $D$ grows polynomially, so we may move the line to $\Re w=-\frac1{10}$. There we apply the functional equation \cite[Theorem~4.15]{IK}
\[
 L(w,\psi)=\frac{\tau(\psi)}{i^{\mathfrak a}\sqrt d}\Big(\frac d\pi\Big)^{\frac12-w}\frac{\Gamma\big(\frac{1-w+\mathfrak a}2\big)}{\Gamma\big(\frac{w+\mathfrak a}2\big)}L(1-w,\bar\psi),\qquad\psi(-1)=(-1)^{\mathfrak a},
\]
to each factor and expand $\prod_iL(1-w-\alpha_i,\bar\psi)=\sum_nd_{-\boldsymbol\alpha}(n)\bar\psi(n)n^{w-1}$, which converges absolutely on this line. The $\psi$-dependent factors are $\tau(\bar\psi)\xi(-1)\cdot\bar\psi(r)\cdot\tau(\psi)^4d^{-2}\cdot\psi(\ell)\bar\psi(n)$, and since $\tau(\psi)\tau(\bar\psi)=\psi(-1)d$ this equals $d^{-1}\tau(\psi)^3\bar\psi(rn\bar\ell)$. Summing over the primitive $\psi$ of parity $\mathfrak a$, by means of the projector $\frac12(1+(-1)^{\mathfrak a}\psi(-1))$, gives $\frac{\varphi(d)}2(\mathscr K_d(v)+(-1)^{\mathfrak a}\mathscr K_d(-v))$ with $v=rn\bar\ell$. The powers combine as $(X/(g\ell))^w(d/\pi)^{2-4w-\sum\alpha_i}n^{w}=d^2\varpi_d\,(n/Y)^w$. Collecting the factor $\mu(r)/\varphi(q')$ and interchanging the absolutely convergent sum and integral gives the formula. Finally $X\asymp qT$ and $1\le\ell\le g^3r^4$, so $Y\ell^{-1}g^3r^4=g^4r^4d^4/(\pi^4X)=q^4/(\pi^4X)\asymp q^3/T=q^{1+\varsigma}$.
\end{proof}

\begin{lemma}[The dual weight]\label{lem:dualweight}
For $j\ge0$ and $A\ge0$, uniformly for $\boldsymbol\alpha$ as above, $(y\partial_y)^jV_{\mathfrak a}(y)\ll_{j,A}\min(y^{9/10},y^{-A})$.
\end{lemma}
\begin{proof}
The gamma quotient $\Gamma(\frac{1-w+\mathfrak a}2)/\Gamma(\frac{w+\mathfrak a}2)$ is holomorphic in $\Re w<1$, and by Stirling's formula it is $\ll(1+|\Im w|)^{\frac12-\Re w}$ uniformly in vertical strips there, and $\asymp(1+|\Im w|)^{\frac12-\Re w}$ once $|\Im w|$ is large (at bounded height it can vanish, for instance at $w=0$ when $\mathfrak a=0$). Only the upper bound is used. As $|\Re\alpha_i|$ is $0$, the integrand of $V_{\mathfrak a}$ is holomorphic in $\Re w<\frac{19}{20}$ and, because $\widetilde\Phi$ decays rapidly, may be integrated on any line $\Re w=c<\frac{19}{20}$. Applying $(y\partial_y)^j$ multiplies the integrand by $w^j$. Take $c=\frac9{10}$ for $y\le1$ and $c=-A$ for $y>1$.
\end{proof}

\subsection{The kernel}
For a prime $p$ and $c\in\F_p^\times$ let $\mathrm{Kl}_3(c;p)=p^{-1}\sum_{x_1x_2x_3=c}e_p(x_1+x_2+x_3)$, the sum over $x_i\in\F_p^\times$, with $e_p(x)=e(x/p)$. By Deligne's theorem $|\mathrm{Kl}_3(c;p)|\le3$: the Kloosterman sheaf has rank $3$ and is pure of weight $0$ \cite[4.1.1]{GKM}.

\begin{lemma}[Local structure of $\mathscr K_d$]\label{lem:Kd}
\begin{enumerate}[label=\textup{(\roman*)}]
\item If $d=\prod_tt$ with pairwise coprime prime powers $t$, then $\mathscr K_d(v)=\prod_t\mathscr K_t(v\,\overline{(d/t)}^{\,3})$, the inverse taken modulo $t$.
\item For a prime $p$ and $(v,p)=1$, $\mathscr K_p(v)=\mathrm{Kl}_3(v;p)+\frac1{p(p-1)}$.
\item For a prime power $t=p^a$ with $a\ge2$, $|\mathscr K_t(v)|\le9$.
\item For a prime $p$, $c\in\F_p^\times$ and $j\in\F_p$, $\sum_{x\in\F_p}\mathscr K_p(cx)e_p(jx)$ vanishes if $j=0$ and has absolute value at most $2\sqrt p+1$ otherwise.
\end{enumerate}
In particular $|\mathscr K_d(v)|\le9^{\omega(d)}\ll_\epsilon d^\epsilon$.
\end{lemma}
\begin{proof}
(i) A character modulo $d$ is a product $\prod_t\psi_t$ of characters modulo the $t$, and it is primitive if and only if every $\psi_t$ is. By the Chinese remainder theorem $\tau(\psi)=\prod_t\psi_t(d/t)\tau(\psi_t)$, so $\tau(\psi)^3\bar\psi(v)=\prod_t\tau(\psi_t)^3\bar\psi_t(v\,\overline{(d/t)}^3)$, and $d\varphi(d)=\prod_tt\varphi(t)$.

(ii) and (iii). Let $t=p^a$ and $(v,p)=1$. Orthogonality of all characters modulo $t$ gives
\[
 \frac1{\varphi(t)}\sum_{\psi\bmod t}\tau(\psi)^3\bar\psi(v)=\sum_{\substack{x_1,x_2,x_3\in(\Z/t\Z)^\times\\x_1x_2x_3=v}}e_t(x_1+x_2+x_3)=:t\,S_t(v).
\]
For $a=1$ the only imprimitive character is the principal one, whose Gauss sum is $-1$; this gives (ii). For $a\ge2$ every imprimitive $\psi$ is periodic modulo $p^{a-1}$, and then $\tau(\psi)=0$, as one sees by summing over the $p$ lifts of each residue modulo $p^{a-1}$. Hence $\mathscr K_t(v)=S_t(v)=t^{-1}\sum_{x,y}e_t(f(x,y))$ with $f(x,y)=x+y+v\,\overline{xy}$, the sum over units $x,y$ modulo $t$.

Let $a=2m$. Write $x=x_0+p^mu$, $y=y_0+p^m\nu$ with $x_0,y_0$ running over units modulo $p^m$ and $u,\nu$ modulo $p^m$. Since $\overline{x_0+p^mu}\equiv\bar x_0-p^mu\bar x_0^2\pmod{p^{2m}}$, we have $f(x,y)\equiv f(x_0,y_0)+p^m(f_xu+f_y\nu)\pmod{p^{2m}}$, where $f_x=1-v\bar x_0^2\bar y_0$ and $f_y=1-v\bar x_0\bar y_0^2$. Summing over $u,\nu$ gives $p^{2m}$ if $f_x\equiv f_y\equiv0\pmod{p^m}$ and $0$ otherwise. These congruences give $x_0\equiv y_0=:z$ and $z^3\equiv v\pmod{p^m}$. The cubing map on $(\Z/p^m\Z)^\times$ has fibres of size at most $3$, the group being cyclic for odd $p$ and of $2$-power order for $p=2$. So $|S_t(v)|\le3$.

Let $a=2m+1$, $m\ge1$. Summing first over translations by multiples of $p^{m+1}$, as above, shows that only pairs $(x,y)\equiv(z,z)\pmod{p^m}$ with $z^3\equiv v\pmod{p^m}$ contribute. For each of the at most three such classes write $x=z+p^mu$, $y=z+p^m\nu$ with $u,\nu$ modulo $p^{m+1}$. Expanding the inverses to second order,
\[
 f(x,y)\equiv f(z,z)+p^m(1-v\bar z^3)(u+\nu)+p^{2m}v\bar z^4(u^2+u\nu+\nu^2)\pmod{p^{2m+1}}.
\]
As $p^m\mid1-v\bar z^3$, the phase depends only on $u,\nu$ modulo $p$, so the class contributes a unimodular factor times $p^{-1}\sum_{u,\nu\in\F_p}e_p(A(u+\nu)+B(u^2+u\nu+\nu^2))$ with $B\ne0$. For $p\ge5$ the quadratic form is non-degenerate and this sum has absolute value $1$; for $p\le3$ it is at most $p\le3$. Hence $|S_t(v)|\le9$.

(iv) Extend $\psi\mapsto\sum_x\bar\psi(x)e_p(jx)$: it is $0$ for $j=0$ and $\psi$ non-principal, and $\psi(j)\tau(\bar\psi)$ for $j\ne0$. Hence the sum vanishes for $j=0$, and for $j\ne0$, using $\tau(\psi)\tau(\bar\psi)=\psi(-1)p$,
\[
 \sum_x\mathscr K_p(cx)e_p(jx)=\frac1{\varphi(p)}\sum_{\psi\ne1}\tau(\psi)^2\bar\psi(-c\bar j)=\sum_{x_1x_2=-c\bar j}e_p(x_1+x_2)-\frac1{p-1},
\]
and Weil's bound for the Kloosterman sum gives the claim. The final assertion follows from (i)--(iii).
\end{proof}

\begin{lemma}[Correlations]\label{lem:corr}
Let $p\ge5$ be prime and $m\in\{2,4\}$. For $1\le i\le m$ let $A_i(x)=a_i(x+b_i)$ with $a_i\in\F_p^\times$, $b_i\in\F_p$, and $\epsilon_i\in\{0,1\}$. Suppose that for all $i\ne j$ either $b_i\ne b_j$ or $a_i\ne\pm a_j$. Then for every $h\in\F_p$
\[
 \Big|\sum_{x\in\F_p}e_p(hx)\prod_{i=1}^m\mathscr K_p\big(A_i(x)\big)^{[\epsilon_i]}\Big|\le C\sqrt p,
\]
where $z^{[0]}=z$, $z^{[1]}=\bar z$, and $C$ is an absolute constant.
\end{lemma}
\begin{proof}
By \cref{lem:Kd}(ii), $\mathscr K_p(A_i(x))=\mathrm{Kl}_3(A_i(x);p)+\frac1{p(p-1)}$ for $x\ne-b_i$, and $0$ for $x=-b_i$. Expanding the product, the terms containing a constant $\frac1{p(p-1)}$ and the at most $m$ excluded points contribute $O(1)$. Since $\overline{\mathrm{Kl}_3(y;p)}=\mathrm{Kl}_3(-y;p)$, the remaining sum is $\sum_xe_p(hx)\prod_i\mathrm{Kl}_3(B_i(x);p)$ over $x\notin\{-b_i\}$, with $B_i=(-1)^{\epsilon_i}A_i$; by hypothesis $B_i\ne\pm B_j$ as affine maps for $i\ne j$. This is \cite[Corollary~3.4]{FKM} for the Kloosterman sheaf $\mathcal{K}\ell_3$, whose special involution is $x\mapsto-x$; that corollary includes the frequency $h=0$. We recall the argument. The function $\mathrm{Kl}_3(\cdot;p)$ is the trace function of a lisse sheaf $\mathcal K\ell_3$ on $\mathbf G_m$, of rank $3$ and pure of weight $0$, tamely ramified at $0$ and wildly ramified at $\infty$ with Swan conductor $1$ \cite[4.1.1]{GKM}, with unipotent local monodromy at $0$ \cite[7.4.3]{GKM}, with dual $[x\mapsto-x]^*\mathcal K\ell_3$ \cite[4.1.3--4.1.4(2)]{GKM}, and with geometric monodromy group $\mathrm{SL}_3$ for $p\ge5$ \cite[11.1]{GKM}. For $B\ne\pm B'$ the sheaves $[B]^*\mathcal K\ell_3$ and $[B']^*\mathcal K\ell_3$ are not geometrically isomorphic, nor is one isomorphic to the dual of the other, after any twist by a rank-one sheaf \cite[\S3.3, Proposition~3.7(2)--(3)]{FKM}: if the roots $-b$ differ, their unipotent ramification occurs at different points, and if the roots agree, $B'=\lambda B$ with $\lambda\ne\pm1$, and the local monodromy at $\infty$ distinguishes them. By the Goursat--Kolchin--Ribet criterion \cite[1.8.2]{ESDE} the geometric monodromy group of $\bigoplus_i[B_i]^*\mathcal K\ell_3$ contains $\mathrm{SL}_3^m$, so $\bigotimes_i[B_i]^*\mathcal K\ell_3$ is geometrically irreducible of rank $3^m>1$ and its twist by the Artin--Schreier sheaf $\mathcal L_{e_p(hx)}$ has no geometric coinvariants. Deligne's theorem \cite{Deligne} and the Grothendieck--Ogg--Shafarevich formula then bound the sum by $C\sqrt p$, with $C$ depending only on $m$ and the ranks and conductors involved.
\end{proof}

\subsection{Three estimates}
Let $q=hk$ be a modulus, $d>1$ a divisor of $q$, and $c$ a unit modulo $d$. Write $d=HK$, where $H$ is the product of the prime powers $p^a\Vert d$ with $a\ge2$, or $p\mid h$, or $p\le3$, and $K$ is the product of the remaining primes; so $K$ is squarefree, $(H,K)=1$, and every prime of $K$ is at least $5$ and divides $k$. By \cref{lem:Kd}(i), $\mathscr K_d(cv)=F_H(v)G_K(v)$, where $F_H$ is periodic modulo $H$ with $|F_H|\le9^{\omega(H)}$, and $G_K(v)=\prod_{p\mid K}\mathscr K_p(c_pv)$ for units $c_p$ modulo $p$. Throughout, $W,W_1,W_2$ denote smooth functions supported in $[\frac12,2]$ with $\|(y\partial_y)^jW\|_\infty\le q^{j\epsilon}$ for $j\le C_0$, a large constant, and implied constants are uniform in them; all estimates are up to factors $q^{O(\epsilon)}$, which we write $q^{o(1)}$.

\begin{lemma}[One smooth variable]\label{lem:onevar}
For $M\ge1$,
\[
 \frac1M\Big|\sum_mW(m/M)\mathscr K_d(cm)\Big|\ll q^{o(1)}\Big(\frac{H\sqrt K}M+\frac1{\sqrt K}\Big).
\]
\end{lemma}
\begin{proof}
By Poisson summation modulo $d$,
\[
 \sum_mW(m/M)\mathscr K_d(cm)=\frac Md\sum_{j\in\Z}\widehat W(jM/d)\widehat{\mathscr K}(j),
\]
where $\widehat W(\xi)=\int W(y)e(-y\xi)\dd y$ and $\widehat{\mathscr K}(j)=\sum_{x\bmod d}\mathscr K_d(cx)e_d(jx)$. By the Chinese remainder theorem $\widehat{\mathscr K}(j)$ factors over the prime powers of $d$; the factor at $t\mid H$ is at most $9t$, and the factor at $p\mid K$ is at most $3\sqrt p$ and vanishes when $p\mid j$, by \cref{lem:Kd}(iv). So $\widehat{\mathscr K}(0)=0$ if $K>1$, and $|\widehat{\mathscr K}(j)|\le q^{o(1)}H\sqrt K$ always. As $\frac Md\sum_{j\ne0}|\widehat W(jM/d)|\ll q^{o(1)}(1+M/d)$, the sum is $\ll q^{o(1)}H\sqrt K(1+M/d)$; and $H\sqrt K\cdot M/d=M/\sqrt K$. If $K=1$ the claim is trivial.
\end{proof}

\begin{lemma}[Bilinear sums]\label{lem:bilinear}
Let $M,N\ge1$, and let $|\alpha_m|,|\beta_n|\le q^{o(1)}$ be supported on $m\in[M,2^8M]$, $n\in[N,2^8N]$. Then
\[
 \Big(\frac1{MN}\Big|\sum_{m,n}\alpha_m\beta_n\mathscr K_d(cmn)\Big|\Big)^2\ll q^{o(1)}\Big(\frac1N+\Big(\frac HM\Big)^{1/2}+\Big(\frac{\sqrt K}M\Big)^{1/2}+K^{-1/4}\Big).
\]
\end{lemma}
\begin{proof}
The left side is at most $q^{o(1)}$ by \cref{lem:Kd}, so we may assume $M\ge4H$. Let $W\ge0$ be smooth, supported in $[\frac12,2^9]$ and equal to $1$ on $[1,2^8]$. By Cauchy's inequality the square of the sum is at most $q^{o(1)}M\sum_{n_1,n_2}|\beta_{n_1}\beta_{n_2}||\mathcal C(n_1,n_2)|$, where $\mathcal C(n_1,n_2)=\sum_mW(m/M)\mathscr K_d(cmn_1)\overline{\mathscr K_d(cmn_2)}$. The terms $n_1=n_2$ give $q^{o(1)}M^2N$, that is, the term $1/N$.

Let $n_1\ne n_2$. Write $\mathscr K_d(cmn_1)\overline{\mathscr K_d(cmn_2)}=\Psi(m)\Gamma(m)$ with $\Psi(m)=F_H(mn_1)\overline{F_H(mn_2)}$, periodic modulo $H$ with $|\Psi|\le q^{o(1)}$, and $\Gamma(m)=G_K(mn_1)\overline{G_K(mn_2)}$, periodic modulo $K$. Put $R=\lfloor M/(4H)\rfloor\ge1$. Replacing $m$ by $m+Hi$ does not change $\Psi$, so
\[
 \mathcal C(n_1,n_2)=\frac1R\sum_m\Psi(m)\sum_{i=1}^RW\Big(\frac{m+Hi}M\Big)\Gamma(m+Hi),
\]
with $m$ in an interval of length $O(M)$. By Cauchy's inequality,
\[
 |\mathcal C|^2\le\frac{q^{o(1)}M}{R^2}\sum_{i_1,i_2\le R}\Big|\sum_mW_{i_1}(m)W_{i_2}(m)\Gamma(m+Hi_1)\overline{\Gamma(m+Hi_2)}\Big|,
\]
where $W_i(m)=W((m+Hi)/M)$.
The terms $i_1=i_2$ contribute $q^{o(1)}M^2/R\ll q^{o(1)}M^2H/M$. For $i_1\ne i_2$, the inner sum is a smooth sum of a function periodic modulo $K$, and completing it costs a factor $\ll(M/K+\log K)$ times the largest complete sum $\sum_{x\bmod K}e_K(jx)\Gamma(x+Hi_1)\overline{\Gamma(x+Hi_2)}$. This factors over $p\mid K$. The local factor is a correlation of four functions $\mathscr K_p(A(x))$, $A\in\{c_pn_1(x+Hi_1),\ c_pn_2(x+Hi_1),\ c_pn_1(x+Hi_2),\ c_pn_2(x+Hi_2)\}$, two of them conjugated; it vanishes unless $p\nmid n_1n_2$. If moreover $p\nmid(i_1-i_2)(n_1^2-n_2^2)$, the hypotheses of \cref{lem:corr} hold (the roots $-Hi_1,-Hi_2$ are distinct, and at each root the slopes $c_pn_1,c_pn_2$ satisfy $n_1\ne\pm n_2$), so the local factor is $O(\sqrt p)$; otherwise it is at most $p\cdot4^4$. Hence the complete sum is $\ll q^{o(1)}\sqrt K\,(K,(i_1-i_2)(n_1^2-n_2^2))^{1/2}$, and
\[
 \frac{|\mathcal C|^2}{M^2}\ll q^{o(1)}\Big(\frac HM+\Big(\frac{\sqrt K}M+\frac1{\sqrt K}\Big)(K,n_1^2-n_2^2)^{1/2}\frac1{R^2}\sum_{i_1\ne i_2}(K,i_1-i_2)^{1/2}\Big).
\]
Since $\sum_{1\le i\le R}(K,i)^{1/2}\le\sum_{v\mid K}v^{1/2}R/v\le q^{o(1)}R$, the last average is $q^{o(1)}$. Taking square roots, $|\mathcal C|/M\ll q^{o(1)}((H/M)^{1/2}+(\sqrt K/M+1/\sqrt K)^{1/2}(K,n_1^2-n_2^2)^{1/4})$. Finally, with $a=n_1-n_2\ne0$ and $b=n_1+n_2$, we have $(K,ab)^{1/4}\le(K,a)^{1/2}+(K,b)^{1/2}$, and $\sum_{n_1\ne n_2}(K,n_1\pm n_2)^{1/2}\le q^{o(1)}N^2$ by the same divisor sum. Inserting these bounds proves the lemma.
\end{proof}

\begin{lemma}[Two smooth variables]\label{lem:twovar}
Let $M\ge N\ge1$. Then
\[
 \Big(\frac1{MN}\Big|\sum_{m,n}W_1(m/M)W_2(n/N)\mathscr K_d(cmn)\Big|\Big)^2\ll q^{o(1)}\Big(\frac{H^2}{MN}+\frac{\sqrt K}M+\frac1{\sqrt K}\Big).
\]
\end{lemma}
\begin{proof}
Let $\Sigma$ be the sum. If $M<4H$ then $H^2/(MN)>\frac1{16}$, as $N\le M$, and the claim follows from the trivial bound; so let $M\ge4H$. Put $I=\lfloor M/(4H)\rfloor\ge1$ and $J=\max(1,\lfloor N/(4H)\rfloor)$. Replacing $(m,n)$ by $(m+Hi,n+Hj)$ does not change $F_H(cmn)$, so
\[
 \begin{split}
 \Sigma=\frac1{IJ}\sum_{m,n}F_H(mn)\sum_{i<I,\,j<J}&W_1\Big(\frac{m+Hi}M\Big)W_2\Big(\frac{n+Hj}N\Big)\\
 &\times G_K\big((m+Hi)(n+Hj)\big),
 \end{split}
\]
where $(m,n)$ runs over a set of $O(MN)$ points. By Cauchy's inequality $|\Sigma|^2\le q^{o(1)}MN(IJ)^{-2}\sum_{i_1,j_1,i_2,j_2}|\Sigma_{i_1j_1i_2j_2}|$, where $\Sigma_{i_1j_1i_2j_2}$ is the sum over $(m,n)$ of the product of the two weighted terms. The $IJ$ diagonal terms give $|\Sigma|^2/(MN)^2\ll q^{o(1)}/(IJ)\ll q^{o(1)}H^2/(MN)$; indeed $IJ\gg MN/H^2$ if $N\ge8H$, and $IJ=I\gg M/H\ge MN/(8H^2)$ otherwise.

If $i_1\ne i_2$, fix $n$ and sum over $m$. At $p\mid K$ the local factor vanishes identically if $p\mid(n+Hj_1)(n+Hj_2)$. Otherwise it is the correlation of $\mathscr K_p(c_p(n+Hj_1)(x+Hi_1))$ with the conjugate of $\mathscr K_p(c_p(n+Hj_2)(x+Hi_2))$, whose roots are distinct when $p\nmid i_1-i_2$; then \cref{lem:corr} with $m=2$ gives $O(\sqrt p)$ for every additive frequency, and otherwise we use the bound $16p$. Completing the $m$-sum as in \cref{lem:bilinear}, $|\Sigma_{i_1j_1i_2j_2}|\ll q^{o(1)}MN(\sqrt K/M+1/\sqrt K)(K,i_1-i_2)^{1/2}$, and the average of $(K,i_1-i_2)^{1/2}$ is $q^{o(1)}$. These terms give $q^{o(1)}(\sqrt K/M+1/\sqrt K)$.

If $i_1=i_2$ and $j_1\ne j_2$, which requires $J\ge2$ and so $N\ge8H$, fix $m$ and complete in $n$ in the same way. This gives $q^{o(1)}MN(\sqrt K/N+1/\sqrt K)(K,j_1-j_2)^{1/2}$ for each such term: at $p\mid(K,j_1-j_2)$ the two roots coincide and we use the bound $16p$. The factor $(K,j_1-j_2)^{1/2}$ averages to $q^{o(1)}$, since $\sum_{j_1\ne j_2<J}(K,j_1-j_2)^{1/2}\le\sum_{v\mid K}v^{1/2}\#\{j_1\ne j_2:\ v\mid j_1-j_2\}\le2J^2\sum_{v\mid K}v^{-1/2}\le2J^2\tau(K)$. These terms are a proportion $1/I\ll H/M$ of all. Since $(H/M)(\sqrt K/N)\le\sqrt K/M$, they contribute $q^{o(1)}(\sqrt K/M+1/\sqrt K)$.
\end{proof}

\subsection{Proof of the non-principal estimate}
\begin{proposition}\label{prop:nonprincipal}
Let $E\subset\mathcal P$ be compact. There is $\eta>0$ such that
\[
 \sum_{\substack{h,k\ge1\\(\log h/L,\log k/L)\in E}}\frac{|\mathcal S^\flat(hk)|}{hk}\ll T^{1-\eta}.
\]
\end{proposition}

Fix a pair $(h,k)$ with $(\log h/L,\log k/L)\in E$ and put $q=hk$. All exponents below are to base $q$: we write $q^\kappa=T^{\kappa\log q/L}$ and so on. Put
\[
 \varsigma=2-\frac{L}{\log q},\qquad\mathfrak b=\frac{\log h}{\log q},\qquad\mathfrak p=\frac{\log P}{\log q},
\]
where $P$ is the largest powerful divisor of $q$ (the product of the $p^a\Vert q$ with $a\ge2$). Then $(\varsigma,\mathfrak b)$ lies in the compact image $E^*\subset\mathcal P^*$ of $E$, and $q^{2-\varsigma}=T$. There is $c_0\in(0,\frac1{100})$, depending on $E$, with
\begin{equation}\label{eq:margins}
 \tfrac12-\tfrac32\varsigma-\mathfrak b\ge c_0,\qquad\mathfrak b-5\varsigma\ge c_0,\qquad\tfrac1{13}-\varsigma\ge c_0\qquad\text{on }E^*.
\end{equation}

\begin{lemma}[Block estimate]\label{lem:block}
There are $\delta>0$ and $\epsilon_0>0$, depending only on $E$, with the following property. Let $(g,d,r)$ and $\ell$ be as in \cref{lem:dualsum}, put $\mathfrak g=\log g/\log q$, $\mathfrak r=\log r/\log q$, let $c$ be a unit modulo $d$, and let $N_1,\dots,N_4\ge1$ with $N_*=N_1N_2N_3N_4\le Yq^{\epsilon_0}$. Put $\mathfrak v=\max(0,\log(Y/N_*)/\log q)$ and
\[
 \tau=\varsigma-2\mathfrak g-3\mathfrak r-\tfrac12\mathfrak p-\tfrac9{10}\mathfrak v .
\]
Then, for weights $w_i$ as in the previous subsection,
\[
 \Sigma:=\frac1{N_*}\Big|\sum_{n_1,\dots,n_4}\prod_{i=1}^4w_i(n_i/N_i)\,\mathscr K_d(cn_1n_2n_3n_4)\Big|\ll q^{-\tau-\delta}.
\]
\end{lemma}
\begin{proof}
Write $N_i=q^{\nu_i}$, $D=\nu_1+\dots+\nu_4$, and order the indices so that $\nu_1\ge\nu_2\ge\nu_3\ge\nu_4\ge0$. Let $d=HK$ be as in the previous subsection. Every prime power in $H$ with exponent at least $2$ divides $P$, so $H\le6Ph$; every prime of $K$ divides $k=q/h$; and $d=q/(gr)$. Hence, with
\[
 \mathfrak h=\mathfrak b+\mathfrak p,\qquad\mathfrak k=\tfrac12(1-\mathfrak b),
\]
we have $H\le q^{\mathfrak h+o(1)}$, $\sqrt K\le q^{\mathfrak k}$, $K\ge q^{1-\mathfrak b-\mathfrak p-\mathfrak g-\mathfrak r-o(1)}$ and $H\sqrt K=\sqrt{Hd}\le q^{(1+\mathfrak b+\mathfrak p-\mathfrak g-\mathfrak r)/2+o(1)}$. Since $\ell\ge1$, \cref{lem:dualsum} gives $Y\ge q^{1+\varsigma-3\mathfrak g-4\mathfrak r-o(1)}$, so
\[
 D\ge\lambda-o(1),\qquad\lambda=1+\varsigma-3\mathfrak g-4\mathfrak r-\mathfrak v .
\]
The trivial bound is $\Sigma\le q^{o(1)}$, by \cref{lem:Kd}. Let $\delta=c_0/100$ and $\tau'=\tau+2\delta$. If $\tau'\le0$ the trivial bound suffices, so assume $\tau'>0$; it is enough to show $\Sigma\ll q^{-\tau'+o(1)}$. Note that $\tau'\le\varsigma+2\delta$. By direct expansion (also checked by computer) we have the identities
\begin{align*}
 \lambda-8\tau&=(1-7\varsigma)+13\mathfrak g+20\mathfrak r+4\mathfrak p+\tfrac{31}5\mathfrak v,\\
 (1-\mathfrak b-\mathfrak p-\mathfrak g-\mathfrak r)-8\tau&=(1-\mathfrak b-8\varsigma)+15\mathfrak g+23\mathfrak r+3\mathfrak p+\tfrac{36}5\mathfrak v,\\
 \lambda-3\tau-\tfrac12(1+\mathfrak b+\mathfrak p-\mathfrak g-\mathfrak r)&=(\tfrac12(1-\mathfrak b)-2\varsigma)+\tfrac72\mathfrak g+\tfrac{11}2\mathfrak r+\mathfrak p+\tfrac{17}{10}\mathfrak v,\\
 \lambda-4\tau-2\mathfrak h&=(1-3\varsigma-2\mathfrak b)+5\mathfrak g+8\mathfrak r+\tfrac{13}5\mathfrak v,\\
 \lambda-6\tau-2\mathfrak k&=(\mathfrak b-5\varsigma)+9\mathfrak g+14\mathfrak r+3\mathfrak p+\tfrac{22}5\mathfrak v.
\end{align*}
By \eqref{eq:margins} the constant terms on the right are at least $c_0$; for instance $1-\mathfrak b-8\varsigma\ge\frac12-\frac{13}2\varsigma\ge c_0$. Replacing $\tau$ by $\tau'$ lowers each left side by at most $16\delta\le c_0/6$. Hence, for $q$ large,
\begin{equation}\label{eq:blockmargins}
 \begin{gathered}
 D>8\tau',\qquad\log K/\log q>8\tau',\qquad\lambda-3\tau'>\tfrac12(1+\mathfrak b+\mathfrak p-\mathfrak g-\mathfrak r),\\
 D-4\tau'>2\mathfrak h,\qquad D-6\tau'>2\mathfrak k,
 \end{gathered}
\end{equation}
each with a margin of at least $c_0/2$. We now consider three cases; in each, the margins in \eqref{eq:blockmargins} absorb the factors $q^{o(1)}$.

\emph{Case 1: a good grouping.} Suppose that some set $\mathcal I\subset\{1,2,3,4\}$ has $\sum_{i\in\mathcal I}\nu_i>2\tau'$ and $\sum_{i\notin\mathcal I}\nu_i>\max(\mathfrak h,\mathfrak k)+4\tau'$. Group the variables $n_i$, $i\notin\mathcal I$, into $m$ and the others into $n$; the coefficients of $m$ and $n$ are bounded by divisor functions, and are supported in ranges $[M,2^8M]$ and $[N,2^8N]$ with $M\asymp\prod_{i\notin\mathcal I}N_i$ and $N\asymp\prod_{i\in\mathcal I}N_i$. In \cref{lem:bilinear} the four terms are then at most $q^{-2\tau'}$, by the two assumed inequalities and $K>q^{8\tau'}$, so $\Sigma^2\ll q^{-2\tau'+o(1)}$.

\emph{Case 2: one long variable.} Suppose that $\nu_2+\nu_3+\nu_4\le2\tau'$. Fix $n_2,n_3,n_4$; then $cn_2n_3n_4$ is a unit modulo $d$ or the inner sum vanishes. \cref{lem:onevar} with $M=N_1\ge q^{D-2\tau'}$ gives $\Sigma\ll q^{o(1)}(H\sqrt K/N_1+1/\sqrt K)$, and by \eqref{eq:blockmargins} both terms are $\ll q^{-\tau'}$, since $D-2\tau'-\tau'\ge\lambda-3\tau'-o(1)$.

\emph{Case 3: two long variables.} Suppose that $\nu_3+\nu_4\le2\tau'$. Fix $n_3,n_4$; again $cn_3n_4$ is a unit modulo $d$ or the inner sum vanishes. Apply \cref{lem:twovar} with $M=N_1\ge N=N_2$. Here $\nu_1+\nu_2\ge D-2\tau'>2\mathfrak h+2\tau'$ and $\nu_1\ge\frac12(D-2\tau')>\mathfrak k+2\tau'$, by \eqref{eq:blockmargins}; so $\Sigma^2\ll q^{-2\tau'+o(1)}$, using $K>q^{8\tau'}$ for the last term.

It remains to show that one of the three cases occurs. Put $\Lambda=\max(\mathfrak h,\mathfrak k)+4\tau'$, $\Delta=\Lambda-D/2$ and $U=D-\Lambda$.

If $\Delta<0$, then $\nu_1\ge D/4>2\tau'$. Either $\nu_2+\nu_3+\nu_4\le2\tau'$, which is Case 2, or $\{2,3,4\}$ and $\{1\}$ both have degree above $2\tau'$, and the larger of the two has degree at least $D/2>\Lambda$, which is Case 1.

Let $\Delta\ge0$. We claim that $3U>D$, $2\Delta<U$ and $4\tau'<U$. Put $\Gamma=\frac{13}2\mathfrak g+10\mathfrak r+\mathfrak p+\frac{31}{10}\mathfrak v\ge0$. By direct expansion (\S2 of the same script),
\[
 \mathfrak h+4\tau-\tfrac\lambda2=\big(\mathfrak b-\tfrac12+\tfrac72\varsigma\big)-\Gamma,\qquad
 \mathfrak k+4\tau-\tfrac\lambda2=\big(\tfrac72\varsigma-\tfrac12\mathfrak b\big)-\Gamma-\mathfrak p .
\]
By \eqref{eq:margins}, $\mathfrak b-\frac12+\frac72\varsigma\le2\varsigma-c_0$ and $\frac72\varsigma-\frac12\mathfrak b\le\varsigma-\frac{c_0}2$. As $\tau'=\tau+2\delta$, $8\delta<\frac{c_0}4$ and $D\ge\lambda-o(1)$, we get, for $q$ large,
\[
 \Lambda-\tfrac\lambda2<2\varsigma-\Gamma-\tfrac{c_0}4,\qquad\text{so}\qquad 0\le\Delta<2\varsigma-\Gamma-\tfrac{c_0}8 .
\]
First, $\Gamma<2\varsigma$ gives $3\mathfrak g+4\mathfrak r+\mathfrak v<\varsigma$, hence $\lambda>1$ and $D>\frac{12}{13}$; so $\Delta<2\varsigma<\frac2{13}<\frac D6$, and $U=\frac D2-\Delta>\frac D3$, which gives $3U>D$ and $2\Delta<\frac D3<U$. Second, $U=D-\Lambda\ge\frac\lambda2-(\Lambda-\frac\lambda2)-o(1)>\frac\lambda2-2\varsigma+\Gamma$, and $\Gamma\ge\frac12(3\mathfrak g+4\mathfrak r+\mathfrak v)=\frac12(1+\varsigma-\lambda)$, so $U>\frac12-\frac32\varsigma$. Since $\tau\le\varsigma$ and $\varsigma<\frac1{13}$,
\[
 U-4\tau'>\tfrac12-\tfrac{11}2\varsigma-8\delta>\tfrac1{13}-8\delta>0 .
\]
This proves the claim.

Now if some $\mathcal I$ has $\sum_{i\in\mathcal I}\nu_i\in(2\tau',U)$, then its complement has degree above $D-U=\Lambda$, which is Case 1. Otherwise every partial sum of the $\nu_i$ lies in $[0,2\tau']\cup[U,D]$. Since $D>8\tau'$, some $\nu_i$, namely $\nu_1$, is at least $U$, and at most two are, as $3U>D$. If only $\nu_1\ge U$, then $\nu_2,\nu_3,\nu_4\le2\tau'$; the partial sums $\nu_2$, $\nu_2+\nu_3$, $\nu_2+\nu_3+\nu_4$ increase by at most $2\tau'$ at each step and cannot enter $(2\tau',U)$, which has length greater than $2\tau'$; so $\nu_2+\nu_3+\nu_4\le2\tau'$, which is Case 2. If $\nu_1,\nu_2\ge U$, then $\nu_3+\nu_4\le D-2U=2\Delta<U$, so $\nu_3+\nu_4\le2\tau'$, which is Case 3.
\end{proof}
\noindent The case analysis has also been tested on random parameters: of $200{,}000$ parameter draws, $30{,}575$ satisfy the selection guards, and for each of these some case gives a saving larger than $\tau$. This is a sanity check; the proof is the argument above.

\begin{proof}[Proof of \cref{prop:nonprincipal}]
Fix $(h,k)$ and $q=hk$ as above, and a term $(g,d,r,\ell,\mathfrak a)$ of \cref{lem:dualsum}. Insert smooth dyadic partitions of unity $1=\sum_{N_i}W(n_i/N_i)$ for the four variables of $d_{-\boldsymbol\alpha}(n)=\sum_{n_1\cdots n_4=n}\prod_in_i^{\alpha_i}$. For a fixed quadruple put $N_*=\prod N_i$; then $n=n_1\cdots n_4\in[N_*/16,16N_*]$, and $n^{-1}V_{\mathfrak a}(n/Y)=\vartheta(n/N_*)$, where $\vartheta(y)=\bar W(y)(yN_*)^{-1}V_{\mathfrak a}(yN_*/Y)$ for a fixed $\bar W\in C_c^\infty((\frac1{32},32))$ equal to $1$ on $[\frac1{16},16]$. By Mellin inversion $\vartheta(y)=\frac1{2\pi}\int\widehat\vartheta(t)y^{-it}\dd t$, and integration by parts and \cref{lem:dualweight} give
\[
 |\widehat\vartheta(t)|\ll_{j,A}(1+|t|)^{-j}N_*^{-1}\min\big((N_*/Y)^{9/10},(N_*/Y)^{-A}\big).
\]
The part $|t|>q^\epsilon$ is negligible by the trivial bound. For $|t|\le q^\epsilon$ the factor $y^{-it}=\prod_i(n_i/N_i)^{-it}$ separates the variables, and the weights $w_i(y)=W(y)(yN_i)^{\alpha_i}y^{-it}$ satisfy the hypotheses of the previous subsection. If $N_*>Yq^{\epsilon_0}$ the factor $(N_*/Y)^{-A}$ makes the quadruple negligible. Otherwise \cref{lem:block} applies, and $\min\big((N_*/Y)^{9/10},(N_*/Y)^{-A}\big)\le q^{-\frac9{10}\mathfrak v}$: if $N_*\le Y$ the first term equals $q^{-\frac9{10}\mathfrak v}$, and if $N_*>Y$ then $\mathfrak v=0$ and the second term is less than $1$. Summing over the $O(L^4)$ quadruples, the $n$-sum in \cref{lem:dualsum} is $\ll q^{-\frac9{10}\mathfrak v-\tau-\delta+o(1)}\le q^{-\varsigma+2\mathfrak g+3\mathfrak r+\frac12\mathfrak p-\delta+o(1)}$.

The prefactor is $d^2\varphi(d)/\varphi(dr)=d^2/\varphi(r)\le q^{2-2\mathfrak g-3\mathfrak r+o(1)}$, since $d=q/(gr)$. The number of triples $(g,d,r)$, the sum $\sum_\ell|b_\ell|$ and $|\varpi_d|$ are $q^{o(1)}$. Hence
\[
 |\mathcal S^\flat(q)|\ll q^{2-\varsigma+\frac12\mathfrak p-\delta+o(1)}=T\,P^{1/2}q^{-\delta+o(1)} .
\]
Finally $q\ge T^{1/2}$, and $\sum_{q\le Z}P(q)^{1/2}q^{-1}\le\sum_{P\le Z}P^{1/2}\cdot\frac{1}{P}\sum_{m\le Z/P}\frac1m\ll(\log Z)^2$, since there are $O(\sqrt x)$ powerful numbers up to $x$. Summing over $(h,k)$, with at most $\tau(q)$ pairs for each $q\le T$, gives $\ll T^{1-\delta/2+o(1)}$.
\end{proof}

\subsection{The factorable row}
\begin{proposition}[Factorable pairing]\label{prop:factorable}
Let $c\in C_c^\infty(\R^3)$ be real with $\operatorname{supp}c\subset\mathcal F^\sharp$, and let $b\in C_c^\infty(\R^3)$. Uniformly for $\Phi$ in bounded subsets of $C_c^\infty$ with a fixed compact support in $(1,2)$,
\[
 \ip{P_b}{P_c}_\Phi=a_3TL^9\Big(\int\Phi\Big)\int_\Om h_bh_c+o(TL^9).
\]
\end{proposition}
\begin{proof}
The reduction to elementary sources in the proof of \cref{prop:balanced} uses only \cref{lem:localiseP}, \cref{lem:sourcenorm}, \cref{lem:sixth} and the fact that $h_c$ vanishes unless $r\in I_c$; so it suffices to prove that for $\boldsymbol u\in(i\R)^3$ with distinct coordinates and $\phi$ as there,
\begin{equation}\label{eq:factelementary}
 \ip{P_{\boldsymbol u,\phi}}{P_c}_\Phi=a_3TL^9\Big(\int\Phi\Big)\int_\Om h_{e^{-\boldsymbol u\cdot a}}h_c+o(TL^9).
\end{equation}
Let $Z(s)=\prod_{k=1}^3\zeta(s+u_k/L)$. By \cref{lem:expansion}, $Z(s)=\sum_nb_{\boldsymbol u}(n)\omega(n/T^8)n^{-s}+O(T^{-A})$, and the coefficients of this polynomial minus $P_{\boldsymbol u,\phi}$ vanish for $\log n/L$ near $I_c$; so $\ip{P_{\boldsymbol u,\phi}}{P_c}_\Phi=\ip Z{P_c}_\Phi+O(T^{-A})$ by \cref{lem:localiseP}.

\emph{Fourier inversion in the third coordinate.} For $\beta\in i\R$ and $x\in\R^2$ let, as in \cref{lem:balkernel},
\[
 \mathcal L_3c(x;\beta)=\int c(x_1,x_2,a_3)e^{\beta(1-a_3)}\dd a_3 .
\]
By Fourier inversion $c(x,a_3)=\frac1{2\pi i}\int_{(0)}\mathcal L_3c(x;\beta)e^{-\beta(1-a_3)}\dd\beta$;
also $|\mathcal L_3c(x;\beta)|\ll_j(1+|\beta|)^{-j}$; moreover the function $F_\beta=\overline{\mathcal L_3c(\cdot;\beta)}$ on $\R^2$ is supported in the compact projection $E$ of $\operatorname{supp}c$, which lies in $\mathcal P$, and its Fourier transform satisfies $\int|\widehat{F_\beta}|\ll_j(1+|\beta|)^{-j}$. On the critical line $\overline{n^{-s}}=n^{-(1-s)}$, and $c$ is real, so
\[
 \overline{P_c(s)}=\frac1{2\pi i}\int_{(0)}e^{-\beta}\,\overline{A_{F_\beta}(s)}\sum_m\omega(m/T^8)m^{-(1-s-\beta/L)}\dd\beta,
\]
the $m$-sum being finite. By \cref{lem:expansion}, whose proof is uniform for shifts of size at most $T^{1/2}$, the $m$-sum is $\zeta(1-s-\beta/L)+O(T^{-A})$ for $|\beta|\le T^{1/2}$, and the rest of the $\beta$-line is negligible. By the functional equation $\zeta(1-s-\beta/L)=\chi(1-s-\beta/L)\zeta(s+\beta/L)$, and by Stirling's formula
\[
 e^{-\beta}\chi(1-s-\beta/L)=\chi(1-s)\big(1+O((1+|\beta|)/L)\big)\qquad(T\le t\le2T,\ |\beta|\le L).
\]
Put $\boldsymbol\sigma(\beta)=(u_1,u_2,u_3,\beta)$. We obtain
\[
 \ip Z{P_c}_\Phi=\frac1{2\pi i}\int_{(0)}\Big(I_{F_\beta}(\boldsymbol\sigma(\beta))+\mathcal E(\beta)\Big)\dd\beta+O(T^{-A}),
\]
where $I_F$ is as in \cref{prop:fourshift} and, by Cauchy--Schwarz, $|\mathcal E(\beta)|\ll\frac{1+|\beta|}L\|Z\|_\Phi\|\zeta(s+\beta/L)A_{F_\beta}\|_\Phi$ for $|\beta|\le L$; for $|\beta|>L$ we use $|e^{-\beta}\chi(1-s-\beta/L)/\chi(1-s)|\ll1$ instead.

\emph{Norms.} By H\"older's inequality, Harper's bound and \cref{lem:sixth}(i) (applied to the Fourier superpositions $A_{F_\beta}=(2\pi)^{-2}\int\widehat{F_\beta}(\xi)A_{g_{\xi_1}}A_{g_{\xi_2}}\dd\xi$ as in \cref{lem:sixth}(ii)), $\|Z\|_\Phi\ll(TL^9)^{1/2}$ and $\|\zeta(s+\beta/L)A_{F_\beta}\|_\Phi\ll_j(TL^9)^{1/2}(1+|\beta|)^{-j}$ for $|\beta|\le LT/4$; the same bounds give $|I_{F_\beta}(\boldsymbol\sigma(\beta))|\ll_j TL^9(1+|\beta|)^{-j}$ there, and the trivial bound with the decay of $\mathcal L_3c$ handles larger $|\beta|$. Hence $\int|\mathcal E|=o(TL^9)$.

\emph{The limit.} For each fixed $\beta\in i\R$, \cref{prop:fourshift} gives $I_{F_\beta}(\boldsymbol\sigma(\beta))=a_3TL^9(\int\Phi)\int\mathcal L_3c(x;\beta)S(\boldsymbol\sigma(\beta);x)\dd x+o(TL^9)$. By dominated convergence (with the majorant $TL^9(1+|\beta|)^{-j}$ from the norm bounds),
\[
 \ip Z{P_c}_\Phi=a_3TL^9\Big(\int\Phi\Big)\frac1{2\pi i}\int_{(0)}\int_{\R^2}\mathcal L_3c(x;\beta)S(\boldsymbol\sigma(\beta);x)\dd x\dd\beta+o(TL^9),
\]
which makes sense also when $\int\Phi=0$.
The integrand is entire in $\beta$. For $0\le\Re\beta\le2r_0$ the function $S(\boldsymbol\sigma(\beta);x)$ is bounded polynomially in $|\beta|$, uniformly for $x$ in a compact set: away from the points $\beta=u_k$ this follows from the explicit formula for $S$, as in the proof of \cref{lem:balkernel}, and near them from the maximum principle. Since $\mathcal L_3c(x;\beta)$ decays rapidly in this strip, we may move the line to $\Re\beta=2r_0$, and the double integral becomes $\Lambda'$ in the proof of \cref{lem:balkernel}. Part (b) of that proof uses only that $c\in C_c^\infty((0,1)^3)$ is real with $a_1+a_2<1$ on its support, which holds here, and gives $\Lambda'=\int_\Om h_{e^{-\boldsymbol u\cdot a}}h_c$. This proves \eqref{eq:factelementary}.
\end{proof}

\subsection{Proof of \Cref{thm:factorable}}
Fix $\Phi\in\mathcal W$. For real $\psi\in C_c^\infty(\mathcal F)$, \cref{lem:synthesis} gives $h_{\mathsf V\psi}=\psi$, \cref{lem:sixth}(ii) gives $\|F_\psi\|_\Phi\ll(TL^9)^{1/2}$, and the coefficients of $F_\psi$ live at $\log N/L\in I_\psi\subset(1,\frac32)$. We go through the proof of \Cref{thm:row}, parts (a)--(d), with $\mathcal R$, $B_\psi$ and \cref{prop:balanced} replaced by $\mathcal F$, $F_\psi$ and \cref{prop:factorable}. Every step applies verbatim: part (a) uses only the pairing of the witness with a compact source; part (b) uses, for the reflected entries, \cref{lem:reflcube}, which requires $\operatorname{supp}\mathsf V\psi\subset(0,1)^3$, and \cref{lem:localiseP} with the disjoint trace ranges $I_\eta$ and $3-I_\psi\subset(\frac32,2)$; part (c) uses \cref{lem:factors} through the same three properties, and for $\mathcal JD_\phi$ the bound $\theta+\max I_\psi<3$; and part (d) uses that $I_\psi\subset(1,2)$, as \cref{prop:mixed} requires.

It remains to treat the cross-pairings between $\mathcal T_{\mathcal R}$ and $\mathcal T_{\mathcal F}$. For real $\psi\in C_c^\infty(\mathcal F)$ and $\eta\in C_c^\infty(\mathcal R)$, \cref{prop:balanced} with $c=\mathsf V\eta$ and $b=\mathsf V\psi$ gives $\mathrm G_T(F_\psi,B_\eta)/\mathcal N_T\to\ip\psi\eta$. By \cref{lem:reflcube}, $\mathrm G_T(\mathcal JF_\psi,B_\eta)=\ip{P_{(\mathsf V\psi)^J}}{B_\eta}_\Phi+o(TL^9)$, and \cref{prop:balanced} with $b=(\mathsf V\psi)^J$ and \cref{lem:reflectedprofile} give $\mathrm G_T(\mathcal JF_\psi,B_\eta)/\mathcal N_T\to\ip{\psi\circ J}\eta$ (which is $0$). Likewise $\mathrm G_T(\mathcal JB_\eta,F_\psi)/\mathcal N_T\to\ip{\eta\circ J}\psi$ by \cref{lem:reflcube} and \cref{prop:factorable} with $b=(\mathsf V\eta)^J$. The remaining entries are complex conjugates of these, by $\ip{\mathcal JH}{\mathcal JK}_\Phi=\overline{\ip HK_\Phi}$ and $\ip H{\mathcal JK}_\Phi=\overline{\ip{\mathcal JH}K_\Phi}$, and all the limits are real. Hence $\mathcal T_6$ is faithful.

\emph{The span.} As in \Cref{thm:row}, $\mathcal H(\mathcal T_6)\supseteq L^2(\mathcal F\cup J\mathcal F)$, and with $v\one_{S_3\cup\mathcal Z\cup\mathcal R\cup J\mathcal R}\in\mathcal H(\mathcal T_5)$ we get $v\one_{U_6}\in\mathcal H(\mathcal T_6)$.

\emph{The measure.} Since $\mathcal F\subset\{1<r<\frac32\}$, we have $\mathcal F\cap JN_0=\emptyset=J\mathcal F\cap N_0$, and as in the proof of \Cref{thm:row} the complement of $U_6$ is, up to a null set, $(N_0\setminus(\mathcal R\cup\mathcal F))\cup J(N_0\setminus(\mathcal R\cup\mathcal F))$. So $\mu(U_6)=C_0+2\mu((N_0\cap\mathcal F)\setminus\mathcal R)$. We describe this set in the coordinates $(k,x,z)$, $k=x+y$, which have Jacobian one. On $k>\frac12$ the inequality $21k+7r<18$ forces $r<\frac{15}{14}$, so outside $\mathcal R$ means $z\ge z_0(k):=\frac{18-28k}7$. Together with $z-x<\frac12$ and the conditions defining $\mathcal F$, the set is
\begin{gather*}
 \tfrac12<k<\tfrac{13}{25},\qquad x_{\min}(k)<x<\tfrac k2-\tfrac32(2k-1),\qquad z_0(k)<z<\tfrac12+x,\\
 x_{\min}(k)=\max\big(5(2k-1),\,z_0(k)-\tfrac12\big).
\end{gather*}
The remaining conditions are implied: $y=k-x<5-9k<z_0(k)<z$, as $k>\frac{17}{35}$; $z<\frac12+x<1$; $r>k+z_0(k)=\frac{18-21k}7>1$, as $k<\frac{11}{21}$; and $r<k+\frac12+x<\frac32$, as $x\le\frac k2$ and $k<\frac23$. The two lower bounds for $x$ cross at $k=\frac{99}{196}$. Since $9!v^2=\frac{9!}4(k-2x)^2(z-x)^2(z-k+x)^2$,
\[
 \begin{split}
 &2\mu\big((N_0\cap\mathcal F)\setminus\mathcal R\big)\\
 &\qquad=\frac{9!}2\int_{1/2}^{13/25}\int_{x_{\min}(k)}^{\frac k2-\frac32(2k-1)}\int_{z_0(k)}^{\frac12+x}(k-2x)^2(z-x)^2(z-k+x)^2\dd z\dd x\dd k ,
 \end{split}
\]
a sum of two polynomial integrals over rational limits, equal to $\frac{840731426179152778718787}{52986177566328125000000000}$. This value is computed by exact polynomial integration and, independently, by an exact quadrature rule; the second computation also checks the inclusions above at the vertices of the two cells. By \cref{prop:projection}, $\liminf_{T\to\infty}M_6(T)/(c_3TL^9)\ge\mu(U_6)$. \qed

\section{The mixed row up to $\kappa<\frac35$}\label{sec:mixedrow}

The region $\mathcal F$ of \Cref{sec:factorable} is limited by \cref{prop:nonprincipal}, which bounds the non-principal part of the four-shift row modulus by modulus. In this section we keep the sum over the moduli $q=hk$ with its signs and its weights. After a double Poisson transform the dual sum of \cref{lem:dualsum} becomes a sum of Kloosterman fractions, and we bound it in four ways: by the bound of Deshouillers and Iwaniec for such sums averaged over the modulus, by the same bound with the two marked factors of $q$ kept as summation variables, by completing a reciprocal frequency with Weil's bound, and by the Kuznetsov formula at the level fixed by the denominators. Each of these bounds is an affine inequality between the degrees of the lengths involved, and the remaining problem is to show that for every configuration one of them saves a power. For the two smallest ranges of $\varsigma$ we do this in prose (\cref{lem:cases}(i)--(ii)), and check the prose by machine; for the three larger bands we give exact rational certificates.

For $\frac12<k<\frac35$ put
\begin{equation}\label{eq:xplus}
 x_+(k)=\begin{cases}
 \frac k2, & \frac12<k<\frac{10}{19},\\[2pt]
 \max\big(\frac{21-39k}2,\frac{7-12k}4\big), & \frac{10}{19}\le k<\frac7{13},\\[2pt]
 \frac{2-3k}4, & \frac7{13}\le k<\frac47,\\[2pt]
 \frac{3-5k}2, & \frac47\le k<\frac35,
 \end{cases}
\end{equation}
and
\begin{gather*}
 \mathcal P^+=\Big\{(a_1,a_2)\in(0,1)^2:\ \tfrac12<k<\tfrac35,\ a_1<x_+(k)\Big\},\qquad k=a_1+a_2,\\
 \mathcal F^+=\{(x,y,z)\in\Om:\ 1<r<\tfrac32,\ (x,y)\in\mathcal P^+\},\\
 \mathcal F^{+\sharp}=\{a\in(0,1)^3:\ 1<|a|<\tfrac32,\ (a_1,a_2)\in\mathcal P^+\}.
\end{gather*}
Since $x_+(k)\le\frac k2$, we have $a_1<a_2$ on $\mathcal P^+$; the set $\mathcal P^+$ is open, because $x_+$ is continuous except at $k=\frac{10}{19}$ and $k=\frac7{13}$, where it jumps down; and $\mathcal P\subset\mathcal P^+$. In the coordinates $\varsigma=2-1/k$, $\mathfrak b=a_1/k$ of \Cref{sec:factorable}, $\mathcal P^+$ becomes
$\mathcal P^{+*}=\{0<\varsigma<\frac13,\ 0<\mathfrak b<\mathfrak b_+(\varsigma)\}$, where
\[
 \mathfrak b_+(\varsigma)=\begin{cases}
 \frac12, & 0<\varsigma<\frac1{10},\\[2pt]
 \max\big(\frac{3-21\varsigma}2,\frac{2-7\varsigma}4\big), & \frac1{10}\le\varsigma<\frac17,\\[2pt]
 \frac{1-2\varsigma}4, & \frac17\le\varsigma<\frac14,\\[2pt]
 \frac{1-3\varsigma}2, & \frac14\le\varsigma<\frac13.
 \end{cases}
\]
The first entry of the maximum is the larger one exactly for $\varsigma\le\frac4{35}$, that is $k\le\frac{35}{66}$. As in \Cref{sec:factorable}, $\mathcal F^+=\mathcal F^{+\sharp}\cap\Om$ and $\mathcal F^+\cap J\mathcal F^+=\emptyset$, and for real $\psi\in C_c^\infty(\mathcal F^+)$ we put $F_\psi=P_{\mathsf V\psi}$.

\begin{theorem}[Mixed row]\label{thm:mixedrow}
Let
\[
 \mathcal T_{\mathcal F^+}=\big\{(F_\psi,\psi),\ (\mathcal JF_\psi,\psi\circ J):\ \psi\in C_c^\infty(\mathcal F^+)\text{ real}\big\},\qquad\mathcal T_7=\mathcal T_5\cup\mathcal T_{\mathcal F^+}.
\]
Then $\mathcal T_7\supseteq\mathcal T_6$ is faithful, $v\one_{U_7}\in\mathcal H(\mathcal T_7)$ with $U_7=S_3\cup\mathcal Z\cup\mathcal R\cup J\mathcal R\cup\mathcal F^+\cup J\mathcal F^+$, and
\begin{align*}
 \mu(U_7)&=C_0+2\mu\big((N_0\cap\mathcal F^+)\setminus\mathcal R\big)\\
 &=\frac{112666832560750258461501594014215737455802061}{2705696476246250148185796472625838239963136}\\
 &=41.64060290940640373015\ldots .
\end{align*}
\end{theorem}

\emph{Overview.} The proof of \Cref{thm:factorable} uses the shape of $\mathcal P$ only through \cref{prop:fourshift}, and \cref{prop:fourshift} uses it only through \cref{prop:nonprincipal}. We prove the corresponding signed estimate on $\mathcal P^+$ (\cref{prop:nonprincipalplus}) and then repeat the remaining steps with $\mathcal F^+$ in place of $\mathcal F$ (\cref{prop:factorableplus} and \Cref{ssec:mixedrowproof}). The proof of \cref{prop:nonprincipalplus} follows that of \cref{prop:nonprincipal} until the dual sum has been separated into four smooth variables, but keeps the sum over the moduli. If one variable is long, or the variables can be split into two long groups, the estimates of \cref{lem:unsignedplus} apply at each modulus. Otherwise two variables are long, and we transform both by Poisson summation (\cref{lem:doublepoisson}). The result is a sum of Kloosterman fractions $e_{u'}(\pm r\,y_0\,\overline{F_0\,e\,d_0\,b})$ over a reduced modulus $u'$, which we bound in the four ways described above (\cref{lem:square,lem:direct,lem:recip,lem:actuallevel}). The results quoted from the literature for these four bounds are collected, in the form used, in \Cref{ssec:quoted}.

\subsection{Quoted estimates}\label{ssec:quoted}
Besides the results already used in \Cref{sec:marked,sec:factorable}, this section uses the following four statements, which we quote in the form and normalisation in which they are applied. Throughout, $\overline x$ denotes an inverse of $x$ modulo the relevant modulus, $e_c(x)=e(x/c)$, and $S(m,n;c)=\sum^*_{x\bmod c}e_c(mx+n\overline x)$.

\emph{(Q1) Sums of Kloosterman fractions} \cite[Theorem~12]{DI}, with the third term as corrected in \cite[Lemma~7.5]{Lichtman}. Let $C,D,N,R,S\ge1$, let $g$ be a smooth function supported in $[C,2C]\times[D,2D]$ with $\partial_c^{\nu_1}\partial_d^{\nu_2}g(c,d)\ll C^{-\nu_1}D^{-\nu_2}$ for $0\le\nu_1,\nu_2\le2$, and let $b_{n,r,s}$ be complex numbers supported on $N\le n\le2N$, $R\le r\le2R$, $S\le s\le2S$. Then
\begin{equation}\label{eq:DI}
 \sum_c\sum_d\sum_{\substack{n,r,s\\(rd,sc)=1}}b_{n,r,s}\,g(c,d)\,e\Big(\frac{n\overline{rd}}{cs}\Big)\ll_\epsilon(CDNRS)^\epsilon\,\mathcal K\,\|b\|_2,
\end{equation}
where $\overline{rd}$ is the inverse modulo $cs$, $\|b\|_2^2=\sum|b_{n,r,s}|^2$, and
\begin{equation}\label{eq:DIK}
 \mathcal K^2=CS(RS+N)(C+DR)+C^2DS\sqrt{(RS+N)R}+D^2NR .
\end{equation}
In \cite[Theorem~12]{DI}, and in its extension to moduli in arithmetic progressions \cite[Theorem~2.1]{Drappeau}, the third term reads $D^2NRS^{-1}$; \cite[Lemma~7.5]{Lichtman} corrects it to $D^2NR$. The two agree for $S=1$, which is the case used in \cref{lem:square}. The right side of \eqref{eq:DI} is nondecreasing in $C,D,N,R,S$; this uses the weaker third term $D^2NR$, since $D^2NRS^{-1}$ is not nondecreasing in $S$. In particular it makes no difference whether the dyadic ranges are taken closed or half-open, and after a dyadic decomposition the same bound holds, with a factor $(\log CDNRS)^3$, when the coefficients are supported in $[1,2N]\times[1,2R]\times[1,2S]$. Negative $n$ (or $d$) are reduced to positive ones by complex conjugation (or by $d\mapsto-d$, $n\mapsto-n$).

\emph{(Q2) The spectral large sieve} \cite[Theorem~2]{DI}, in the normalisation of \eqref{eq:kuznetsov}. Let $A\ge1$, let $\mathfrak c$ be a cusp of $\Gamma_0(A)$ equivalent to $\infty$ or to $1/m$ with $m\mid A$ and $(m,A/m)=1$ (so that $\mu(\mathfrak c)=1/A$ in the notation of \cite{DI}), let $U,N\ge1$, and let $(a_n)_{N<n\le2N}$ be complex. Then
\begin{equation}\label{eq:LS}
 \sum_{|t_j|\le U}\Big|\sum_na_n\rho_{j\mathfrak c}(n)\Big|^2\ll_\epsilon(ANU)^\epsilon\Big(U^2+\frac NA\Big)\sum_n|a_n|^2,
\end{equation}
where $\rho_{j\mathfrak c}=\rho^{\rm raw}_{j\mathfrak c}/\sqrt{\cosh(\pi t_j)}$ for an orthonormal basis $\{u_j\}$ of Maass cusp forms of $\Gamma_0(A)$ (Petersson product not divided by the volume), the exceptional forms included; and the same bound holds for the holomorphic forms of weight at most $U$ with the coefficients $\psi_{jk\mathfrak c}$ of \eqref{eq:kuznetsov}, and for the Eisenstein series with $\frac1{4\pi}\sum_{\mathfrak c'}\int_{-U}^U$. In \cite[Theorem~2]{DI} the raw coefficients carry the weight $1/\cosh(\pi t_j)$, which is our normalisation; at an exceptional parameter $t_j=i\tau_j$ this weight lies in $[1,1.07]$ by the Kim--Sarnak bound $\tau_j\le\frac7{64}$ \cite{KimSarnak}. The same inequality, with the same proof, was used in the proof of \cref{prop:moments}. For indices $-2N\le n<-N$ apply \eqref{eq:LS} to the basis $u_j(-\bar z)$, which is again orthonormal; the reflection maps the cusps $\infty$ and $1/m$ to equivalent cusps and changes the scaling matrices by a translation, which multiplies $\rho_{j\mathfrak c}(n)$ by a unimodular factor depending only on $n$.

\emph{(Q3) Exceptional spectrum with a weight}, from \cite[Lemma~3.9 and Corollary~3.10]{Pascadi}. For $Y\ge1$ and $(a_n)_{N<n\le2N}$, with the normalisation of (Q2) at the cusp $\infty$ of $\Gamma_0(A)$,
\begin{equation}\label{eq:EX}
 \sum_{t_j=i\tau_j\ \mathrm{exc}}Y^{2\tau_j}\Big|\sum_na_n\rho_{j\infty}(n)\Big|^2\ll_\epsilon(ANY)^\epsilon\Big(1+\frac NA+\frac{N^{3/2}Y^{1/2}}A\Big)\sum_n|a_n|^2 .
\end{equation}
This is not stated as such in \cite{Pascadi}; it follows from \cite[Corollary~3.10]{Pascadi} by the argument sketched in \cite[\S2]{Pascadi}. With $X=Y$ and the cusp $\infty$ of $\Gamma_0(A)$, for which $\mu(\infty)=1/A$, that corollary bounds the left side of \eqref{eq:EX} by a constant times
\[
 \sum_{c\equiv0\,(A)}\frac1c\sum_{m,n}|a_ma_n|\,|S(m,n;c)|\,\Phi\Big(\frac{Y\sqrt{mn}}c\Big)+O_\epsilon\Big(\Big(1+\frac{N^{1+\epsilon}}A\Big)\|a\|_2^2\Big),
\]
where $\Phi\ge0$ is a fixed smooth function supported on $t\asymp1$. Here the moduli of the cusp $\infty$ are the $c\equiv0\pmod A$ \cite[Lemma~3.2]{Pascadi}; the regular spectrum, which gives the second term, is bounded by \cite[Lemma~3.9]{Pascadi}, that is by \cite[Theorem~2]{DI}; and the test function in the proof of the corollary is $\phi(y)=\Phi(yY/4\pi)$, so that only the moduli $c\asymp NY$ occur. (The printed statement of \cite[Corollary~3.10]{Pascadi} has $\sqrt{mn}/c$ as the argument of $\Phi$; the factor $Y$ comes from this choice of $\phi$, and the bound below is the one obtained in \cite[\S2]{Pascadi}.) The coefficients in \cite{Pascadi} are those of \cite{DI}, without the factor $\cosh(\pi t_j)^{-1/2}$, which at an exceptional parameter lies in $[1,1.07]$ as noted in (Q2). For each modulus, Weil's bound (Q4) and $\sum_{m,n}|a_ma_n|(m,c)^{1/2}\ll c^\epsilon N\|a\|_2^2$ bound the double sum of Kloosterman sums by $c^{1/2+\epsilon}N\|a\|_2^2$, and summing over the $O(NY/A)$ moduli $c\equiv0\pmod A$ with $c\asymp NY$, each with the factor $c^{-1}$, gives $(ANY)^\epsilon N^{3/2}Y^{1/2}A^{-1}\|a\|_2^2$. We use \eqref{eq:EX} only with $Y=A^2/N^3\ge1$, where its right side is $\ll(ANY)^\epsilon\|a\|_2^2$. As in (Q2), it holds also at a cusp $1/m$ with $m\mid A$, $(m,A/m)=1$, since by \cref{lem:ALcusp} the coefficients of $u_j$ at this cusp are those of $u_j|W_{A/m}$ at $\infty$, and $\{u_j|W_{A/m}\}$ is an orthonormal basis of the same eigenspaces.

\emph{(Q4) Weil's bound} \cite[Corollary~11.12]{IK}: $|S(m,n;c)|\le\tau(c)(m,n,c)^{1/2}c^{1/2}$ for all integers $m,n$ and $c\ge1$; for $m=0$ this includes the Ramanujan sum.

We also use the Kuznetsov formula for the cusp pair $(\infty,\mathfrak a)$ exactly as in \eqref{eq:kuznetsov}, with arbitrary nonzero integers $m',n'$ in place of $1,\pm n$:
\begin{equation}\label{eq:kuzgeneral}
 \begin{split}
 \sum_\gamma\frac{S_{\infty\mathfrak a}(m',n';\gamma)}\gamma\,\phi\Big(\frac{4\pi\sqrt{|m'n'|}}\gamma\Big)={}&\sum_j\overline{\rho_{j\infty}(m')}\rho_{j\mathfrak a}(n')\widetilde\phi(t_j)\\
 &+\one_{m'n'>0}\sum_{k,j}\widetilde\phi(k)\overline{\psi_{jk\infty}(m')}\psi_{jk\mathfrak a}(n')\\
 &+\frac1{4\pi}\sum_{\mathfrak c}\int\overline{\rho_{\mathfrak c\infty}(m',t)}\rho_{\mathfrak c\mathfrak a}(n',t)\widetilde\phi(t)\dd t,
 \end{split}
\end{equation}
where $\gamma$ runs over the moduli of the pair, the transforms are those of \cite[Theorem~1]{DI} for the sign of $m'n'$, and for $\phi$ supported on $\xi\asymp\xi_0\le1$ with $\xi^i\phi^{(i)}(\xi)\ll_i\Phi_0$ they satisfy \eqref{eq:transforms} with $X^\epsilon Q\sqrt A$ replaced by $(\Phi_0/\xi_0)^\epsilon\Phi_0$ and $K_B$ by $\xi_0^{-1}$. In this direction (Kloosterman sums to spectrum) the formula has no diagonal term, whether or not the two cusps are equivalent. The proof of \cref{lem:ALcusp} gives, for all integers $m',n'$,
\begin{equation}\label{eq:ALgeneral}
 S_{\infty\mathfrak a}(m',n';c\sqrt a)=S(m',n'\overline a;c)\qquad(m\mid c,\ (c,a)=1),
\end{equation}
since there the Kloosterman phase is $e((m'x+n'y)/c)$ with $y\equiv\overline a\,\overline x\pmod c$ and $x$ running once over $(\Z/c\Z)^*$.

\subsection{The four-shift row on $\mathcal P^+$}

\begin{proposition}[Four-shift row]\label{prop:fourshiftplus}
The conclusion of \cref{prop:fourshift} holds for every $F\in C_c^\infty(\R^2)$ with support in a compact subset of $\mathcal P^+$.
\end{proposition}
\begin{proof}
We follow the proof of \cref{prop:fourshift}, with $\kappa_+=\max\{a_1+a_2:\ a\in\operatorname{supp}F\}<\frac35$ in place of $\max\{a_1+a_2\}<\frac{13}{25}$. The stationary-phase formula \eqref{eq:saddle} uses only $\kappa_+<1$, so it holds with error $O(T^{\kappa_++\epsilon})$, which is $O(T^{\frac35})$ for $\epsilon<\frac35-\kappa_+$, and its main term is $2\pi\sum_{h,k}\overline{F(\log h/L,\log k/L)}(hk)^{-1}\mathcal S(hk)$. In the character expansion \eqref{eq:charexp} we treat the three parts as there. For the principal characters off the poles, the new line contributes $\ll T^{\frac12+\delta}N^{2\epsilon}$ for each $N=hk$, as shown in the proof of \cref{prop:fourshift}; with the weight $1/N$ and $N\le T^{\kappa_+}$ the total is $O(T^{\frac12+2\delta})$, whatever $\kappa_+<1$. The residue computation does not use the size of $\kappa_+$. The non-principal characters contribute
\[
 2\pi\sum_{h,k\ge1}\frac{\overline{F(\log h/L,\log k/L)}}{hk}\,\mathcal S^\flat(hk)\ll T^{1-\eta}
\]
by \cref{prop:nonprincipalplus}.
\end{proof}

\begin{proposition}\label{prop:nonprincipalplus}
Let $F\in C_c^\infty(\R^2)$ have support in a compact set $E\subset\mathcal P^+$. There is $\eta>0$ such that, uniformly for $\boldsymbol\sigma$ in compact subsets of $(i\R)^4$ and $\Phi$ in bounded subsets of $\mathcal W$,
\[
 \Sigma_F:=\sum_{h,k\ge1}\frac{\overline{F(\log h/L,\log k/L)}}{hk}\,\mathcal S^\flat(hk)\ll T^{1-\eta}.
\]
\end{proposition}

Unlike \cref{prop:nonprincipal}, this bounds a signed sum: in the cases of \cref{lem:direct,lem:actuallevel} the cancellation comes from the sum over the moduli, with the weight $F$ evaluated at the factorisation $q=hk$. The proof occupies the rest of this subsection and is completed after \cref{lem:cases}.

\subsubsection*{The signed sum and its conductor labels}
Insert \cref{lem:dualsum} for each $q=hk$, and note that $d^2\varphi(d)/\varphi(dr)=d^2/\varphi(r)$ because $(d,r)=1$. We first record how the terms depend on the primitive conductor $d$.

\begin{lemma}[Conductor labels]\label{lem:labels}
In \cref{lem:dualsum}, $b_\ell(g,d,r)=b^\circ_\ell(g,r)\one_{(\ell,d)=1}$, where $b^\circ_\ell(g,r)$ depends only on $g$, $r$ and $\boldsymbol\alpha$, vanishes unless $\ell\mid g^3r^4$, and satisfies $\sum_\ell|b^\circ_\ell(g,r)|\ll_\epsilon(gr)^\epsilon$. With $q=gdr$ and $X=qT/(2\pi)$,
\[
 Y=\frac{g\ell d^4}{\pi^4X}=\frac{2\ell d^3}{\pi^3rT},\qquad\Big|\frac{d^2\varpi_d}{\varphi(r)}\Big|=\frac{d^2}{\pi^2\varphi(r)},\qquad\varpi_d=\pi^{-2+\sum\alpha_i}d^{-\sum\alpha_i}.
\]
\end{lemma}
\begin{proof}
In the proof of \cref{lem:dualsum}, for each prime $p\mid gr$ let $v=v_p(g)$ and put
\[
 P_p(X)=\begin{cases}d_{\boldsymbol\alpha}(p^v)\prod_{i=1}^4(1-p^{-\alpha_i}X), & p\mid r,\\[2pt]
 \prod_{i=1}^4(1-p^{-\alpha_i}X)\sum_{j\ge0}d_{\boldsymbol\alpha}(p^{v+j})X^j, & p\nmid r,\end{cases}
\]
a polynomial of degree at most $4$ (if $p\mid r$) or $3$ (if $p\nmid r$), whose coefficients depend only on $p$, $v$ and $\boldsymbol\alpha$. The local factor of $D(w)$ divided by that of $\prod_iL(w+\alpha_i,\psi)$ is $P_p(\psi(p)p^{-w})$ for every $p\mid gr$: for $p\mid r$ and for $p\mid g$, $p\nmid q'$ this is the computation in that proof, and for $p\mid g$, $p\mid d$ we have $\psi(p)=0$, the local factor of the $L$-functions is $1$, and $P_p(0)=d_{\boldsymbol\alpha}(p^v)$ is the constant local factor of $D(w)$ found there. At $p\nmid gr$ the two local factors agree. Hence $\sum_\ell b_\ell\psi(\ell)\ell^{-w}=\prod_{p\mid gr}P_p(\psi(p)p^{-w})=\sum_\ell b^\circ_\ell(g,r)\psi(\ell)\ell^{-w}$, where $b^\circ_\ell(g,r)$ are the coefficients of the expanded product, which is independent of $d$; and $\psi(\ell)=0$ unless $(\ell,d)=1$. The support and the bound follow as in \cref{lem:dualsum}. The formula for $Y$ is $X=gdrT/(2\pi)$ substituted, and $|\varpi_d|=\pi^{-2}$ because the $\alpha_i$ are imaginary.
\end{proof}

Consequently, writing $\epsilon_{\mathfrak a}(\pm)=\frac12(\pm1)^{\mathfrak a}$ for the parity projector,
\begin{equation}\label{eq:SigmaF}
 \begin{split}
 \Sigma_F={}&\sum_{g,r}\frac{\mu(r)}{\varphi(r)}\sum_\ell b^\circ_\ell(g,r)\sum_{\mathfrak a,\pm}\epsilon_{\mathfrak a}(\pm)\sum_{\substack{d>1\\(d,r\ell)=1}}d^2\varpi_d\\
 &\times\sum_{hk=gdr}\frac{\overline{F(\log h/L,\log k/L)}}{hk}\sum_{n\ge1}\frac{d_{-\boldsymbol\alpha}(n)}n\,V_{\mathfrak a}\Big(\frac n{Y}\Big)\mathscr K_d\big(\pm rn\overline\ell\big),
 \end{split}
\end{equation}
with $r$ squarefree. At fixed labels $(g,r,\ell,\mathfrak a,\pm)$ the summand depends on $d$ only through the power $d^{2-\sum\alpha_i}$, the cut $(d,r\ell)=1$, the smooth weight $V_{\mathfrak a}(n/Y)$ with $Y\propto d^3$, the kernel $\mathscr K_d$ with the unit dilation $\pm r\overline\ell$, and the modulus weight at the factorisations $hk=gdr$. No coefficient depends on $d$ in any other way; this is what allows us to sum over $d$ with the signs kept.

\emph{Separation of variables.} Insert smooth dyadic partitions of unity in $h$, $k$, $d$ and in each of the four variables of $d_{-\boldsymbol\alpha}(n)=\sum_{n_1n_2n_3n_4=n}\prod_in_i^{\alpha_i}$, and dyadic ranges for the labels $g,r,\ell$. On a box $h\asymp H$, $k\asymp K$ the function $\overline{F(\log h/L,\log k/L)}W(h/H)W(k/K)$ is a smooth function of $(\log h,\log k)$ supported in a fixed square, with all derivatives bounded (each derivative of $F(\cdot/L)$ gains $1/L$); by Fourier inversion in $(\log h,\log k)$ it is $\iint G(\theta_1,\theta_2)h^{i\theta_1}k^{i\theta_2}\dd\theta$ with $\iint|G|\ll1$. By the definition of $V_{\mathfrak a}$,
\[
 V_{\mathfrak a}\Big(\frac nY\Big)=\frac1{2\pi i}\int_{(c)}\widetilde\Phi(w)\,G_{\mathfrak a}(w)\Big(\frac{\pi^3rT}{2\ell}\Big)^{w}d^{-3w}n^w\dd w,
\]
with $G_{\mathfrak a}$ the gamma quotient, on any line $c<\frac{19}{20}$ (proof of \cref{lem:dualweight}). On a box $n_i\asymp N_i$, $d\asymp D_1$ we take $c=\frac9{10}$ if $N_*:=\prod N_i\le16Y_1$, where $Y_1$ is the value of $Y$ at $d=D_1$, and $c=-A$ otherwise; the part with $|\Im w|>Q^\epsilon$ is negligible. After these insertions every weight is a product of functions of the single variables $h,k,d,n_1,\dots,n_4$, each of the form $W(y/Y_0)\,y^{\xi}$ with $W$ a fixed smooth bump and $|\Im\xi|\le Q^\epsilon$, so with rescaled derivatives $\ll Q^{j\epsilon}$; the integrals over $\theta$ and $w$ cost $Q^{O(\epsilon)}$. Here and below $Q=T^{\kappa}$ with $\kappa$ the value of $\log q/L$ at a point of $E$ near which $F$ is supported; by a partition of unity we may assume that $\operatorname{supp}F$ lies in a small ball around that point, so that $\log q/\log Q=1+O(\epsilon)$ for all $q=hk$ in the support. We write $Q^{o(1)}$ for factors $Q^{O(\epsilon)}$ and $(\log Q)^{O(1)}$.

\emph{Degrees.} All degrees are to base $Q$: $[x]=\log x/\log Q$, taken at the dyadic scale. For a dyadic datum $\mathfrak D$ (the boxes of $h,k,g,r,\ell,d,n_1,\dots,n_4$) put $\mathfrak g=[g]$, $\mathfrak r=[r]$, $\mathfrak i=[\ell]$, let $\varsigma=2-L/\log q$ and $\mathfrak b=\log h/\log q$ at the box, and let $\mathfrak v=\max(0,[Y_1/N_*])$ be the tail degree. Put
\begin{gather*}
 \sigma=1-\mathfrak g-\mathfrak r=[d],\qquad
 \lambda=1+\varsigma-3\mathfrak g-4\mathfrak r-\mathfrak v,\\
 \tau=\varsigma-2\mathfrak g-3\mathfrak r-\tfrac9{10}\mathfrak v,\qquad\mathsf B_*=\tfrac12\sigma+2\tau .
\end{gather*}
Order the four dual degrees $[N_i]$, let $\mathsf m\ge\mathsf n$ be the two largest and $\mathsf j$ the sum of the other two, and put $\mathsf D=\mathsf m+\mathsf n+\mathsf j$. For a modulus $q$ let $P(q)$ be its powerful part and $\mathfrak p=[P(q)]$, and put $\tau_{\mathfrak p}=\tau-\frac12\mathfrak p$; in \cref{prop:nonprincipal} the letter $\tau$ denoted $\tau_{\mathfrak p}$.

\begin{lemma}[Degrees and normalisation]\label{lem:degrees}
\textup{(i)} Up to an error that tends to $0$ with $\epsilon$,
\[
 \mathsf D=\lambda+\mathfrak i=1+\varsigma-3\mathfrak g-4\mathfrak r-\mathfrak v+\mathfrak i,\qquad0\le\mathfrak i\le3\mathfrak g+4\mathfrak r,
\]
so $\lambda\le\mathsf D\le1+\varsigma$; boxes with $N_*>Y_1Q^{\epsilon}$ contribute $O(T^{-A})$.

\textup{(ii)} Let $\mathscr B_{\mathfrak D}(q)$ be the contribution of the datum $\mathfrak D$ to $\mathcal S^\flat(q)$, and $\mathscr B_{\mathfrak D}=\sum_{h,k}\overline{F(\cdot)}(hk)^{-1}\mathscr B_{\mathfrak D}(hk)$ its contribution to $\Sigma_F$. Normalise the sum over $n$ by $N_*$ and the sum over the moduli by their number. Then the trivial bound is $|\mathscr B_{\mathfrak D}|\le Q^{2-2\mathfrak g-3\mathfrak r-\frac9{10}\mathfrak v+o(1)}$, and $\mathscr B_{\mathfrak D}\ll T^{1-\delta}$ holds as soon as
\begin{itemize}
\item[(a)] the normalised sum over $d$ and $n$ (all moduli of the datum together) is $\ll Q^{-\tau-\delta'}$, or
\item[(b)] for every modulus $q$ of the datum, the normalised sum over $n$ at that modulus is $\ll Q^{-\tau_{\mathfrak p(q)}-\delta'}$,
\end{itemize}
for some $\delta'>0$, where in (b) $\delta$ depends on $\delta'$ only.
\end{lemma}
\begin{proof}
(i) By \cref{lem:labels}, $[Y_1]=\mathfrak i+3\sigma-\mathfrak r-(2-\varsigma)=1+\varsigma-3\mathfrak g-4\mathfrak r+\mathfrak i$, since $T=q^{2-\varsigma}$. If $N_*\le16Y_1$ then $[N_*]=[Y_1]-\mathfrak v+o(1)$. If $N_*>16Y_1$ then $\mathfrak v=0$, the line $c=-A$ gives the factor $(N_*/Y_1)^{-A}$, so $N_*\le Y_1Q^\epsilon$ up to $O(T^{-A})$, and $[N_*]=[Y_1]+O(\epsilon)$. The bound $\mathfrak i\le3\mathfrak g+4\mathfrak r$ is $\ell\mid g^3r^4$, and $\mathsf D\le1+\varsigma$ follows.

(ii) On the box, $|n^{-1}V_{\mathfrak a}(n/Y)|\ll N_*^{-1}(N_*/Y_1)^{9/10}\le N_*^{-1}Q^{-\frac9{10}\mathfrak v}$ for $c=\frac9{10}$, and $\ll N_*^{-1}$ for $c=-A$ (where $\mathfrak v=0$); this is \cref{lem:dualweight}. With $|\mathscr K_d|\le q^{o(1)}$ (\cref{lem:Kd}) the $n$-sum is at most $Q^{-\frac9{10}\mathfrak v+o(1)}$ times its normalisation, and $d^2/\varphi(r)\le Q^{2\sigma-\mathfrak r+o(1)}=Q^{2-2\mathfrak g-3\mathfrak r+o(1)}$. Summing over the $q^{o(1)}$ labels and factorisations of each $q$ with the weight $1/q$, over $q\le T$, gives the trivial bound, and $Q^{2-\varsigma}=T$ gives (a). For (b), the hypothesis gives $|\mathscr B_{\mathfrak D}(q)|\ll Q^{2-\varsigma+\frac12\mathfrak p-\delta'+o(1)}=TP(q)^{1/2}q^{-\delta'+o(1)}$, and, as at the end of the proof of \cref{prop:nonprincipal}, $\sum_{q\le T}P(q)^{1/2}q^{-1}\ll(\log T)^2$, the at most $\tau(q)=T^{o(1)}$ pairs $(h,k)$ for each $q$, and $q\ge T^{1/2}$ give $\mathscr B_{\mathfrak D}\ll T^{1-\delta'/2+o(1)}$.
\end{proof}

The number of data $\mathfrak D$ is $(\log T)^{O(1)}$, so it suffices to prove (a) or (b) for each $\mathfrak D$ with a $\delta'$ that is uniform in $\mathfrak D$. Whether (a) or (b) is used, and which estimate below proves it, is decided by the limiting degrees of $\mathfrak D$ alone; in particular the decision does not depend on $\mathfrak p$, nor on the modulus within $\mathfrak D$. If $\tau<0$ (strictly, at the limiting tuple) the trivial bound gives (a).

\subsubsection*{Estimates at each modulus}

\begin{lemma}[Unsigned estimates]\label{lem:unsignedplus}
Let $d>1$, let $c$ be a unit modulo $d$, and let $P$ be the powerful part of $d$. Let $W$ be smooth as in \Cref{sec:factorable}.
\begin{enumerate}[label=\textup{(\roman*)}]
\item For $M\ge1$, $\frac1M\big|\sum_mW(m/M)\mathscr K_d(cm)\big|\ll q^{o(1)}\big(\sqrt{dP}/M+\sqrt{P/d}\big)$.
\item For $M,N\ge1$ and coefficients $|\alpha_m|,|\beta_n|\le q^{o(1)}$ supported on $m\in[M,2^8M]$, $n\in[N,2^8N]$,
\[
 \Big(\frac1{MN}\Big|\sum_{m,n}\alpha_m\beta_n\mathscr K_d(cmn)\Big|\Big)^2\ll q^{o(1)}\Big(\frac1N+\frac{\sqrt{dP}}M+\sqrt{\frac Pd}\Big).
\]
\end{enumerate}
Consequently the normalised dual sum of a datum with $\tau\ge0$ is $\ll Q^{-\tau_{\mathfrak p}-\delta'}$ at every modulus, which is (b) of \cref{lem:degrees}, if either $\mathsf m>\frac12\sigma+\tau$, or $\sigma>4\tau$ and the four dual variables split into two groups of degrees $\mathsf s\le\mathsf l$ with $\mathsf s>2\tau$ and $\mathsf l>\mathsf B_*$ (each strictly, at the limiting tuple).
\end{lemma}
\begin{proof}
(i) Let $H'$ be the product of the prime powers $p^a\Vert d$ with $a\ge2$, and $K'=d/H'$, which is squarefree and coprime to $H'$; so $H'\le P$. By Poisson summation modulo $d$,
\[
 \sum_mW(m/M)\mathscr K_d(cm)=\frac Md\sum_{j\in\Z}\widehat W(jM/d)\widehat{\mathscr K}(j),
\]
where $\widehat{\mathscr K}(j)=\sum_{x\bmod d}\mathscr K_d(cx)e_d(jx)$.
By \cref{lem:Kd}(i) and the Chinese remainder theorem, $\widehat{\mathscr K}(j)$ is a product of complete sums modulo the prime powers of $d$. The factor at $t\mid H'$ is at most $9t$ by \cref{lem:Kd}(iii). The factor at $p\mid K'$ is at most $2\sqrt p+1$, and vanishes if $p\mid j$, by \cref{lem:Kd}(iv), which holds at every prime. Hence $\widehat{\mathscr K}(0)=0$ if $K'>1$, and $|\widehat{\mathscr K}(j)|\le q^{o(1)}H'\sqrt{K'}$. Since $\frac Md\sum_{j\ne0}|\widehat W(jM/d)|\ll q^{o(1)}(1+M/d)$, the left side of (i) is $\ll q^{o(1)}(H'\sqrt{K'}/M+H'\sqrt{K'}/d)$, and $H'\sqrt{K'}=\sqrt{dH'}\le\sqrt{dP}$, $H'\sqrt{K'}/d=\sqrt{H'/d}\le\sqrt{P/d}$. If $K'=1$ then $d\le P$ and (i) is trivial.

(ii) By Cauchy's inequality in $m$, with a smooth $W\ge0$ equal to $1$ on $[1,2^8]$, the square of the sum is at most $q^{o(1)}M\sum_{n_1,n_2}|\beta_{n_1}\beta_{n_2}||\mathcal C(n_1,n_2)|$ with $\mathcal C(n_1,n_2)=\sum_mW(m/M)\mathscr K_d(cmn_1)\overline{\mathscr K_d(cmn_2)}$. The terms $n_1=n_2$ contribute $q^{o(1)}M^2N$, which is the term $1/N$. For $n_1\ne n_2$ complete the $m$-sum modulo $d$ as in (i): $\mathcal C=\frac Md\sum_j\widehat W(jM/d)\widehat{\mathcal C}(j)$, where $\widehat{\mathcal C}(j)$ factors over the prime powers $t\Vert d$ into the complete sums $\sum_{x\bmod t}\mathscr K_t(c_txn_1)\overline{\mathscr K_t(c_txn_2)}e_t(j_tx)$ with units $c_t$. If $p\mid n_1n_2$ for some $p\mid d$ the summand vanishes identically. Let $K''$ be the product of the primes $p\ge5$ with $p\Vert d$, and $H''=d/K''\le6P$. At $p\mid K''$ with $p\nmid n_1n_2(n_1^2-n_2^2)$, \cref{lem:corr} with $m=2$, $A_1(x)=c_pn_1x$, $A_2(x)=c_pn_2x$ (equal roots, slopes $c_pn_1\ne\pm c_pn_2$) bounds the factor by $C\sqrt p$; at the other primes of $K''$ we use the bound $16p$, and at $t\mid H''$ the bound $81t$. So $|\widehat{\mathcal C}(j)|\le q^{o(1)}H''\sqrt{K''}(K'',n_1^2-n_2^2)^{1/2}$, and $H''\sqrt{K''}=\sqrt{dH''}\le\sqrt{6dP}$. Hence
\[
 \frac{|\mathcal C(n_1,n_2)|}M\ll q^{o(1)}\Big(\frac{\sqrt{dP}}M+\sqrt{\frac Pd}\Big)(K'',n_1-n_2)^{1/2}(K'',n_1+n_2)^{1/2},
\]
using that $K''$ is squarefree. By $xy\le\frac12(x^2+y^2)$ and $\sum_{0<|a|\le A}(K'',a)\le\tau(K'')\cdot2A$, the sum of the last two factors over $n_1\ne n_2$ is $\ll q^{o(1)}N^2$. Collecting the terms gives (ii). No condition on the sizes of the factors of $q$ is used.

For the consequences, let $Q^{\mathfrak p}$ be the powerful part of the modulus; the powerful part of $d$ is at most that of $q$. In (i), with $M=Q^{\mathsf m}$ and the other three variables fixed (the product with the unit $c$ is a unit or the inner sum vanishes), the saving is $\min(\mathsf m-\frac12(\sigma+\mathfrak p),\frac12(\sigma-\mathfrak p))$, which exceeds $\tau_{\mathfrak p}=\tau-\frac12\mathfrak p$ if $\mathsf m>\frac12\sigma+\tau$ and $\sigma>2\tau$. Here $\sigma-2\tau=1-2\varsigma+3\mathfrak g+5\mathfrak r+\frac95\mathfrak v>0$ for $\varsigma<\frac12$. In (ii), grouping the variables, with divisor-bounded coefficients, into $m$ of degree $\mathsf l$ and $n$ of degree $\mathsf s$, the squared saving is $\min(\mathsf s,\mathsf l-\frac12(\sigma+\mathfrak p),\frac12(\sigma-\mathfrak p))$, which exceeds $2\tau_{\mathfrak p}=2\tau-\mathfrak p$ if $\mathsf s>2\tau$, $\mathsf l>\mathsf B_*$ and $\sigma>4\tau$, for every $\mathfrak p\ge0$. The identity $\sigma-4\tau=1-4\varsigma+7\mathfrak g+11\mathfrak r+\frac{18}5\mathfrak v$ shows $\sigma>4\tau$ for $\varsigma<\frac14$; for $\frac14\le\varsigma<\frac13$ the failure $4\tau\ge\sigma$ is one of the alternatives in \cref{lem:cases}(iii).
\end{proof}

\emph{Remark.} In \cref{lem:onevar,lem:bilinear} the primes of $h$ were put into $H$ in order to shift by multiples of $H$; the form above completes directly modulo $d$, puts only the powerful part (and in (ii) the primes $2,3$) into $H'$, $H''$, and needs no condition on the sizes of the factors of $q$. The thresholds in the last paragraph do not depend on $\mathfrak p$; the density of powerful moduli enters only through \cref{lem:degrees}(ii)(b).

\subsubsection*{Double Poisson}

In the remaining cases $\mathsf m\le\frac12\sigma+\tau$ and no split as in \cref{lem:unsignedplus} succeeds. Let $n_1,n_2$ be the variables of degrees $\mathsf m\ge\mathsf n$, of lengths $U_m=Q^{\mathsf m}$, $U_n=Q^{\mathsf n}$, and keep the other two in their product $t=n_3n_4\asymp Q^{\mathsf j}$, whose coefficient $\gamma_t$ (a convolution of two separated weights) is divisor-bounded and independent of $d$.

\begin{lemma}[Double Poisson]\label{lem:doublepoisson}
Let $d>1$, let $d_{(1)}$ be the product of the primes $p$ with $p\Vert d$, and let $c,t$ be units modulo $d$. For integers $a,b$,
\[
 \frac1d\sum_{x_1,x_2\bmod d}\mathscr K_d(c\,t\,x_1x_2)\,e_d(bx_1+ax_2)=\one_{(ab,d)=1}\sum_{l\mid d_{(1)}}\frac1{\varphi(l)}\,e_{d/l}\big(c\,t\,\overline l\,\overline{ab}\big).
\]
Consequently
\begin{multline*}
 \sum_{n_1,n_2}W_1\Big(\frac{n_1}{U_m}\Big)W_2\Big(\frac{n_2}{U_n}\Big)\mathscr K_d(ctn_1n_2)\\
 =\frac{U_mU_n}d\sum_{(ab,d)=1}\widehat W_2\Big(\frac{aU_n}d\Big)\widehat W_1\Big(\frac{bU_m}d\Big)\sum_{l\mid d_{(1)}}\frac{e_{d/l}(ct\overline l\,\overline{ab})}{\varphi(l)},
\end{multline*}
the dual variables have effective lengths of degrees $\mathsf a=\sigma-\mathsf n$ and $\mathsf b=\sigma-\mathsf m$, both positive, and the trivial bound of the term $l$, normalised as in \cref{lem:degrees}, exceeds the required $Q^{-\tau}$ by $Q^{\mathsf j+\mathfrak v/10-\mathfrak i-\mathsf z+o(1)}$, where $\mathsf z=[l]$.
\end{lemma}
\begin{proof}
By \cref{lem:Kd}(i) and the Chinese remainder theorem, both sides factor over the prime powers $t'\Vert d$ (with the dilations $\overline{(d/t')}$, which combine as $\prod_{t'}e_{t'}(y\overline{(d/t')})=e_d(y)$), so it suffices to take $d=p^k$. For a primitive $\psi$ modulo $p^k$ and an integer $a$, $\sum_{x}\bar\psi(x)e_{p^k}(ax)$ is $\psi(a)\tau(\bar\psi)$ if $p\nmid a$ and $0$ if $p\mid a$. Hence the left side is $0$ unless $p\nmid ab$, and then equals
\[
 \frac1{p^k\cdot p^k\varphi(p^k)}\sum_\psi^*\tau(\psi)^3\tau(\bar\psi)^2\bar\psi(ct)\psi(ab)=\frac1{\varphi(p^k)}\sum_\psi^*\tau(\psi)\bar\psi(y),\qquad y=ct\overline{ab},
\]
using $\tau(\psi)\tau(\bar\psi)=\psi(-1)p^k$. Since $\sum_{\psi\bmod p^k}\tau(\psi)\bar\psi(y)=\varphi(p^k)e_{p^k}(y)$, and the imprimitive characters have $\tau(\psi)=0$ if $k\ge2$ (proof of \cref{lem:Kd}) while for $k=1$ the principal character has $\tau=-1$, the right side is $e_{p^k}(y)$ if $k\ge2$ and $e_p(y)+\frac1{p-1}$ if $k=1$. Multiplying over $p$ and expanding the product gives the sum over $l\mid d_{(1)}$, with $e_{d/l}(y\overline l)$ as the product of the remaining additive characters. The second display follows by Poisson summation modulo $d$ in $n_1$ and $n_2$ (with $\widehat W(\xi)=\int W(y)e(-y\xi)\dd y$). The transforms $\widehat W_i$ are Schwartz, so $|a|\le Q^{\epsilon}d/U_n$, $|b|\le Q^\epsilon d/U_m$ up to $O(T^{-A})$. Here $\mathsf a\ge\mathsf b=\sigma-\mathsf m\ge\frac12\sigma-\tau>0$. After division by $U_mU_n$, the term $l$ is trivially at most $Q^{o(1)}d^{-1}\varphi(l)^{-1}(d/U_n)(d/U_m)=Q^{\sigma-\mathsf m-\mathsf n-\mathsf z+o(1)}$, and $\tau+\sigma-\mathsf m-\mathsf n=\mathsf j+\frac1{10}\mathfrak v-\mathfrak i$ by \cref{lem:degrees}(i).
\end{proof}

\emph{Reduction of the modulus.} Write $d=lu$ with $l\mid d_{(1)}$, so that $l$ is squarefree and $(l,u)=1$, and fix $l$. The transform requires $(t,d)=1$ (otherwise $\mathscr K_d(ct\,\cdot)\equiv0$), which is $(t,l)=1$ and $(t,u)=1$. Expand the second condition by M\"obius inversion \emph{before} any absolute value is taken: $\one_{(t,u)=1}=\sum_{f\mid(t,u)}\mu(f)$, and write $u=fu'$, $t=fy_0$. Since $t=fy_0$ and $u=fu'$, every term satisfies $e_u(ct\overline l\,\overline{ab})=e_{u'}(cy_0\overline l\,\overline{ab})$, the inverses now taken modulo $u'$ (an inverse modulo $u$ is also one modulo $u'$). The conditions $(ab,d)=1$ become $(ab,u')=1$ and $(ab,lf)=1$. Expand $(a,lf)=1$ by M\"obius inversion as well: $a=ed_0$ with $e\mid lf$. Keep $(b,lf)=1$ as a cut in the weight of $b$, and $(fy_0,l)=1$ in that of $y_0$. With $c=\pm r\overline\ell$ and $F_0=l\ell$ the phase becomes
\begin{equation}\label{eq:reducedphase}
 e_{u'}\big(\pm r\,y_0\,\overline{F_0\,e\,d_0\,b}\big)\qquad\text{on}\quad(rF_0ed_0b,u')=1,
\end{equation}
where $(r\ell,u')=1$ comes from $(d,r\ell)=1$, $(l,u')=1$ from $(l,u)=1$, and $(ed_0b,u')=1$ from $(ab,u')=1$. Here $y_0\ge1$, since $t\ge1$. Put $\mathsf x=[f]$, $\mathsf E=[e]$, $\mathsf Z=\mathsf z+\mathsf x$, $\mathsf y=\mathsf j-\mathsf x$ and $\mathsf C=\sigma-\mathsf Z=[u']$; then $\mathsf E\le\min(\mathsf Z,\mathsf a)$, $d_0$ has effective length of degree $\mathsf a-\mathsf E$, and $b$ of degree $\mathsf b$. For fixed $(l,f,e)$ let $\Sigma(l,f,e)$ be the corresponding part of the normalised sum of \cref{lem:degrees}(ii)(a), all moduli $d=lfu'$ of the datum together. Its terms are
\begin{equation}\label{eq:reducedsum}
 \Sigma(l,f,e)=\frac{\mu(f)\mu(e)}{\varphi(l)}\sum_{u'}\sum_{y_0}\sum_{d_0}\sum_b\varpi(u',y_0,d_0,b)\,e_{u'}\big(\pm r\,y_0\,\overline{F_0ed_0b}\big),
\end{equation}
summed under \eqref{eq:reducedphase}, where the weight $\varpi$, of size $Q^{o(1)}$ times a normalising constant, is the product of: the modulus weight $\Theta(u')=\sum_{hk=grlfu'}\omega_h(h)\omega_k(k)$ with the separated mark weights $\omega_h,\omega_k$; the separated weights of $d=lfu'$ and the cut $(d,r\ell)=1$; the coefficient $\gamma_{fy_0}\one_{(fy_0,l)=1}$; $\widehat W_2(ed_0U_n/d)$; and $\widehat W_1(bU_m/d)\one_{(b,lf)=1}$. On dyads of $u',|d_0|,|b|$ the last two factors separate into products of functions of single variables by Fourier inversion in the logarithms, with the rescaled derivative bounds of the Schwartz seminorms (also as the arguments tend to $0$), at a cost $Q^{o(1)}$; we do this in \cref{lem:square,lem:direct,lem:recip} but not in \cref{lem:actuallevel}.

\emph{Normalisation.} The factor $1/\varphi(l)$ saves $\mathsf z$; for fixed $f$ the variable $y_0$ has $Q^{\mathsf j-\mathsf x}$ values instead of $Q^{\mathsf j}$, so the $f$-dyad saves $\mathsf x$ after its $Q^{\mathsf x}$ terms are summed; and the substitution $a=ed_0$ saves $\mathsf E$. Summed over the dyads of $l$ and $f$ (which have harmonic weights in the average over the moduli, as $d=lfu'$), the trivial bounds of the $\Sigma(l,f,e)$ add up to $Q^{\sigma-\mathsf m-\mathsf n-\mathsf Z-\mathsf E+o(1)}$. Hence (a) of \cref{lem:degrees} holds for the datum if, for every choice of the dyads of $l,f$ and of $e$, the sum of $\Sigma(l,f,e)$ over these dyads is smaller than its trivial bound by $Q^{\mathsf u+\delta'}$, where
\[
 \mathsf u=\tau+\sigma-\mathsf m-\mathsf n-\mathsf Z-\mathsf E=\mathsf j+\tfrac1{10}\mathfrak v-\mathfrak i-\mathsf Z-\mathsf E .
\]
If $\mathsf u<0$ the trivial bound suffices. The trivial bound of the raw sum \eqref{eq:reducedsum} at fixed $(l,f,e)$, with $\varpi$ replaced by its size, is $Q^{\mathsf C+\mathsf y+(\mathsf a-\mathsf E)+\mathsf b+o(1)}$, so the requirement is that the raw sum be $\ll Q^{\mathsf B-\delta'}$ with
\[
 \mathsf B=\mathsf C+\mathsf y+\mathsf a-\mathsf E+\mathsf b-\mathsf u=2\sigma+\mathsf y-\tau .
\]

\subsubsection*{Four estimates for the Kloosterman fractions}

Below, \emph{margin} means the allowed exponent minus the exponent of the bound; an estimate saves a power when its margin is positive at the limiting tuple. All exponents are computed with the largest dual scales $Q^{\mathsf a-\mathsf E}$, $Q^{\mathsf b}$; every bound used is nondecreasing in these scales, so smaller dyads of $d_0$ and $b$ are covered. Put $\mathsf c=\mathsf C-\mathsf w$, $\mathsf d=\mathsf a-\mathsf E$, $\mathsf R=\mathsf b+\mathfrak r+\mathsf z+\mathfrak i+\mathsf E$, $\mathsf N=2\mathfrak r+\mathsf y$, where $\mathsf w=[S_0]$ is defined in \cref{lem:direct}. The identities between affine forms used below have also been checked by computer expansion.

\begin{lemma}[Spectral square]\label{lem:square}
Suppose $\mathsf b>\mathsf j+\mathfrak r$. Put $\mathcal A=\sigma+\mathsf z-\mathsf x+\mathsf a+2\mathsf b-\frac15\mathfrak v+\mathsf E$ and
\[
 \begin{aligned}
 \mathsf P_0&=\mathsf b-\mathsf j-\tfrac15\mathfrak v+2\mathsf z+\mathsf x+2\mathsf E,\\
 \mathsf P_1&=\mathsf m-\tfrac12\mathsf n-\mathsf j-\mathfrak i-\tfrac15\mathfrak v+\tfrac12(\mathsf Z+\mathsf E),\\
 \mathsf P_2&=\tfrac12\mathsf a-\mathsf j-\tfrac12\mathfrak i-\tfrac15\mathfrak v+\tfrac32\mathsf z+\mathsf x+\mathsf E,\\
 \mathsf P_3&=\tfrac12(\sigma+\mathsf m)-\tfrac32\mathsf j-\tfrac12(\mathfrak r+\mathfrak i)-\tfrac15\mathfrak v+\tfrac12\mathsf z+\tfrac12\mathsf x+\tfrac32\mathsf E .
 \end{aligned}
\]
If $\mathsf P_0,\mathsf P_2,\mathsf P_3>0$, and $\mathsf P_1>0$ or $\mathsf z+\mathfrak i+2\mathsf b+\mathsf a<\mathsf C$, then the sum of \eqref{eq:reducedsum} over the dyads of $l$ and $f$ saves $Q^{\mathsf u+\delta'}$ for every $e$.
\end{lemma}
\begin{proof}
Fix $l,f,e$ and one sign; separate the weights on dyads $u'\asymp U'$, $d_0\asymp A_0\le Q^{\mathsf a-\mathsf E+o(1)}$ (with a sign), $b\asymp B_0\le Q^{\mathsf b+o(1)}$ (with a sign), $y_0\asymp Y_0=Q^{\mathsf y+o(1)}$. The sum is then $\sum_{u',d_0}\alpha(u',d_0)\sum_{y_0,b}\beta(y_0)\gamma(b)e_{u'}(\pm ry_0\overline{F_0ed_0b})$ with $|\alpha|,|\beta|,|\gamma|\le Q^{o(1)}$, where $\alpha$ contains $\Theta(u')$ and all the dependence on the factorisations $hk$, and the inner sum is over $(b,u')=1$. By Cauchy's inequality in $(u',d_0)$, majorising the dyads by smooth nonnegative $W(u'/U')V(d_0/A_0)$, dropping the conditions $(r,u')=1$ and (when $e$ is not expanded, as in \cref{lem:cases}(i)) $(d_0,lf)=1$ in this nonnegative majorant, and keeping $(F_0ed_0,u')=1$,
\[
 |\Sigma|^2\le Q^{o(1)}U'A_0\sum_{y_0,b,y_0',b'}\beta\gamma\overline{\beta'\gamma'}\sum_{(R_0d_0,u')=1}W\Big(\frac{u'}{U'}\Big)V\Big(\frac{d_0}{A_0}\Big)e_{u'}\big(k\,\overline{R_0d_0}\big),
\]
with $k=\pm r(y_0b'-y_0'b)$ and $R_0=F_0ebb'$, since $y_0\overline b-y_0'\overline{b'}\equiv(y_0b'-y_0'b)\overline{bb'}$.

\emph{Diagonal.} If $k=0$ then $y_0b'=y_0'b$; for each value $v\le4Y_0B_0$ of this product there are $\tau(v)^2$ quadruples, so these terms contribute $Q^{o(1)}U'A_0\cdot U'A_0Y_0B_0$.

\emph{Off-diagonal.} Group the terms with $k\ne0$ by the sign of $k$ and by $(|k|,R_0)$, with coefficient $b_{k,R_0}=\sum\beta\gamma\overline{\beta'\gamma'}$ over the quadruples with these values. At fixed $R_0$ there are divisor-many pairs $(b,b')$; for one such pair with $v=(b,b')$, the solutions $(y_0,y_0')$ of $y_0b'-y_0'b=k/(\pm r)$ differ by multiples of $(b/v,b'/v)$, so there are $O(1+Y_0v/B_0)$ of them. Bounding one factor of $\sum|b_{k,R_0}|^2$ by the largest fibre and using $\sum_{b,b'\asymp B_0}(b,b')\ll B_0^2\log B_0$,
\[
 \|b_{k,R_0}\|_2\ll Q^{o(1)}Y_0B_0\Big(1+\frac{Y_0}{B_0}\Big)^{1/2}\le Q^{\mathsf y+\mathsf b+o(1)},
\]
as $\mathsf y\le\mathsf j<\mathsf b$. Apply (Q1) with $S=1$, $s=1$, smooth variables $c=u'$, $d=d_0$, numerators $n=|k|\le K_0=Q^{\mathfrak r+\mathsf y+\mathsf b+o(1)}$ and $r=R_0\asymp Q^{\mathsf z+\mathfrak i+\mathsf E+2\mathsf b+o(1)}$; the phase is $e(n\overline{rd}/c)$ or its conjugate, and the coprimality is exactly $(R_0d_0,u')=1$. As $\mathfrak r+\mathsf y\le\mathsf j+\mathfrak r<\mathsf b$, the numerator scale is below the denominator scale, $R+K_0\asymp R$, and \eqref{eq:DIK} has the three terms of degrees
\[
 \mathsf C+\mathsf R'+\max(\mathsf C,\mathsf R'+\mathsf d),\qquad2\mathsf C+\mathsf d+\mathsf R',\qquad2\mathsf d+\mathsf K_0+\mathsf R',
\]
with $\mathsf R'=\mathsf z+\mathfrak i+\mathsf E+2\mathsf b$ and $\mathsf K_0=\mathfrak r+\mathsf y+\mathsf b$. The off-diagonal part of the expanded sum is therefore $\ll Q^{(\mathsf y+\mathsf b)+\frac12\max(\cdots)+o(1)}$.

\emph{Budget.} We need $|\Sigma|\ll Q^{\mathsf C+\mathsf d+\mathsf y+\mathsf b-\mathsf u-\delta'}$, that is, the expanded sum must be $\ll Q^{\mathsf C+\mathsf d+2\mathsf y+2\mathsf b-2\mathsf u-2\delta'}$. Using $\mathsf u\le\mathsf j+\frac1{10}\mathfrak v-\mathsf Z-\mathsf E$ (we do not use the term $-\mathfrak i$), this exponent is at least $\mathcal A$. Subtracting the four bounds from $\mathcal A$ gives the margins: $\mathsf P_0$ for the diagonal; $\mathsf P_1$ for the first term when $\mathsf R'+\mathsf d\ge\mathsf C$, and $\mathsf P_2+\frac12\mathsf d>\mathsf P_2$ when $\mathsf R'+\mathsf d<\mathsf C$; $\mathsf P_2$ and $\mathsf P_3$ for the other two. Summing over $f$ and $l$ in their dyads, and over the $O(Q^{o(1)})$ dyads and signs, proves the claim.
\end{proof}

\begin{lemma}[Direct estimate with the marks kept]\label{lem:direct}
Fix the labels $g,r,\ell$, the dyads of $l,f$ and of the other lengths, and $V_0=grlf$, $V_1=(h,V_0)$, $V_2=V_0/V_1$. The map
\[
 (h,k)\longmapsto(V_1,S_0,C_0),\qquad h=V_1S_0,\quad k=V_2C_0,\quad u'=S_0C_0,
\]
is a bijection from the factorisations $hk=V_0u'$ onto the triples with $V_1\mid V_0$, $S_0C_0=u'$ and $(S_0,V_2)=1$. Let $\mathsf w=[S_0]$, so that $\mathsf w\le\min(\mathsf C,\mathfrak b)$. Then the sum of \eqref{eq:reducedsum} over all $l,f,e,V_1$ in the given dyads, at fixed dyad of $S_0$, has raw exponent at most
\[
 \mathsf Z+\tfrac12\big(\mathsf y+\mathsf b+\mathsf w-\mathsf z-(\mathsf E-\mathsf z)_+ +\max\mathcal K\big),
\]
where $\mathcal K$ consists of the seven affine forms
\begin{equation}\label{eq:Kset}
 \begin{gathered}
 \mathsf c+\mathsf w+\mathsf p+\mathsf q\quad(\mathsf p\in\{\mathsf R+\mathsf w,\mathsf N\},\ \mathsf q\in\{\mathsf c,\mathsf R+\mathsf d\}),\\
 2\mathsf c+\mathsf d+\mathsf w+\tfrac12(\mathsf p+\mathsf R)\quad(\mathsf p\in\{\mathsf R+\mathsf w,\mathsf N\}),\qquad
 2\mathsf d+\mathsf N+\mathsf R .
 \end{gathered}
\end{equation}
Since there are $Q^{\mathsf Z+o(1)}$ pairs $(l,f)$, each with the budget $Q^{\mathsf B}$, the estimate saves if $\frac12(\mathsf y+\mathsf b+\mathsf w-\mathsf z-(\mathsf E-\mathsf z)_++\max\mathcal K)<\mathsf B$. It saves without the gain $(\mathsf E-\mathsf z)_+$ under the same condition with this term omitted.
\end{lemma}
\begin{proof}
\emph{The bijection.} Given $(h,k)$, $V_1=(h,V_0)$, $S_0=h/V_1$ and $V_2=V_0/V_1$ are coprime, and $S_0k=V_2u'$, so $S_0\mid u'$; put $C_0=u'/S_0$, then $k=V_2C_0$. Conversely $(V_1S_0,V_0)=V_1(S_0,V_2)=V_1$. The condition $(S_0,V_2)=1$ involves $S_0$ only, and $S_0\le h$, $S_0\mid u'$ give $\mathsf w\le\min(\mathsf C,\mathfrak b)$.

\emph{The sum.} Insert $r$ into numerator and denominator: by \eqref{eq:reducedphase} the phase is $e_{S_0C_0}(\pm N_0\overline{R_0d_0})$ with
\[
 N_0=r^2y_0,\qquad R_0=rF_0eb=r\ell\,l\,e\,b,
\]
and every condition that involves $C_0$ or $d_0$ is the single coprimality $(R_0d_0,S_0C_0)=1$: indeed $(r\ell l,u')=1$ and $(ebd_0,u')=1$ are \eqref{eq:reducedphase}, and no other condition on $d=lfS_0C_0$ remains after the expansions in $f$ and $e$. The mark weight at $k=V_2C_0$ and the weights of $d=lfS_0C_0$ depend on $C_0$ and on the representation; on common dyads they separate by Fourier inversion in the logarithms of $V_2,l,f,S_0,C_0,e,d_0$, so that the weights of $C_0$ and $d_0$ become fixed smooth functions $g(C_0,d_0)$ as in (Q1), and all the rest enters the coefficient $b_{N_0,R_0,S_0}$, which collects the terms of all $(l,f,e,V_1,y_0,b)$ in the dyads with the given $(N_0,R_0,S_0)$: the mark weight at $h=V_1S_0$, the cut $(S_0,V_2)=1$, the coefficient of $y_0$ and the cuts $(fy_0,l)=1$, $(b,lf)=1$. The sign of $R_0$ is moved into $N_0$. In particular no coefficient depending on the modulus is attached to $C_0$, and the sum over $C_0$ keeps its sign.

\emph{The coefficient norm.} At fixed $(N_0,R_0,S_0)$ the integer $y_0$ is fixed and $(l,e,b)$ is a factorisation of $R_0/(r\ell)$, with divisor-many choices; then $e\mid lf$ forces $e/(e,l)\mid f$, and $e/(e,l)\ge e/l$ has degree at least $(\mathsf E-\mathsf z)_+$, so $f$ has $Q^{\mathsf x-(\mathsf E-\mathsf z)_++o(1)}$ choices in its dyad; and $V_1\mid grlf$ has divisor-many values. The number of all tuples is $Q^{\mathsf z+\mathsf x+\mathsf y+\mathsf b+\mathsf w+o(1)}$. Bounding the sum of the squares by the largest fibre times the number of tuples,
\[
 \|b\|_2^2\le Q^{\mathsf z+2\mathsf x+\mathsf y+\mathsf b+\mathsf w-(\mathsf E-\mathsf z)_++o(1)}
\]
after extracting the common size of the weights.

\emph{The spectral factor.} Apply (Q1) with $c=C_0$ (degree $\mathsf c=\mathsf C-\mathsf w$), $d=d_0$ ($\mathsf d$), $n=N_0$ ($\mathsf N$), $r=R_0$ ($\mathsf R$), $s=S_0$ ($\mathsf w$). The three terms of $\mathcal K^2$ in \eqref{eq:DIK} have the degrees
\begin{gather*}
 \mathsf c+\mathsf w+\max(\mathsf R+\mathsf w,\mathsf N)+\max(\mathsf c,\mathsf R+\mathsf d),\\
 2\mathsf c+\mathsf d+\mathsf w+\tfrac12\big(\max(\mathsf R+\mathsf w,\mathsf N)+\mathsf R\big),\qquad2\mathsf d+\mathsf N+\mathsf R,
\end{gather*}
whose maximum is $\max\mathcal K$. So the sum is
\[
 \ll Q^{\frac12(\mathsf z+2\mathsf x+\mathsf y+\mathsf b+\mathsf w-(\mathsf E-\mathsf z)_++\max\mathcal K)+o(1)},
\] which is the raw exponent stated, as $\frac12(\mathsf z+2\mathsf x)=\mathsf Z-\frac12\mathsf z$. Without pooling $f$ (fixing $f$ and summing $Q^{\mathsf x}$ bounds) the coefficient norm loses the term $-(\mathsf E-\mathsf z)_+$ and the bound is the same with this term omitted.
\end{proof}

\begin{lemma}[Reciprocal completion]\label{lem:recip}
If $\mathsf u<\min(\frac12\mathsf C,\ \mathsf b-\frac12\mathsf C)$, then \eqref{eq:reducedsum} saves $Q^{\mathsf u+\delta'}$ at each modulus $u'$, for every $l,f,e$.
\end{lemma}
\begin{proof}
Fix $l,f,e,u',y_0,d_0$ and a dyad $|b|\asymp B_0\le Q^{\mathsf b+o(1)}$, and expand $(b,lf)=1$ as $\sum_{\nu\mid(b,lf)}\mu(\nu)$; terms with $(\nu,u')>1$ vanish. For $b=\nu b_1$ the weight of $b_1$ is a smooth function $W(\nu b_1/B_0)$ with bounded rescaled derivatives (after separation), and the summand is $e_{u'}(w_0\overline{b_1})\one_{(b_1,u')=1}$ with $w_0\equiv\pm ry_0\overline{F_0ed_0\nu}$. Poisson summation modulo $u'$ gives
\[
 \sum_{b_1}W\Big(\frac{\nu b_1}{B_0}\Big)e_{u'}(w_0\overline{b_1})=\frac{B_0}{\nu u'}\sum_{h'\in\Z}\widehat W\Big(\frac{h'B_0}{\nu u'}\Big)S(h',w_0;u').
\]
Since $r,F_0,e,d_0,\nu$ are units modulo $u'$, $(h',w_0,u')\le(y_0,u')$, and (Q4) gives $|S(h',w_0;u')|\le Q^{o(1)}\sqrt{u'}(y_0,u')^{1/2}$ for every $h'$, including $h'=0$. So the $b$-sum is $\ll Q^{o(1)}(y_0,u')^{1/2}(\sqrt{u'}+B_0/\sqrt{u'})$. As $y_0\ge1$, $\sum_{y_0\le Y_0}(y_0,u')^{1/2}\le\sum_{v\mid u'}v^{1/2}Y_0/v\le\tau(u')Y_0$. Compared with the trivial bound $Q^{\mathsf y+\mathsf b}$ of the sum over $(y_0,b)$, the saving is $\min(\frac12\mathsf C,\mathsf b-\frac12\mathsf C)$ for every $u',d_0$; the divisor sums over $\nu$ cost $Q^{o(1)}$. Since this is an estimate at each modulus, the dyads of $l,f$ are summed trivially.
\end{proof}

\begin{lemma}[Kuznetsov at the represented level]\label{lem:actuallevel}
Fix $l,f,e$ and put $\mathsf J=\mathsf C-\mathsf d=\mathsf n-\mathsf Z+\mathsf E$ and $\mathsf A_*=\mathsf R+\mathsf w$.
\begin{enumerate}[label=\textup{(\roman*)}]
\item If $\mathsf J<0$, the sum \eqref{eq:reducedsum} vanishes for all sufficiently large $T$.
\item If $\mathsf J\ge0$, the sum of \eqref{eq:reducedsum} over $V_1$ and over one dyad of $S_0$, in the parametrisation of \cref{lem:direct}, has raw exponent at most
\begin{equation}\label{eq:Fstar}
 \begin{split}
 \mathsf F_*={}&\mathsf d+\tfrac12\mathsf R+\mathsf b+\mathsf w+\tfrac12(\mathsf J+\mathsf y)\\
 &+\tfrac12\big((\mathsf J-\mathsf A_*)_++(\mathsf N-\mathsf A_*)_+\big)
 +\tfrac7{32}\big(\mathsf C+\mathsf J+\mathsf N-\tfrac32\mathsf R-2\mathsf w\big)_+ .
 \end{split}
\end{equation}
\end{enumerate}
So the estimate saves if $\mathsf J<0$, or if $\mathsf F_*<\mathsf B$.
\end{lemma}
\begin{proof}
Fix the labels $g,r,\ell$, the sign, $l,f,e$, the dyads, $V_1$ (hence $V_2$), $S_0$ and $b$. As in the proof of \cref{lem:direct}, the phase is $e_{S_0C_0}(\pm N_0\overline{R_0d_0})$ with $N_0=r^2y_0$ and $R_0=r\ell\,l\,e\,b$, and every condition on $C_0$ or $d_0$ is $(R_0d_0,S_0C_0)=1$. Here we do not separate the weight of $d_0$. With $d=lfS_0C_0$ it is $\widehat W_2(ed_0U_n/d)=\widehat W_2(d_0/D_0)$, $D_0=lfS_0C_0/(eU_n)$, and $d_0$ runs over all integers coprime to $q'=S_0C_0$ (the expansion $a=ed_0$ leaves no other condition on $d_0$). Put $T_0=q'/D_0=eU_n/(lf)$; it does not depend on $C_0$, and $[T_0]=\mathsf E+\mathsf n-\mathsf z-\mathsf x=\mathsf J$.

\emph{Poisson summation in $d_0$.} Splitting $d_0$ into classes $x$ modulo $q'$ and applying Poisson summation to $k\mapsto\widehat W_2((x+kq')/D_0)$, whose Fourier transform is $W_2(-\,\cdot)$ by Fourier inversion,
\[
 \sum_{(d_0,q')=1}\widehat W_2\Big(\frac{d_0}{D_0}\Big)e_{q'}\big(\pm N_0\overline{R_0d_0}\big)=\frac1{T_0}\sum_{t\in\Z}W_2\Big(-\frac t{T_0}\Big)S\big(t,\pm N_0\overline{R_0};q'\big).
\]
$W_2$ is supported in a fixed compact subset of $(0,\infty)$, so only $t\ne0$ with $-t\asymp T_0$ occur. If $\mathsf J<0$ then $T_0\to0$, no such $t$ exists for large $T$, and the sum vanishes. This proves (i). Assume now $\mathsf J\ge0$.

\emph{The Kloosterman sums as sums over a cusp pair.} Put $A=|R_0|S_0$. Since $(R_0,S_0)=1$, \cref{lem:ALcusp} applies to $\Gamma_0(A)$ with the exact divisor $|R_0|$ in place of $a$ and $S_0$ in place of $m$; the moduli of the cusp pair $(\infty,\mathfrak a)$, $\mathfrak a=1/S_0$, are $\gamma=c\sqrt{|R_0|}$ with $S_0\mid c$ and $(c,R_0)=1$, that is $c=S_0C_0$ with $(C_0,R_0)=1$, which is exactly our range of $C_0$. By \eqref{eq:ALgeneral}, with $n'=\pm\operatorname{sgn}(R_0)\,r^2y_0$,
\[
 S\big(t,\pm N_0\overline{R_0};S_0C_0\big)=S\big(t,n'\,\overline{|R_0|};S_0C_0\big)=S_{\infty\mathfrak a}\big(t,n';S_0C_0\sqrt{|R_0|}\big).
\]
All remaining weights of $C_0$ (the mark weight at $k=V_2C_0$, the separated weights of $d=lfS_0C_0$, including the factor $d^{-3w}$ from $V_{\mathfrak a}$, and $\widehat W_1(bU_m/d)$) form a smooth function $\Psi(C_0)$ supported on $C_0\asymp Q^{\mathsf c+o(1)}$, with $(C_0\partial_{C_0})^i\Psi\ll_iQ^{o(1)}\|\Psi\|_\infty$; the weights of $y_0$ form a coefficient $\beta(y_0)\ll Q^{o(1)}$. Write $\Psi(C_0)=\gamma^{-1}\phi(4\pi\sqrt{|tn'|}/\gamma)$ with $\gamma=S_0C_0\sqrt{|R_0|}$, and remove the dependence of $\phi$ on $|tn'|$ by Mellin inversion in $|t|/T_0$ and $|n'|/(r^2Y_0)$, where $Y_0=Q^{\mathsf y+o(1)}$ is the scale of $y_0$, on fixed enlarged intervals, as after \cref{lem:ALcusp}; the resulting functions $\phi$ are supported on $\xi\asymp\xi_0$ with
\[
 [\xi_0]=\tfrac12(\mathsf J+\mathsf N)-\big(\mathsf C+\tfrac12[R_0]\big),\qquad\Phi_0:=\sup|\phi|=Q^{\mathsf C+\frac12[R_0]+o(1)}\|\Psi\|_\infty,
\]
and $\xi^i\phi^{(i)}\ll_iQ^{o(1)}\Phi_0$, and the Mellin factors $(|t|/T_0)^{i\upsilon_1}$, $(|n'|/r^2Y_0)^{i\upsilon_2}$ go into the coefficients $\alpha_t=W_2(-t/T_0)(|t|/T_0)^{i\upsilon_1}$ and $\beta'_{n'}=\beta(y_0)(|n'|/r^2Y_0)^{i\upsilon_2}$. Put $\mathsf X=\mathsf C+\frac12\mathsf R-\frac12(\mathsf J+\mathsf N)$; at the full $b$-scale $[R_0]=\mathsf R$ and $-[\xi_0]=\mathsf X$. By the identity $2\mathsf X=2-\varsigma+\mathfrak v$ (\S6), and since in general $[R_0]\ge\mathsf R-\mathsf b\ge\mathsf R-1$, we have $-[\xi_0]\ge\mathsf X-\frac12\ge\frac13$, so $\xi_0\le1$ for large $T$.

\emph{Spectral expansion.} By \eqref{eq:kuzgeneral} with $m'=t$,
\[
 \begin{split}
 \sum_{C_0}\Psi(C_0)S_{\infty\mathfrak a}(t,n';\gamma)&=\sum_j\overline{\rho_{j\infty}(t)}\rho_{j\mathfrak a}(n')\widetilde\phi(t_j)\\
 &\qquad+(\text{holomorphic terms})+(\text{Eisenstein terms}),
 \end{split}
\]
and the transforms satisfy $\widetilde\phi(t)\ll_AQ^{o(1)}\Phi_0(1+|t|)^{-A}$ for real $t$, $\widetilde\phi(k)\ll_AQ^{o(1)}\Phi_0k^{-A}$, and $\widetilde\phi(i\tau)\ll Q^{o(1)}\Phi_0\xi_0^{-2\tau}$ for $0<\tau\le\theta=\frac7{64}$. Sum over $t$ and $y_0$ with the coefficients $\alpha_t$, $\beta'_{n'}$. For the regular spectrum we use Cauchy's inequality on each dyadic range of $|t_j|$ (or of $k$) and \eqref{eq:LS} at level $A$ for the cusps $\infty$ (indices $|t|\asymp T_0$) and $\mathfrak a$ (indices $|n'|\asymp Q^{\mathsf N}$, supported on multiples of $r^2$); with the rapid decay of the transforms this gives
\[
 \ll Q^{o(1)}\Phi_0\Big(1+\frac{T_0}A\Big)^{1/2}\Big(1+\frac{Q^{\mathsf N}}A\Big)^{1/2}\|\alpha\|_2\|\beta'\|_2 .
\]
For the exceptional spectrum put $Y_1=\max(1,A^2/T_0^3)$ and $Y_2=\max(1,A^2Q^{-3\mathsf N})$, and write $\xi_0^{-2\tau_j}=(\xi_0^{-2}/Y_1Y_2)^{\tau_j}\,Y_1^{\tau_j}Y_2^{\tau_j}$. The first factor is at most $\max(1,\xi_0^{-2}/Y_1Y_2)^{\theta}$. By Cauchy's inequality, and by \eqref{eq:EX} with $Y=Y_i$ when $Y_i>1$ (its right side is then $\ll Q^{o(1)}$ times the norm, as noted after \eqref{eq:EX}) or \eqref{eq:LS} with $U=1$ when $Y_i=1$, the exceptional terms are bounded by the same expression times $\max(1,\xi_0^{-2}/Y_1Y_2)^{\theta}$. Here
\[
 \Big[\frac{\xi_0^{-2}}{Y_1Y_2}\Big]\le2\mathsf X-(2\mathsf A_*-3\mathsf J)-(2\mathsf A_*-3\mathsf N)=2\big(\mathsf C+\mathsf J+\mathsf N-\tfrac32\mathsf R-2\mathsf w\big)
\]
at the full $b$-scale, where $[A]=\mathsf A_*$ (the last equality is an identity checked in \S6; for smaller $|b|$ see below), and $2\theta=\frac7{32}$.

\emph{Collecting.} $\|\alpha\|_2^2=Q^{\mathsf J+o(1)}$, $\|\beta'\|_2^2\le Q^{\mathsf y+o(1)}$, $[A]=\mathsf A_*$ and $[\Phi_0]=\mathsf C+\frac12\mathsf R$ at the full $b$-scale. Multiplying by $T_0^{-1}=Q^{-\mathsf J}$ and using $\mathsf C-\mathsf J=\mathsf d$, the sum at fixed $(V_1,S_0,b)$ is $\ll Q^{\mathsf F_*-\mathsf b-\mathsf w+o(1)}$. There are divisor-many $V_1$, $Q^{\mathsf w}$ values of $S_0$ in its dyad and $Q^{\mathsf b}$ values of $b$, which gives \eqref{eq:Fstar}. For $|b|$ in a smaller dyad, of degree $\mathsf b'<\mathsf b$, the same computation gives $\mathsf F_*$ with $\mathsf b$ replaced by $\mathsf b'$ in $\mathsf b$ and in $\mathsf R$; its derivative in $\mathsf b'$ is at least $\frac32-\frac12-\frac12-\frac32\cdot\frac7{32}=\frac{11}{64}>0$, so the full dyad is the worst case, and the $O(\log Q)$ dyads cost $Q^{o(1)}$. No bound depending on $C_0$ individually was used, and the sum over $C_0$ was kept with its signs and weights.
\end{proof}

\emph{Checks of the normalisation.} We record the two points of \cref{lem:actuallevel} that depend on conventions. (1) Two cusps. In \eqref{eq:kuzgeneral} both Maass coefficients carry $\cosh(\pi t_j)^{-1/2}$, as in \eqref{eq:kuznetsov}; \eqref{eq:LS} and \eqref{eq:EX} are stated for exactly these coefficients, at the cusps $\infty$ and $\mathfrak a=1/S_0$ of $\Gamma_0(|R_0|S_0)$, for which $\mu(\mathfrak c)=1/A$; the geometric modulus is $S_0C_0\sqrt{|R_0|}$, so the scalar in front of $\phi$ is $\gamma$, not $S_0C_0$, and the Bessel argument is $4\pi\sqrt{|tn'|}/\gamma$; this is the source of $\frac12\mathsf R$ in $\mathsf F_*$ and of $\mathsf X$. The direction Kloosterman sums $\to$ spectrum has no diagonal term. (2) Conductor labels. By \cref{lem:labels} the only dependence of the summand on $d$ is the one listed after \eqref{eq:SigmaF}; after the bijection of \cref{lem:direct} it is a smooth function of $C_0$, and the coprimality conditions are those of \eqref{eq:reducedphase}, which are exactly the conditions $(c,R_0)=1$, $S_0\mid c$ on the moduli of the cusp pair.

\subsubsection*{The case analysis}

\begin{lemma}[Covering]\label{lem:cases}
Let $(\varsigma,\mathfrak b)\in\mathcal P^{+*}$, and consider a limiting tuple of the degrees of a datum, together with the degrees $\mathsf z,\mathsf x,\mathsf E,\mathsf w$ of $l,f,e,S_0$. Then one of the following has positive margin at the tuple: the trivial bound ($\tau<0$, or $\mathsf u<0$), \cref{lem:unsignedplus}, \cref{lem:square}, \cref{lem:direct}, \cref{lem:recip} and \cref{lem:actuallevel}. More precisely:
\begin{enumerate}[label=\textup{(\roman*)}]
\item for $0<\varsigma<\frac1{10}$ and every $\mathfrak b$, the trivial bound, \cref{lem:unsignedplus} or \cref{lem:square} saves;
\item for $\frac1{10}\le\varsigma<\frac18$ and $\mathfrak b<\frac{3-21\varsigma}2$, the trivial bound, \cref{lem:unsignedplus}, \cref{lem:square} or \cref{lem:direct} (without the gain $(\mathsf E-\mathsf z)_+$) saves;
\item on each of the bands $\frac1{10}\le\varsigma<\frac17$, $\frac17\le\varsigma<\frac14$, $\frac14\le\varsigma<\frac13$, with the respective caps $\mathfrak b<f(\varsigma)=\frac{2-7\varsigma}4$, $\frac{1-2\varsigma}4$, $\frac{1-3\varsigma}2$, the trivial bound, \cref{lem:unsignedplus}, \cref{lem:recip}, \cref{lem:direct} (with or without the gain) or \cref{lem:actuallevel} saves.
\end{enumerate}
\end{lemma}
Since $\frac{3-21\varsigma}2\ge\frac{2-7\varsigma}4$ exactly for $\varsigma\le\frac4{35}<\frac18$, the three parts cover $\mathcal P^{+*}$.

\begin{proof}
Assume $\tau\ge0$ and $\mathsf u\ge0$ at the tuple (otherwise the trivial bound saves), and that \cref{lem:unsignedplus} does not apply: $\mathsf m\le\frac12\sigma+\tau$, and no split succeeds. We use $\lambda\le\mathsf D\le1+\varsigma$ and $0\le\mathfrak i\le3\mathfrak g+4\mathfrak r$ (\cref{lem:degrees}), $3\mathfrak r\le2\mathfrak g+3\mathfrak r\le\varsigma$ (from $\tau\ge0$), and the identities below, all checked by expansion in \S6:
\[
 \begin{gathered}
 \sigma-4\tau=1-4\varsigma+7\mathfrak g+11\mathfrak r+\tfrac{18}5\mathfrak v,\qquad
 \tfrac23\lambda-\mathsf B_*=\tfrac16-\tfrac43\varsigma+\tfrac52\mathfrak g+\tfrac{23}6\mathfrak r+\tfrac{17}{15}\mathfrak v,\\
 \lambda-\mathsf B_*-4\tau=\tfrac12-5\varsigma+\tfrac{19}2\mathfrak g+\tfrac{29}2\mathfrak r+\tfrac{22}5\mathfrak v .
 \end{gathered}
\]
Let $d_1\ge d_2\ge d_3\ge d_4$ be the four dual degrees, so $\mathsf m=d_1$, $\mathsf n=d_2$, $\mathsf j=d_3+d_4$, and put $\mathsf U=\mathsf D-\mathsf B_*$.

\emph{Step 1: gaps.} Let $\varsigma<\frac18$. Then $\sigma>4\tau$, and every set $G$ of the dual variables has degree $\deg G\le2\tau$ or $\deg G\ge\mathsf U$. Indeed, if $2\tau<\deg G<\mathsf U$, then the complement has degree $\mathsf D-\deg G>\mathsf B_*\ge2\tau$; if $\deg G$ is the smaller degree the split $(G,G^c)$ has $\mathsf s=\deg G>2\tau$ and $\mathsf l>\mathsf B_*$, and otherwise $\mathsf s=\mathsf D-\deg G>\mathsf B_*\ge2\tau$ and $\mathsf l=\deg G>\mathsf s>\mathsf B_*$; either way the split succeeds. Moreover $3\mathsf U>\mathsf D$, since $\frac23\mathsf D-\mathsf B_*\ge\frac23\lambda-\mathsf B_*>0$; so at most two variables have degree $\ge\mathsf U$.

\emph{Step 2: the square margins when $\mathsf j\le2\tau$.} Then $\mathsf m-\frac12\mathsf n\ge\frac14(\mathsf m+\mathsf n)=\frac14(\mathsf D-\mathsf j)$, $\mathsf n\le\frac12(\mathsf D-\mathsf j)$ and $\mathsf b=\sigma-\mathsf m\ge\frac12\sigma-\tau$, and, with $\mathsf D-4\mathfrak i=\lambda-3\mathfrak i\ge\lambda-3(3\mathfrak g+4\mathfrak r)\ge\lambda-4(3\mathfrak g+4\mathfrak r)$ and $\mathsf D\le1+\varsigma$,
\[
 \begin{aligned}
 \mathsf b-\mathsf j-2\mathfrak r-\tfrac15\mathfrak v&\ge\tfrac12\sigma-3\tau-2\mathfrak r-\tfrac15\mathfrak v=\tfrac12-3\varsigma+\tfrac{11}2\mathfrak g+\tfrac{13}2\mathfrak r+\tfrac52\mathfrak v,\\
 4\mathsf P_1&\ge1-9\varsigma+5\mathfrak g+10\mathfrak r+\tfrac{36}5\mathfrak v+2(\mathsf Z+\mathsf E),\\
 4\mathsf P_2&\ge1-7\varsigma+4\mathfrak g+8\mathfrak r+\tfrac{23}5\mathfrak v,\\
 \mathsf P_3&\ge\tfrac12-3\varsigma+4\mathfrak g+6\mathfrak r+\tfrac52\mathfrak v .
 \end{aligned}
\]
For $\varsigma<\frac18$ the first line gives $\mathsf b>\mathsf j+\mathfrak r$ and $\mathsf P_0\ge\mathsf b-\mathsf j-\frac15\mathfrak v>0$, and $\mathsf P_2,\mathsf P_3>0$; if moreover $\varsigma<\frac19$, then $\mathsf P_1>0$ and \cref{lem:square} saves.

(i) Here $\varsigma<\frac1{10}$, so Step~1 applies and $\mathsf U-4\tau\ge\lambda-\mathsf B_*-4\tau>0$. If $d_3\ge\mathsf U$, then $\mathsf D\ge3\mathsf U$, which is impossible; so $d_3,d_4\le2\tau$, and $\deg\{d_3,d_4\}=\mathsf j\le4\tau<\mathsf U$ forces $\mathsf j\le2\tau$ by Step~1. By Step~2, \cref{lem:square} saves.

(ii) Here $\frac1{10}\le\varsigma<\frac18$ and $\mathsf w\le\mathfrak b<\frac{3-21\varsigma}2$. We use the margins $\mathsf M_{\mathsf k}=\mathsf B-\frac12(\mathsf y+\mathsf b+\mathsf w-\mathsf z+\mathsf k)$ of \cref{lem:direct} without the gain, one for each $\mathsf k\in\mathcal K$. Expanding (\S6),
\[
 \begin{aligned}
 2\mathsf M_{(\mathsf R+\mathsf w,\mathsf c)}&=1-\varsigma+\mathfrak g+\mathfrak r+\tfrac45\mathfrak v+(\mathsf m-\mathsf n)+(\mathsf Z-\mathsf E)+\mathsf z-\mathsf w,\\
 2\mathsf M_{(\mathsf N,\mathsf c)}&=1-2\varsigma+3\mathfrak g+3\mathfrak r+\tfrac95\mathfrak v+\mathsf m+2\mathsf x+3\mathsf z,\\
 2\mathsf M_{(\mathsf N,\mathsf R+\mathsf d)}&=2\mathsf m+\mathsf n-2\varsigma+4\mathfrak g+3\mathfrak r-\mathfrak i+\tfrac95\mathfrak v+\mathsf x+\mathsf z-\mathsf w,\\
 2\mathsf M_{\mathsf N}'&=\mathsf m+\tfrac12\mathsf n-\tfrac32\varsigma+3\mathfrak g+3\mathfrak r+\tfrac{13}{10}\mathfrak v+\tfrac32\mathsf x+\tfrac52\mathsf z+\tfrac12\mathsf E,\\
 \mathsf M_1:=\mathsf M_{(\mathsf R+\mathsf w,\mathsf R+\mathsf d)}&=\mathsf m-\tfrac12\varsigma+\mathfrak g+\tfrac12\mathfrak r+\tfrac25\mathfrak v-\mathsf w-\tfrac12\mathfrak i-\tfrac12\mathsf E,\\
 2\mathsf M_2:=2\mathsf M_{\mathsf R+\mathsf w}'&=\mathsf m-\varsigma+2\mathfrak g+2\mathfrak r+\tfrac45\mathfrak v-\tfrac12\mathsf w+\mathsf Z+\mathsf z,\\
 2\mathsf M_3:=2\mathsf M''&=2(\mathsf m+\mathsf n)-2\tau-3\mathfrak r-\mathfrak i+\mathsf E-\mathsf w,
 \end{aligned}
\]
where the subscript $(\mathsf p,\mathsf q)$ refers to the first line of \eqref{eq:Kset}, $\mathsf M'_{\mathsf p}$ to the second and $\mathsf M''$ to the third. Below we always have $\mathsf j\le4\tau$, so $\mathsf m+\mathsf n\ge\lambda-4\tau\ge1-3\varsigma$. With $\mathsf m\ge\mathsf n$, $\mathsf E\le\mathsf Z$, $\mathsf w\le\frac12$, $4\mathfrak g+3\mathfrak r-\mathfrak i\ge-\mathfrak r\ge-\frac13\varsigma$, $2\mathsf m+\mathsf n\ge\frac32(\mathsf m+\mathsf n)$ and $\mathsf m+\frac12\mathsf n\ge\frac34(\mathsf m+\mathsf n)$, the first four margins are at least $\frac12-\varsigma$, $1-2\varsigma$, $1-\frac{41}6\varsigma$ and $\frac34-\frac{15}4\varsigma$, all positive. It remains to treat $\mathsf M_1,\mathsf M_2,\mathsf M_3$.

\emph{Case 1: $\mathsf j\le2\tau$.} By Step~2, $\mathsf P_0,\mathsf P_2,\mathsf P_3>0$ and $\mathsf b>\mathsf j+\mathfrak r$, and \cref{lem:square} saves unless $\mathsf P_1\le0$. Then $\varsigma\ge\frac19$ and $\mathsf E\le\mathsf Z+\mathsf E\le\frac12(9\varsigma-1-5\mathfrak g-10\mathfrak r-\frac{36}5\mathfrak v)$. With $\mathsf m\ge\frac12(\mathsf m+\mathsf n)\ge\frac12(\lambda-2\tau)$,
\[
 \begin{aligned}
 \mathsf M_1&\ge\tfrac12(\lambda-2\tau)-\tfrac12\varsigma+\mathfrak g+\tfrac12\mathfrak r+\tfrac25\mathfrak v-\tfrac12(3\mathfrak g+4\mathfrak r)-\mathsf w-\tfrac12\mathsf E\\
 &\ge\tfrac34-\tfrac{13}4\varsigma+\tfrac54\mathfrak g+2\mathfrak r+\tfrac{13}5\mathfrak v-\mathsf w\\
 &>\tfrac34-\tfrac{13}4\varsigma-\tfrac{3-21\varsigma}2=\tfrac{29\varsigma-3}4>0,\\
 2\mathsf M_2&\ge\tfrac12(\lambda-2\tau)-\varsigma-\tfrac12\mathsf w\ge\tfrac12-\tfrac32\varsigma-\tfrac14>0,\\
 2\mathsf M_3&\ge2(\lambda-2\tau)-2\tau-3\mathfrak g-7\mathfrak r-\mathsf w=2-4\varsigma+3\mathfrak g+3\mathfrak r+\tfrac{17}5\mathfrak v-\mathsf w>0 .
 \end{aligned}
\]
So \cref{lem:direct} saves.

\emph{Case 2: $\mathsf j>2\tau$.} By Step~1, $\mathsf j=\deg\{d_3,d_4\}\ge\mathsf U$. If $\mathsf n>2\tau$, then also $\mathsf n\ge\mathsf U$ and $\mathsf m\ge\mathsf U$, and $\mathsf D\ge3\mathsf U$, which is impossible; so $\mathsf n\le2\tau$ and $\mathsf j\le4\tau$. The sets $\{d_2,d_4\}$ and $\{d_2,d_3\}$ have degree at least $\mathsf j>2\tau$, hence at least $\mathsf U$, so $d_3,d_4\ge\mathsf U-\mathsf n$ and $\mathsf n\ge d_3\ge\mathsf U-2\tau$. Hence
\[
 \mathsf m=\mathsf D-\mathsf n-\mathsf j\ge\lambda-6\tau,\qquad\mathsf b-\mathsf j=\sigma-\mathsf D+\mathsf n\ge\sigma-\mathsf B_*-2\tau=\tfrac12\sigma-4\tau,
\]
and $\mathsf b-\mathsf j-2\mathfrak r-\frac15\mathfrak v\ge\frac12-4\varsigma+\frac{15}2\mathfrak g+\frac{19}2\mathfrak r+\frac{17}5\mathfrak v>0$, which gives $\mathsf b>\mathsf j+\mathfrak r$ and $\mathsf P_0>0$. With $\mathsf a=\sigma-\mathsf n\ge\sigma-2\tau$,
\[
 \begin{aligned}
 \mathsf P_1&\ge1-10\varsigma+16\mathfrak g+25\mathfrak r+\tfrac{87}{10}\mathfrak v+\tfrac12(\mathsf Z+\mathsf E),\\
 \mathsf P_2&\ge\tfrac12-5\varsigma+8\mathfrak g+\tfrac{25}2\mathfrak r+\tfrac{43}{10}\mathfrak v+\tfrac32\mathsf z+\mathsf x+\mathsf E,\\
 \mathsf P_3&\ge1-\tfrac{17}2\varsigma+\tfrac{29}2\mathfrak g+22\mathfrak r+\tfrac{37}5\mathfrak v+\tfrac12\mathsf z+\tfrac12\mathsf x+\tfrac32\mathsf E .
 \end{aligned}
\]
If these are positive, \cref{lem:square} saves. Otherwise, using $\mathsf E\le\mathsf Z$, one of $\mathsf E\le10\varsigma-1$, $\mathsf E\le5\varsigma-\frac12$, $\mathsf E\le\frac13(17\varsigma-2)$ holds. Now
\[
 \begin{aligned}
 \mathsf M_1&\ge(\lambda-6\tau)-\tfrac12\varsigma+\mathfrak g+\tfrac12\mathfrak r+\tfrac25\mathfrak v-\tfrac12(3\mathfrak g+4\mathfrak r)-\mathsf w-\tfrac12\mathsf E\\
 &\ge1-\tfrac{11}2\varsigma-\mathsf w-\tfrac12\mathsf E>5\varsigma-\tfrac12-\tfrac12\mathsf E,
 \end{aligned}
\]
which is $\ge0$, $\ge\frac12(5\varsigma-\frac12)\ge0$ and $\ge\frac16(13\varsigma-1)>0$ in the three cases; and $2\mathsf M_2\ge\lambda-6\tau-\varsigma-\frac12\mathsf w\ge\frac34-6\varsigma>0$, $2\mathsf M_3\ge2\lambda-10\tau-3\mathfrak r-\mathsf w=2-8\varsigma+14\mathfrak g+19\mathfrak r+7\mathfrak v-\mathsf w>0$. So \cref{lem:direct} saves.

The affine reductions in (i) and (ii) are also checked by machine, independently of Steps~1 and~2: with the four dual degrees kept separately, all seven splits of the four variables allowed to fail and the powerful-part degree free, the failure of every estimate named in (i), respectively (ii), gives $13122$, respectively $76545$, linear systems, and a computer verifies exact certificates of the kind used in (iii) at $22$, respectively $69$, nodes of the corresponding disjunction tree and checks that these nodes cover every system. In (i) the program fixes $\mathsf x=0$. This is not a monotone reduction for the alternative $\mathsf z+\mathfrak i+2\mathsf b+\mathsf a\ge\mathsf C$ of \cref{lem:square}, because $\mathsf C$ decreases in $\mathsf x$; that alternative is covered for every $\mathsf x$ by the lower bound for $4\mathsf P_1$ in Step~2.

(iii) If $\mathsf J<0$, \cref{lem:actuallevel}(i) saves; so assume $\mathsf J\ge0$. Since $\mathsf u\ge0$ gives $\mathsf Z\le\mathsf j+\frac1{10}\mathfrak v$ and $\mathsf j\le\frac12\mathsf D\le\frac12(1+\varsigma)$, the identity $\sigma-\frac12(1+\varsigma)-\frac1{10}\mathfrak v=\frac12-\frac12\varsigma-\mathfrak g-\mathfrak r-\frac1{10}\mathfrak v$ (\S6) and $\mathfrak g+\mathfrak r+\frac1{10}\mathfrak v\le\frac12\varsigma$ give $\mathsf C\ge\frac12-\varsigma>0$; so all moduli $u'$ are large. The hypotheses of an uncovered tuple are: all degrees nonnegative, $\tau\ge0$, $\mathsf m\ge\mathsf n$, $2\mathsf n\ge\mathsf j$, $\mathsf y\ge0$, $\mathsf E\le\min(\mathsf Z,\mathsf a)$, $\mathsf w\le\min(\mathsf C,\frac12)$, $\mathsf m\le\frac12\sigma+\tau$, $\mathsf u\ge0$, $\mathfrak i\le3\mathfrak g+4\mathfrak r$ and $\mathsf J\ge0$. Failure of the split $\{d_1,d_2\}\,|\,\{d_3,d_4\}$ and of the split $\{d_1\}\,|\,\{d_2,d_3,d_4\}$, of \cref{lem:recip}, of \cref{lem:direct} without and with the gain $(\mathsf E-\mathsf z)_+$, and of \cref{lem:actuallevel}(ii) means that at least one entry in each of the following rows is nonnegative:
\begin{equation}\label{eq:rows}
 \begin{array}{l}
 2\tau-\mathsf j,\quad \mathsf B_*-\mathsf m-\mathsf n;\\
 2\tau-\mathsf m,\quad 2\tau-(\mathsf D-\mathsf m),\quad \mathsf B_*-(\mathsf D-\mathsf m);\\
 \mathsf u-\tfrac12\mathsf C,\quad \mathsf u+\tfrac12\mathsf C-\mathsf b;\\
 \text{for one }\mathsf k\in\mathcal K,\ \text{both }L_{\mathsf k}=\mathsf y+\mathsf b+\mathsf w-\mathsf z+\mathsf k-2\mathsf B\ \text{and}\ L_{\mathsf k}-\mathsf E+\mathsf z;\\
 \text{one of the eight affine expansions of }2\mathsf F_*-2\mathsf B .
 \end{array}
\end{equation}
(For the fourth row: failure with the gain means $L_{\mathsf k}\ge(\mathsf E-\mathsf z)_+$ for some $\mathsf k$, which implies both entries; for the fifth, each positive part in $\mathsf F_*$ is replaced by $0$ or by its argument.) On the third band the first two rows may be replaced by the single alternative $4\tau-\sigma\ge0$, under which \cref{lem:unsignedplus} does not apply. For a band $[\varsigma_0,\varsigma_1)$ with cap $f$, append $\varsigma-\varsigma_0\ge0$, $\varsigma_1-\varsigma-g_*\ge0$ and $f(\varsigma)-\mathsf w-g_*\ge0$; a tuple inside the band would have $g_*>0$, since $\mathsf w\le\mathfrak b<f(\varsigma)$. This gives $2\cdot3\cdot2\cdot7\cdot8=672$ linear systems on each of the first two bands and $7\cdot2\cdot7\cdot8=784$ on the third. For each of them a computer verifies, in exact rational arithmetic, a nonnegative combination of the inequalities and of the constant $1$ equal to $-g_*$ or to $-1$; so $g_*\le0$, a contradiction.
\end{proof}

\begin{remark}[Without the represented level]\label{rem:fallback}
\Cref{lem:actuallevel} is the only place where \eqref{eq:LS}, \eqref{eq:EX} and \eqref{eq:kuzgeneral} are used. Without it, (i) and (ii) of \cref{lem:cases} are unchanged, and in (iii) the hypothesis $\mathsf J\ge0$, the second entry of the fourth row of \eqref{eq:rows} and the fifth row are dropped; the same computation verifies the $2\cdot3\cdot2\cdot7=84$ resulting systems on each of the bands $\frac1{10}\le\varsigma<\frac17$ and $\frac17\le\varsigma<\frac14$ with the caps $\mathfrak b<\frac{13-47\varsigma}{30}$ and $\mathfrak b<\frac{7-19\varsigma}{30}$. With these caps in place of those of (iii), the argument below gives the coefficient
\[
 \frac{1767870158617164811674095218440187374694865783}{42458728008803622772455544218525820350000000}=41.6373792038\ldots
\]
(computed in \S6), which then depends only on (Q1) and (Q4) among the quoted estimates.
\end{remark}

\begin{proof}[Proof of \cref{prop:nonprincipalplus}]
By \cref{lem:degrees} it suffices to prove (a) or (b) for every datum $\mathfrak D$, with a $\delta'>0$ that does not depend on $\mathfrak D$ or $T$. The degrees of all data lie in a compact set $\mathcal X$: $(\varsigma,\mathfrak b)$ lies in the compact image $E^*\subset\mathcal P^{+*}$ of $E$, and all other degrees lie in $[0,2]$. Adjoin the degrees $\mathsf z,\mathsf x,\mathsf E,\mathsf w$ of $l,f,e,S_0$, which also lie in $[0,2]$. Each estimate of \cref{lem:cases} has as margin a minimum of finitely many affine functions of these degrees, and the margins of \cref{lem:direct,lem:actuallevel} (a maximum over the cases of \eqref{eq:Kset} or over the positive parts) are continuous; so if one has a positive margin at a point it has margin at least half that value on a neighbourhood. By \cref{lem:cases} every point has such an estimate, and by compactness there are $\delta_0>0$ and a finite assignment of estimates to the sets of a finite open cover such that at every point the assigned estimate has margin at least $\delta_0$. Since the degrees of a datum at height $T$ differ from points of $\mathcal X$ by $O(\epsilon)+o(1)$, we may take $\epsilon$ small and $T$ large so that all dyadic, divisor, truncation and Mellin losses, all $Q^{o(1)}$ factors, and the deviation of the actual degrees from the cover are below $\frac12\delta_0$.

For a datum whose limiting degrees lie in a set assigned to \cref{lem:unsignedplus}, or with $\tau<0$, we prove (b), respectively (a), directly; the choice does not depend on $\mathfrak p$, nor on the modulus within the datum. Otherwise we apply \cref{lem:doublepoisson} at every modulus of the datum, reduce the modulus as in \eqref{eq:reducedsum}, and split the result into the $Q^{o(1)}$ sums over the dyads of $l,f,S_0$ and the divisors $e\mid lf$; for each of these sums the assigned estimate (\cref{lem:square,lem:direct,lem:recip,lem:actuallevel}, or the trivial bound if $\mathsf u<0$) gives the saving $Q^{\mathsf u+\delta_0/2}$, which is (a) for the datum. In \cref{lem:direct,lem:actuallevel} the choice is made on whole dyads of the labels and of $S_0$, so no coefficient depending on the modulus is attached to the smooth variable $C_0$, and the signed sum over $C_0$ is kept with the weight $\overline F$ at $q=hk$. All estimates hold uniformly for $\boldsymbol\sigma$ in compact subsets of $(i\R)^4$ and $\Phi$ in bounded subsets of $\mathcal W$, because they use only the size and the rescaled derivatives of the weights. Summing over the $(\log T)^{O(1)}$ data gives $\Sigma_F\ll T^{1-\eta}$.
\end{proof}

\subsection{The mixed pairing}

\begin{proposition}[Mixed pairing]\label{prop:factorableplus}
Let $c\in C_c^\infty(\R^3)$ be real with $\operatorname{supp}c\subset\mathcal F^{+\sharp}$, and let $b\in C_c^\infty(\R^3)$. Uniformly for $\Phi$ in bounded subsets of $C_c^\infty$ with a fixed compact support in $(1,2)$,
\[
 \ip{P_b}{P_c}_\Phi=a_3TL^9\Big(\int\Phi\Big)\int_\Om h_bh_c+o(TL^9).
\]
\end{proposition}
\begin{proof}
We go through the proof of \cref{prop:factorable} and record, for each step, what it uses.

\emph{Reduction to elementary sources.} This step of the proof of \cref{prop:balanced} uses only \cref{lem:localiseP}, \cref{lem:sourcenorm}, \cref{lem:sixth} and the fact that $h_c$ vanishes unless $r\in I_c$. These hold for every real $c\in C_c^\infty((0,1)^3)$ supported in $1<|a|<2$, in particular on $\mathcal F^{+\sharp}$. So it suffices to prove \eqref{eq:factelementary} for $\boldsymbol u\in(i\R)^3$ with distinct coordinates, and as there $\ip{P_{\boldsymbol u,\phi}}{P_c}_\Phi=\ip Z{P_c}_\Phi+O(T^{-A})$ with $Z(s)=\prod_k\zeta(s+u_k/L)$.

\emph{Fourier inversion in the third coordinate.} Put
\[
 \mathcal L_3c(x;\beta)=\int c(x_1,x_2,a_3)e^{\beta(1-a_3)}\dd a_3,\qquad F_\beta=\overline{\mathcal L_3c(\cdot;\beta)} .
\]
The function $F_\beta$ is supported in the compact projection $E$ of $\operatorname{supp}c$ to the first two coordinates, which lies in $\mathcal P^+$ because $\operatorname{supp}c\subset\mathcal F^{+\sharp}$; and $\int|\widehat{F_\beta}|\ll_j(1+|\beta|)^{-j}$, $|\mathcal L_3c|\ll_j(1+|\beta|)^{-j}$, by integration by parts in $a_3$, as there. The representation of $\overline{P_c(s)}$ as an integral over $\beta$ of $e^{-\beta}\overline{A_{F_\beta}(s)}$ times a finite $m$-sum, \cref{lem:expansion} for the $m$-sum, the functional equation and Stirling's formula are used exactly as in the proof of \cref{prop:factorable}; none of them depends on the support of $c$. This gives
\[
 \ip Z{P_c}_\Phi=\frac1{2\pi i}\int_{(0)}\Big(I_{F_\beta}(\boldsymbol\sigma(\beta))+\mathcal E(\beta)\Big)\dd\beta+O(T^{-A}),
\]
where $\boldsymbol\sigma(\beta)=(u_1,u_2,u_3,\beta)$, with $|\mathcal E(\beta)|\ll\frac{1+|\beta|}L\|Z\|_\Phi\|\zeta(s+\beta/L)A_{F_\beta}\|_\Phi$ for $|\beta|\le L$.

\emph{Norms.} By H\"older's inequality, Harper's bound \cite{Harper} (which rests on the Riemann hypothesis, a theorem of Alp\"oge \cite{Alpoge}) and \cref{lem:sixth}(i), applied to the Fourier superpositions of $A_{F_\beta}$, we have $\|Z\|_\Phi\ll(TL^9)^{1/2}$, $\|\zeta(s+\beta/L)A_{F_\beta}\|_\Phi\ll_j(TL^9)^{1/2}(1+|\beta|)^{-j}$ and $|I_{F_\beta}(\boldsymbol\sigma(\beta))|\ll_jTL^9(1+|\beta|)^{-j}$ for $|\beta|\le LT/4$; these bounds use only $\int|\widehat{F_\beta}|$ and not the support. Hence $\int|\mathcal E|=o(TL^9)$, as there.

\emph{The limit.} For each fixed $\beta\in i\R$, \cref{prop:fourshiftplus}, applied to $F_\beta$ (supported in the compact set $E\subset\mathcal P^+$) and $\boldsymbol\sigma(\beta)$, gives $I_{F_\beta}(\boldsymbol\sigma(\beta))=a_3TL^9(\int\Phi)\int\mathcal L_3c(x;\beta)S(\boldsymbol\sigma(\beta);x)\dd x+o(TL^9)$; this is the only place where $\mathcal P^+$ enters. By dominated convergence, with the majorant $TL^9(1+|\beta|)^{-j}$,
\[
 \ip Z{P_c}_\Phi=a_3TL^9\Big(\int\Phi\Big)\frac1{2\pi i}\int_{(0)}\int_{\R^2}\mathcal L_3c(x;\beta)S(\boldsymbol\sigma(\beta);x)\dd x\dd\beta+o(TL^9).
\]
The contour shift to $\Re\beta=2r_0$ uses only the polynomial bound for $S(\boldsymbol\sigma(\beta);x)$ on compact sets of $x$ and the decay of $\mathcal L_3c$, and identifies the double integral with $\Lambda'$ in the proof of \cref{lem:balkernel}. Part (b) of that proof needs only that $c\in C_c^\infty((0,1)^3)$ is real with $a_1+a_2<1$ on its support; on $\mathcal F^{+\sharp}$ we have $a_1+a_2<\frac35$. Hence $\Lambda'=\int_\Om h_{e^{-\boldsymbol u\cdot a}}h_c$, which is \eqref{eq:factelementary}.
\end{proof}

\subsection{Proof of \Cref{thm:mixedrow}}\label{ssec:mixedrowproof}
Fix $\Phi\in\mathcal W$, and let $\psi,\eta$ be real. For $\psi\in C_c^\infty(\mathcal F^+)$, \cref{lem:synthesis} gives $h_{\mathsf V\psi}=\psi$, and \cref{lem:sixth}(ii) gives $\|F_\psi\|_\Phi\ll(TL^9)^{1/2}$. The coefficients of $F_\psi$ live at $\log N/L\in I_\psi\subset(1,\frac32)$, and those of $P_{(\mathsf V\psi)^J}$ at $\log N/L\in3-I_\psi\subset(\frac32,2)$. We prove the limits required by \cref{def:faithful} for all pairs involving $\mathcal T_{\mathcal F^+}$; pairs within $\mathcal T_5$ are covered by \Cref{thm:row}. Since $\mathcal F\subset\mathcal F^+$, the family $\mathcal T_{\mathcal F^+}$ contains $\mathcal T_{\mathcal F}$, and $\mathcal T_7\supseteq\mathcal T_6$.

\emph{(a) Numerators.} As in part (a) of the proof of \Cref{thm:row}, $\mathrm N_T(F_\psi)=\ip{P_{c_0}}{F_\psi}_\Phi+O(T^{-A})$ with $c_0(a)=\varphi_1(|a|)\omega_0(a)$, where $\varphi_1\in C_c^\infty((1,2))$ is $1$ near $I_\psi$. By \cref{prop:factorableplus} with $c=\mathsf V\psi$ and $b=c_0$, $\mathrm N_T(F_\psi)/\mathcal N_T\to\ip{h_{c_0}}\psi=\ip{\varphi_1(r)v}\psi=\ip v\psi$. Further $\mathrm N_T(\mathcal JF_\psi)=\overline{\mathrm N_T(F_\psi)}$ has the limit $\ip v\psi=\ip v{\psi\circ J}$, as $v\circ J=v$.

\emph{(b) Gram entries within $\mathcal T_{\mathcal F^+}$, and the witness norm.} For real $\psi,\eta\in C_c^\infty(\mathcal F^+)$, \cref{prop:factorableplus} with $c=\mathsf V\psi$ and $b=\mathsf V\eta$ gives $\mathrm G_T(F_\eta,F_\psi)/\mathcal N_T\to\ip{h_{\mathsf V\eta}}{h_{\mathsf V\psi}}=\ip\eta\psi$. The case $\eta=\psi$ is the norm of the witness, $\|F_\psi\|_\Phi^2/\mathcal N_T\to\|\psi\|^2$; it needs no separate argument, because \cref{prop:factorableplus} places no condition on the support of $b$ beyond smoothness and compactness. Next $\mathrm G_T(\mathcal JF_\psi,\mathcal JF_\eta)=\overline{\mathrm G_T(F_\psi,F_\eta)}$ has the limit $\ip\psi\eta=\ip{\psi\circ J}{\eta\circ J}$. For the cross terms, \cref{lem:reflcube} and the norm bound give $\mathrm G_T(\mathcal JF_\eta,F_\psi)=\ip{P_{(\mathsf V\eta)^J}}{F_\psi}_\Phi+o(TL^9)$, and this is $O(T^{-A})$ by \cref{lem:localiseP}, since $3-I_\eta\subset(\frac32,2)$ and $I_\psi\subset(1,\frac32)$ are disjoint. This agrees with $\ip{\eta\circ J}\psi=0$, since $J\mathcal F^+\cap\mathcal F^+=\emptyset$. The entry $\mathrm G_T(F_\psi,\mathcal JF_\eta)$ is the complex conjugate.

\emph{(c) Cross-pairings with $\mathcal T_3^\infty$.} As in part (c) of the proof of \Cref{thm:row}, it suffices to show $\mathrm G_T(K,F_\psi)/\mathcal N_T\to\ip{h_{F_K}}\psi$ and $\mathrm G_T(\mathcal JK,F_\psi)/\mathcal N_T\to\ip{h_{F_K}\circ J}\psi$ for $K\in\{U,V,D_\phi\}$. For $U$ and $V$ this is \cref{lem:factors} with $W_\psi$, $I_\psi$, $\psi^\flat$ replaced by $F_\psi$, $I_\psi$, $\psi$; its proof uses $\|F_\psi\|_\Phi\ll(TL^9)^{1/2}$ (\cref{lem:sixth}(ii)), \cref{lem:localiseP}, the pairing of $F_\psi$ with a smooth compact source (\cref{prop:factorableplus}), and in Step~4 that $\psi$ vanishes unless $r\in I_\psi$. For $K=D_\phi$, with $\operatorname{supp}\phi\subset(0,\theta)$ and $\theta<1$, the profiles $\phi(r)v$ and $\phi(3-r)v$ vanish on $1<r<2$; $\mathrm G_T(D_\phi,F_\psi)=O(T^{-A})$ by \cref{lem:localiseP}, and $\mathrm G_T(\mathcal JD_\phi,F_\psi)=O(T^{-A})$ by the integration by parts in part (c) of \Cref{sec:markedproof} with $\rho=\max I_\psi<\frac32$, since $\theta+\rho<3$.

\emph{(d) Cross-pairings with marked tests.} Let $\eta\in C_c^\infty(\mathcal Y)$ be real. Every $a\in\operatorname{supp}\mathsf V\psi$ has $|a|\in I_\psi\subset(1,2)$, so \cref{prop:mixed} gives $\mathrm G_T(F_\psi,W_\eta)/\mathcal N_T\to\ip\psi{\eta^\flat}$, and $\mathrm G_T(W_\eta,F_\psi)$ is the conjugate. By \cref{lem:reflcube} and $\|W_\eta\|_\Phi\ll(TL^9)^{1/2}$, $\mathrm G_T(F_\psi,\mathcal JW_\eta)=\overline{\ip{P_{(\mathsf V\psi)^J}}{W_\eta}_\Phi}+o(TL^9)$; every $a\in\operatorname{supp}(\mathsf V\psi)^J$ has $|a|\in3-I_\psi\subset(1,2)$, so \cref{prop:mixed} and \cref{lem:reflectedprofile} give the limit $\ip{\psi\circ J}{\eta^\flat}=\ip\psi{\eta^\flat\circ J}$. The entries with $\mathcal JF_\psi$ are the complex conjugates.

\emph{(e) Cross-pairings with $\mathcal T_{\mathcal R}$.} For real $\eta\in C_c^\infty(\mathcal R)$, \cref{prop:balanced} with $c=\mathsf V\eta$ and $b=\mathsf V\psi$ gives $\mathrm G_T(F_\psi,B_\eta)/\mathcal N_T\to\ip\psi\eta$ (the supports of $\psi$ and $\eta$ may overlap). By \cref{lem:reflcube}, $\mathrm G_T(\mathcal JF_\psi,B_\eta)=\ip{P_{(\mathsf V\psi)^J}}{B_\eta}_\Phi+o(TL^9)$, and \cref{prop:balanced} with $b=(\mathsf V\psi)^J$ and \cref{lem:reflectedprofile} give the limit $\ip{\psi\circ J}\eta$ (which is $0$). Likewise $\mathrm G_T(\mathcal JB_\eta,F_\psi)/\mathcal N_T\to\ip{\eta\circ J}\psi$ by \cref{lem:reflcube} and \cref{prop:factorableplus} with $c=\mathsf V\psi$ and $b=(\mathsf V\eta)^J$. The remaining entries are complex conjugates of these, by $\ip{\mathcal JH}{\mathcal JK}_\Phi=\overline{\ip HK_\Phi}$ and $\ip H{\mathcal JK}_\Phi=\overline{\ip{\mathcal JH}K_\Phi}$, and all the limits are real.

All limits in (a)--(e) hold for every $\Phi\in\mathcal W$. Hence $\mathcal T_7$ is faithful.

\emph{The span.} The real $\psi\in C_c^\infty(\mathcal F^+)$ are dense in $L^2(\mathcal F^+)$ and the $\psi\circ J$ in $L^2(J\mathcal F^+)$, so $\mathcal H(\mathcal T_7)\supseteq L^2(\mathcal F^+\cup J\mathcal F^+)$; with $v\one_{S_3\cup\mathcal Z\cup\mathcal R\cup J\mathcal R}\in\mathcal H(\mathcal T_5)$ we get $v\one_{U_7}=v\one_{S_3\cup\mathcal Z\cup\mathcal R\cup J\mathcal R}+v\one_{U_7\setminus(S_3\cup\mathcal Z\cup\mathcal R\cup J\mathcal R)}\in\mathcal H(\mathcal T_7)$.

\emph{The measure.} Since $\mathcal F^+\subset\{1<r<\frac32\}$, as in the proof of \Cref{thm:factorable} we have $\mu(U_7)=C_0+2\mu((N_0\cap\mathcal F^+)\setminus\mathcal R)$. In the coordinates $(k,x,z)$ the set $(N_0\cap\mathcal F^+)\setminus\mathcal R$ lies in $\frac12<k<\frac7{12}$: for $k\ge\frac7{12}$ the cap $x_+(k)=\frac{3-5k}2$ is at most the lower face $\frac12(k-\frac12)$ that $N_0$ imposes on $x$ (from $z>\max(k-x,1-k)$ and $z<\frac12+x$). On $N_0$ the upper face for $z$ is $\frac12+x$ and the lower face is $\max(k-x,1-k)$; outside $\mathcal R$ it is also $z\ge z_0(k)=\frac{18-28k}7$, which matters only for $k<\frac{11}{21}$, since $z_0(k)-(1-k)=\frac{11-21k}7$. The gain over $\mu(U_6)$ consists of three parts, integrated with the reflected density $\frac{9!}2(k-2x)^2(z-x)^2(z-k+x)^2$: the part of $N_0\setminus(\mathcal R\cup\mathcal F)$ with $\frac12<k<\frac{10}{19}$, where every mark is admitted; the two cells with $\frac{10}{19}<k<\frac8{15}$ below the cap $a_0(k)=\frac{21-39k}2$; and the seven cells of \cref{tab:sevencells}.

\begin{table}[ht]
\centering
\small
\caption{The seven cells of $(N_0\cap\mathcal F^+)\setminus\mathcal R$ with $\frac{35}{66}<k<\frac7{12}$ beyond the cap $a_0(k)=\frac{21-39k}2$, with $a_1=\frac{7-12k}4$, $a_2=\frac{2-3k}4$, $a_3=\frac{3-5k}2$. The upper face for $z$ is $\frac12+x$; the masses are $2\mu$ of the cell (both reflected halves).}\label{tab:sevencells}
\begin{tabular}{@{}lllll@{}}
\toprule
Range of $k$ & Lower $x$ & Upper $x$ & Lower $z$ & $2\mu$ \\
\midrule
$(\frac{35}{66},\frac8{15})$ & $a_0(k)$ & $a_1(k)$ & $1-k$ & $0.000809878521\ldots$ \\
$(\frac8{15},\frac7{13})$ & $\frac12(k-\frac12)$ & $2k-1$ & $k-x$ & $0.000206853802\ldots$ \\
$(\frac8{15},\frac7{13})$ & $2k-1$ & $a_1(k)$ & $1-k$ & $0.002612978444\ldots$ \\
$(\frac7{13},\frac6{11})$ & $\frac12(k-\frac12)$ & $2k-1$ & $k-x$ & $0.000484475814\ldots$ \\
$(\frac7{13},\frac6{11})$ & $2k-1$ & $a_2(k)$ & $1-k$ & $0.000292151447\ldots$ \\
$(\frac6{11},\frac47)$ & $\frac12(k-\frac12)$ & $a_2(k)$ & $k-x$ & $0.001041043286\ldots$ \\
$(\frac47,\frac7{12})$ & $\frac12(k-\frac12)$ & $a_3(k)$ & $k-x$ & $0.000022888772\ldots$ \\
\midrule
Total & & & & $0.005470270089\ldots$ \\
\bottomrule
\end{tabular}
\end{table}

The cap $a_0$ meets $a_1$ at $k=\frac{35}{66}$, the cap $a_2$ meets $2k-1$ at $k=\frac6{11}$, and $a_3$ meets $\frac12(k-\frac12)$ at $k=\frac7{12}$. The first part of the gain has mass $0.017403320153\ldots$, of which the strip $\frac{11}{21}<k<\frac{10}{19}$ contributes $\frac{6048414814819537375}{1666766405108292973184}=0.003628831728\ldots$; the second has mass $\frac{1592111573351043355967}{264644248141002375000000}=0.006016044499\ldots$; the third is the total of \cref{tab:sevencells}. Exact integration gives
$2\mu((N_0\cap\mathcal F^+)\setminus\mathcal R)=C-C_0$ with
\[
 C=\frac{112666832560750258461501594014215737455802061}{2705696476246250148185796472625838239963136},
\]
and the seven cells of \cref{tab:sevencells} add exactly
\[
 \frac{268755110285287452500182503673016801}{49130135420790450741077959596000000000}=0.00547027008949\ldots .
\]
These values are computed by exact iterated polynomial integration and, independently, by exact triangulation of each cell. The same computation checks that each cell lies in $N_0$ and below the cap of the estimate that covers it, that the cells meet neither $\mathcal R$ nor each other in positive measure, and that they exhaust the region defined directly by \eqref{eq:xplus}.

\emph{The passage to the liminf.} \Cref{prop:projection} turns faithfulness into a bound for the sharp moment: for each $\epsilon>0$ it chooses a finite combination of tests, applies the smoothed inequality \eqref{eq:smoothbound}, which holds for every $\Phi\in\mathcal W$ and all large $T$, with a $\Phi\le\one_{(1,2)}$ of integral at least $1-\epsilon$, at the heights $2^{-k}T$, $1\le k\le K$, and sums. It uses no bound for $M_6(T)$; Harper's bound enters the proof of \Cref{thm:mixedrow} only through the norm bounds of \cref{lem:sixth,lem:sourcenorm} in (a)--(e) and in \cref{prop:factorableplus}. Its second part uses only $v\one_{U_7}\in\mathcal H(\mathcal T_7)$ and $U_7=JU_7$. Hence $\liminf_{T\to\infty}M_6(T)/(c_3TL^9)\ge\mu(U_7)=C$. \qed

\emph{Computation.} The certificates of (iii) and the computations of this section are checked by one program, which uses only SymPy and exact rational arithmetic and runs in under a minute. It also checks the case analyses of \cref{lem:cases}(i)--(ii) and every affine identity used in \cref{lem:degrees,lem:unsignedplus,lem:doublepoisson,lem:square,lem:direct,lem:actuallevel,lem:cases}, computes the coefficient of the remark after \cref{lem:cases}, and recomputes the ceiling of \cref{prop:ceiling} with $\mathcal F^+$ in its list. The certificates were found by floating-point linear programming and are checked exactly.

\section{Proof of \Cref{thm:main}}\label{sec:main}

By \Cref{thm:mixedrow} the family $\mathcal T_7$ is faithful and $v\one_{U_7}\in\mathcal H(\mathcal T_7)$. Faithfulness means that the limits of \cref{def:faithful} hold for every $\Phi\in\mathcal W$. The first part of \cref{prop:projection} turns these smoothed limits into a bound for the sharp moment: it applies \eqref{eq:smoothbound}, with a $\Phi\le\one_{(1,2)}$ of integral at least $1-\epsilon$, at the heights $2^{-k}T$ for $1\le k\le K$, sums, and lets $\epsilon\to0$ and $K\to\infty$. No upper bound for $M_6(T)$ is used in this step; Harper's bound, and with it the Riemann hypothesis \cite{Alpoge}, enters only through the norm bounds of \cref{lem:sixth,lem:sourcenorm} in the proofs of faithfulness. The second part of \cref{prop:projection} uses only $v\one_{U_7}\in\mathcal H$ and $JU_7=U_7$, so
\[
 \liminf_{T\to\infty}\frac{M_6(T)}{c_3T(\log T)^9}\ \ge\ 9!\,\|P_{\mathcal H(\mathcal T_7)}v\|^2\ \ge\ \mu(U_7)=C .
\]
This proves \Cref{thm:main}. \qed

For the record, the regions are nested as follows, and so are the closed spans of the families. (In passing from $\mathcal T_3$ to $\mathcal T_4$ the family $\mathcal T_3$ is replaced by its smooth part $\mathcal T_3^\infty$, which has the same closed span by \Cref{thm:marked}; the assembly uses only the regions and the spans.)
\begin{center}
\small
\begin{tabular}{@{}lllr@{}}
\toprule
Family & Region & Added by & $\mu$ \\
\midrule
$\mathcal T_3$ & $S_3$ & \Cref{thm:third} & $41.3163069650$ \\
$\mathcal T_4=\mathcal T_3^\infty\cup\mathcal T_{\mathcal M}$ & $S_3\cup\mathcal Z$ & \Cref{thm:marked} & $41.5917968750$ \\
$\mathcal T_5=\mathcal T_4\cup\mathcal T_{\mathcal R}$ & $\cup\,\mathcal R\cup J\mathcal R$ & \Cref{thm:row} & $41.5958462794$ \\
$\mathcal T_6=\mathcal T_5\cup\mathcal T_{\mathcal F}$ & $\cup\,\mathcal F\cup J\mathcal F$ & \Cref{thm:factorable} & $41.6117132746$ \\
$\mathcal T_7=\mathcal T_5\cup\mathcal T_{\mathcal F^+}$ & $\cup\,\mathcal F^+\cup J\mathcal F^+$ & \Cref{thm:mixedrow} & $41.6406029094$ \\
\bottomrule
\end{tabular}
\end{center}
The missing mass $42-C=0.35939709059359626984\ldots$ is that of $(N_0\setminus(\mathcal R\cup\mathcal F^+))\cup J(N_0\setminus(\mathcal R\cup\mathcal F^+))$. Inside $N_0$ the region $\mathcal R$ is limited by the trace condition $21k+7r<18$, which comes from the exponents in Heath-Brown's bounds (\cref{lem:exponents}), and $\mathcal F^+$ by the mark caps \eqref{eq:xplus}, which are the conditions under which the case analysis of \cref{lem:cases} closes. Both regions meet $N_0$ only in $\{k<\frac7{12}\}$; on $N_0\cap\{k>\frac7{12}\}$, of mass $\frac{75263}{524288}=0.1435527801\ldots$, no pairing beyond $\mathcal T_4$ is available.

\begin{remark}[A variance form of the bound]\label{rem:variance}
Let $D(x)=\sum_{n\le x}d_3(n)$ and
\[
 V(T)=\frac{T}{2c_3(\log T)^9}\int_T^{T^{3/2}}\big[D(e^{2/T}x)-2D(e^{1/T}x)+D(x)\big]^2\frac{\dd x}{x^2}.
\]
In separate work in preparation, the identity $m_*(T)=2+2V(T)+o(1)$ is derived, for a fixed smooth unit-mass time average $m_*(T)$ of $|\zeta(\frac12+iTy)|^6/(c_3(\log T)^9)$; it uses the Riemann hypothesis \cite{Alpoge} and a central-polynomial estimate that is not proved in this paper. Granting it, \Cref{thm:main} gives
\begin{multline*}
 \liminf_{T\to\infty}V(T)\ \ge\ \frac{C-2}2\\
 =\frac{107255439608257758165130001068964060975875789}{5411392952492500296371592945251676479926272}\\
 =19.8203014547\ldots .
\end{multline*}
The conjecture $M_6(T)\sim42\,c_3T(\log T)^9$ of Conrey and Ghosh corresponds to $V(T)\to20$. Neither $42$ nor $20$ is proved here.
\end{remark}

\section{The limit of support enlargement: proof of \Cref{prop:ceiling}}\label{sec:ceiling}

Let $\mathcal U$ be the union of the sets listed in \Cref{prop:ceiling}, that is, of $\mathcal D_+=\{x+y<z\}$, $\{r<1\}$, $\{x<\frac14\}$, $\mathcal R$ and $\mathcal F^+$, and of their images under $J$. If every test in a faithful family has its profile in $L^2(\mathcal U)$, then $\mathcal H\subseteq L^2(\mathcal U)$, and
\[
 9!\,\|P_{\mathcal H}v\|^2\le9!\,\|P_{L^2(\mathcal U)}v\|^2=9!\,\|v\one_{\mathcal U}\|^2=\mu(\mathcal U).
\]
So it suffices to show that $\mu(\mathcal U)$ is the stated number. The set $U_7$ of \Cref{thm:mixedrow} is contained in $\mathcal U$: by definition $\mathcal D_{\rm acc}\subseteq\{x<\frac14\}\cup J\{x<\frac14\}$ (if $x\ge\frac14$ and $\varrho\ge\frac14$ then $3\min(x,\varrho)+5\max(x,\varrho)\ge2$; and $x+y<\frac12$ forces $x<\frac14$), $\mathcal A=\{r<1\}\cup J\{r<1\}$, and $\mathcal Z\subseteq\mathcal D_+\cup J\mathcal D_+$ (\Cref{sec:marked}). Thus \Cref{prop:ceiling} bounds every extension of the present argument that adds tests only with profiles in these supports.

\emph{The complement.} Since $J\mathcal D_+=\{y+z-x>1\}$, $J\{r<1\}=\{r>2\}$ and $J\{x<\frac14\}=\{z>\frac34\}$, the complement in $\Om$ of the first three sets and their reflections is, up to a null set,
\[
 B=\{(x,y,z)\in\Om:\ 1<r<2,\ x>\tfrac14,\ z<\tfrac34,\ x+y>z,\ y+z-x<1\},
\]
which is $J$-invariant. On $\mathcal F$ we have $x<\frac k2-\frac32(2k-1)=\frac32-\frac52k<\frac14$, as $k>\frac12$; hence $\mathcal F\cap B=\emptyset$, and $J\mathcal F\cap B=J(\mathcal F\cap B)=\emptyset$.

The larger region $\mathcal F^+$ does meet $B$, in the part of $\{\frac12<k<\frac{10}{19}\}$ with $x>\frac14$ (where every mark is admitted), and only there: for $k\ge\frac{10}{19}$ the cap \eqref{eq:xplus} is below $\frac9{38}<\frac14$. Since $\mathcal F^+\subset\{r<\frac32\}$, $J\mathcal F^+\subset\{r>\frac32\}$ and $\mathcal R\subset\{r<\frac{15}{14}\}$, adding $\mathcal F^+$ to the list increases $\mu(\mathcal U)$ by
\[
 \begin{split}
 2\mu\big((\mathcal F^+\cap B)\setminus\mathcal R\big)&=\frac{196813181035147418047099}{474637777079666241191850000000}\\
 &=4.1465974800\ldots\cdot10^{-7}.
 \end{split}
\]
This increment, and the value of $\mu(\mathcal U)$ found below, have also been checked by computer, together with $\mu(B)$ and $2\mu(B\cap\mathcal R)$.

As $\mathcal R\cap J\mathcal R=\emptyset$ and $J$ preserves $\mu$ and $B$,
\[
 \mu(\mathcal U)=42-\mu(B)+2\mu(B\cap\mathcal R)+2\mu\big((B\cap\mathcal F^+)\setminus\mathcal R\big).
\]

\emph{The mass of $B$.} Let $A(a,b,c)=(\frac{a+b}2,\frac{a+c}2,\frac{b+c}2)$, with inverse $(x,y,z)\mapsto(x+y-z,\,x-y+z,\,y+z-x)$. It preserves the trace, has Jacobian $\frac14$, and halves the three pairwise differences, so $v\circ A=\frac18v$ (as $v=\frac12\Delta$) and $\mu(A(S))=2^{-8}\mu(S)$ for measurable $S\subseteq\Om$. The conditions $x<y<z$, $x+y>z$ and $y+z-x<1$ on $(x,y,z)=A(a,b,c)$ say that $0<a<b<c<1$, and $x>\frac14$, $z<\frac34$ say that $a+b>\frac12$, $b+c<\frac32$. Hence $B=A(B')$ with
\[
 \begin{gathered}
 B'=\{(a,b,c)\in\Om:\ 1<a+b+c<2\}\setminus(E'\cup JE'),\\
 E'=\{(a,b,c)\in\Om:\ a+b\le\tfrac12,\ a+b+c>1\},
 \end{gathered}
\]
because $J\{a+b\le\frac12\}=\{b+c\ge\frac32\}$, a point of $E'$ has trace below $\frac32$, and $E'\cap JE'=\emptyset$ (on $E'$, $b<\frac12$, so $b+c<\frac32$). By \Cref{sec:third}, $\mu(\{r<1\})=\mu(\{r>2\})=1$, so the central band has mass $40$. On $E'$ put $k=a+b$; then $0<k<\frac12$, $0<a<\frac k2$ and $1-k<c<1$ (the conditions $b<c$ and $c<1$ being implied), and $9!v^2=\frac{9!}4(k-2a)^2(c-a)^2(c-k+a)^2$, so
\[
 \mu(E')=\frac{9!}{4}\int_0^{1/2}\int_{1-k}^1\int_0^{k/2}(k-2a)^2(c-a)^2(c-k+a)^2\dd a\dd c\dd k=\frac{105}{8}.
\]
Thus $\mu(B')=40-\frac{105}4=\frac{55}4$ and $\mu(B)=\frac{55}{1024}$.

\emph{The mass of $B\cap\mathcal R$.} On $\mathcal R$ we have $k<\frac{18-7r}{21}$, and $x+y>z$ is $k>\frac r2$; these are compatible only for $r<\frac{36}{35}$. For $1<r<\frac{36}{35}$ and $\frac r2<k<\frac{18-7r}{21}$ the remaining conditions of $B$ and $\Om$ reduce to $\frac14<x<\frac k2$: indeed $z=r-k<\frac r2<\frac34$; $y<z$ is $x>2k-r$, and $2k-r<\frac{36-35r}{21}<\frac14$; and $y+z-x<1$ is $x>\frac{r-1}2$, with $\frac{r-1}2<\frac14$. In the coordinates $(r,k,x)$ the $x$-integral is given by the formula in the proof of \Cref{thm:row} with $t=k-\frac12$, and
\[
 2\mu(B\cap\mathcal R)=\frac{134173355669}{1470888975150000000}.
\]
Hence $42-\mu(B)+2\mu(B\cap\mathcal R)=42-\frac{55}{1024}+\frac{134173355669}{1470888975150000000}=\frac{493586674117277079727}{11767111801200000000}$, and with the increment from $\mathcal F^+$ found above, $\mu(\mathcal U)=\frac{13860404849860091}{330432202525696}$. Both masses are computed by exact rational triangulation. This computation integrates $E'$ and $B$ directly and checks $\mu(E')=\frac{105}8$, $\mu(B)=\frac{55}{1024}$ and the value of $2\mu(B\cap\mathcal R)$. A second computation recomputes $2\mu(B\cap\mathcal R)$ by iterated polynomial integration, taking $\mu(B)=\frac{55}{1024}$ as input. Any open subset of $B\setminus(\mathcal R\cup J\mathcal R\cup\mathcal F^+\cup J\mathcal F^+)$, for instance a small ball about $(0.30,0.45,0.50)$, which has trace $1.25\in(\frac{15}{14},\frac32)$ and $k=0.75>\frac35$, carries positive mass that no profile in $L^2(\mathcal U)$ can see. \qed

\begin{remark}
A similar computation shows that the marked construction alone, granting every dominant-region pairing in both orientations, has ceiling $1339/32=41.84375$.
\end{remark}

\section{Formal verification in Lean}\label{sec:lean}
The finite parts of the proof have been checked in the Lean~4 proof assistant \cite{Lean4} with the Mathlib library \cite{Mathlib}: the masses $\mu(\Om)=42$ and $\mu(U_7)=C$, \Cref{prop:ceiling}, the passage from faithful pairings to \Cref{thm:main}, and the covering \cref{lem:cases}. The analytic theorems that produce the pairings are not formalised. They enter as a single explicit hypothesis. The development uses version~4.29.1 of Lean and the matching release of Mathlib, and is available at \coderepo.

\emph{Statements.} The formal statements define $M_6(T)$ with Mathlib's Riemann zeta function, the factor $a_3$ as a product over primes, and $C$ as the rational number of \Cref{thm:main}, and state \Cref{thm:main} as the assertion that $M_6(T)\ge c\,c_3T(\log T)^9$ for all large $T$, for every $c<C$. They define the chamber $\Om$, the function $v$, the measure $\mu$, the region $U_7$ and the region of \Cref{prop:ceiling}, following \Cref{sec:schur,sec:third,sec:marked,sec:balanced,sec:factorable,sec:mixedrow,sec:ceiling}. They also define tests, admissible weights, the smoothed numerators and Gram pairings, faithful families (\cref{def:faithful}) and closed spans. The hypothesis asserts that some faithful family has $v\one_{U_7}$ in its closed span. This is the conclusion of \Cref{thm:mixedrow}, which rests on \Cref{thm:third,thm:marked,thm:row,thm:factorable}. Lean proves the following.
\begin{itemize}\small
\item $\mu(\Om)=42$ and $\mu(U_7)=C$, exactly.
\item The region of \Cref{prop:ceiling} has mass $13860404849860091/330432202525696<42$, and $9!\|P_{\mathcal H}v\|^2$ is at most this mass for every family whose profiles vanish outside it. This is \Cref{prop:ceiling}.
\item \Cref{thm:main}, from the hypothesis above. Its proof formalises \cref{prop:projection} and \Cref{sec:main}.
\item \Cref{lem:cases} and the covering of the mark region that follows from it, for the affine exponent constraints of \Cref{sec:mixedrow}.
\end{itemize}
The formal statements were checked against the paper before the proofs were written. One correction was made afterwards: they had required the weights to be real analytic, which with their support condition left no nonzero weight, and they now require them to be infinitely differentiable, as in \Cref{sec:schur}.

\emph{Proof.} Lean proves the integration formula for polynomials over cells bounded by affine faces, and computes $\mu(\Om)$, $\mu(U_7)$ and the mass of the region of \Cref{prop:ceiling} by cutting each region into such cells, up to null affine boundaries. The seven cells of \Cref{tab:sevencells} are among them, and their masses agree exactly with the rational values computed by the programs. For \cref{lem:cases}, Lean checks the affine identities of \Cref{sec:mixedrow} by kernel computation, and the branches of part (iii) by Farkas certificates with nonnegative rational multipliers. These are the certificates used in the proof of \cref{lem:cases}, and a checker proved correct in Lean verifies them.

\emph{Trust.} Except for \cref{lem:cases}, all of the above depends only on the three standard axioms of Lean, namely propositional extensionality, the axiom of choice and the soundness of quotients, with every computation carried out in the kernel. The proof of \cref{lem:cases} also uses three evaluations by Lean's \texttt{native\_decide} tactic, which runs compiled code, so accepting it means trusting the Lean compiler for those evaluations, which take a few seconds, as well as the Lean kernel. The development contains no unproved steps and declares no further axioms.

\begin{thebibliography}{99}\small
\bibitem{CG} J.\,B.~Conrey and A.~Ghosh, \emph{A conjecture for the sixth power moment of the Riemann zeta-function}, Int.\ Math.\ Res.\ Not.\ (1998), 775--780.
\bibitem{KS} J.\,P.~Keating and N.\,C.~Snaith, \emph{Random matrix theory and $\zeta(1/2+it)$}, Comm.\ Math.\ Phys.\ \textbf{214} (2000), 57--89.
\bibitem{CFKRS} J.\,B.~Conrey, D.\,W.~Farmer, J.\,P.~Keating, M.\,O.~Rubinstein and N.\,C.~Snaith, \emph{Integral moments of $L$-functions}, Proc.\ London Math.\ Soc.\ \textbf{91} (2005), 33--104.
\bibitem{RS} M.~Radziwi\l\l\ and K.~Soundararajan, \emph{Continuous lower bounds for moments of zeta and $L$-functions}, Mathematika \textbf{59} (2013), 119--128.
\bibitem{HS} W.~Heap and K.~Soundararajan, \emph{Lower bounds for moments of zeta and $L$-functions revisited}, Mathematika \textbf{68} (2022), 1--14.
\bibitem{Harper} A.\,J.~Harper, \emph{Sharp conditional bounds for moments of the Riemann zeta function}, arXiv:1305.4618.
\bibitem{Ng} N.~Ng, \emph{The sixth moment of the Riemann zeta function and ternary additive divisor sums}, Discrete Anal.\ (2021); arXiv:1610.04977.
\bibitem{DP} B.~Durkan and T.~Page, \emph{Amplified moments of the Riemann zeta function}, arXiv:2606.27323v2, 2026.
\bibitem{Fraccaro} N.~Fraccaro, hexagonmath.org/2610.00634v3; Lean artifacts at github.com/Niklas-Fraccaro/sixth-moment-lean.
\bibitem{Mac} I.\,G.~Macdonald, \emph{Symmetric functions and Hall polynomials}, 2nd ed., Oxford Univ.\ Press, 1995.
\bibitem{St} R.\,P.~Stanley, \emph{Enumerative combinatorics}, Vol.~2, Cambridge Univ.\ Press, 1999.
\bibitem{BG} D.~Bump and A.~Gamburd, \emph{On the averages of characteristic polynomials from classical groups}, Comm.\ Math.\ Phys.\ \textbf{265} (2006), 227--274.
\bibitem{Anderson} G.\,W.~Anderson, \emph{A short proof of Selberg's generalized beta formula}, Forum Math.\ \textbf{3} (1991), 415--417.
\bibitem{Forrester} P.\,J.~Forrester, \emph{Log-gases and random matrices}, Princeton Univ.\ Press, 2010.
\bibitem{HB} D.\,R.~Heath-Brown, \emph{The divisor function $d_3(n)$ in arithmetic progressions}, Acta Arith.\ \textbf{47} (1986), 29--56.
\bibitem{DI} J.-M.~Deshouillers and H.~Iwaniec, \emph{Kloosterman sums and Fourier coefficients of cusp forms}, Invent.\ Math.\ \textbf{70} (1982), 219--288.
\bibitem{Drappeau} S.~Drappeau, \emph{Sums of Kloosterman sums in arithmetic progressions, and the error term in the dispersion method}, Proc. Lond. Math. Soc. (3) \textbf{114} (2017), 684--732; arXiv:1504.05549.
\bibitem{Watt} N.~Watt, \emph{Kloosterman sums and a mean value for Dirichlet polynomials}, J.\ Number Theory \textbf{53} (1995), doi:10.1006/jnth.1995.1086.
\bibitem{Alpoge} L.~Alp\"oge, \emph{Remark}, 20 November 2026, \url{https://alpo.ge/remark.pdf}. Formal Lean proof of the Riemann hypothesis and of GRH for Dirichlet $L$-functions: \url{https://github.com/leanprover-community/riemann-hypothesis}. Independent verification: A.\,V.~Sutherland, \url{https://math.mit.edu/~drew/grh/}; M.~Rubinstein, \url{https://www.math.uwaterloo.ca/~mrubinst/alpoge_claude_grh/}.
\bibitem{Apostol} T.\,M.~Apostol, \emph{Introduction to analytic number theory}, Springer, 1976.
\bibitem{AtkinLehner} A.\,O.\,L.~Atkin and J.~Lehner, \emph{Hecke operators on $\Gamma_0(m)$}, Math.\ Ann.\ \textbf{185} (1970), 134--160.
\bibitem{BM} V.~Blomer and D.~Mili\'cevi\'c, \emph{Kloosterman sums in residue classes}, J.\ Eur.\ Math.\ Soc.\ \textbf{17} (2015), 51--69.
\bibitem{Humphries} P.~Humphries, \emph{Density theorems for exceptional eigenvalues for congruence subgroups}, Algebra Number Theory \textbf{12} (2018), 1581--1610.
\bibitem{IwSpec} H.~Iwaniec, \emph{Spectral methods of automorphic forms}, 2nd ed., Grad.\ Stud.\ Math.\ \textbf{53}, Amer.\ Math.\ Soc., 2002.
\bibitem{IK} H.~Iwaniec and E.~Kowalski, \emph{Analytic number theory}, Amer.\ Math.\ Soc.\ Colloq.\ Publ.\ \textbf{53}, 2004.
\bibitem{KimSarnak} H.\,H.~Kim, \emph{Functoriality for the exterior square of $\mathrm{GL}_4$ and the symmetric fourth of $\mathrm{GL}_2$}, with Appendix~1 by D.~Ramakrishnan and Appendix~2 by H.\,H.~Kim and P.~Sarnak, J.\ Amer.\ Math.\ Soc.\ \textbf{16} (2003), 139--183.
\bibitem{KY} E.\,M.~K\i ral and M.\,P.~Young, \emph{The fifth moment of modular $L$-functions}, J.\ Eur.\ Math.\ Soc.\ \textbf{23} (2021), 237--314.
\bibitem{Lichtman} J.~D.~Lichtman, \emph{Primes in arithmetic progressions to large moduli and shifted primes without large prime factors}, arXiv:2211.09641.
\bibitem{Pascadi} A.~Pascadi, \emph{Large sieve inequalities for exceptional Maass forms and the greatest prime factor of $n^2+1$}, Forum Math.\ Pi \textbf{14} (2026), e8, doi:10.1017/fmp.2026.10025; arXiv:2404.04239.
\bibitem{PascadiExp} A.~Pascadi, \emph{On the exponents of distribution of primes and smooth numbers}, arXiv:2505.00653v1.
\bibitem{SPY} R.~Schulze-Pillot and A.~Yenirce, \emph{Petersson products of bases of spaces of cusp forms and estimates for Fourier coefficients}, Int.\ J.\ Number Theory \textbf{14} (2018), 2277--2290.
\bibitem{Deligne} P.~Deligne, \emph{La conjecture de Weil.~II}, Publ.\ Math.\ IH\'ES \textbf{52} (1980), 137--252.
\bibitem{GKM} N.\,M.~Katz, \emph{Gauss sums, Kloosterman sums, and monodromy groups}, Ann.\ of Math.\ Stud.\ \textbf{116}, Princeton Univ.\ Press, 1988.
\bibitem{ESDE} N.\,M.~Katz, \emph{Exponential sums and differential equations}, Ann.\ of Math.\ Stud.\ \textbf{124}, Princeton Univ.\ Press, 1990.
\bibitem{FKM} \'E.~Fouvry, E.~Kowalski and Ph.~Michel, \emph{A study in sums of products}, Philos.\ Trans.\ Roy.\ Soc.\ A \textbf{373} (2015), 20140309.
\bibitem{Titchmarsh} E.\,C.~Titchmarsh, \emph{The theory of the Riemann zeta-function}, 2nd ed., revised by D.\,R.~Heath-Brown, Oxford Univ.\ Press, 1986.
\bibitem{Lean4} L.~de~Moura and S.~Ullrich, \emph{The Lean~4 theorem prover and programming language}, in: Automated Deduction -- CADE~28, Lecture Notes in Comput.\ Sci.\ \textbf{12699}, Springer, 2021, 625--635.
\bibitem{Mathlib} The mathlib Community, \emph{The Lean mathematical library}, in: Proceedings of the 9th ACM SIGPLAN International Conference on Certified Programs and Proofs (CPP 2020), ACM, 2020, 367--381.
\end{thebibliography}

\end{document}
